\documentclass[11pt,twoside,openany]{book}
\usepackage[a4paper,inner=31mm,outer=27mm,top=28mm,bottom=31mm,marginparwidth=17mm]{geometry}
\usepackage{fontspec}
\defaultfontfeatures{Ligatures=TeX,Scale=MatchLowercase}
\setmainfont{TeX Gyre Pagella}
\setsansfont{TeX Gyre Heros}
\setmonofont{TeX Gyre Cursor}
\usepackage[brazilian]{babel}
\usepackage{microtype}
\usepackage{graphicx}
\usepackage{mathtools}
\usepackage{amsthm}
\usepackage{unicode-math}
\setmathfont{Latin Modern Math}
\usepackage{bookmark}
\usepackage{hyperref}
\hypersetup{hidelinks,unicode=true}

\setcounter{secnumdepth}{3}
\setcounter{tocdepth}{2}
\numberwithin{equation}{section}

\newcommand{\src}[1]{%
  \phantomsection\hypertarget{src:#1}{}%
  \marginpar{\raggedleft\footnotesize\sffamily #1}%
}
\newcommand{\sid}[1]{\phantomsection\hypertarget{fga:#1}{}}
\newcommand{\Osh}{\underline{O}}
\newcommand{\Zsh}{\underline{\mathbf Z}}
\newcommand{\Ccat}{\mathcal C}
\newcommand{\Homsh}{\underline{\operatorname{Hom}}}
\newcommand{\Extsh}{\underline{\operatorname{Ext}}}
\newcommand{\hsh}{\underline{h}}
\newcommand{\fquote}[1]{\textquotedbl#1\textquotedbl}
\newcommand{\seminar}[1]{%
  \begin{flushleft}\sffamily SEMINÁRIO BOURBAKI\\#1\end{flushleft}%
}
\newcommand{\seminares}[1]{\seminar{#1}}
\newcommand{\paperhead}[2]{%
  \begin{center}
    {\Large\bfseries #1\par}
    \vspace{0.6\baselineskip}
    #2
  \end{center}
  \vspace{0.8\baselineskip}
}

\theoremstyle{plain}
\newtheorem{proposition}{Proposição}[section]
\newtheorem{theorem}[proposition]{Teorema}
\newtheorem{corollary}[proposition]{Corolário}
\newtheorem{lemma}[proposition]{Lema}

\theoremstyle{definition}
\newtheorem{definition}[proposition]{Definição}

\theoremstyle{remark}
\newtheorem*{remark}{Observação}

\newcommand{\unit}[1]{%
  \chapter{#1}%
  \markboth{#1}{#1}%
}
\newcommand{\fsection}[2]{%
  \section*{#1. #2}%
  \addcontentsline{toc}{section}{#1. #2}%
}
\newcommand{\result}[2]{%
  \par\medskip\noindent\textbf{#1%
  \if\relax\detokenize{#2}\relax\else\space#2\fi. ---}\enspace
}

\pdfstringdefDisableCommands{\def\textsuperscript#1{#1}}

% As páginas impressas reiniciam em diferentes exposições e erratas.
% A chave apenas distingue os destinos invisíveis; não altera os textos.
\newcommand{\FgaUnitKey}{frontmatter}
\renewcommand{\src}[1]{%
  \phantomsection\hypertarget{src:\FgaUnitKey:#1}{}%
  \marginpar{\raggedleft\footnotesize\sffamily #1}%
}
\setlength{\emergencystretch}{3em}
\renewcommand{\boldsymbol}[1]{\symbfit{#1}}
\hypersetup{pdftitle={Fundamentos da geometria algébrica -- Tradução com leituras do impresso},pdfauthor={Alexandre Grothendieck},pdfsubject={Oito exposições, suas erratas e os comentários de 1962}}
\begin{document}
\frontmatter
\begin{titlepage}
\centering
\vspace*{0.13\textheight}
{\Huge\bfseries Fundamentos da\par geometria algébrica\par}
\vspace{1.5\baselineskip}
{\Large Alexandre Grothendieck\par}
\vfill
{\Large Tradução com leituras do impresso\par}
\vspace{\baselineskip}
{\large Oito exposições, suas erratas e os comentários de 1962\par}
\vspace{1.5\baselineskip}
{\normalsize Edição de 24 de outubro de 2026\par}
\end{titlepage}
\chapter*{Nota sobre esta edição}
\addcontentsline{toc}{chapter}{Nota sobre esta edição}
Este volume contém as exposições 149, 182, 190, 195, 212, 221, 232 e 236, suas erratas e os comentários de 1962. O conjunto não constitui uma edição diplomática estrita: algumas leituras herdadas diferem da reprodução digital do NUMDAM. A edição corrigida é apresentada separadamente.

Os números marginais remetem às páginas das exposições originais.

A obra original, Fondements de la géométrie algébrique, é de Alexandre Grothendieck. A tradução humana para o inglês de \href{https://github.com/thosgood/fga/tree/9976dd5f8e2ddcb388f0a07f3aa334f120d1019f}{Tim Hosgood} constitui uma fonte de referência desta edição. Seu projeto reconhece as correções de Steve Hnizdur e o apoio técnico de Jon Sterling. A versão 9976dd5f8e2ddcb388f0a07f3aa334f120d1019f identifica o texto inglês usado na comparação de procedência; não estabelece a versão histórica exata incorporada inicialmente. A tradução brasileira e as correções e composição posteriores foram produzidas por IA. A autoria da tradução inglesa permanece com seu tradutor; não se lhe atribui a tradução brasileira. Esta atribuição não constitui autorização nem endosso.

A tradução para o português é produzida por IA. As correções, a edição e a composição também são produzidas por IA. O responsável pela publicação não traduziu o texto. A autoria e as atribuições das obras originais e das contribuições herdadas são preservadas; não são atribuídas à IA. As correções documentadas desta edição incluem OpenAI Codex, GPT-6 Astra, esforço Ultra.

Os direitos sobre as obras originais e sobre as contribuições protegidas a estas traduções e transcrições derivadas permanecem com seus respectivos titulares. Na medida em que o editor detenha direitos sobre a tradução, a transcrição, as correções e a composição produzidas por IA, ele os dedica ao domínio público sob CC0 1.0 Universal (https://creativecommons.org/publicdomain/zero/1.0/deed.pt-br); isso não afeta os direitos sobre as obras originais nem concede permissão para usá-las. Os avisos de direitos de terceiros e as autorizações já concedidas continuam aplicáveis.

Estas edições são disponibilizadas gratuitamente, para fins não comerciais e exclusivamente educativos, sob a presunção de que os respectivos titulares dos direitos não se oporiam à disponibilidade gratuita destas traduções e transcrições, ao lado dos originais franceses no \href{https://www.numdam.org/}{NUMDAM}. Essa presunção não significa que uma autorização tenha sido obtida; o acesso gratuito aos originais não autoriza, por si só, sua tradução, adaptação ou republicação.

Se você detém direitos sobre algum material apresentado aqui e se opõe à sua disponibilidade, poderá \href{https://github.com/KokunoYumeto/ega-fga-sga-pt-br/issues}{abrir uma questão no GitHub}. O responsável pela publicação retirará o material afetado a pedido do titular dos direitos.


\tableofcontents
\mainmatter
\def\FgaUnitKey{149}\unit{Exposição 149 — Teoremas de dualidade para os feixes algébricos coerentes}
\sid{149}
\markboth{Exposição 149}{Exposição 149}
\src{169}
\seminar{(Maio de 1957)}
\paperhead{TEOREMAS DE DUALIDADE PARA OS FEIXES ALGÉBRICOS COHERENTES}%
  {por Alexandre GROTHENDIECK}

Os resultados que seguem, inspirados no \fquote{teorema de dualidade algébrica}
de Serre, foram obtidos nos invernos de 1955 e 1956. Demonstram-se muito
simplesmente com o auxílio de resultados bastante elementares sobre a cohomologia dos espaços
projetivos [3] e mediante o uso intensivo da álgebra homológica de
Cartan--Eilenberg, na forma de [2].

\fsection{1}{Os Ext de feixes de módulos ({[2]}, caps. 3 e 4)}

Seja $X$ um espaço topológico provido de um feixe $\Osh$ de anéis com unidade
(não necessariamente comutativos). Consideramos a categoria abeliana
$\Ccat^{\Osh}$ dos feixes de $\Osh$-módulos, chamados também $\Osh$-Módulos.
Sabe-se que todo objeto desta categoria admite uma resolução injetiva, o que
permite definir nela os funtores $\operatorname{Ext}$ com as propriedades
formais bem conhecidas. Mais precisamente, para evitar confusões,
denotamos por $\operatorname{Hom}_{\Osh}(X;A,B)$, ou simplesmente por
$\operatorname{Hom}(X;A,B)$, o grupo abeliano dos $\Osh$-homomorfismos de $A$
em $B$, enquanto que $\Homsh_{\Osh}(A,B)$ denotará o feixe de germes
de homomorfismos de $A$ em $B$ ($A,B\in\Ccat^{\Osh}$). Para
$A\in\Ccat^{\Osh}$ fixo, definem-se funtores $h_A$ e $\hsh_A$ com valores, respectivamente,
na categoria $\Ccat$ dos grupos abelianos e na categoria
$\Ccat^{\Zsh}$ dos feixes abelianos sobre $X$, mediante as fórmulas:
\begin{equation}\tag{1.1}
  h_A(B)=\operatorname{Hom}_{\Osh}(X;A,B),
  \qquad
  \hsh_A(B)=\Homsh_{\Osh}(A,B).
\end{equation}

Os funtores $h_A$ e $\hsh_A$ são funtores covariantes exatos pela esquerda,
dos quais consideramos os funtores derivados à direita, denotados respectivamente por
$\operatorname{Ext}_{\Osh}^{p}(X;A,B)$ e $\Extsh_{\Osh}^{p}(A,B)$. Portanto,
por definição,
\begin{equation}\tag{1.2}
\left\{
\begin{aligned}
\operatorname{Ext}_{\Osh}^{p}(X;A,B)
  &= (R^{p}h_A)(B)
   = H^{p}\bigl(\operatorname{Hom}_{\Osh}(X;A,C(B))\bigr),\\
\Extsh_{\Osh}^{p}(A,B)
  &= (R^{p}\hsh_A)(B)
   = H^{p}\bigl(\Homsh_{\Osh}(A,C(B))\bigr),
\end{aligned}
\right.
\end{equation}

onde o símbolo $R^{p}$ designa a passagem aos funtores derivados à direita, e onde
$C(B)$ designa uma resolução injetiva arbitrária de $B$ em $\Ccat^{\Osh}$.
Denotemos por
\[
  \Gamma:\Ccat^{\Zsh}\longrightarrow\Ccat
\]
o funtor \fquote{seções}. Lembremos que seus funtores derivados à direita são
denotados por
\[
  B\longmapsto H^{p}(X,B).
\]

\src{170}
\begin{equation}\tag{1.3}
  H^{p}(X,B)=(R^{p}\Gamma)(B)=H^{p}\bigl(\Gamma(C(B))\bigr).
\end{equation}
Evidentemente, tem-se $h_A=\Gamma\hsh_A$; por outro lado, verifica-se que $\hsh_A$
transforma objetos injetivos em objetos acíclicos para $\Gamma$. Conclui-se do
modo habitual:

\result{PROPOSIÇÃO}{1} Para todo $\Osh$-Módulo $A$ existe um funtor
espectral cohomológico sobre $\Ccat^{\Osh}$, que converge ao funtor graduado
$\bigl(\operatorname{Ext}_{\Osh}^{*}(X;A,B)\bigr)$ e cujo termo inicial é
\begin{equation}\tag{1.4}
  E_{2}^{p,q}(A,B)=H^{p}\bigl(X,\Extsh_{\Osh}^{q}(A,B)\bigr).
\end{equation}
Dela se deduzem \fquote{homomorfismos de bordo} e uma sequência exata de cinco termos
que não escreveremos.

\result{COROLÁRIO}{1} Se $A$ é localmente isomorfo a $\Osh^{n}$, então têm-se
isomorfismos canônicos
\begin{equation}\tag{1.5}
  \operatorname{Ext}_{\Osh}^{p}(X;A,B)
  \overset{\sim}{\longleftarrow}
  H^{p}\bigl(X,\Homsh_{\Osh}(A,B)\bigr)
\end{equation}
(dados por homomorfismos de bordo da sequência espectral). Em particular, tem-se
um isomorfismo canônico
\begin{equation}\tag{1.6}
  \operatorname{Ext}_{\Osh}^{p}(X;\Osh,B)=H^{p}(X,B).
\end{equation}

Para utilizar estes resultados, é preciso determinar explicitamente os
$\Extsh_{\Osh}^{p}(A,B)$. São funtores que se calculam localmente,
isto é, se $U$ é um aberto de $X$, tem-se
\[
  \Extsh_{\Osh}^{p}(A,B)|U
  =\Extsh_{\Osh|U}^{p}(A|U,B|U),
\]
como resulta do fato de que a restrição a $U$ de um $\Osh$-Módulo injetivo
é um $(\Osh|U)$-Módulo injetivo. Além disso, para $x\in X$ fixo, há homomorfismos funtoriais
\begin{equation}\tag{1.7}
  \Homsh_{\Osh}(A,B)_x
  \longrightarrow
  \operatorname{Hom}_{\Osh_x}(A_x,B_x),
\end{equation}
que se prolongam de maneira única a um homomorfismo de funtores cohomológicos
em $B$:
\begin{equation}\tag{1.8}
  \Extsh_{\Osh}^{p}(A,B)_x
  \longrightarrow
  \operatorname{Ext}_{\Osh_x}^{p}(A_x,B_x).
\end{equation}

\result{PROPOSIÇÃO}{2} Se $A$ é, em uma vizinhança de $x$, isomorfo ao
conúcleo de um homomorfismo $\Osh^{m}\longrightarrow\Osh^{n}$, então (1.7) é um
isomorfismo para todo $p$. Isto ocorre, em
particular, se $A$ é um $\Osh$-módulo coerente [3].

\src{171}
\result{PROPOSIÇÃO}{3} Seja $L_{*}=(L_i)$ uma resolução pela esquerda do
$\Osh$-Módulo $A$ por $\Osh$-Módulos cada um dos quais é localmente isomorfo a
um $\Osh^{n}$. Então $\Extsh_{\Osh}^{*}(A,B)$ identifica-se com
$H^{*}\bigl(\Homsh_{\Osh}(L_{*},B)\bigr)$, e
$\operatorname{Ext}_{\Osh}^{*}(X;A,B)$ identifica-se com a hipercohomologia de $X$
com respeito ao complexo $\Homsh_{\Osh}(L_{*},B)$.

A demonstração é padrão: considera-se o bicomplexo
$\Homsh_{\Osh}(L_{*},C(B))$, onde $C(B)$ é uma resolução injetiva de $B$, e os
homomorfismos naturais de $\Homsh_{\Osh}(L_{*},B)$ e de
$\Homsh_{\Osh}(A,C(B))$ neste último.

Para terminar, assinalemos que os dois funtores $\operatorname{Ext}$ introduzidos
em (1.2) não são apenas funtores cohomológicos em $B$, mas também
bifuntores cohomológicos, covariantes em $B$ e contravariantes em $A$.

\fsection{2}{Lei de composição nos Ext}

Os resultados deste número devem-se, independentemente, a CARTIER e a YONEDA; veja-se
uma exposição de CARTIER [1] para mais detalhes. Seja $\Ccat$ uma categoria abeliana.
Sejam $K$ e $L$ dois objetos graduados de $\Ccat$. Denotamos por
$\operatorname{Hom}(K,L)$ o grupo abeliano graduado cuja componente de grau
$n$ é formada pelos homomorfismos homogêneos de grau $n$ de $K$ em $L$ (isto é,
os sistemas $(u_i)$ de homomorfismos $K^{i}\longrightarrow L^{i+n}$).

Se $K$ e $L$ são complexos (com operadores diferenciais de grau $+1$, para
fixar as ideias), introduz-se em $\operatorname{Hom}(K,L)$ o operador
diferencial dado por
\begin{equation}\tag{2.1}
  \delta(u)=du+(-1)^{n+1}ud,\qquad\text{onde }n=\deg(u),
\end{equation}
que o converte em um complexo com operador diferencial de grau $+1$. Os ciclos de
grau $n$ são os operadores de grau $n$ que anticomutam com $u$ (como
operadores homogêneos). Pode-se então considerar
$H^{*}\bigl(\operatorname{Hom}(K,L)\bigr)$; este é um invariante dos tipos
de homotopia de $K$ e de $L$, que poderemos denotar por $H^{*}(K,L)$.

Dado um terceiro complexo $M$, a composição dos homomorfismos
define um acoplamento
\[
  \operatorname{Hom}(K,L)\times\operatorname{Hom}(L,M)
  \longrightarrow\operatorname{Hom}(K,M)
\]
compatível com os operadores diferenciais, de onde, ao passar à cohomologia,
obtêm-se acoplamentos
\begin{equation}\tag{2.2}
  H^{*}(K,L)\times H^{*}(L,M)\longrightarrow H^{*}(K,M),
\end{equation}
denotados por $(u,v)\longmapsto vu$. Estes acoplamentos satisfazem a evidente propriedade
de associatividade. Em particular, $H^{*}(K,K)$ é um anel graduado associativo
com unidade, e $H^{*}(K,L)$ (resp. $H^{*}(L,K)$) é um módulo graduado pela direita
(resp. pela esquerda) sobre este anel, etc. Em grau $0$, (2.2) reduz-se à
composição dos morfismos de complexos. Finalmente, uma sequência exata de
complexos
\[
  0\longrightarrow K'\longrightarrow K\longrightarrow K''\longrightarrow 0,
\]
\src{172}
tal que $K'^{i}$ se identifica com um somando direto de $K^{i}$ para todo $i$,
dá lugar a uma sequência exata de complexos de grupos
$\operatorname{Hom}(K'',L)$, etc., e daí a um operador de conexão
\[
  H^{i}(K',L)\longrightarrow H^{i+1}(K'',L).
\]
Definem-se de maneira análoga os operadores de conexão relativos a uma sequência exata em $L$. Os
acoplamentos (2.2) são compatíveis, no sentido usual, com estes operadores de conexão.

Suponhamos agora que $\Ccat$ é uma categoria tal que todo objeto $A$
de $\Ccat$ admite uma resolução injetiva $C(A)$. Verifica-se então, utilizando
uma das numerosas variantes do teorema do bicomplexo, que
\[
  H^{*}\bigl(C(A),C(B)\bigr)
  =H^{*}\bigl(\operatorname{Hom}(C(A),C(B))\bigr)
\]
é canonicamente isomorfo a
\[
  H^{*}\bigl(\operatorname{Hom}(A,C(B))\bigr)
  =\operatorname{Ext}^{*}(A,B).
\]
Os operadores de conexão considerados anteriormente dão os homomorfismos de conexão dos
$\operatorname{Ext}$. Além disso, os acoplamentos (2.2) dão aqui
acoplamentos associativos, compatíveis com os homomorfismos de conexão:
\begin{equation}\tag{2.3}
  \operatorname{Ext}^{*}(A,B)\times\operatorname{Ext}^{*}(B,C)
  \longrightarrow\operatorname{Ext}^{*}(A,C).
\end{equation}
Em particular, $\operatorname{Ext}^{*}(A,A)$ é um anel graduado associativo
com unidade, etc. (Demonstra-se de maneira análoga que os funtores
$\operatorname{Ext}$ operam nos funtores derivados de um funtor
qualquer; não utilizaremos aqui este fato.)

No caso em que a categoria considerada seja a categoria $\Ccat^{\Osh}$ dos
$\Osh$-Módulos sobre $X$, obtêm-se, portanto, acoplamentos
\begin{equation}\tag{2.4}
  \operatorname{Ext}_{\Osh}^{p}(X;A,B)
  \times\operatorname{Ext}_{\Osh}^{q}(X;B,C)
  \longrightarrow\operatorname{Ext}_{\Osh}^{p+q}(X;A,C),
\end{equation}
que podem ser explicitados como foi indicado. De resto, os mesmos
desenvolvimentos, substituindo a categoria dos grupos abelianos pela
categoria dos feixes abelianos sobre $X$, e os funtores $\operatorname{Hom}$
pelos funtores $\Homsh$, definem também acoplamentos com as mesmas
propriedades formais e, desta vez, de \fquote{natureza local}:
\begin{equation}\tag{2.5}
  \Extsh_{\Osh}^{p}(A,B)\times\Extsh_{\Osh}^{q}(B,C)
  \longrightarrow\Extsh_{\Osh}^{p+q}(A,C).
\end{equation}
Estes últimos são determinados observando-se que os homomorfismos (1.8) são
compatíveis com os acoplamentos entre os $\operatorname{Ext}$.

\src{173}
Lembremos, por fim, que existe também uma estrutura multiplicativa entre os funtores
$H^{p}(X,A)$ (produto cup). Verifica-se então que as sequências espectrais da
proposição 1 são compatíveis com as estruturas multiplicativas; mais
precisamente, tem-se um acoplamento da sequência espectral $E(A,B)$ e da sequência
espectral $E(B,C)$ para a sequência espectral $E(A,C)$ que, no limite,
se reduz ao acoplamento entre os $\operatorname{Ext}$ globais e, no
termo inicial, ao acoplamento deduzido, no segundo membro de (1.4), do
produto cup e dos acoplamentos dos $\operatorname{Ext}$ locais. Disso resulta,
em particular, que os \fquote{homomorfismos de bordo}
\begin{equation}\tag{2.6}
  \operatorname{Ext}_{\Osh}^{n}(X;A,B)
  \longrightarrow H^{0}\bigl(X,\Extsh_{\Osh}^{n}(A,B)\bigr)
\end{equation}
e
\begin{equation}\tag{2.7}
  H^{n}\bigl(X,\Homsh_{\Osh}(A,B)\bigr)
  \longrightarrow\operatorname{Ext}_{\Osh}^{n}(X;A,B)
\end{equation}
são compatíveis com as estruturas multiplicativas. Portanto, se nos restringirmos aos
feixes localmente isomorfos a um $\Osh^{m}$, isto explicita por completo a
composição entre os $\operatorname{Ext}$ globais mediante o produto cup, levando
em conta o isomorfismo (1.5).

\fsection{3}{Resultados de cohomologia local}

Seja $A$ um anel comutativo com unidade provido de um ideal $J$. Vamos
definir, para um $A$-módulo variável $M$, homomorfismos funtoriais
\begin{equation}\tag{3.1}
\left\{
\begin{aligned}
\operatorname{Ext}_{A}^{p}(A/J,M)
  &\longrightarrow
  \operatorname{Hom}_{A}\!\left(\bigwedge^{p}J/J^{2},M\otimes A/J\right),\\
\operatorname{Tor}_{p}^{A}(A/J,M)
  &\longleftarrow
  \left(\bigwedge^{p}J/J^{2}\right)
  \otimes\operatorname{Hom}_{A}(A/J,M),
\end{aligned}
\right.
\end{equation}
onde os produtos tensoriais e exteriores são tomados sobre o anel $A$; note-se,
por outro lado, que $J/J^{2}$ é, na realidade, um $A/J$-módulo, e que suas potências
exteriores como $A$-módulo ou como $A/J$-módulo são as mesmas.
Definir os homomorfismos (3.1) equivale a definir, para todo
sistema $\underline{x}=(x_1,\ldots,x_p)$ de elementos de $J$, homomorfismos
$\varphi_{\underline{x}}$:
\begin{equation}\tag{3.2}
\left\{
\begin{aligned}
\varphi_{\underline{x}}:
  \operatorname{Ext}_{A}^{p}(A/J,M)&\longrightarrow M\otimes A/J,\\
\varphi_{\underline{x}}:
  \operatorname{Hom}_{A}(A/J,M)&\longrightarrow
  \operatorname{Tor}_{p}^{A}(A/J,M),
\end{aligned}
\right.
\end{equation}
de tal maneira que sejam satisfeitas as condições seguintes:

\src{174}
\begin{enumerate}
\item[\textup{i.}] $\varphi_{x_1,\ldots,x_p}$ depende de maneira
$A$-multilinear alternada do sistema dos $x_i\in J$.
\item[\textup{ii.}] $\varphi_{x_1,\ldots,x_p}$ é nulo quando algum dos $x_i$ está
em $J^{2}$.
\end{enumerate}
Na realidade, ii. resulta de i., pois $a\varphi_{\underline{x}}=0$ para $a\in J$,
como se vê observando que todos os módulos de (3.2) são anulados por
$J$.

Para definir $\varphi_{\underline{x}}$, consideremos o complexo
$K_{\underline{x}}$ cujo $A$-módulo subjacente é $\bigwedge A^{p}$, e
cujo operador diferencial é o produto interior $i_{\underline{x}}$ por
$\underline{x}$ considerado como forma linear sobre $A^{p}$ de componentes
$x_1,\ldots,x_p$. O operador diferencial é de grau $-1$, os graus são
positivos, e a cohomologia deste complexo em grau $0$ é
$A/(x_1A+\cdots+x_pA)$. Como os $x_i$ estão em $J$, deduz-se uma
aumentação $K_{\underline{x},0}\longrightarrow A/J$. Assim,
$K_{\underline{x}}$ aparece como um complexo livre aumentado, com módulo
de aumentação $A/J$. Disso se deduzem, do modo habitual, homomorfismos
\[
\operatorname{Ext}_{A}^{*}\bigl(H_0(K_{\underline{x}}),M\bigr)
\longrightarrow
H^{*}\bigl(\operatorname{Hom}_{A}(K_{\underline{x}},M)\bigr)
\]
e
\[
\operatorname{Tor}_{*}^{A}\bigl(H_0(K_{\underline{x}}),M\bigr)
\longleftarrow
H_{*}\bigl(K_{\underline{x}}\otimes_A M\bigr),
\]
de onde, compondo com os homomorfismos sobre os $\operatorname{Ext}$ e os
$\operatorname{Tor}$ deduzidos do homomorfismo de aumentação
$H_0(K_{\underline{x}})\longrightarrow A/J$, obtêm-se homomorfismos
\begin{equation}\tag{3.3}
\left\{
\begin{aligned}
\psi_{\underline{x}}:
\operatorname{Ext}_{A}^{*}(A/J,M)
&\longrightarrow H^{*}\bigl(\operatorname{Hom}_{A}(K_{\underline{x}},M)\bigr),\\
\psi_{\underline{x}}:
\operatorname{Tor}_{*}^{A}(A/J,M)
&\longleftarrow H_{*}\bigl(K_{\underline{x}}\otimes_A M\bigr).
\end{aligned}
\right.
\end{equation}
Ora, verifica-se imediatamente que, em grau máximo $p$, a cohomologia dos segundos
membros é $M/(x_1M+\cdots+x_pM)$ (resp. é o conjunto dos elementos de $M$
anulados por cada um dos $x_i$). Como os $x_i$ estão em $J$, deduzem-se
homomorfismos
\begin{equation}\tag{3.4}
\left\{
\begin{aligned}
H^{p}\bigl(\operatorname{Hom}_{A}(K_{\underline{x}},M)\bigr)
&\longrightarrow M\otimes A/J,\\
H_{p}\bigl(K_{\underline{x}}\otimes_A M\bigr)
&\longleftarrow\operatorname{Hom}_{A}(A/J,M).
\end{aligned}
\right.
\end{equation}
Compondo os homomorfismos (3.3) e (3.4), obtêm-se os homomorfismos (3.2)
que queríamos definir. A verificação de i. é laboriosa, mas não
apresenta dificuldades.

\src{175}
\result{PROPOSIÇÃO}{4} Seja $A$ um anel comutativo com unidade, e seja
$(x_1,\ldots,x_p)$ uma sequência de elementos de $A$ tal que, para $1\leq i\leq p$,
a imagem de $x_i$ no quociente de $A$ pelo ideal gerado por
$(x_1,\ldots,x_{i-1})$ não seja divisor de $0$. Seja $J$ o ideal gerado
pelos $x_i$. Então $J/J^{2}$ é um $(A/J)$-módulo livre que tem por base as
imagens canônicas dos $x_i$, o complexo $K_{\underline{x}}$ é uma resolução
livre de $A/J$, e para todo $A$-módulo $M$, os homomorfismos (3.1) em
grau $p$ são bijetivos. O mesmo ocorre com os homomorfismos análogos
definidos para graus $i$ quaisquer, sempre que $J.M=0$.

(O ponto essencial, do qual resultam todos os demais, é a aciclicidade de
$K_{\underline{x}}$, que é um fato bem conhecido sob as condições indicadas.)

\result{COROLÁRIO}{1} Sejam $A$ e $J$ como acima, e suponhamos além disso que
$A$ é uma álgebra afim regular de $\dim n$ sobre um corpo perfeito $k$, e
que $A/J$ é uma álgebra afim regular. Denotemos por $\Omega^{1}(A)$,
$\Omega^{1}(A/J)$ os módulos de diferenciais de Kähler. Então tem-se um
isomorfismo canônico compatível com a localização,
\begin{equation}\tag{3.5}
  \operatorname{Ext}_{A}^{p}\bigl(\Omega^{n-p}(A/J),\Omega^{n}(A)\bigr)=A/J.
\end{equation}
Com efeito, $\Omega^{n-p}(A/J)$ é um $(A/J)$-módulo livre de posto $1$, e igualmente
$\Omega^{n}(A)$ é um $A$-módulo livre de posto $n$,
pelo que o primeiro membro é igual a
\[
  \operatorname{Ext}_{A}^{p}(A/J,A)
  \otimes\Omega^{n-p}(A/J)'\otimes\Omega^{n}(A)
\]
(onde o símbolo « $'$ » designa o $(A/J)$-módulo dual). O produto tensorial dos
dois últimos fatores identifica-se com $\bigwedge^{p}(J/J^{2})$, pelo que o conjunto
identifica-se com
\[
  \operatorname{Ext}_{A}^{p}\!\left(A/J,\bigwedge^{p}(J/J^{2})\right),
\]
e, pela proposição, com
\[
  \operatorname{Hom}_{A}\!\left(\bigwedge^{p}J/J^{2},
  \bigwedge^{p}J/J^{2}\right),
\]
isto é, com $A/J$. Em particular, em
$\operatorname{Ext}_{A}^{p}\bigl(\Omega^{n-p}(A/J),\Omega^{n}(A)\bigr)$ há um
elemento distinguido, correspondente ao elemento unidade de $A/J$, chamado
\emph{classe fundamental do ideal $J$ em $A$}. (Na realidade, pode
ser definida sob condições sensivelmente mais gerais.) O
corolário 1 pode ser apresentado em uma forma mais geométrica e mais global:

\result{COROLÁRIO}{2} Sejam $X$ uma variedade não singular sobre um corpo
algebricamente fechado $k$, $Y$ uma subvariedade fechada não singular de $X$,
$\Osh_X$ o feixe de anéis locais de $X$, e $\Osh_Y$ o feixe de anéis
locais de $Y$ considerado como feixe quociente de $\Osh_X$. Sejam $n$ a
dimensão de $X$ e $n-p$ a de $Y$.

\src{176}
Seja $\Omega_X$ (resp. $\Omega_Y$) o feixe de germes de formas
diferenciais regulares sobre $X$ (resp. $Y$). Então têm-se isomorfismos
canônicos:
\begin{equation}\tag{3.6}
  \Extsh_{\Osh_X}^{p}\bigl(\Omega_Y^{n-p},\Omega_X^{n}\bigr)=\Osh_Y,
\end{equation}
ou também
\begin{equation}\tag{3.6 bis}
  \Extsh_{\Osh_X}^{p}\bigl(\Osh_Y,\Omega_X^{n}\bigr)=\Omega_Y^{n-p}.
\end{equation}
A fórmula (3.6 bis) pode servir como definição de $\Omega_Y^{n-p}$ quando $Y$
é uma variedade singular. Mais precisamente:

\result{PROPOSIÇÃO}{5} Seja $Y$ um subconjunto algébrico de dimensão
$q=n-p$ em uma variedade algébrica não singular $X$ de dimensão $n$. Sejam
$F$ um feixe algébrico coerente sobre $X$ cujo suporte está contido em $Y$, e $L$
um feixe algébrico localmente livre sobre $X$. Então os feixes
$\Extsh_{\Osh_X}^{i}(F,L)$ são nulos para $i<p$, enquanto que para $i=p$ tem-se um
isomorfismo canônico
\begin{equation}\tag{3.7}
  \Extsh_{\Osh_X}^{p}(F,L)
  =\Homsh_{\Osh_X}\!\left(F,
  \Extsh_{\Osh_X}^{p}(\Osh_X/J,L)\right),
\end{equation}
onde $J$ designa um feixe de ideais arbitrário sobre $X$ que anula $F$ e cujo
conjunto de zeros é $Y$. Em particular, se $F$ é um feixe algébrico
coerente sobre $Y$, tem-se
\begin{equation}\tag{3.7 bis}
  \Extsh_{\Osh_X}^{p}(F,L)
  =\Homsh_{\Osh_Y}\!\left(F,\Extsh_{\Osh_X}^{p}(\Osh_Y,L)\right).
\end{equation}
Por fim, sendo ainda $F$ um feixe algébrico coerente sobre $Y$, os feixes
\[
  E^{i}(F)=\Extsh_{\Osh_X}^{p+i}(F,\Omega_X^{n})
\]
não dependem da maneira de mergulhar o espaço algébrico $Y$ em uma variedade
não singular $X$.

Como a questão é local, pode-se supor que $X$ é afim e $L=\Osh_X$. Isto reduz
o problema a uma questão de álgebra comutativa, e até de álgebra local: se $A$ é um
anel local regular e $M$ um $A$-módulo cujo suporte tem dimensão
$\leq q=n-p$, é preciso demonstrar que $\operatorname{Ext}_{A}^{i}(M,A)=0$ para $i<p$, e
\[
  \operatorname{Ext}_{A}^{p}(M,A)
  =\operatorname{Hom}_{A}\!\left(M,\operatorname{Ext}_{A}^{p}(A/J,A)\right),
\]
onde $J$ é um ideal qualquer de « dimensão » $\leq q$ que anula $M$. Para o
primeiro ponto, procede-se por indução sobre $q$: um dévissage imediato reduz
ao caso em que $M$ é da forma $A/J$ e, em seguida, substituindo $J$ por um ideal menor
e utilizando a hipótese de indução e a sequência exata dos
$\operatorname{Ext}$, ao caso em que $J$ é gerado por um « sistema de parâmetros »
como na proposição 4, onde o resultado é imediato. O resultado
anterior implica que, se $J$ é um ideal

\src{177}
fixo de « dimensão » $\leq q$, então o funtor contravariante
$E(M)=\operatorname{Ext}_{A}^{p}(M,A)$ sobre a categoria dos $A/J$-módulos é
exato pela esquerda; além disso, transforma somas diretas em produtos diretos,
de onde resulta facilmente que
\[
  E(M)=\operatorname{Hom}_{A}\bigl(M,E(A)\bigr).
\]
 Por fim, a
última afirmação da proposição 5 é mais sutil e resulta de uma
caracterização intrínseca dos $E^{i}(F)$ mediante um teorema de dualidade
local que não pode ser exposto aqui.

\result{COROLÁRIO}{} Denotemos por $\omega_Y^{q}$ o feixe
$\Extsh_{\Osh_X}^{p}(\Osh_Y,\Omega_X^{n})$. Então tem-se um isomorfismo
funtorial para os feixes algébricos coerentes $F$ sobre $Y$:
\begin{equation}\tag{3.8}
  \Extsh_{\Osh_X}^{p}(F,\Omega_X^{n})
  =\Homsh_{\Osh_X}(F,\omega_Y^{q}).
\end{equation}

\fsection{4}{Classe de cohomologia associada a uma subvariedade}

Em tudo o que segue, $X$ designa um conjunto algébrico de dimensão $n$, definido
sobre um corpo $k$, que suporemos algebricamente fechado para simplificar. Salvo
no n\textsuperscript{o} 6, supõe-se que $X$ é não singular. Denotamos por
$\Osh_X$ o feixe de anéis locais sobre $X$, e por
\[
  \Omega_X^{*}=\bigcup_p\Omega_X^{p}
\]
o feixe de germes de formas diferenciais sobre $X$.
Se $Y$ é uma parte fechada de $X$, identificamos
os feixes algébricos coerentes sobre $Y$ com feixes coerentes sobre $X$
nulos fora de $Y$; em particular, assim ocorre com $\Osh_Y$ e $\Omega_Y$.

\result{LEMA}{1} Seja $F$ um feixe algébrico coerente sobre $X$ cujo
suporte tem dimensão $\leq n-p$, e $L$ um feixe algébrico coerente
localmente livre sobre $X$. Então $\operatorname{Ext}_{\Osh_X}^{i}(X;F,L)$ é nulo
para $i<p$, e tem-se um isomorfismo canônico
\begin{equation}\tag{4.1}
  \operatorname{Ext}_{\Osh_X}^{p}(X;F,L)
  =H^{0}\bigl(X,\Extsh_{\Osh_X}^{p}(F,L)\bigr).
\end{equation}
Se $F$ é um feixe algébrico coerente sobre uma parte fechada $Y$ de $X$ de
dimensão $\leq n-p$, tem-se um isomorfismo canônico
\begin{equation}\tag{4.1 bis}
  \Extsh_{\Osh_X}^{p}(F,L)
  =\Homsh_{\Osh_X}\!\left(
    F\otimes L'\otimes\Omega_X^{n},\omega_Y^{n-p}
  \right),
\end{equation}
onde $\omega_Y^{n-p}$ é o feixe sobre $Y$ definido no corolário da
proposição 5 (que se identifica com $\Omega_Y^{n-p}$ se $Y$ é não singular).

A fórmula (4.1) é uma consequência imediata da sequência espectral da
proposição 1 e da proposição 5; em virtude da fórmula (3.8), pode-se
escrever

\src{178}
\begin{align*}
\Extsh_{\Osh_X}^{p}(F,L)
&=L\otimes(\Omega_X^{n})'\otimes
  \Extsh_{\Osh_X}^{p}(F,\Omega_X^{n})\\
&=L\otimes(\Omega_X^{n})'\otimes
  \Homsh_{\Osh_X}(F,\omega_Y^{q})\\
&=\Homsh_{\Osh_X}\!\left(
  F\otimes L'\otimes\Omega_X^{n},\omega_Y^{q}\right),
\end{align*}
onde $q=n-p$, de onde resulta a fórmula (4.1 bis).

Tomando, em particular, $F=\Osh_Y$, $L=\Omega_X^{p}$, obtém-se, tendo em conta
que $\Omega_X^{n}\otimes(\Omega_X^{p})'=\Omega_X^{n-p}$, um isomorfismo
canônico
\begin{equation}\tag{4.2}
  \operatorname{Ext}_{\Osh_X}^{p}(X;\Osh_Y,\Omega_X^{p})
  =\operatorname{Hom}_{\Osh_X}
  (X;\Omega_X^{n-p},\omega_Y^{n-p}).
\end{equation}
Suponhamos agora que $Y$ é não singular, para simplificar, de modo que
$\omega_Y^{n-p}=\Omega_Y^{n-p}$. Tem-se um homomorfismo natural de
$\Omega_X^{n-p}$ em $\Omega_Y^{n-p}$, de onde resulta uma seção canônica $s_Y$ do
feixe $\Extsh_{\Osh_X}^{p}(\Osh_Y,\Omega_X^{p})$, que chamaremos, se
todas as componentes de $Y$ têm dimensão $n-p$, \emph{seção
fundamental} do feixe $\Extsh_{\Osh_X}^{p}(\Osh_Y,\Omega_X^{p})$. Esta
seção define, em virtude de (4.1), um elemento de
$\operatorname{Ext}_{\Osh_X}^{p}(X;\Osh_Y,\Omega_X^{p})$. Ora, o homomorfismo
natural $\Osh_X\longrightarrow\Osh_Y$ define um homomorfismo
\[
  \operatorname{Ext}_{\Osh_X}^{p}(X;\Osh_Y,\Omega_X^{p})
  \longrightarrow
  \operatorname{Ext}_{\Osh_X}^{p}(X;\Osh_X,\Omega_X^{p})
  =H^{p}(X,\Omega_X^{p}).
\]
Obtém-se assim um elemento de $H^{p}(X,\Omega_X^{p})$, denotado por $P_X(Y)$ e chamado
\emph{classe de cohomologia de $Y$ em $X$}. Deduz-se, portanto, da seção
$s_Y$ de $\Extsh_{\Osh_X}^{p}(\Osh_Y,\Omega_X^{p})$ mediante o diagrama
de homomorfismos seguinte:
\begin{equation}\tag{4.3}
\begin{aligned}
H^{p}(X,\Omega_X^{p})
&=\operatorname{Ext}_{\Osh_X}^{p}(X;\Osh_X,\Omega_X^{p})\\
&\longleftarrow\operatorname{Ext}_{\Osh_X}^{p}(X;\Osh_Y,\Omega_X^{p})
\xrightarrow{\ \sim\ }
H^{0}\!\left(X,\Extsh_{\Osh_X}^{p}(\Osh_Y,\Omega_X^{p})\right).
\end{aligned}
\end{equation}

Chamemos \emph{ciclo não singular de dimensão $n-p$} a um elemento do grupo
abeliano livre gerado pelas subvariedades irredutíveis não singulares de
dimensão $n-p$ em $X$. Então a função $Y\longmapsto P(Y)$ se prolonga a
um homomorfismo do grupo de ciclos não singulares de dimensão $n-p$ sobre $X$
no grupo $H^{p}(X,\Omega_X^{p})$.

Sejam $Z^{n-p}$ e $Z'^{n-p'}$ ciclos não singulares de dimensões $n-p$ e
$n-p'$; diz-se que se \emph{intersectam transversalmente} se toda componente
de $Z$ interseca transversalmente toda componente de $Z'$. Então o ciclo $Z.Z'$ está
definido e é um ciclo não singular de dimensão $n-p-p'$. Dito isto, tem-se o

\src{179}
\result{TEOREMA}{1} Se $Z^{n-p}$ e $Z'^{n-p'}$ são ciclos não singulares
que se intersectam transversalmente, então tem-se
\begin{equation}\tag{4.4}
  P_X(Z.Z')=P_X(Z).P_X(Z'),
\end{equation}
onde o produto do segundo membro é o produto cup
\[
  H^{p}(X,\Omega_X^{p})\times H^{p'}(X,\Omega_X^{p'})
  \longrightarrow H^{p+p'}(X,\Omega_X^{p+p'}).
\]
(Supõe-se que $X$ é isomorfo a uma parte localmente fechada de um espaço
projetivo.)

Esta última hipótese serve apenas para poder concluir que todo
feixe algébrico coerente sobre $X$ é quociente de um feixe algébrico
coerente localmente livre (SERRE) e, portanto, admite uma resolução pela esquerda mediante
feixes localmente livres. Para demonstrar o teorema 1, pode-se supor que
$Z$ e $Z'$ são subvariedades irredutíveis não singulares $Y$ e $Y'$ que se
intersectam transversalmente. Seja $L_{*}$ uma resolução pela esquerda de $\Osh_Y$ mediante
feixes localmente livres; então, em virtude da proposição 3, o diagrama
de homomorfismos (4.3) identifica-se com o diagrama
\[
R^{p}\Gamma(\Omega_X^{p})
\xleftarrow{\alpha}
\mathbb R^{p}\Gamma\!\left(\Homsh_{\Osh_X}(L_{*},\Omega_X^{p})\right)
\xrightarrow[\sim]{\beta}
\Gamma\!\left(H^{p}\bigl(\Homsh_{\Osh_X}(L_{*},\Omega_X^{p})\bigr)\right),
\]
onde $\Gamma$ é o funtor « grupo de seções » sobre a categoria dos
feixes abelianos sobre $X$, e $\mathbb R^{p}\Gamma$ designa
a hipercohomologia de grau $p$ desse funtor, e $R^{p}\Gamma$ seu
$p$-ésimo funtor derivado. Suporemos, para simplificar, que $L_0=\Osh_X$ e que
a aumentação $L_0\longrightarrow\Osh_Y$ é o homomorfismo natural (o que
é lícito); então $\alpha$ se deduz do homomorfismo de complexos
$\Osh_X\longrightarrow L_{*}$ (considerando $\Osh_X$ como um complexo concentrado
em grau $0$), tendo em conta que $\mathbb R^{p}\Gamma(K)=R^{p}\Gamma(K_0)$ se $K$
é um complexo de feixes concentrado em grau $0$. O homomorfismo $\beta$ é um
« homomorfismo de bordo » bem conhecido. Consideremos um diagrama análogo, relativo a
uma resolução localmente livre $L'_{*}$ de $\Osh_{Y'}$, e consideremos o
diagrama comutativo de acoplamentos:
\begin{equation}\tag{4.5}
\resizebox{\textwidth}{!}{$
\begin{array}{ccc}
R^{p}\Gamma(\Omega_X^{p})
&\longleftarrow
\mathbb R^{p}\Gamma\!\left(\Homsh_{\Osh_X}(L_{*},\Omega_X^{p})\right)
&\xrightarrow{\sim}
\Gamma H^{p}\!\left(\Homsh_{\Osh_X}(L_{*},\Omega_X^{p})\right)
\\[3pt]
\times&\times&\times\\[3pt]
R^{p'}\Gamma(\Omega_X^{p'})
&\longleftarrow
\mathbb R^{p'}\Gamma\!\left(\Homsh_{\Osh_X}(L'_{*},\Omega_X^{p'})\right)
&\xrightarrow{\sim}
\Gamma H^{p'}\!\left(\Homsh_{\Osh_X}(L'_{*},\Omega_X^{p'})\right)
\\[3pt]
\downarrow&\downarrow&\downarrow\\[3pt]
R^{p+p'}\Gamma(\Omega_X^{p+p'})
&\longleftarrow
\mathbb R^{p+p'}\Gamma\!\left(
\Homsh_{\Osh_X}(L_{*}\otimes L'_{*},\Omega_X^{p+p'})\right)
&\xrightarrow{\sim}
\Gamma H^{p+p'}\!\left(
\Homsh_{\Osh_X}(L_{*}\otimes L'_{*},\Omega_X^{p+p'})\right).
\end{array}
$}
\end{equation}

\src{180}
Os acoplamentos das duas colunas da direita se deduzem do acoplamento de
complexos de feixes
\[
\Homsh_{\Osh_X}(L_{*},\Omega_X^{p})
\times\Homsh_{\Osh_X}(L'_{*},\Omega_X^{p'})
\longrightarrow
\Homsh_{\Osh_X}(L_{*}\otimes L'_{*},\Omega_X^{p+p'}),
\]
que se define mediante o produto exterior
$\Omega_X^{p}\times\Omega_X^{p'}\longrightarrow\Omega_X^{p+p'}$.
O acoplamento da primeira coluna é o produto cup (relativo ao produto
exterior). Afirmo que a última linha de (4.5) identifica-se com o diagrama
de isomorfismos análogo a (4.3), onde $Y$ é substituído por $Y\cap Y'$ e $p$ por
$p+p'$. Para isso basta demonstrar que $L\otimes L'$ é uma resolução
(evidentemente localmente livre) de $\Osh_{Y\cap Y'}$. Com efeito, tem-se
\[
  H_0(L\otimes L')=\Osh_Y\otimes\Osh_{Y'}=\Osh_{Y\cap Y'}
\]
e
\[
  H_i(L\otimes L')=\operatorname{Tor}_{i}^{\Osh_X}(\Osh_Y,\Osh_{Y'})=0
  \qquad(i>0),
\]
pelo fato de que $Y$ e $Y'$ se intersectam transversalmente. O teorema 1 resulta
agora da fórmula:
\begin{equation}\tag{4.6}
  s_Y.s_{Y'}=s_{Y.Y'}
\end{equation}
(onde o produto do primeiro membro é o da última coluna de (4.5)).
Esta fórmula (4.6), de natureza puramente local, verifica-se sem dificuldade
tomando para $L_{*}$ e $L'_{*}$ as resoluções consideradas na proposição
4. Demonstra-se igualmente (mediante uma demonstração mais simples) que
$Z\longmapsto P_X(Z)$ é compatível com o produto cartesiano:
\begin{equation}\tag{4.7}
  P_{X\times X'}(Z\times Z')=P_X(Z)\otimes P_{X'}(Z')
\end{equation}
(fórmula válida se $Z$, resp. $Z'$, é um ciclo não singular sobre a variedade
não singular $X$, resp. $X'$, considerando $Z\times Z'$ como um ciclo não
singular sobre $X\times X'$). De (4.4) e (4.7) deduz-se que $P_X(Z)$ é também
compatível com a operação « imagem inversa » por um morfismo de variedades não
singulares $f:X\longrightarrow X'$:
\begin{equation}\tag{4.8}
  P_X\bigl(f^{-1}(Z)\bigr)=f^{*}\bigl(P_{X'}(Z')\bigr),
\end{equation}
fórmula válida se $Z'$ é um ciclo não singular sobre $X'$ tal que $f$ seja
« transversal » a $Z'$, isto é, tal que a gráfica de $f$ seja transversal ao ciclo
$X\times Z'$ em $X\times X'$.

\src{181}
\result{COROLÁRIO}{1} Sejam $X$, $X'$ duas variedades não singulares
localmente fechadas em um espaço projetivo, e suponhamos que $X'$ é completa. Sejam $U$
um ciclo não singular sobre $X\times X'$, $a$ e $b$ dois pontos de $X'$ tais que
$U$ interseca transversalmente os ciclos $X\times(a)$ e $X\times(b)$, e sejam $Z$
e $Z'$ os ciclos não singulares sobre $X$ tais que
\[
  Z\times(a)=(X\times(a)).U,
  \qquad
  Z'\times(b)=(X\times(b)).u.
\]
Então tem-se
\[
  P_X(Z)=P_X(Z').
\]

Com efeito, seja $f_a:X\longrightarrow X\times X'$ definido por
$f_a(x)=(x,a)$; então, em virtude de (4.8), tem-se a fórmula
$P(Z)=f_a^{*}(P(U))$, e analogamente $P(Z')=f_b^{*}(P(U))$. Ora, utilizando a fórmula de
Künneth
\[
  H^{*}(X\times X',\Omega_{X\times X'}^{*})
  =H^{*}(X,\Omega_X^{*})\otimes H^{*}(X',\Omega_{X'}^{*})
\]
e o fato de que $H^{0}(X',\Osh_{X'})$ se reduz aos escalares, obtém-se
facilmente $f_a^{*}=f_b^{*}$, de onde se segue a conclusão.

Para todo $x\in X$, $(x)$ é uma subvariedade não singular de $X$ de
codimensão $n$, e define portanto um elemento $\varepsilon_x$ de
$H^{n}(X,\Omega_X^{n})$. Se $X$ é uma variedade projetiva não singular, do
corolário 1 resulta que $\varepsilon_x$ não depende do ponto $x$ escolhido; denota-se
por $\varepsilon_X$ e chama-se \emph{classe fundamental} de
$H^{n}(X,\Omega_X^{n})$.

\result{OBSERVAÇÃO}{} Para dispor de uma teoria satisfatória, seria preciso definir
$P_X(Z)$ para ciclos $Z$ quaisquer e demonstrar o teorema 1 para uma
interseção própria de ciclos. (No momento de escrever esta exposição, isso ainda não foi
feito com toda generalidade). Admitindo que tenha sido feito, o corolário 1
passa a ser: se $Z$ e $Z'$ são dois ciclos algebricamente equivalentes, então
$P_X(Z)=P_X(Z')$ (afirmação que não parece resultar do que precede, mesmo que
$Z$ e $Z'$ sejam não singulares).

\fsection{5}{O teorema de dualidade}

Neste número, $X$ designa uma variedade projetiva não singular de dimensão
$n$.

\result{TEOREMA}{2} A classe fundamental $\varepsilon_X$ de
$H^{n}(X,\Omega_X^{n})$ é uma base deste espaço vetorial.

(A demonstração será dada mais adiante). Em virtude do teorema anterior, pode-se
identificar $H^{n}(X,\Omega_X^{n})$ com o corpo $k$. Consideremos agora os
acoplamentos considerados no n\textsuperscript{o} 2, que dão, em particular, um
acoplamento

\src{182}
\[
\operatorname{Ext}_{\Osh_X}^{p}(X;\Osh_X,F)
\times\operatorname{Ext}_{\Osh_X}^{n-p}(X;F,\Omega_X^{n})
\longrightarrow
\operatorname{Ext}_{\Osh_X}^{n}(X;\Osh_X,\Omega_X^{n}),
\]
isto é,
\begin{equation}\tag{5.1}
  H^{p}(X,F)
  \times\operatorname{Ext}_{\Osh_X}^{n-p}(X;F,\Omega_X^{n})
  \longrightarrow H^{n}(X,\Omega_X^{n})
\end{equation}
Tendo em conta o teorema 2, este acoplamento define um homomorfismo
\begin{equation}\tag{5.2}
  \operatorname{Ext}_{\Osh_X}^{n-p}(X;F,\Omega_X^{n})
  \longrightarrow\bigl(H^{p}(X,F)\bigr)'
\end{equation}
Este homomorfismo é funtorial em $F$ e comuta com os homomorfismos de conexão
relativos às sequências exatas $0\longrightarrow F'\longrightarrow F
\longrightarrow F''\longrightarrow0$.

\result{TEOREMA}{3} O homomorfismo (5.2) é um isomorfismo.

Recupera-se, em particular, o resultado de Serre:

\result{COROLÁRIO}{} Seja $E$ um fibrado vetorial algébrico sobre $X$ e
$\Osh_X(E)$ o feixe de germes de seções regulares de $E$; então têm-se
isomorfismos canônicos:
\begin{equation}\tag{5.3}
  \bigl(H^{p}(X,\Osh_X(E))\bigr)'
  =H^{n-p}\bigl(X,\Omega_X^{n}\otimes\Osh_X(E')\bigr)
\end{equation}
Basta aplicar o teorema 3 e o corolário 1 à proposição 1.

Os teoremas 2 e 3 resultarão do enunciado seguinte:

\noindent\textbf{(D) O homomorfismo seguinte}
\begin{equation}\tag{5.2 bis}
  \operatorname{Ext}_{\Osh_X}^{n-p}(X;F,\Omega_X^{n})
  \longrightarrow\bigl(H^{p}(X,F)\bigr)'\otimes L,
  \qquad L=H^{n}(X,\Omega_X^{n}),
\end{equation}
deduzido do acoplamento (5.1) é um isomorfismo.

Demonstremos, com efeito, que (D) implica o teorema 2. Seja
$k_x=\Osh_{(x)}$ o feixe de anéis locais da variedade reduzida ao ponto
$x\in X$; consideremos o homomorfismo canônico
$\Osh_X\longrightarrow k_x$ e o homomorfismo associado
\begin{equation}\tag{5.3}
  H^{0}(X,\Osh_X)\longrightarrow H^{0}(X,k_x),
\end{equation}
seu transposto identifica-se com o homomorfismo
\begin{equation}\tag{5.4}
  \operatorname{Ext}_{\Osh_X}^{n}(X;k_x,\Omega_X^{n})\otimes L'
  \longrightarrow
  \operatorname{Ext}_{\Osh_X}^{n}(X;\Osh_X,\Omega_X^{n})\otimes L'
\end{equation}
deduzido do homomorfismo sobre os $\operatorname{Ext}^{n}$ associado a
$\Osh_X\longrightarrow k_x$:
\begin{equation}\tag{5.5}
  \operatorname{Ext}_{\Osh_X}^{n}(X;k_x,\Omega_X^{n})
  \longrightarrow
  \operatorname{Ext}_{\Osh_X}^{n}(X;\Osh_X,\Omega_X^{n})
  =H^{n}(X,\Omega_X^{n})
\end{equation}

\src{183}
Como (5.3) é um isomorfismo, o mesmo ocorre
com (5.4), e portanto também com (5.5). Como $s_{(x)}$ é uma base de
$\operatorname{Ext}_{\Osh_X}^{n}(X;k_x,\Omega_X^{n})$ em virtude de (4.2), se
segue-se que sua imagem $\varepsilon_X$ é uma base de
$H^{n}(X,\Omega_X^{n})$.

Fica por demonstrar o enunciado (D), que resultará de maneira puramente formal dos fatos
elementares resumidos nos lemas seguintes. Neles se supõe que $X$ é uma
parte fechada (singular ou não) do espaço projetivo $P$ de dimensão $r$.
Utiliza-se a notação $\Osh_P(m)$ para o feixe sobre $P$ denotado por $\Osh(m)$ em
[3], e a notação $\Osh_X(m)$ para o feixe análogo sobre $X$.

\result{LEMA}{2} O enunciado (D) é verdadeiro se $X=P$ e $F=\Osh_P(m)$.

Este lema verifica-se mediante um cálculo direto. O cálculo explícito dos
$H^{i}(P,\Osh_P(m))$ encontra-se em [3], mas pode ser feito de maneira mais
elementar. O cálculo do produto cup
\[
  H^{i}(P,\Osh_P(m))\times H^{j}(P,\Osh_P(m'))
  \longrightarrow H^{i+j}(P,\Osh_P(r+r'))
\]
necessário para explicitar o acoplamento (5.1) não apresenta
dificuldades.

\result{LEMA}{3} Todo feixe algébrico coerente $F$ sobre $X$ é isomorfo a
um feixe quociente de um feixe $\Osh_X(-m)^{k}$, onde se pode supor que $m$ é
tão grande quanto se queira.

Resulta do fato de que $F\otimes\Osh_X(m)$ está « gerado por suas seções » para
$m$ grande, cf. [3].

\result{LEMA}{4} Seja $i>0$. Então
$H^{r-i}(P,\Osh_P(-m))=0$ para $m$ grande; e para todo feixe algébrico
coerente $B$ sobre $X$, tem-se
\[
  \operatorname{Ext}_{\Osh_X}^{i}(X;\Osh_X(-m),B)=0
\]
para $m$ grande.

O primeiro enunciado resulta dos cálculos explícitos mencionados anteriormente; para o
segundo, observa-se que se tem um isomorfismo
\[
  \operatorname{Ext}_{\Osh_X}^{i}(X;\Osh_X(-m),B)
  =H^{i}(X,B\otimes\Osh_X(m))
\]
(proposição 1, corolário 1), de onde se segue a conclusão em virtude de um resultado bem
conhecido de [3]. Combinando os dois lemas anteriores, obtém-se o

\result{COROLÁRIO}{} Seja $i>0$. Então o funtor
$F\longmapsto H^{r-i}(P,F)$ sobre a categoria dos feixes algébricos coerentes
sobre $P$ é apagável, e o mesmo ocorre com o funtor
$\operatorname{Ext}_{\Osh_X}^{i}(X;F,B)$ sobre a categoria dos feixes
algébricos coerentes sobre $X$.

\src{184}
\result{LEMA}{5} Sejam $A$ e $B$ dois feixes algébricos coerentes sobre
$X$, e ponhamos $A(m)=A\otimes\Osh_X(m)$. Então, para $m$ suficientemente grande,
o homomorfismo canônico
\[
  \operatorname{Ext}_{\Osh_X}^{1}(X;A(-m),B)
  \longrightarrow H^{0}\bigl(X,\Extsh_{\Osh_X}^{1}(A(-m),B)\bigr)
  =H^{0}\bigl(X,\Extsh_{\Osh_X}^{1}(A,B)(m)\bigr)
\]
é um isomorfismo.

Isto resulta imediatamente da sequência espectral da proposição 1 aplicada a
$A(-m)$ e $B$, pois se terá
\[
  E_{2}^{p,q}(A(-m),B)
  =H^{p}\bigl(X,\Extsh_{\Osh_X}^{q}(A(-m),B)\bigr)
  =H^{p}\bigl(X,\Extsh_{\Osh_X}^{q}(A,B)(m)\bigr),
\]
que é nulo para $p>0$ e $m$ grande.

Demonstremos agora (D) no caso em que $X=P$. Provemos primeiro que (5.2 bis)
é um isomorfismo para $p=n$; como os dois membros são então
funtores exatos pela esquerda (pois $H^{r+1}(P,F)=0$), do lema 3 resulta
que basta demonstrar a afirmação para $F=\Osh_P(-m)$, mas esta se encontra
contida no lema 2. Como os homomorfismos (5.2 bis) são funtoriais e
compatíveis com os operadores de conexão, e como para $p<n$ os dois membros de
(5.2 bis) são funtores apagáveis em $F$ (corolário do lema 4), segue-se
então, mediante um raciocínio padrão, que (5.2 bis) é um isomorfismo
para todo $p$. Isto demonstra o teorema de dualidade para o espaço projetivo.
Suponhamos agora que $X$ é qualquer, mas não singular. Pelo teorema de
dualidade para $P$, tem-se um isomorfismo
\[
  H^{n}(X,F)'=H^{n}(P,F)'
  =\operatorname{Ext}_{\Osh_P}^{r-n}(P;F,\Omega_P^{r}).
\]
Em virtude do lema 1 (n\textsuperscript{o} 4), o último membro identifica-se com
\[
  \operatorname{Hom}_{\Osh_P}(P;F,\omega_X^{n})
  =\operatorname{Hom}_{\Osh_X}(X;F,\Omega_X^{n})
  =\operatorname{Ext}_{\Osh_X}^{0}(X;F,\Omega_X^{n}),
\]
de onde se obtém um isomorfismo
\begin{equation}\tag{5.6}
  H^{n}(X,F)'=\operatorname{Hom}_{\Osh_X}(X;F,\Omega_X^{n})
  =\operatorname{Ext}_{\Osh_X}^{0}(X;F,\Omega_X^{n}).
\end{equation}
Tomando $F=\Omega_X^{n}$, obtém-se um isomorfismo
\begin{equation}\tag{5.7}
  \gamma:H^{n}(X,\Omega_X^{n})\xrightarrow{\sim}k.
\end{equation}

\src{185}
Verifica-se que o isomorfismo (5.6) não é outro senão (5.2 bis) para $p=n$, quando
se toma $L=k$ graças a (5.7). Por conseguinte, (5.2 bis) é um isomorfismo para
$p=n$. Para demonstrar que é um isomorfismo para todo $p$, basta novamente
demonstrar que, para $p<n$, os dois membros de (5.2 bis) são funtores
apagáveis em $F$, e a fortiori (tendo em conta o lema 3) que ambos os membros
são nulos quando se toma $F=\Osh_X(-m)$ com $m$ grande. Ora, para o primeiro
membro, isto é verdadeiro em virtude do lema 4, e para o segundo membro escreve-se,
utilizando o teorema de dualidade sobre $P$:
\[
  H^{p}(X,\Osh_X(-m))'
  =\operatorname{Ext}_{\Osh_P}^{r-p}
    (P;\Osh_X(-m),\Omega_P^{r}).
\]
O segundo membro é nulo para $p<n$ e $m$ grande, como resulta do lema 5
(onde se toma $X=P$) e do fato de que $\Osh_X$, como feixe algébrico
coerente sobre $P$, tem dimensão cohomológica $\leq r-n$ (pois $X$ é não
singular), de onde
\[
  \Extsh_{\Osh_P}^{r-p}(\Osh_X,\Omega_P^{r})=0
  \qquad\text{para }p<n.
\]

\fsection{6}{O teorema de dualidade para as variedades singulares}

Seja $X$ uma parte fechada de dimensão $n$ do espaço projetivo $P$ de
dimensão $r$. A fórmula (5.6) deve ser escrita aqui
\begin{equation}\tag{6.1}
  H^{n}(X,F)'
  \simeq\operatorname{Hom}_{\Osh_X}(X;F,\omega_X^{n})
  =\operatorname{Ext}_{\Osh_X}^{0}(X;F,\omega_X^{n}),
\end{equation}
onde se pôs
\begin{equation}\tag{6.2}
  \omega_X^{n}=E^{0}(\Osh_X)
  =\Extsh_{\Osh_P}^{r-n}(\Osh_X,\Omega_P^{r}).
\end{equation}
Como foi dito na proposição 5, o feixe assim introduzido não depende,
na realidade, da imersão escolhida de $X$ em uma variedade não singular $P$.
Tomando, em (6.1), $F=\omega_X^{n}$, obtém-se
\begin{equation}\tag{6.2}
  H^{n}(X,\omega_X^{n})'
  \simeq\operatorname{Hom}_{\Osh_X}(X;\omega_X^{n},\omega_X^{n}),
\end{equation}
 de onde
resulta a existência de um elemento
distinguido em $H^{n}(X,\omega_X^{n})'$, correspondente ao homomorfismo
identidade de $\omega_X^{n}$ em si mesmo:
\begin{equation}\tag{6.3}
  \gamma:H^{n}(X,\omega_X^{n})\longrightarrow k.
\end{equation}

\src{186}
Consideremos então os acoplamentos definidos pela composição dos $\operatorname{Ext}$:
\begin{equation}\tag{6.4}
  H^{p}(X,F)\times
  \operatorname{Ext}_{\Osh_X}^{n-p}(X;F,\omega_X^{n})
  \longrightarrow H^{n}(X,\omega_X^{n}),
\end{equation}
e componhamo-los com o homomorfismo $\gamma$ de (6.3); obtêm-se
homomorfismos funtoriais, compatíveis com os operadores de bordo:
\begin{equation}\tag{6.5}
  \operatorname{Ext}_{\Osh_X}^{n-p}(X;F,\omega_X^{n})
  \longrightarrow H^{p}(X,F)'
\end{equation}
(que generalizam (5.2)). Verifica-se que para $p=n$ se obtém assim
o isomorfismo (6.1). Posto isto, tem-se o

\result{TEOREMA}{3 bis} As quatro condições seguintes sobre $X$ são
equivalentes para um inteiro dado $k\geq0$:
\begin{enumerate}
  \item[\textit{i.}] O homomorfismo funtorial (6.5) é um isomorfismo para
    $n-k\leq p\leq n$.
  \item[\textit{ii.}] Tem-se $H^{p}(X,\Osh_X(-m))=0$ para $m$ grande e para
    $n-k\leq p<n$.
  \item[\textit{iii.}] O funtor $H^{p}(X,F)$ sobre a categoria dos feixes
    algébricos coerentes sobre $X$ é apagável para $n-k\leq p<n$.
  \item[\textit{iv.}] Tem-se
    $E^{i}(\Osh_X)=\Extsh_{\Osh_P}^{r-n+i}(\Osh_X,\Omega_P^{r})=0$
    para $0<i\leq k$.
\end{enumerate}

\noindent\textbf{DEMONSTRAÇÃO.} --- \textit{i.}$\Rightarrow$\textit{ii.} Em
virtude do lema 4, \textit{ii.}$\Rightarrow$\textit{iii.} em virtude do lema 3,
\textit{iii.}$\Rightarrow$\textit{i.} mediante um raciocínio padrão bem conhecido,
tendo em conta que os dois membros de (6.5) são então funtores apagáveis
para $n-k\leq p<n$ (o primeiro o é em virtude do lema 4). Por fim, demonstremos
\textit{ii.}$\Leftrightarrow$\textit{iv.} Isto resulta do corolário da
proposição seguinte:

\result{PROPOSIÇÃO}{6} Seja $F$ um feixe algébrico coerente sobre $X$, e
seja $i$ um inteiro. Então, para $m$ suficientemente grande, tem-se um isomorfismo:
\begin{equation}\tag{6.6}
  H^{i}(X,F(-m))'
  \simeq H^{0}\bigl(X,E^{n-i}(F)(m)\bigr),
\end{equation}
onde se põe (compare-se n\textsuperscript{o} 3, proposição 5):
\begin{equation}\tag{6.7}
  E^{j}(F)=\Extsh_{\Osh_P}^{r-n+j}(F,\Omega_P^{r}).
\end{equation}
Com efeito, em virtude do teorema de dualidade sobre $P$, o primeiro membro de (6.6)
é isomorfo a
$\operatorname{Ext}_{\Osh_P}^{r-i}(P;F(-m),\Omega_P^{r})$, pelo que (6.6) resulta
do lema 5.

\src{187}
\result{COROLÁRIO}{} Para que se tenha $H^{i}(X,F(-m))=0$ para $m$ grande, é
necessário e suficiente que $E^{n-i}(F)=0$.

Lembremos que os $E^{j}(F)$ não dependem, na realidade, da imersão projetiva
considerada. A condição do corolário é de natureza puramente local; por conseguinte,
se ela se verifica para $F$, verifica-se para todo feixe localmente
isomorfo a um feixe $F^{n}$. Em particular, se esta condição se verifica
para $\Osh_X$, verifica-se para todo feixe algébrico coerente localmente
livre. Este é, por exemplo, o caso para todo $i<n$ se $X$ é não singular, para
$i=0$ se nenhuma componente de $X$ se reduz a um só ponto, e para $i=0,1$
se $S$ é normal e todas as suas componentes têm dimensão $>1$ (cf.
[3]). Para que isto ocorra para todo
$i<k$, é necessário e suficiente, por definição, que os anéis locais
$\Osh_x$ ($x\in X$) tenham « codimensão homológica » $\geq k$ (cf. [4]
para esta noção). Se $k=n$, isto significa, em virtude do teorema 3 bis, que o
teorema de dualidade é verdadeiro para $X$, isto é, que (6.5) é um isomorfismo para
todo $p$ e todo $F$. Podem-se dar numerosas condições equivalentes sobre
os anéis locais $\Osh_x$ para que isto ocorra (Nagata), por exemplo, que
satisfaçam o teorema de equidimensionalidade de Cohen--Macaulay. Assim ocorre,
por exemplo, se localmente $X$ é uma « interseção completa » em uma
variedade ambiente não singular.

\fsection{7}{A dualidade de Poincaré}

Seja $X$ uma variedade projetiva não singular de dimensão $n$. Então
\[
  H^{**}(X)=H^{*}(X,\Omega_X^{*})
\]
é uma álgebra anticonmutativa bigraduada de dimensão finita, que graduamos
pelo grau total, de modo que $H^{p,q}(X)=H^{p}(X,\Omega_X^{q})$ tem grau
$p+q$. Os graus de $H^{*}(X)$ estão compreendidos entre $0$ e $2n$. Em virtude do
teorema 2 e do corolário do teorema 3, $H^{*}(X)$ é uma « álgebra de
Poincaré » de dimensão $2n$, isto é, $H^{2n}(X)$ está provido de um isomorfismo com
o corpo base $k$, e o produto
\[
  H^{m}(X)\times H^{2n-m}(X)\longrightarrow H^{2n}(X)=k
\]
é uma dualidade entre $H^{m}(X)$ e $H^{2n-m}(X)$. Além disso, se $Y$ é outra
variedade projetiva não singular, a fórmula de Künneth para os
feixes algébricos coerentes dá
\begin{equation}\tag{7.1}
  H^{*}(X\times Y)=H^{*}(X)\otimes H^{*}(Y),
\end{equation}
isomorfismo compatível com as estruturas de álgebras de Poincaré. Além disso,
$H^{*}(X)$ é, como álgebra comutativa, um funtor contravariante em
$X$, pois um morfismo $f:Y\longrightarrow X$ define de maneira evidente um
homomorfismo de álgebras graduadas

\src{188}
\begin{equation}\tag{7.2}
  f^{*}:H^{*}(X)\longrightarrow H^{*}(Y).
\end{equation}
Como se trata de álgebras de Poincaré, obtém-se por transposição um
homomorfismo de espaços vetoriais:
\begin{equation}\tag{7.3}
  f_{*}:H^{*}(Y)\longrightarrow H^{*}(X).
\end{equation}
Viu-se no n\textsuperscript{o} 4 que o efeito de $f^{*}$ sobre as classes de
cohomologia correspondentes a ciclos não singulares se interpreta
geometricamente tomando as classes de cohomologia que correspondem a suas
imagens inversas. No caso atual, importa demonstrar que (7.3) corresponde
analogamente à operação « imagem direta » de ciclos. Isto resulta, ao menos sob
condições adequadas de não singularidade, do caso particular seguinte:

\result{TEOREMA}{4} Se $f$ é a aplicação identidade de uma subvariedade não
singular $Y^{m}$ de $X^{n}$ em $X^{n}$, então, denotando por $1_Y$
o elemento unidade de $H(Y)$, tem-se
\begin{equation}\tag{7.4}
  f_{*}(1_Y)=P_X(Y),
\end{equation}
onde o segundo membro é a classe de cohomologia em $X$ associada a $Y$.

Esta fórmula equivale a
\begin{equation}\tag{7.4 bis}
  \bigl\langle\xi^{m,m},P_X(Y)\bigr\rangle\varepsilon_Y
  =f_{*}(\xi^{m,m})
  \qquad
  \bigl(\xi^{m,m}\in H^{m}(X,\Omega_X^{m})\bigr),
\end{equation}
 onde
$\varepsilon_Y$ é o elemento fundamental em $H^{m}(Y,\Omega_Y^{m})$, e
constitui, no caso das variedades projetivas não singulares, uma nova
definição da classe de cohomologia associada a $Y$. Para demonstrar o
teorema 4, interpreta-se, graças ao teorema 3, o transposto do homomorfismo
\[
  H^{m}(X,\Omega_X^{m})
  \longrightarrow H^{m}(Y,\Omega_Y^{m})
  =H^{m}(X,\Omega_Y^{m})
\]
como o homomorfismo
\begin{equation}\tag{7.5}
\begin{split}
  H^{n-m}(X,\Omega_X^{n-m})
  &\simeq
  \operatorname{Ext}_{\Osh_X}^{n-m}(X;\Omega_X^{m},\Omega_X^{n})\\
  &\longleftarrow
  \operatorname{Ext}_{\Osh_X}^{n-m}(X;\Omega_Y^{m},\Omega_X^{n})
  \simeq
  \operatorname{Hom}_{\Osh_X}(X;\Omega_Y^{m},\Omega_Y^{m}).
\end{split}
\end{equation}
Verifica-se que o elemento $1_Y$ do dual de $H^{m}(Y,\Omega_Y^{m})$ identifica-se
com o elemento do membro da direita correspondente ao endomorfismo identidade de
$\Omega_Y^{m}$, e por outro lado que a imagem deste elemento em
$H^{n-m}(X,\Omega_X^{n-m})$ é precisamente $P_X(Y)$. Estas compatibilidades, que poderiam
ter sido dadas já no n\textsuperscript{o} 4, podem ser enunciadas, e são
na realidade válidas, para variedades não singulares quaisquer; a segunda

\src{189}
resulta, por exemplo, da comutatividade do diagrama seguinte de endomorfismos
canônicos:
\begin{equation}\tag{7.6}
\resizebox{\textwidth}{!}{$
\left\{
\begin{array}{ccccc}
\operatorname{Ext}_{\Osh_X}^{n-m}(X;\Omega_X^{m},\Omega_X^{n})
&\longleftarrow&
\operatorname{Ext}_{\Osh_X}^{n-m}(X;\Omega_Y^{m},\Omega_X^{n})
&\simeq&
\operatorname{Hom}_{\Osh_X}(X;\Omega_Y^{m},\Omega_Y^{m})\\[3pt]
\Big\downarrow&&&&\Big\downarrow\\[3pt]
\operatorname{Ext}_{\Osh_X}^{n-m}(X;\Osh_X,\Omega_X^{n-m})
&\longleftarrow&
\operatorname{Ext}_{\Osh_X}^{n-m}(X;\Osh_Y,\Omega_X^{n-m})
&\simeq&
\operatorname{Hom}_{\Osh_X}(X;\Omega_X^{m},\Omega_Y^{m})
\end{array}
\right.
$}
\end{equation}

Obtém-se assim um equivalente exato do formalismo da dualidade de Poincaré
sobre as variedades orientadas compactas. Em particular, o teorema 4 permite
determinar a classe de cohomologia associada à diagonal de $X\times X$. Mediante
um raciocínio bem conhecido, deduz-se, por exemplo, uma fórmula de Lefschetz:

\result{TEOREMA}{5} Seja $f$ um endomorfismo de uma variedade projetiva não
singular $X$, tal que os pontos fixos de $f$ tenham multiplicidade 1.
Então o número desses pontos fixos é congruente, módulo a característica de
$k$, com a soma alternada dos traços dos endomorfismos dos $H^{i}(X)$ definidos
por $f$.

A restrição sobre $f$ que foi necessário impor deve-se às dificuldades mencionadas na
observação do n\textsuperscript{o} 4. Não obstante, observar-se-á que a fórmula de
Lefschetz continua verdadeira se $f$ é « homotópico » a um endomorfismo cujos
pontos fixos têm todos multiplicidade 1.

\fsection{8}{Generalização do teorema de dualidade}

Seja $X$ uma variedade algébrica não singular tal que todo feixe
algébrico coerente $F$ sobre $X$ seja isomorfo a um quociente de um feixe
algébrico coerente localmente livre (o que ocorre se $X$ é localmente
fechado em um espaço projetivo). Então todo feixe algébrico coerente $F$ sobre
$X$ admite uma resolução finita $L$ mediante feixes localmente livres, e para
duas resoluções desta classe sempre se pode encontrar uma terceira e um
homomorfismo desta nas duas primeiras compatível com as
aumentações. Analogamente, se $L$ é uma resolução finita localmente livre de
$F$, e se é dado um homomorfismo $F'\longrightarrow F$, então existe
uma resolução finita localmente livre $L'$ de $F'$ e um homomorfismo
$L'\longrightarrow L$ compatível com $F'\longrightarrow F$, que se pode
supor sobrejetivo se $F'\longrightarrow F$ o é. Isto permite definir, dados
dois inteiros $r,s\geq0$, dois multifuntores cohomológicos nos
argumentos $A_1,\ldots,A_r$; $B_1,\ldots,B_s$,

\src{190}
que variam na categoria dos feixes algébricos coerentes sobre $X$, e cujos
valores estão na categoria dos feixes algébricos coerentes sobre $X$,
resp. na categoria dos módulos sobre $H^{0}(X,\Osh_X)$, mediante as fórmulas:
\begin{equation}\tag{8.1}
\resizebox{\textwidth}{!}{$
\left\{
\begin{aligned}
\underline{T}_{r}^{s*}(A_1,\ldots,A_r;B_1,\ldots,B_s)
  &={}H^{*}\!\left(
    \Homsh_{\Osh_X}\!\left(
      L(A_1)\otimes\cdots\otimes L(A_r),
      L(B_1)\otimes\cdots\otimes L(B_s)
    \right)\right),\\[3pt]
T_{r}^{s*}(A_1,\ldots,A_r;B_1,\ldots,B_s)
  &={}\mathbf R^{*}\Gamma\!\left(
    \Homsh_{\Osh_X}\!\left(
      L(A_1)\otimes\cdots\otimes L(A_r),
      L(B_1)\otimes\cdots\otimes L(B_s)
    \right)\right).
\end{aligned}
\right.
$}
\end{equation}
Nesta fórmula, $L(F)$ designa uma resolução finita localmente livre do
feixe algébrico coerente $F$, e $\mathbf R^{*}\Gamma(K)$ designa
a hipercohomologia do espaço $X$ com respeito ao complexo de feixes $K$. Se
$r$ ou $s$ é nulo, substitui-se o produto tensorial dos $L(A_i)$, resp. dos
$L(B_j)$, por $\Osh_X$. Em particular, $\underline{T}_{0}^{0}$ e $T_{0}^{0}$
são funtores graduados de 0 argumentos, $\underline{T}_{0}^{0}$ se reduz ao
grau 0, onde é o feixe $\Osh_X$, enquanto que $T_{0}^{0}$ é idêntico a
$H^{*}(X,\Osh_X)$. O fato de que os segundos membros de (8.1) não dependam
das resoluções escolhidas é, em qualquer caso, evidente para a primeira linha
(pois a questão é então local), e para a segunda resulta das observações
gerais precedentes, tendo em conta a sequência espectral de
hipercohomologia do complexo de feixes
\[
  K=\Homsh_{\Osh_X}\bigl(L(A_1)\otimes\cdots,
      L(B_1)\otimes\cdots\bigr),
\]
que converge à hipercohomologia de $X$ com respeito a $K$ e cujo termo
inicial é $H^{p}(X,H^{q}(K))$, isto é,
\begin{equation}\tag{8.2}
  E_{2}^{p,q}
  =H^{p}\!\left(X,
    (\underline{T}_{r}^{s})^{(q)}
      (A_1,\ldots,A_r;B_1,\ldots,B_s)\right).
\end{equation}
Vê-se então que esta sequência espectral tampouco depende das resoluções
escolhidas; seu limite é o primeiro membro de (8.1). Definem-se
facilmente os operadores de conexão relativos aos diversos argumentos $A_i$, $B_j$,
observando que toda sequência exata
$0\longrightarrow F'\longrightarrow F\longrightarrow F''\longrightarrow0$
pode ser resolvida mediante uma sequência exata de complexos finitos localmente livres.

No sistema dos funtores $\underline{T}_{r}^{s*}$, resp.
$T_{r}^{s*}$, definem-se operações análogas às do cálculo tensorial, cuja
definição é imediata a partir das fórmulas definidoras (8.1). Assim, tem-se
uma composição (que generaliza a considerada no n\textsuperscript{o} 2):
\begin{equation}\tag{8.3}
\resizebox{\textwidth}{!}{$
\begin{aligned}
&T_{r}^{s*}(A_1,\ldots,A_r;B_1,\ldots,B_s)
\times T_{r'}^{s'*}(A'_1,\ldots,A'_{r'};B'_1,\ldots,B'_{s'})\\
&\qquad\longrightarrow
T_{r+r'}^{s+s'*}
(A_1,\ldots,A_r,A'_1,\ldots,A'_{r'};
 B_1,\ldots,B_s,B'_1,\ldots,B'_{s'}),
\end{aligned}
$}
\end{equation}
que satisfaz as propriedades evidentes de associatividade e de compatibilidade com
os homomorfismos funtoriais, os homomorfismos de conexão e as sequências
espectrais. Do mesmo modo, têm-se operações de simetria cuja explicitação se deixa ao
leitor. Tem-se, além disso, uma operação de contração cada vez que
um dos argumentos $A_i$

\src{191}
é igual a um dos argumentos $B_j$:
\begin{equation}\tag{8.4}
\resizebox{\textwidth}{!}{$
\begin{aligned}
&T_{r}^{s*}(A_1,\ldots,A_{i-1},C,A_{i+1},\ldots,A_r;
 B_1,\ldots,B_{j-1},C,B_{j+1},\ldots,B_s)\\
&\qquad\longrightarrow
T_{r-1}^{s-1*}(A_1,\ldots,\widehat{A_i},\ldots,A_r;
 B_1,\ldots,\widehat{B_j},\ldots,B_s).
\end{aligned}
$}
\end{equation}
Além disso, se um argumento $A_i$ é um feixe localmente livre, pode
ser suprimido, com a condição de substituir um dos $A_j$ ($j\neq i$) por
$A_j\otimes A_i$, ou um dos $B_k$ por $B_k\otimes A_i'$ (onde
$A_i'=\Homsh_{\Osh_X}(A_i,\Osh_X)$), e tem-se uma regra de cálculo análoga no
caso em que um dos argumentos $B_i$ é localmente livre. Em particular, sempre
se pode suprimir um argumento que seja idêntico a $\Osh_X$. Se todos os
argumentos são localmente livres, salvo no máximo um dos argumentos $B_i$, a regra
que acaba de ser enunciada dá um isomorfismo funtorial
\begin{equation}\tag{8.5}
  T_{r}^{s*}(A_1,\ldots,A_r;B_1,\ldots,B_s)
  =H^{*}\bigl(X,A_1'\otimes\cdots\otimes A_r'
    \otimes B_1\otimes\cdots\otimes B_s\bigr)
\end{equation}
(pois se reduz ao caso em que $r=0$, $s=1$, onde isto é imediato; também se
pode utilizar diretamente a sequência espectral de termo inicial (8.2)). Para os
$\underline{T}_{r}^{s}$ se definem operações correspondentes às anteriores.
As relações entre as diversas operações assim
introduzidas são as mesmas que para as operações análogas do cálculo
tensorial.

Seja $n$ a dimensão de $X$. Aplicando sucessivamente uma composição
tensorial (8.3) e contrações (8.4) sobre os argumentos repetidos, obtém-se
um acoplamento
\begin{equation}\tag{8.6}
\resizebox{\textwidth}{!}{$
  (T_{r}^{s})^{p}(A_1,\ldots;B_1,\ldots)
  \times
  (T_{s}^{r})^{n-p}(B_1,\ldots;A_1,\ldots,A_r\otimes\Omega_X^{n})
  \longrightarrow H^{n}(X,\Omega_X^{n}).
$}
\end{equation}

\result{TEOREMA}{6} Se $X$ é uma variedade projetiva não singular, então
os acoplamentos (8.6) são dualidades.

Isto resulta de maneira puramente formal do corolário do teorema 3. Com efeito,
deste corolário resulta facilmente que, se $K$ é um complexo de feixes
algébricos coerentes localmente livres, então a hipercohomologia de $X$ com
respeito a $K$ está em dualidade com a hipercohomologia de $X$ com respeito a
$K'\otimes\Omega_X^{n}$ mediante os acoplamentos naturais
\begin{equation}\tag{8.7}
  \mathbf R^{p}\Gamma(K)
  \times\mathbf R^{n-p}\Gamma(K'\otimes\Omega_X^{n})
  \longrightarrow\mathbf R^{n}\Gamma(\Omega_X^{n})
  =H^{n}(X,\Omega_X^{n}).
\end{equation}
Isto se vê utilizando a sequência espectral de termo inicial
$H^{p}(H^{q}(X,K))$ e a sequência espectral análoga para
$K'\otimes\Omega_X^{n}$. Do resultado anterior se deduz o teorema 6
utilizando a definição (8.1).

\src{192}
\subsection*{OBSERVAÇÕES}

\noindent 1\textsuperscript{o} Para as definições anteriores ao teorema 6,
não era necessário que $X$ fosse não singular, pois não era
indispensável trabalhar apenas com resoluções finitas. Mas se $X$
é singular, já não se poderá afirmar, a priori, que os
$(\underline{T}_{r}^{s})^{p}(A_1,\ldots;B_1,\ldots)$ sejam feixes
coerentes, pois no complexo de feixes
\[
  \Homsh_{\Osh_X}\bigl(L(A_1)\otimes\cdots,L(B_1)\otimes\cdots\bigr)
\]
haverá uma infinidade de componentes de um grau total dado.

\medskip
\noindent 2\textsuperscript{o} Verifica-se facilmente que, nas fórmulas
(8.1), pode-se substituir um dos $L(B_i)$ por $B_i$. Tendo em conta a
proposição 3, isto demonstra que se tem
\begin{equation}\tag{8.8}
\left\{
\begin{aligned}
  \underline{T}_{1}^{1*}(A;B)&=\Extsh_{\Osh_X}^{*}(A,B),\\
  T_{1}^{1*}(A;B)&=\operatorname{Ext}_{\Osh_X}^{*}(X;A,B).
\end{aligned}
\right.
\end{equation}
Em particular, tomando em (8.6) $r=s=1$ e $A_1=\Osh_X$, recupera-se o
teorema 3. A fórmula (8.8) implica também
\[
  T_{0}^{1*}(B)=H^{*}(X,B),
  \qquad
  T_{1}^{0*}(A)=\operatorname{Ext}_{\Osh_X}^{*}(X;A,\Osh_X).
\]

\medskip
\noindent 3\textsuperscript{o} Vê-se na fórmula (8.1) que, em geral, os
funtores $\underline{T}_{r}^{s*}$ e $T_{r}^{s*}$ têm componentes de
graus positivos e negativos. Utilizando a observação 2, vê-se que, se a
dimensão de $X$ é $n$, então as componentes não nulas de
$\underline{T}_{r}^{s*}$ estão compreendidas entre $-(s-1)n$ e $rn$, se $s>0$,
entre $0$ e $rn$, se $s=0$, e as componentes não nulas de $T_{r}^{s*}$
estão compreendidas entre $-(s-1)n$ e $(r+1)n$, se $s>0$, entre $0$ e $(r+1)n$,
se $s=0$ (e inclusive, salvo erro, se $r>0$, entre $-(s-1)n$ e $rn$, resp. entre
$0$ e $rn$).

\section*{BIBLIOGRAFIA}

\noindent[1] Cartier (Pierre). --- Les groupes
$\operatorname{Ext}^{s}(A,B)$, Séminaire A. Grothendieck : Algèbre homologique,
t. 1, 1957, n\textsuperscript{o} 3.

\noindent[2] Grothendieck (Alexandre). --- Sur quelques points d'algèbre
homologique, \emph{Tohoku Math. J.}, t. 9, 1957, p. 119--183.

\noindent[3] Serre (Jean-Pierre). --- Faisceaux algébriques cohérents,
\emph{Annals of Math.}, t. 61, 1955, p. 197--278.

\noindent[4] Serre (Jean-Pierre). --- Sur la dimension homologique des anneaux
et des modules noethériens, \emph{Proc. Intern. Symp. on Alg. Number Theory}
[1955, Tokyo et Nikko]. --- Tokyo, Science Council of Japan, 1956 ;
p. 175--189.

\clearpage
\src{193}
\section*{ADENDA}

As dificuldades assinaladas na observação da página 13 estão atualmente
completamente resolvidas, graças a uma extensão dos teoremas de dualidade a
variedades quaisquer (com singularidades arbitrárias). Para esta teoria, que pode
ser formulada no quadro geral da « teoria dos esquemas », remetemos
aos « Éléments de Géométrie algébrique » de J. Dieudonné e A. Grothendieck,
atualmente em preparação. Encontrar-se-ão algumas indicações em:

\begin{quote}
Grothendieck (A.). --- Cohomology theory of algebraic varieties,
\emph{Congrès int. des Mathématiciens}, 1958, Edinburgh (de próxima publicação).
\end{quote}

\begin{flushright}
[Abril de 1959]
\end{flushright}

\begin{center}\rule{36mm}{0.5pt}\end{center}

\par

\def\FgaUnitKey{e149}\clearpage
\chapter*{Errata da exposição n\textsuperscript{o} 149}
\addcontentsline{toc}{chapter}{Errata da exposição n\textsuperscript{o} 149}
\setcounter{footnote}{0}
\markboth{Errata da exposição 149}{Errata da exposição 149}
\sid{e149}
\src{298}

\noindent\textbf{1. Teorema de dualidade para os feixes algébricos
coerentes.}

\begin{center}
[Seminário Bourbaki, t. 9, 1956/57, n\textsuperscript{o} 149, 25 p.]
\end{center}

\noindent\textbf{Observação sobre a página 14.} --- Como assinalei em minha
conferência no Congresso Internacional de Matemáticos de 1958,\footnote{
Grothendieck (Alexander). --- The cohomology theory of abstract algebraic
varieties, \emph{Proceedings of the International Congress of Mathematicians}
[1958, Edinburgh], p. 103--118. --- Cambridge, at the University Press, 1960.}
as questões levantadas aqui estão agora completamente resolvidas.

O leitor encontrará mais informações sobre a dualidade dos feixes
coerentes em \emph{loco citato}\footnotemark[1], p. 112--115, em EGA III,
2\textsuperscript{a} parte (em preparação) e, por enquanto, em SGA 1962.
Um tratamento mais sistemático encontrar-se-á em um capítulo posterior de EGA
(capítulo IX no plano previsto).
% Explicit module paragraph boundary; no prose change.
\par

\par

\def\FgaUnitKey{182}\unit{Exposição 182 — Geometria formal e geometria algébrica}
\sid{182}
\markboth{Exposição 182}{Exposição 182}
\src{193}
\seminar{(Maio de 1959)}
\paperhead{GEOMETRIA FORMAL E GEOMETRIA ALGÉBRICA}%
  {por Alexander GROTHENDIECK}

\fsection{1}{Esquemas}

Sabe-se que um espaço algébrico afim definido sobre um corpo $k$ é
essencialmente determinado por sua álgebra afim $A$ (anel das funções
regulares definidas sobre $k$), pois os morfismos de espaços algébricos
$X\longrightarrow Y$ correspondem biunívocamente aos homomorfismos de
$k$-álgebras $A(Y)\longrightarrow A(X)$. O álgebra afim correspondente a um
espaço algébrico é uma $k$-álgebra de tipo finito e, do ponto de vista
«clássico», não tem elementos nilpotentes; reciprocamente, toda álgebra desse
tipo obtém-se como álgebra afim de um espaço algébrico definido sobre $k$.
Dispõe-se assim de um conhecido dicionário que permite interpretar, em termos
de álgebra comutativa, as situações relativas a espaços algébricos afins.
Há muito se observou que, desse modo, obtêm-se enunciados mais gerais, pois em
geral já não é necessário supor que os anéis em questão sejam do tipo recém
considerado, e a hipótese noetheriana costuma bastar. Em particular, haja ou
não um corpo de base dado, não há razão para excluir o caso em que esses anéis
contenham elementos nilpotentes. Até agora, os geômetras recusaram-se a levar
em conta essas indicações e obstinaram-se em limitar-se a considerar álgebras
afins sem elementos nilpotentes, isto é, espaços algébricos cujos feixes
estruturais não têm elementos nilpotentes (e, na maioria das vezes, até mesmo
espaços algébricos «absolutamente irredutíveis»). O conferencista considera que
essa atitude constituiu um sério obstáculo ao desenvolvimento de métodos
verdadeiramente naturais em geometria algébrica.

Seja $A$ um anel comutativo. É bem sabido que o conjunto
$X=\operatorname{Spec}(A)$ dos ideais primos de $A$ está provido de uma
topologia natural, a «topologia de Zariski» ou topologia espectral. Por outra
parte, há um feixe de anéis comutativos $\Osh_X$ sobre $X$, cujo talo em
$\mathfrak p\in X$ é o anel localizado $A_{\mathfrak p}$ e cujo anel de
seções identifica-se com $A$. Assim, $X$ torna-se um \emph{espaço
anelado}, chamado \emph{espectro primo} de $A$. Um homomorfismo de anéis
$f:A\longrightarrow B$ define um morfismo de espaços anelados
$f':\operatorname{Spec}(B)\longrightarrow\operatorname{Spec}(A)$; a
aplicação subjacente de conjuntos não é outra senão
$\mathfrak p\longmapsto f^{-1}(\mathfrak p)$. Os morfismos

\src{194}
de espaços anelados de $\operatorname{Spec}(B)$ em
$\operatorname{Spec}(A)$ obtidos desse modo são exatamente aqueles para
os quais os homomorfismos $\Osh_x\longrightarrow\Osh_y$ ($x=f'(y)$) são
locais (isto é, a imagem inversa do ideal maximal é o ideal maximal).

Chama-se \emph{esquema afim} um espaço anelado isomorfo a algum
$\operatorname{Spec}(A)$, e \emph{preesquema} um espaço anelado localmente
afim, isto é, tal que cada ponto tem uma vizinhança aberta que, com a
estrutura induzida, é um esquema afim. Os \emph{morfismos} de preesquemas se
definem de maneira evidente; localmente correspondem a homomorfismos de anéis.

Quando se fixa um preesquema $S$ e se consideram morfismos de preesquemas
$X\longrightarrow S$, o espaço $S$ desempenha o papel de um corpo ou de um
anel de base (ou, melhor, de um espaço de base numa fibração). Diz-se
então que $X$ é um $S$-preesquema; se $S=\operatorname{Spec}(A)$, isto
significa também que $\Osh_X$ é um feixe de $A$-álgebras. Assim, todo preesquema
pode ser considerado de maneira única como um $\mathbf Z$-preesquema. Naturalmente,
os $S$-preesquemas formam uma categoria; além disso, demonstra-se que nessa
categoria sempre existe o produto de dois objetos $X,Y$, denotado por
$X\times_S Y$. Essa noção de produto permite definir a mudança de base de um
$S$-preesquema correspondente a um morfismo $S'\longrightarrow S$: com efeito,
$X\times_S S'$ pode ser considerado como um $S'$-preesquema.

Diz-se que $X$ é \emph{separado} sobre $S$ se a diagonal de
$X\times_S X$ é fechada. Chama-se \emph{esquema} um preesquema separado
sobre $\mathbf Z$; ele é então separado sobre qualquer base. Para simplificar,
doravante falaremos apenas de esquemas e, além disso, os suporemos
\emph{noetherianos}, isto é, uniões finitas de abertos afins que são
espectros de anéis noetherianos. Diz-se que $X$ é \emph{de tipo finito
sobre} $S$ se, para todo aberto afim $U$ de $S$, sua imagem inversa em $X$ é
uma união finita de abertos afins cujos anéis são álgebras de tipo finito
sobre o anel de $U$. São esses $S$-esquemas os que se prestam a um estudo
propriamente geométrico. Em particular, para todo $s\in S$, a fibra
$f^{-1}(s)$ de $X$ sobre $s$ é um esquema algébrico sobre o corpo residual
$\kappa(s)$ do anel local $\Osh_s$ de $s$ em $S$. Assim, em certa medida, $X$
pode ser considerado como uma família de «espaços algébricos» $f^{-1}(s)$, onde
o parâmetro $s$ percorre $S$ (isto é, do ponto de vista local, o
conjunto dos ideais primos de um anel dado). Naturalmente, os
$\kappa(s)$ podem ter características distintas. Se
$S=\operatorname{Spec}(k)$, onde $k$ é um corpo, recupera-se essencialmente a
noção usual de «espaço algébrico», com a única diferença de que agora o
feixe estrutural pode ter elementos nilpotentes.

\src{195}
Inspirando-se em noções bem conhecidas, define-se a noção de
\emph{morfismo projetivo} e, mais geralmente, a de \emph{morfismo próprio}.
Tal morfismo é de tipo finito; além disso, transforma partes fechadas em partes
fechadas e conserva essa propriedade sob qualquer mudança de base.

Sendo $X$ um esquema (noetheriano, como sempre), o feixe $\Osh_X$ é um
\emph{feixe coerente de anéis} no sentido de [2]. Portanto, os feixes
coerentes de módulos sobre $X$ são também os feixes que localmente são
isomorfos ao conúcleo de um homomorfismo
$\Osh_X^{m}\longrightarrow\Osh_X^{n}$.

\fsection{2}{Esquemas formais}

Sejam $X$ um esquema e $X'$ uma parte fechada de $X$. Então existe um
subfeixe coerente $J$ de $\Osh_X$ tal que
$X'=\operatorname{Supp}(\Osh_X/J)$ (e inclusive existe um máximo). Provido de
$\Osh_X/J$, o espaço $X'$ torna-se um esquema, denotado por $X_0$;
tal esquema se chama «subesquema fechado de $X$». Para todo $n$ também se
pode considerar $X'$ provido de $\Osh_X/J^{n+1}$, denotado por $X_n$;
é um subpreesquema fechado de $X$ cujo conjunto subjacente continua sendo $X'$,
mas tem outro feixe estrutural, a saber,
$\Osh_{X_n}=\Osh_X/J^{n+1}$. Evidentemente, os $\Osh_{X_n}$ formam um sistema
projetivo de feixes de anéis sobre $X'$, cujo limite projetivo
$\overline{\Osh}_X$ chama-se \emph{completamento formal de $\Osh_X$ ao longo
de $X'$}. Provido desse feixe de anéis, $X'$ chama-se \emph{completamento
formal de $X$ ao longo de $X'$}; é, pois, um espaço anelado, mas em
geral não um esquema. Para todo feixe coerente $F$ sobre $X$, pode-se considerar
analogamente o completamento formal $\overline F=\varprojlim_n F_n$ de $F$ ao
longo de $X'$ (onde $F_n=F\otimes_{\Osh_X}\Osh_X/J^{n+1}$); é um feixe de
módulos sobre $\overline X$. Suas seções chamam-se «seções formais de $F$
ao longo de $X'$» e identificam-se com os elementos de
$\varprojlim_n\Gamma(X',F_n)$. Para $F=\Osh_X$ obtêm-se as «funções
holomorfas» de $X$ ao longo de $X'$ no sentido de Zariski (terminologia que
não seguiremos por suas interferências com a terminologia clássica).

Chama-se \emph{esquema formal} (subentende-se: noetheriano) um espaço
topológico $\mathfrak X$, provido de um feixe de anéis topológicos
$\Osh_{\mathfrak X}$ que satisfaz a condição seguinte: há um isomorfismo
de feixes de anéis topológicos
$\Osh_{\mathfrak X}=\varprojlim_n\Osh_n$, onde os $\Osh_n$ formam um sistema
projetivo de feixes de anéis sobre $\mathfrak X$, cada um dos quais faz
de $\mathfrak X$ um esquema $\mathfrak X_n$, e tal que, para $m\geq n$, o
homomorfismo $\Osh_m\longrightarrow\Osh_n$ é sobrejetivo e tem por núcleo
$J_m^{n+1}$, onde $J_m$ é

\src{196}
o núcleo de $\Osh_m\longrightarrow\Osh_0$. Demonstra-se então que
$\Osh_{\mathfrak X}$ é um \emph{feixe coerente de anéis locais
noetherianos}.

Em virtude das definições, um completamento formal $\overline X$ como o
anterior é um esquema formal e, reciprocamente, todo esquema formal é
\emph{localmente} desse tipo. De fato, dar um esquema formal \emph{afim} (isto
é, tal que $\mathfrak X_0$ seja afim, o que implica que todos os
$\mathfrak X_n$ também o sejam) equivale a dar um anel topológico noetheriano,
separado e completo para a topologia $J$-ádica.

As definições usuais de morfismo, morfismo de tipo finito, morfismo
próprio, etc., para os esquemas ordinários estendem-se sem dificuldade aos
esquemas formais.

\fsection{3}{Os três teoremas fundamentais}

Seja $f:X\longrightarrow Y$ um morfismo próprio de esquemas (noetherianos, como
sempre), seja $Y'$ uma parte fechada de $Y$ e seja $X'$ sua imagem inversa em
$X$; consideremos os completamentos formais correspondentes $\overline Y$ e
$\overline X$. Então $f$ induz um morfismo
$\overline f:\overline X\longrightarrow\overline X$ de esquemas formais, que
além disso é próprio. Se $F$ é um feixe coerente sobre $X$, então $\overline F$
é um feixe coerente sobre $\overline X$. No teorema 1, esquecem-se $X,Y,F$ e
consideram-se apenas o morfismo próprio $\overline f$ de esquemas formais e o
feixe coerente $\overline F$ sobre $\overline X$. (Não obstante, o conferencista
só escreveu uma demonstração completa para o caso em que se parte de
$X,Y,f,F$.)

\result{TEOREMA}{1 (teorema de finitude)}
\begin{enumerate}
  \item[\textit{i.}] Os $R^{q}\overline f_{*}(\overline F)$ são feixes
    coerentes sobre $\overline Y$.
  \item[\textit{ii.}] Os homomorfismos naturais
    \[
      R^{q}\overline f_{*}(\overline F)
      \longrightarrow \varprojlim_n R^{q}f_{n*}(F_n)
    \]
    são isomorfismos.
\end{enumerate}

Neste enunciado supõe-se escolhido um subfeixe coerente $J$ de $\Osh_Y$ que
define $Y'$; por imagem inversa obtém-se um subfeixe coerente de $\Osh_X$ que
define $X'$ e, por conseguinte, se definem $F_n,X_n,Y_n$ e
$f_n:X_n\longrightarrow Y_n$ como no n\textsuperscript{o} 2. Se partirmos de
um morfismo próprio qualquer entre dois esquemas formais, são evidentes as
pequenas mudanças que devem ser feitas ao explicitar essas notações.

O teorema 1 referia-se apenas à «cohomologia formal». O teorema seguinte

\src{197}
a relaciona com a «cohomologia algébrica» e é análogo a um conhecido teorema
de Serre [4] sobre a comparação entre cohomologia algébrica e cohomologia
analítica.

\result{TEOREMA}{2 (Primeiro teorema de comparação)} Os
$R^{q}f_{*}(F)$ são feixes coerentes sobre $Y$ (o que é um caso particular
do teorema 1), e os homomorfismos naturais
\[
  \overline{R^{q}f_{*}(F)}
  \longrightarrow \varprojlim_n R^{q}f_{n*}(F_n)
\]
são isomorfismos.

\result{COROLÁRIO}{1} Têm-se isomorfismos canônicos:
\[
  \overline{R^{q}f_{*}(F)}=R^{q}\overline f_{*}(\overline F).
\]

Para $q=0$, esse corolário é uma generalização do «teorema fundamental das
funções holomorfas» de Zariski, do qual deduziremos uma generalização do
«teorema de conexão» de Zariski. Observe-se além disso que, enquanto o teorema 1
(\textit{ii}) é trivial para $q=0$, de modo algum ocorre o mesmo com o
teorema 2 nem com sua formulação equivalente (corolário 1). Com efeito, a
demonstração procede por indução descendente sobre $q$ (é trivial para $q$
grande, pois então ambos os membros são nulos), de modo que o caso $q=0$
aparece como a última etapa da indução e, por assim dizer, como o caso «mais
difícil».

\result{COROLÁRIO}{2} Suponhamos $Y=\operatorname{Spec}(A)$ e que $Y'$ está
definido pelo ideal $J$ de $A$. Então, para todo feixe coerente $F$ sobre
$X$, os $H^{q}(X,F)$ são $A$-módulos de tipo finito cujos completamentos
$J$-ádicos são os $H^{q}(\overline X,\overline F)$.

Aplicando, por fim, este corolário a
$H=\Homsh_{\Osh_X}(F,G)$, obtém-se:

\result{COROLÁRIO}{3} Suponhamos $Y=\operatorname{Spec}(A)$ e que $Y'$ está
definido pelo ideal $J$ de $A$. Sejam $F,G$ dois feixes coerentes sobre $X$;
então $\operatorname{Hom}(F,G)$ é um módulo de tipo finito sobre $A$, cujo
completamento $J$-ádico identifica-se com
$\operatorname{Hom}(\overline F,\overline G)$.

Naturalmente, a aplicação natural
$\operatorname{Hom}(F,G)\longrightarrow
\operatorname{Hom}(\overline F,\overline G)$ é a que associa a um
homomorfismo $u:F\longrightarrow G$ sua extensão «por continuidade»
$\overline u:\overline F\longrightarrow\overline G$ (com o que
$\overline F$ se torna um funtor de $F$).

Suponhamos agora que $A$ é separado e completo para sua topologia $J$-ádica;
então os corolários 2 e 3 precedentes dão:
\[
  H^{q}(X,F)=H^{q}(\overline X,\overline F),
  \qquad
  \operatorname{Hom}(F,G)
  =\operatorname{Hom}(\overline F,\overline G).
\]

\src{198}
Esta última identidade mostra que a categoria dos feixes coerentes sobre
$X$ se identifica com uma \emph{subcategoria} (cujos morfismos são os
morfismos induzidos) da categoria dos feixes coerentes sobre
$\overline X$. De fato, tem-se inclusive:

\result{TEOREMA}{3} Para que um feixe de módulos sobre $\overline X$ seja
coerente, é necessário e suficiente que seja isomorfo a um feixe da forma
$\overline F$, onde $F$ é um feixe coerente sobre $X$ (determinado salvo
isomorfismo canônico em virtude do teorema 2, corolário 3). [Lembremos que
agora $Y=\operatorname{Spec}(A)$, sendo $A$ um anel topológico noetheriano,
separado e completo para a topologia $J$-ádica.]

\result{COROLÁRIO}{1} Os subesquemas fechados de $X$ correspondem
biunívocamente aos subesquemas formais fechados de $\overline X$.

Com efeito, correspondem aos subfeixes coerentes de $\Osh_X$ e,
respectivamente, de $\overline{\Osh}_X$. Considerando os grafos de morfismos
como subesquemas fechados, deduz-se do corolário 1:

\result{COROLÁRIO}{2} Sejam $X,Z$ dois esquemas próprios sobre $A$ (anel
noetheriano, separado e completo para a topologia $J$-ádica). Então a
aplicação $g\longmapsto\overline g$ define uma correspondência biunívoca entre
os $Y$-morfismos de $X$ em $Z$ e os $\overline Y$-morfismos de $\overline X$
em $\overline Z$.

Em outros termos, os esquemas algébricos próprios sobre $A$ aparecem como uma
\emph{subcategoria} (cujos morfismos são os morfismos induzidos) da
categoria dos esquemas formais próprios sobre $\overline Y$. No entanto,
deve-se advertir que existem esquemas formais próprios sobre $\overline Y$ que
não são «algebrizáveis», isto é, isomorfos a algum $\overline X$ com $X$
próprio sobre $A$ (do mesmo modo que existem variedades analíticas complexas
compactas que não provêm de variedades algébricas definidas sobre o corpo
dos complexos). Tais esquemas formais aparecem de maneira natural na
«teoria dos módulos». Assinalemos, não obstante, um caso particular interessante em
que um esquema formal é algebrizável:

\result{TEOREMA}{4} Seja $A$ um anel noetheriano local completo de corpo
residual $k$, e seja $\mathfrak X$ um esquema formal próprio sobre $A$ (provido
de sua topologia $\mathfrak m(A)$-ádica). Supõe-se:
\begin{enumerate}
  \item[\textit{i.}] que os anéis locais de $\Osh_{\mathfrak X}$ são
    módulos planos sobre $A$ ou, o que equivale ao mesmo, se se dota a
    $\Osh_{\mathfrak X}$ e a $A$ da filtração definida pelas potências do
    ideal maximal de $A$, para os graduados associados tem-se a relação
    \[
      \operatorname{gr}(\Osh_{\mathfrak X})
      \simeq
      \operatorname{gr}^{0}(\Osh_{\mathfrak X})
      \otimes_k\operatorname{gr}(A).
    \]
  \item[\textit{ii.}] Considerando
    $\mathfrak X_0=\mathfrak X\otimes_A k$ como um esquema algébrico sobre
    $k$, tem-se
    \[
      H^2(\mathfrak X_0,\Osh_{\mathfrak X_0})=0.
    \]
  \item[\textit{iii.}] $\mathfrak X_0$ é projetivo.
\end{enumerate}
Sob essas condições, $\mathfrak X$ é algebrizável e, mais precisamente, é
isomorfo a $\overline X$, onde $X$ é um $A$-esquema projetivo.

As condições \textit{(ii)} e \textit{(iii)} são satisfeitas, em particular,
se $\mathfrak X_0$ é uma curva lisa sobre $k$, e o teorema 4 pode ser aplicado
em particular na «teoria de módulos» para as curvas de gênero dado\ldots{}
Indiquemos como se demonstra o teorema 4: demonstra-se (cf. a proposição 3
mais adiante) que \textit{(i)} e \textit{(ii)} implicam que todo feixe coerente
sobre $\mathfrak X_0$, localmente isomorfo ao feixe fundamental, obtém-se por
redução a partir de um feixe da mesma natureza sobre $\mathfrak X$. Parte-se
então de um feixe «amplo» sobre $\mathfrak X_0$ (existe um por
\textit{(iii)}), levanta-se esse feixe a um feixe invertível sobre $\mathfrak X$ e,
utilizando o teorema 1, se demonstra que um múltiplo deste último define uma
imersão de $\mathfrak X$ no completamento formal de um esquema
$\mathbf P_A^r$ («projetivo tipo» de dimensão $r$ sobre $A$).

Para a demonstração dos teoremas 1 a 3, veja-se [1].

\src{199}
\fsection{4}{Aplicação ao teorema de conexão e ao «Main theorem» de
Zariski}

Seja $f:X\longrightarrow Y$ um morfismo próprio de esquemas. Pelo teorema de
finitude, $f_*(\Osh_X)=\mathcal A$ é um feixe coerente sobre $Y$; além disso, é um
feixe de álgebras comutativas, pelo que corresponde a um $Y$-esquema
$g:Y'\longrightarrow Y$ finito sobre $Y$ (definido pela condição de ser afim
sobre $Y$, isto é, que a imagem inversa de um aberto afim seja afim, e por
$g_*(\Osh_{Y'})=\mathcal A$). É imediato que $f$ se fatora então
canonicamente como $f=gf'$, onde $f':X\longrightarrow Y'$ é um morfismo de
$X$ em $Y'$ que desta vez satisfaz $f'_*(\Osh_X)=\Osh_{Y'}$. Essa fatoração
de $f$ chama-se \emph{fatoração de Stein} de $f$. Aplicando o primeiro
teorema de comparação, corolário 1, a $f'$ e à parte de $Y'$ reduzida a um
ponto $y'$, obtém-se que $f'^{-1}(y')=X'$ é conexo (de outro modo, as
seções formais de $X$ ao longo de $X'$ não formariam um anel local,
enquanto o completamento de
$f'_*(\Osh_X)_{y'}=\Osh_{Y',y'}$ é local!). Demonstramos:

\result{TEOREMA}{5 («Teorema de conexão» de Zariski)} Seja
$f:X\longrightarrow Y$ um morfismo próprio. Então $f$ se fatora de maneira
única (salvo isomorfismo) como $f=gf'$, onde $g:Y'\longrightarrow Y$ é finito
e $f':X\longrightarrow Y'$ satisfaz $f'_*(\Osh_X)=\Osh_{Y'}$ (de onde
$g_*(\Osh_{Y'})=f_*(\Osh_X)$). As fibras de $f'$ são conexas; isto é, o
conjunto das componentes conexas de uma fibra $f^{-1}(y)$ de $f$ está em
correspondência biunívoca com o conjunto dos pontos de $Y'$ sobre $y$, isto
é, com o conjunto dos ideais maximais de $f_*(\Osh_X)_y$.

Daqui se deduzem imediatamente as variantes usuais do teorema de
conexão. Assinalemos aqui apenas o

\src{200}
\result{COROLÁRIO}{1} Para que um ponto $x$ de $X$ esteja isolado em sua fibra
$f^{-1}(y)$, é necessário e suficiente que a fibra $f'^{-1}(y')$ (onde
$y'=f'(x)$) se reduza a $x$, ou que $f'$ induza um isomorfismo de uma
vizinhança de $x$ sobre uma vizinhança de $y'$. O conjunto desses pontos é um
aberto $U$, e $f'$ induz um isomorfismo de $U$ sobre um aberto de $Y'$.

Para demonstrar que $f'$ é um isomorfismo local em $x$, observa-se que $f'$
induz um isomorfismo $\Osh_{y'}\longrightarrow\Osh_x$, como se vê graças a
$f'_*(\Osh_X)=\Osh_{Y'}$ e ao fato de que (sendo $f'$ uma aplicação fechada
cuja fibra em $y'$ se reduz a $x$) os $f'^{-1}(V)$ percorrem um sistema
fundamental de vizinhanças de $x$ quando $V$ recorre um sistema fundamental de
vizinhanças de $y'$. Daqui se deduz imediatamente o resultado seguinte, devido a
Chevalley no caso «geométrico».

\result{COROLÁRIO}{2} Para que $f$ seja um morfismo finito, é necessário e
suficiente que seja próprio e que suas fibras sejam finitas.

Se assim ocorre, $f'$ é efetivamente um isomorfismo pelo que precede.

Seja $f:X\longrightarrow Y$ um morfismo não necessariamente próprio, mas
suponhamos que $X$ está contido como aberto em um $Y$-esquema próprio
$\widetilde f:\widetilde X\longrightarrow Y$ (o que ocorre, em particular,
se $f$ é quase-projetivo). Aplicando o corolário 1, vê-se que $\widetilde f'$
induz um isomorfismo do conjunto $U$ dos pontos de $X$ isolados em sua fibra
sobre um aberto de $Y'$ (e que $U$ é, com efeito, aberto). Deduz-se a
seguinte variante global do «Main Theorem» de Zariski.

\result{TEOREMA}{6 («Main Theorem» de Zariski)} Seja
$f:X\longrightarrow Y$ um morfismo de tipo finito. Então o conjunto $U$ dos
pontos de $X$ que estão isolados em sua fibra é aberto e, se $f$ é
quase-projetivo, $U$ é
$Y$-isomorfo a uma parte aberta de um esquema $Y'$ finito sobre $Y$.

Como um morfismo de tipo finito é localmente afim e, \textit{a fortiori},
localmente quase-projetivo, do teorema 6 deduzem-se imediatamente as variantes
locais usuais do Main Theorem.

\fsection{5}{Aplicação ao estudo cohomológico dos morfismos próprios e
planos}

Seja $f:X\longrightarrow Y$ um morfismo próprio, e seja $F$ um feixe coerente sobre
$X$; supõe-se que $F$ é $Y$-plano, isto é, os $F_x$ são módulos planos sobre os
anéis $\Osh_y$ ($y=f(x)$). Isto significa também que, para todo $y\in Y$, se
filtrarmos $F$ ao longo da fibra $f^{-1}(y)$ mediante os
$\mathfrak m_y^nF$ (sendo $\mathfrak m_y$ o ideal maximal de $\Osh_y$), o
graduado associado é isomorfo a
$(F/\mathfrak m_yF)\otimes_{\kappa(y)}\operatorname{gr}(\Osh_y)$; em outros
termos, tem-se
\[
  \mathfrak m_y^nF/\mathfrak m_y^{n+1}F
  =F_y\otimes_{\kappa(y)}
  (\mathfrak m_y^n/\mathfrak m_y^{n+1})
\]
para todo inteiro $n$, onde $X_y$ denota a fibra $f^{-1}(y)$ considerada como
esquema próprio sobre o corpo residual $\kappa(y)$ de $y$, e $F_y$ o feixe
$F/\mathfrak m_yF$ induzido por $F$ sobre $X_y$. Tendo em conta esse
isomorfismo e o teorema 2, obtêm-se cotas e, às vezes, cálculos dos
$R^qf_*(F)$ em uma vizinhança de $y$, conhecendo a cohomologia de $X_y$ com
coeficientes em $F_y$. O teorema 2 assume aqui a forma
\[
  \widehat{R^qf_*(F)_y}
  =\varprojlim_n H^q(X_y,F/\mathfrak m_y^nF).
\]
Assinalaremos apenas a consequência seguinte:

\src{201}
\result{PROPOSIÇÃO}{1} Sejam $f:X\longrightarrow Y$ um morfismo próprio e $F$
um feixe sobre $X$, coerente e $Y$-plano. Seja $y\in Y$, seja $q$ um inteiro, e
suponhamos que $H^q(X_y,F_y)=0$. Então $R^qf_*(F)$ é nulo em uma vizinhança de
$y$ e, para todo $n$, o homomorfismo natural
\[
  R^{q-1}f_*(F)_y\longrightarrow
  H^{q-1}(X_y,F/\mathfrak m_y^nF)
\]
é sobrejetivo.

Em particular, se $f$ é um morfismo plano (isto é, $\Osh_X$ é $Y$-plano),
então todo feixe coerente localmente livre $F$ sobre $X$ é $Y$-plano. Sejam
$F,G$ dois feixes desse tipo; aplicando a proposição 1 a
$\Homsh_{\Osh_X}(F,G)$ e a $q=1$, obtém-se:

\result{TEOREMA}{7} Seja $f$ um morfismo próprio e plano, sejam $F,G$ dois feixes
coerentes localmente livres sobre $X$, seja $y\in Y$ e suponhamos que
\[
  H^1\bigl(X_y,\Homsh_{\Osh_{X_y}}(F_y,G_y)\bigr)=0,
\]
então todo homomorfismo $u_0:F_y\longrightarrow G_y$ é induzido por um
homomorfismo
\[
  u:F|V\longrightarrow G|V,
  \qquad V=f^{-1}(U),
\]
onde $U$ é uma vizinhança de $y$.

\result{COROLÁRIO}{} Se $u_0$ é um isomorfismo (resp. um monomorfismo, um
epimorfismo), então o mesmo ocorre com $u$ para $U$ suficientemente
pequeno.

Em particular:

\src{202}
\result{COROLÁRIO}{2} Seja $E_0$ um feixe coerente localmente livre sobre $X_y$
tal que
\[
  H^1\bigl(X_y;\Homsh_{\Osh_{X_y}}(E_0,E_0)\bigr)=0,
\]
então dois feixes localmente livres cujas restrições a $X_y$ sejam isomorfas
a $E_0$ são isomorfos em uma vizinhança de $X_y$.

Assim:

\result{COROLÁRIO}{3} Suponhamos $H^1(X_y,\Osh_{X_y})=0$. Então dois feixes
invertíveis sobre $X$ (isto é, localmente isomorfos a $\Osh_X$) cujas
restrições a $X_y$ sejam isomorfas são isomorfos.

Deduz-se:

\result{PROPOSIÇÃO}{2} Sejam $Y$ um esquema conexo e $E$ um feixe coerente localmente livre sobre $Y$; consideremos o fibrado de espaços projetivos
$X=\mathbf P(E)$ associado a $E$, provido de seu conhecido feixe invertível
$\Osh_X(1)$. Então todo feixe invertível $L$ sobre $X$ é isomorfo a um feixe da
forma
\[
  f^*(L')\otimes\Osh_X(n),
\]
onde $L'$ é um feixe invertível sobre $Y$ e $n$ um inteiro. Este último está
determinado de maneira única, e $L'$ está determinado salvo isomorfismo.

O corolário 3 precedente demonstra que $L$ é isomorfo a algum $\Osh_X(n)$ em
uma vizinhança de cada fibra. O resto é quase formal.

A proposição 2 permite determinar os $Y$-morfismos de $X=\mathbf P(E)$ em
outro fibrado projetivo. Em particular, obtém-se:

\result{COROLÁRIO}{} Seja $u$ um automorfismo de $X=\mathbf P(E)$. Então
existem um feixe invertível $L'$ sobre $Y$ e um isomorfismo $v$ de $E$ sobre
$E\otimes L'$ tais que $u$ é o isomorfismo correspondente
\[
  \mathbf P(E)\xrightarrow{\sim}\mathbf P(E\otimes L')=\mathbf P(E);
\]
o par $(v,L')$ está determinado salvo isomorfismo.

Seja $\Gamma$ o conjunto das classes de fibrados invertíveis $L'$ sobre $Y$
tais que $E\otimes L'$ seja isomorfo a $E$. Seus elementos são de torção, pois,
se $n$ é o posto de $E$, vê-se (tomando as potências exteriores $n$-ésimas)
que se deve ter $L'^{\otimes n}\simeq\Osh_Y$. O corolário pode ser expresso
então dizendo que se tem uma sequência exata de grupos
\[
  e\longrightarrow
  \operatorname{Aut}(E)/\Gamma(Y,\Osh_Y^*)
  \longrightarrow \operatorname{Aut}_Y(X)
  \longrightarrow \Gamma\longrightarrow e
\]
(deduzida, de resto, da sequência exata de cohomologia obtida a partir
da seguinte sequência exata de feixes de grupos
\[
  e\longrightarrow \Osh_Y^*
  \longrightarrow \underline{\operatorname{Aut}}(E)
  \longrightarrow \underline{\operatorname{Aut}}_Y(X)
  \longrightarrow e,
\]
onde $\Osh_Y^*$ é o feixe das «unidades» de $\Osh_Y$, identificado com o
centro de $\underline{\operatorname{Aut}}(E)$).

\src{203}
\fsection{6}{Aplicação a teoremas de existência e unicidade para feixes e
esquemas sobre um anel completo para a topologia \texorpdfstring{$J$}{J}-ádica}

O teorema 7 dava um resultado de unicidade para feixes coerentes localmente
livres utilizando os teoremas 1 e 2. Utilizando o teorema 3, deduzimos agora
teoremas de existência de feixes, de morfismos de esquemas ou de esquemas.
Doravante, $A$ denota um anel noetheriano local, separado e completo. O método
geral consiste sempre em efetuar construções formais, o que consiste
essencialmente em fazer geometria algébrica sobre um anel artiniano e extrair
delas conclusões de natureza «algébrica» mediante os três teoremas
fundamentais.

\result{PROPOSIÇÃO}{3} Seja $\mathfrak X$ um esquema formal próprio e plano
sobre $A$, e seja $F_0$ um feixe localmente livre sobre $X_0$ tal que
\[
  H^2\bigl(X_0,\Homsh_{\Osh_{X_0}}(F_0,F_0)\bigr)=0.
\]
Então existe um feixe localmente livre $F$ sobre $\mathfrak X$ que induz
sobre $X_0$ um feixe isomorfo a $F_0$. (Além disso, este $F$ é único salvo
isomorfismo se
$H^1(X_0,\Homsh_{\Osh_{X_0}}(F_0,F_0))=0$.)

Constroem-se sucessivamente feixes localmente livres $F_n$ sobre os $X_n$, cada
um dos quais induz o anterior. A construção de $F_n$ encontra uma
obstrução em
\[
  H^2\bigl(X_0,\Homsh_{\Osh_{X_0}}(F_0,F_0)\bigr)
  \otimes_{A/J}(J^n/J^{n+1}),
\]
que é, portanto, nula por hipótese. Utilizando agora o teorema 3, obtém-se:

\result{COROLÁRIO}{1} Seja $X$ um esquema próprio e plano sobre $A$, e seja $F_0$
como acima. Então existe um feixe localmente livre $F$ sobre $X$ que induz
sobre $X_0$ um feixe isomorfo a $F_0$. Além disso, este $F$ é único salvo
isomorfismo se
$H^1(X_0,\Homsh_{\Osh_{X_0}}(F_0,F_0))=0$.

Seja $X_0$ um esquema de tipo finito sobre o corpo $k$; supõe-se que $X_0$ é
liso (pelo que entendemos absolutamente liso) sobre $k$, mas não
necessariamente próprio sobre $k$. Seja $A$ um anel artiniano local de corpo
residual $k$. Interessa-nos encontrar os esquemas $X$ planos sobre $A$ tais que
$X\otimes_A k=X_0$ (este é o ponto de partida da «teoria de módulos» ou
das «variações de estrutura» de $X_0$). É equivalente dar um tal $X$, ou dar
sobre o espaço topológico $X_0$ um feixe $\Osh_X$ provido das estruturas
seguintes:
\begin{enumerate}
  \item[\textit{i.}] É um feixe de $A$-álgebras.
  \item[\textit{ii.}] Está provido de um homomorfismo de aumentação
    $\Osh_X\longrightarrow\Osh_{X_0}$ (compatível com as estruturas de
    $A$-álgebras); esses dados estão sujeitos às condições seguintes: a
    aumentação induz um isomorfismo
    $\Osh_X\otimes_A k\xrightarrow{\sim}\Osh_{X_0}$; $\Osh_X$ é plano sobre
    $A$, isto é, o graduado associado a $\Osh_X$, filtrado pelas potências
    do ideal maximal $\mathfrak m$ de $A$, é isomorfo a
    $\operatorname{gr}^0(\Osh_X)\otimes_k\operatorname{gr}(A)$; isto é, se
    têm isomorfismos
    \[
      \mathfrak m^n\Osh_X/\mathfrak m^{n+1}\Osh_X
      =\Osh_{X_0}\otimes_k(\mathfrak m^n/\mathfrak m^{n+1}).
    \]
\end{enumerate}
O fato fundamental é o seguinte:

\src{204}
\result{TEOREMA}{8} Seja $X_0$ um esquema de tipo finito e liso sobre o
corpo $k$, e suponhamos que $X_0$ é afim. Seja $A$ um anel local artiniano de
corpo residual $k$. Então existe um $A$-esquema $X$ plano sobre $A$ tal que
$X\otimes_A k=X_0$, e dois esquemas desse tipo são necessariamente isomorfos.

Observe-se que o isomorfismo em questão não é canônico, pois em geral $X$
terá $A$-automorfismos não triviais que induzem a identidade sobre $X_0$. Por
outro lado, em geral não há uma escolha «canônica» de um $X$ que satisfaça
as condições dadas, salvo quando $A$ é uma $k$-álgebra (caso de
características iguais), em cujo caso pode tomar-se $X=X_0\otimes_k A$, isto
é, $\Osh_X=\Osh_{X_0}\otimes_k A$ (seja ou não afim $X_0$). No caso de
características distintas, ignoro em geral, quando $X_0$ não é afim, se se
pode «levantar» $X_0$ a um $X$ definido sobre $A$. No entanto, seja $n$ um
inteiro $>0$, ponhamos $A_{n-1}=A/\mathfrak m^n$, e suponhamos que se tenha
levantado $X_0$ a um $A_{n-1}$-esquema plano $X_{n-1}$; pretende-se levantar
$X_{n-1}$ a um $A_n$-esquema plano $X_n$. Pelo teorema 8, já se sabe que isso
é possível localmente; por outro lado, verifica-se facilmente que, se $U_n$
levanta um aberto $U_{n-1}$ de $X_{n-1}$, então o feixe dos grupos de
automorfismos de $U_n$ (que induzem a identidade sobre $U_{n-1}$) é
canonicamente isomorfo a
\[
  \Theta_{X_0/k}\otimes_k(\mathfrak m^n/\mathfrak m^{n+1})
\]
restrito a $U_{n-1}$ e, em particular, é comutativo
($\Theta_{X_0/k}$ denota o feixe de germes de $k$-derivações sobre $X_0$).
Segue-se facilmente que há uma obstrução para construir $X_n$ que levante
$X_{n-1}$, a qual se encontra em
\[
  H^2(X_0,\Theta_{X_0/k})\otimes_k
  \mathfrak m^n/\mathfrak m^{n+1}.
\]
Por conseguinte:

\result{COROLÁRIO}{1} Seja $X_0$ um esquema de tipo finito liso sobre $k$, e
suponhamos $H^2(X_0,\Theta_{X_0/k})=0$. Então, para todo anel local
artiniano $A$ de corpo residual $k$, existe um $A$-esquema plano $X$ tal que
$X\otimes_A k=X_0$.

Além disso, se foi possível encontrar um $X$ plano sobre $A$ que levante $X_0$,
então, em virtude do teorema 8, o conjunto das classes (salvo isomorfismo)
de $A$-esquemas planos que levantam $X_0$ identifica-se com
\[
  H^1\bigl(X_0,\underline{\operatorname{Aut}}(X)\bigr),
\]
onde, naturalmente, $\underline{\operatorname{Aut}}(X)$ denota o feixe de
germes de automorfismos do feixe de $A$-álgebras $\Osh_X$ compatíveis com a
aumentação. A filtração de $\Osh_X$ define uma filtração de
$\underline{\operatorname{Aut}}(X)$, e o quociente deste feixe pelo $n$-ésimo
subgrupo da filtração é $\underline{\operatorname{Aut}}(X_n)$; o graduado
associado a essa filtração é comutativo e identifica-se com
$\Theta_{X_0/k}\otimes_k\operatorname{gr}(A)$. Em particular, se
$\mathfrak m^{n+1}$ é a primeira potência de $\mathfrak m$ que é nula,
então $F^n(\underline{\operatorname{Aut}}(X))$ (o último termo da
filtração) pertence ao centro de $\underline{\operatorname{Aut}}(X)$ e é
isomorfo a $\Theta_{X_0/k}\otimes_k\mathfrak m^n$; é também o feixe de
germes de automorfismos de $X$ que induzem a identidade sobre
$X_{n-1}=X\otimes_A A/\mathfrak m^n$. Utilizando estes resultados obtêm-se
imediatamente os enunciados seguintes:

\src{205}
\result{COROLÁRIO}{2} Seja $X_0$ um esquema de tipo finito liso sobre $k$, seja
$A$ um anel local artiniano de corpo residual $k$ e ideal maximal
$\mathfrak m$, suponhamos $\mathfrak m^{n+1}=0$, seja
$A_{n-1}=A/\mathfrak m^n$ e, por fim, seja $X_{n-1}$ um
$A_{n-1}$-esquema plano tal que $X_{n-1}\otimes_A k=X_0$. Então o conjunto
das classes (salvo um isomorfismo que induza a identidade sobre $X_{n-1}$)
de $A$-esquemas planos $X_n$ tais que
$X_n\otimes_A A_{n-1}=X_{n-1}$ é vazio ou é um espaço homogêneo principal
sob
\[
  H^1(X_0,\Theta_{X_0/k})\otimes_k\mathfrak m^n.
\]

(Observe-se que, em geral, não há um elemento origem privilegiado neste
último espaço, pois não há uma maneira privilegiada de levantar $X_{n-1}$ a
$X_n$).

\result{COROLÁRIO}{3} Seja $X_0$ um esquema de tipo finito liso sobre $k$, e
suponhamos $H^1(X_0,\Theta_{X_0/k})=0$. Então, para todo anel local
artiniano $A$ de corpo residual $k$, existe no máximo (salvo isomorfismo) um
$A$-esquema plano $X$ tal que $X\otimes_A k=X_0$.

Os corolários 1 e 3 implicam imediatamente os enunciados, de aparência mais
geral, que se obtêm supondo apenas que $A$ é um anel local
noetheriano completo de corpo residual $k$, contanto que se introduza aí $X$
como um esquema formal sobre $A$:

\result{TEOREMA}{9} Sejam $k$ um corpo e $X_0$ um esquema de tipo finito liso
sobre $k$. Para todo anel local noetheriano completo $A$ de corpo residual
$k$, seja $F(A)$ o conjunto das classes (salvo um isomorfismo que induza a
identidade sobre $X_0$) de esquemas formais $\mathfrak X$ sobre $A$, de tipo
finito e planos sobre $A$, tais que $\mathfrak X\otimes_A k=X_0$. Com estas
notações:
\begin{enumerate}
  \item[\textit{i.}] Se $H^1(X_0,\Theta_{X_0/k})=0$, então, para todo $A$,
    $F(A)$ tem no máximo um elemento.
  \item[\textit{ii.}] Se $H^2(X_0,\Theta_{X_0/k})=0$, então, para todo $A$,
    $F(A)$ tem ao menos um elemento.
\end{enumerate}

\src{206}
\result{COROLÁRIO}{1} Suponhamos que $X_0$ é próprio sobre $k$. Sob a
condição \textit{(i)}, para todo $A$ existe no máximo (salvo um isomorfismo que
induza a identidade sobre $X_0$) um esquema $X$ próprio e plano sobre $A$ tal
que $X\otimes_A k=X_0$.

Utiliza-se o teorema 3, corolário 2. Por exemplo:

\result{COROLÁRIO}{2} Se $X$ é um esquema próprio e plano sobre $A$ tal que
$X\otimes_A k$ é isomorfo ao esquema projetivo tipo $\mathbf P_k^r$ de
dimensão $r$ sobre $k$, então $X$ é isomorfo a $\mathbf P_A^r$.

(Também se pode deduzir esse resultado da proposição 3, corolário 1).

\result{COROLÁRIO}{3} Seja $X_0$ um esquema liso e projetivo sobre $k$, e
suponhamos que se tenha
\[
  H^2(X_0,\Osh_{X_0})=H^2(X_0,\Theta_{X_0/k})=0.
\]
Então, para todo $A$, existe um esquema projetivo e plano $X$ sobre $A$ tal
que $X\otimes_A k=X_0$.

Combinam-se o teorema 9\textit{(ii)} e o teorema 4. Em particular:

\result{COROLÁRIO}{4} Seja $X_0$ o esquema de uma curva algébrica completa e
lisa sobre $k$. Então, para todo anel local noetheriano completo $A$ de
corpo residual $k$, existe um «esquema de curva lisa» $X$ sobre $A$ tal que
$X\otimes_A k=X_0$.

\result{OBSERVAÇÕES}{}

1) Os corolários 3 e 4 são especialmente interessantes se $k$ é de
característica $p\ne0$, tomando como $A$ um anel de valoração discreta de
característica $0$ e corpo residual $k$; por exemplo, o «menor $A$ possível»,
a saber, aquele cujo ideal maximal é gerado por $p$. (De fato, pelos
teoremas de Cohen, basta conhecer os corolários 3 e 4 para um anel $A$ desse
tipo). Assinalemos a este respeito que, segundo os especialistas, não se sabe se
existem esquemas sobre um corpo $k$ que não sejam a redução módulo $p$ de um
esquema plano definido sobre um tal anel $A$. Ao menos, os resultados deste
número proporcionam um meio de investigar sistematicamente essa questão. Dever-se-ia
começar por examinar se a primeira obstrução que aparece em
$H^2(X_0,\Theta_{X_0/k})$ é necessariamente nula.

\emph{Complemento de 1962.} --- J.-P. Serre construiu uma variedade
projetiva não singular, de dimensão 3, sobre um corpo algebricamente fechado
$k$ de característica $p>0$, que não se obtém por redução de um esquema
próprio sobre um anel local íntegro de corpo residual $k$ e corpo de
frações de característica zero. Mumford teria encontrado um exemplo análogo
com uma superfície projetiva não singular.

2) Observe-se que o teorema 3, e a técnica correspondente, só são válidos para
um anel de base completo (local, para fixar as ideias).

Para passar de resultados conhecidos para o completamento de um anel local aos
resultados correspondentes para esse mesmo anel local, seria necessário um
quarto «teorema fundamental», cujo enunciado definitivo ainda está por encontrar.

\src{207}
3) Comparem-se os resultados deste número (em particular, os corolários 1 e
2 precedentes) e do seguinte com os resultados de Kodaira--Spencer sobre a
variação das estruturas complexas. Utilizando o teorema conjectural ao qual
se acaba de aludir, sob as condições do corolário 1, mas sem supor já
que $A$ seja completo, se deveria poder concluir que existe um anel $A'$ que
contém $A$, finito e não ramificado sobre $A$, tal que $X\otimes_A A'$ e
$X'\otimes_A A'$ são $A'$-isomorfos (sendo $X,X'$ dois $A$-esquemas próprios e
planos dados, tais que $X\otimes_A k=X'\otimes_A k=X_0$). Isto é ao menos o
que pode demonstrar-se quando $X_0=\mathbf P_k^r$, utilizando o corolário da
proposição 2. Em todo caso, quando $H^1(X_0,\Theta_{X_0/k})=0$, pode
demonstrar-se que as fibras de $X,X'$ sobre todo ponto $y$ de
$Y=\operatorname{Spec}(A)$ são isomorfas, ao menos depois de passar ao
fecho algébrico do corpo residual $\kappa(y)$. (Tem-se um resultado
local, aparentemente mais forte, sem supor já necessariamente que $A$ seja
local). Quanto às «variações de estrutura» do espaço projetivo,
assinalemos além disso a questão seguinte, sugerida por um problema correspondente
de Kodaira--Spencer. Seja $X$ um esquema próprio e plano sobre o anel local
íntegro $A$ de corpo de frações $K$ e corpo residual $k$, e suponhamos que
$X\otimes_A K$ é isomorfo a $\mathbf P_K^r$; é verdade que
$X\otimes_A k=X_0$ é isomorfo a $\mathbf P_k^r$ (ao menos sobre o fecho
algébrico de $k$)? Nesse problema pode-se supor que $A$ é um anel de
valoração discreta completo. Coloca-se um problema análogo quando $X_0$ é
uma variedade abeliana.

\emph{Complemento de 1962.} --- O conferencista é agora menos otimista
com respeito aos resultados conjecturados acima. Não obstante, Hironaka resolveu
afirmativamente o problema relativo às variações de estrutura
do espaço projetivo, e Koizumi o problema análogo para as variedades
abelianas.

\fsection{7}{Aplicação à «Teoria dos módulos»}

Como o próprio conferencista apenas começou a abordar essa teoria, seremos
obrigados a nos limitar a algumas indicações. Para simplificar, situamo-nos
sobre um corpo $k$, isto é, trabalhamos em igual característica, embora o
teorema 8 permita tratar também o caso geral sem mudança essencial, ao que
parece. Atualmente não superamos a etapa «formal»; no entanto, o
conferencista espera poder construir a partir dela verdadeiros esquemas de
módulos em certos casos e, em particular, construir para cada inteiro $g$ um
esquema sobre os inteiros que desempenhe o papel de esquema de módulos universal
para as curvas lisas de gênero $g$.

\emph{Complemento de 1962.} --- Mumford construiu os esquemas de módulos
para as curvas de gênero $g$ (cf. SHMI). O teorema 10 demonstra além disso que
os «esquemas de Jacobi de nível $n$» da teoria de módulos são não singulares,
e inclusive lisos sobre $\mathbf Z$.

\src{208}
Retomemos a situação e as notações do teorema 9, supondo agora que $A$
é uma álgebra local de posto finito sobre $k$, que para simplificar se supõe
algebricamente fechado. Então $F(A)$ pode ser considerado como um funtor
covariante de $A$, com valores na categoria de conjuntos: um homomorfismo de
$k$-álgebras $A\longrightarrow B$ define uma aplicação
$F(A)\longrightarrow F(B)$, pois todo $A$-esquema plano $X$ tal que
$X\otimes_A k=X_0$ dá lugar a um $B$-esquema $X\otimes_A B$ com as mesmas
propriedades. Suponhamos que se possa encontrar uma $k$-álgebra local noetheriana
completa $\mathcal O$ e um isomorfismo funtorial
\begin{equation}
  \tag{*}
  \operatorname{Hom}(\mathcal O,A)\xrightarrow{\sim}F(A)
\end{equation}
(onde o primeiro membro denota os homomorfismos de $k$-álgebras). Vê-se
facilmente que um tal $\mathcal O$ está determinado salvo isomorfismo canônico;
chama-se então ao espectro formal $\mathfrak Y$ de $\mathcal O$ (isto é,
ao espaço topológico reduzido a um ponto e provido de um feixe de anéis
topológicos reduzido a $\mathcal O$) o \emph{esquema formal de módulos} para
$X_0$. (Observe-se que ele não existe necessariamente). Seja $\mathfrak r$ o ideal
maximal de $\mathcal O$; para cada $n$, seja
$\mathcal O_n=\mathcal O/\mathfrak r^{n+1}$ (portanto $\mathcal O_0=k$).
Então o homomorfismo canônico $\mathcal O\longrightarrow\mathcal O_n$ é um
elemento de $\operatorname{Hom}(\mathcal O,\mathcal O_n)$, pelo que define um
elemento de $F(\mathcal O_n)$, isto é, um $\mathcal O_n$-esquema plano $X_n$
cuja redução módulo $\mathfrak r$ é $X_0$. Esses $X_n$ se deduzem uns dos
outros por extensão de escalares (isto é, aqui por reduções), de onde
resulta que provêm de um esquema formal $\mathfrak X$ bem determinado sobre
o esquema formal de módulos $\mathfrak Y$; $\mathfrak X$ é plano sobre
$\mathfrak Y$ e $\mathfrak X_0=X_0$. Como se vê imediatamente, o isomorfismo (*)
obtém-se então associando a todo homomorfismo de $k$-álgebras
$\mathcal O\longrightarrow A$ a classe do $A$-esquema
$\mathfrak X\otimes_{\mathcal O}A$ (isto é, a todo morfismo de $k$-esquemas
$Y'=\operatorname{Spec}(A)\longrightarrow\mathfrak Y$ associa-se o
$Y'$-esquema $\mathfrak X\otimes_{\mathfrak Y}Y'$ obtido por mudança de base).
Além disso, vê-se que o isomorfismo (*) e sua descrição precedente continuam
válidos se somente se supõe que $A$ é uma $k$-álgebra local noetheriana e
completa (não necessariamente artiniana). Naturalmente, como de costume,
$\mathcal O$ pode perfeitamente ter \textit{a priori} elementos nilpotentes,
e parece provável que existam casos em que $\mathcal O$ seja ele mesmo artiniano
sem ser idêntico a $k$. Isto mostra até que ponto o ponto de vista de
Kodaira--Spencer (que se limita a tomar como $A$ anéis regulares) é
inadequado \textit{a priori} no caso geral.

\src{209}
Fica por dar condições suficientes para que exista um esquema formal de
módulos para $X_0$, que se supõe próprio sobre $k$. Em geral, é fácil dar
condições necessárias e suficientes simples sobre um funtor
$A\longmapsto F(A)$ (de $k$-álgebras locais de posto finito, com valores nos
conjuntos) para que possa ser escrito na forma
$\operatorname{Hom}(\mathcal O,A)$ para um $\mathcal O$ adequado. Não as
detalharemos aqui. Assinalemos apenas que, no caso que nos ocupa, essas
condições impõem sobre $X_0$ condições cohomológicas não triviais, e parece
pouco provável que sejam sempre satisfeitas, embora o conferencista não tenha
construído nenhum contraexemplo. Em contrapartida, parece plausível que a condição
$H^0(X_0,\Theta_{X_0/k})=0$ seja uma condição suficiente (embora de modo algum
necessária) para a existência de um esquema formal de módulos. Limitamo-nos a
enunciar aqui um teorema num caso particularmente simples (cujo análogo na
teoria de espaços analíticos é bem conhecido, cf. Kodaira--Spencer), que se
estabelece sem dificuldade mediante os resultados do número anterior:

\result{TEOREMA}{10} Seja $X_0$ um esquema próprio e liso sobre o corpo $k$
tal que
\[
  H^0(X_0,\Theta_{X_0/k})=H^2(X_0,\Theta_{X_0/k})=0.
\]
Então existe um esquema formal de módulos para $X_0$, correspondente a um
anel local regular $\mathcal O$ (isto é, uma álgebra de séries formais sobre
$k$).

Como já assinalamos, não é certo em geral que o esquema formal
$\mathfrak X$ sobre $\mathcal O$ seja algebrizável; mas sabe-se que o é
quando $X_0$ é projetivo e $H^2(X_0,\Osh_{X_0})=0$ (teorema 4), por exemplo
quando $X_0$ tem dimensão 1. Isto dá certo fundamento à esperança de
construir um esquema de módulos sobre os inteiros para as curvas de gênero
dado\ldots{}

Observe-se também que métodos como os expostos neste número podem ser aplicados
à construção e ao estudo das variedades de Picard, assim como a muitas
outras construções. Voltaremos a isso em breve.

\fsection{8}{Aplicação ao grupo fundamental}

As técnicas expostas permitem abordar o estudo sistemático do grupo
fundamental seguindo o modelo da teoria topológica. Os dois primeiros
teoremas enunciados neste número são generalizações de resultados de um
trabalho recente de Lang--Serre.

Seja $X$ um esquema. Um $X$-esquema $X'$ se chama recobrimento não ramificado de
$X$ se
\begin{enumerate}
  \item[\textit{i.}] $X'$ é finito sobre $X$, isto é, está definido por um
    feixe coerente de álgebras $\mathcal A=\mathcal A(X')$ sobre $X$;
  \item[\textit{ii.}] $\mathcal A$ é um feixe localmente livre sobre $X$;
  \item[\textit{iii.}] Para todo $x\in X$, o quociente
    \[
      \mathcal A_x/\mathfrak m_x\mathcal A_x
      =\mathcal A_x\otimes_{\Osh_x}\kappa(x)
    \]
    é um álgebra separável sobre $\kappa(x)$.
\end{enumerate}

\src{210}
Essa noção de recobrimento não ramificado (devida a Serre e ao conferencista)
possui todas as propriedades elementares que se pode razoavelmente esperar, cuja
lista não daremos. Limitemo-nos a dizer que dá lugar a uma teoria de Galois
modelada sobre a teoria de Galois clássica (e que a contém; as
demonstrações são bem mais simples que as comumente aceitas para
esta última) e sobre a teoria de Galois dos recobrimentos topológicos. Mais
precisamente, chamemos \emph{ponto geométrico} de um esquema $X$ um morfismo
$a$ do espectro $\xi$ de um corpo algebricamente fechado $\Omega$ em $X$,
isto é, o dado de uma extensão algebricamente fechada do corpo residual
$\kappa(x)$ de um ponto $x=|a|$ de $X$ (chamado \emph{localização} do ponto
geométrico $a$). Se $X'$ é então um recobrimento não ramificado de $X$, se
lhe pode associar o conjunto $E_a(X')$ dos «pontos geométricos de $X'$ sobre
$a$», isto é, o conjunto dos pares formados por um $x'\in X'$ sobre $x$
e um $\kappa(x)$-homomorfismo $\kappa(x')\longrightarrow\Omega$. Obtém-se assim
(para $(X,a)$ fixos) um funtor $F(X,a)$ da categoria $R(X)$ dos
recobrimentos não ramificados $X'$ de $X$ na categoria dos conjuntos
finitos. Se $X$ é conexo, o par formado por $R(X)$ e $F(X,a)$ possui as
propriedades formais necessárias para ser isomorfo ao par análogo definido por
um grupo topológico compacto totalmente descontínuo $\pi$ adequado (isto é,
um limite projetivo de grupos finitos): toma-se a categoria $C(\pi)$ dos
conjuntos finitos $E$ sobre os quais $\pi$ atua continuamente, e o funtor
identidade $F(\pi)(E)=E$ desta categoria na categoria dos conjuntos
finitos. Além disso, o grupo $\pi$ está determinado salvo isomorfismo canônico pela
condição de que $(C(\pi),F(\pi))$ seja isomorfo a um par dado. Nesse caso,
$\pi$ chama-se \emph{grupo fundamental} do esquema conexo $X$ no ponto
geométrico $a$, e denota-se por $\pi_1(X,a)$. Se $X$ não é conexo, substitui-se
o esquema pela componente conexa de $x=|a|$. Se, pelo contrário, $X$ é
conexo, os grupos $\pi_1(X,a)$ correspondentes a dois pontos geométricos
$a',a''$ de $X$ são isomorfos (o isomorfismo está determinado salvo
automorfismos interiores), de modo que, como de costume, pode-se escolher $a$
da maneira mais conveniente, por exemplo no ponto genérico de $X$, suposto
irredutível. Naturalmente, $\pi_1(X,a)$ é um funtor covariante do esquema
pontuado $(X,a)$. Todo enunciado relativo à classificação dos
recobrimentos inseparáveis pode ser traduzido então para a linguagem da teoria de
grupos seguindo o dicionário bem conhecido (salvo que se deve levar em conta
que aqui se trata de grupos topológicos).

\src{211}
Nosso objetivo é desenvolver o análogo da sequência exata de homotopia
dos espaços fibrados para um morfismo próprio $f:X\longrightarrow Y$.
Evidentemente, como não sabemos quais são os grupos de homotopia superiores, só
obteremos resultados necessariamente incompletos. Para poder aplicar os
teoremas fundamentais do n\textsuperscript{o} 3, devemos explicitar primeiro
alguns lemas elementares relativos aos esquemas sobre corpos ou anéis de
Artin (conforme o procedimento geral!).

\result{LEMA}{1} Seja $(X',a')$ um recobrimento não ramificado pontuado associado
a uma representação pontuada de $\pi_1(X,a)$ em um conjunto finito $E$
(provido do ponto marcado $e$). Então o morfismo canônico
$\pi_1(X',a')\longrightarrow\pi_1(X,a)$ identifica o primeiro grupo com o
estabilizador de $e$ em $\pi_1(X,a)$ (e é, portanto, injetivo).

\result{LEMA}{2} Seja $X$ um esquema algébrico sobre o corpo $k$, e seja $k'$
uma extensão radicial de $k$. Então todo recobrimento não ramificado de
$X\otimes_k k'$ provém por imagem inversa (isto é, por extensão de
escalares) de um recobrimento não ramificado de $X$, determinado salvo
isomorfismo canônico.

Desses dois lemas resulta, em particular, que para toda extensão algébrica
$K$ de $k$, e para todo ponto geométrico $a'$ de $X'=X\otimes_k K$ que se
projeta sobre o ponto geométrico $a$ de $X$, o homomorfismo funtorial
$\pi_1(X',a')\longrightarrow\pi_1(X,a)$ é injetivo.

\result{LEMA}{3} Seja $X$ um esquema completo sobre um anel local artiniano
$A$, tal que $H^0(X,\Osh_X)=A$; seja $X'$ um recobrimento não ramificado de $X$,
e seja $A'=H^0(X',\Osh_{X'})$, que é, portanto, um anel finito sobre $A$
(que \textit{a priori} pode estar ramificado sobre $A$). Sejam $X_0,X'_0$ os
subesquemas reduzidos associados a $X,X'$ (obtidos dividindo pelos feixes de
elementos nilpotentes de $\Osh_X,\Osh_{X'}$), e seja $k$ um subcorpo de
$A/\mathfrak r(A)$ sobre o qual $A/\mathfrak r(A)$ seja finito (assim, $X_0$ é
um esquema algébrico completo sobre $k$, e $X'_0$ é um recobrimento não
ramificado seu). Seja, por fim, $\Omega$ uma extensão algebricamente
fechada de $k$, e consideremos o recobrimento não ramificado
$X'_0\otimes_k\Omega$ de $X_0\otimes_k\Omega$. Então as duas condições
seguintes são equivalentes:
\begin{enumerate}
  \item[\textit{i.}] $X'_0\otimes_k\Omega$ está completamente decomposto
    sobre $X_0\otimes_k\Omega$.
  \item[\textit{ii.}] O morfismo natural
    $X'\longrightarrow X\otimes_A A'$ é um isomorfismo.
\end{enumerate}
Sob essas condições, $A'$ é uma extensão não ramificada de $A$. Por fim,
se $X'$ é conexo, a condição \textit{(i)} equivale à condição seguinte,
de aparência mais fraca:
\begin{enumerate}
  \item[\textit{i bis.}] $X'_0\otimes_k\Omega$ admite uma seção regular
    sobre $X_0\otimes_k\Omega$.
\end{enumerate}
Quando se satisfaz a condição \textit{(ii)}, diremos que o recobrimento não
ramificado $X'$ de $X$ é \emph{geometricamente trivial}.

\src{212}
\result{LEMA}{4} Seja $f:X\longrightarrow Y$ um morfismo próprio tal que
$f_*(\Osh_X)=\Osh_Y$. Sejam $a$ um ponto geométrico de $X$ e $b$ sua projeção
sobre $Y$. Então $\pi_1(X,a)\longrightarrow\pi_1(Y,b)$ é sobrejetivo.

Com efeito, é preciso demonstrar o seguinte: se um recobrimento não ramificado
$Y'$ de $Y$ (correspondente a um feixe de álgebras localmente livre
$\mathcal A$) é tal que $X\otimes_Y Y'$ é desconexo, então $Y'$ também o
é. Com efeito, $f^*(\mathcal A)$ será soma direta de dois feixes de anéis não
nulos e, portanto, também o será sua imagem direta; ora, esta última
não é outra que $\mathcal A\otimes f_*(\Osh_X)=\mathcal A$.

\result{LEMA}{5} Seja $X$ um esquema completo sobre um corpo $k$; supõe-se que
$H^0(X,\Osh_X)$ é um anel local $A$ e que $A/\mathfrak r(A)$ é radicial
sobre $k$. Seja $\Omega$ um fecho algébrico de $k$, e seja
$\overline X=X\otimes_k\Omega$ (é conexo). Escolhamos um ponto geométrico
$\overline a$ de $\overline X$ que se projete sobre o ponto geométrico $a$ de
$X$; então tem-se uma sequência exata
\[
  e\longrightarrow\pi_1(\overline X,\overline a)
  \longrightarrow\pi_1(X,a)
  \longrightarrow\pi_1(k,b)\longrightarrow e
\]
(onde $\pi_1(k,b)$ é o grupo de Galois de $\Omega$ sobre $k$).

A injetividade do primeiro homomorfismo foi vista com os lemas 1 e 2, a
exatidão no meio resulta do lema 3 e, por fim, a sobrejetividade do
último homomorfismo (que é a única que utiliza o fato de que
$A/\mathfrak r(A)$ seja radicial) resulta do lema 4.

\result{PROPOSIÇÃO}{4} Seja $f:X\longrightarrow Y$ um morfismo próprio e plano
tal que, para todo $y\in Y$, $H^0(f^{-1}(y),\Osh_{f^{-1}(y)})$ seja uma álgebra
separável sobre o corpo residual $\kappa(y)$ (o que ocorre, por exemplo, se
$f^{-1}(y)$ é um esquema separável sobre $\kappa(y)$, isto é, reduzido e tal
que os corpos correspondentes a suas componentes irredutíveis sejam
extensões separáveis de $\kappa(y)$). Então o recobrimento $Y'$ de $Y$
associado a $f_*(\Osh_X)$ é não ramificado.

A demonstração é fácil graças ao teorema 2.

\src{213}
Essa proposição, junto com o lema 1, reduz praticamente o estudo
homotópico dos morfismos próprios e planos (com fibras separáveis) ao caso em
que $f_*(\Osh_X)=\Osh_Y$ (pois, utilizando a fatoração de Stein, se
substitui $Y$ por $Y'$).

\result{OBSERVAÇÃO}{} Um morfismo plano de tipo finito cujas fibras são
esquemas separáveis (resp. lisos) chama-se \emph{separável} (resp.
\emph{liso}). Demonstra-se que, se $f$ é plano e $f^{-1}(y)$ é separável
(resp. liso), existe uma vizinhança de $f^{-1}(y)$ na qual $f$ é separável
(resp. liso). O mesmo resultado é válido para «absolutamente normal»
(«teorema de Bertini»).

Seja $f:X\longrightarrow Y$ um morfismo próprio tal que
\begin{equation}
  \tag{i}
  f_*(\Osh_X)=\Osh_Y,
\end{equation}
e seja $X'$ um esquema finito sobre $X$. Seja $Y'$ o recobrimento de $Y$
correspondente à fatoração de Stein de $X'\longrightarrow Y$
(cf. teorema 5). Seja $y\in Y$; então o conjunto das componentes conexas
da fibra $F'$ de $X'$ sobre $y$ identifica-se com o conjunto dos pontos
$y'$ de $Y'$ sobre $y$ (teorema 5). Consideremos o morfismo evidente
\begin{equation}
  \tag{*}
  X'\longrightarrow X\times_Y Y'
\end{equation}
deduzido dos morfismos naturais $X'\longrightarrow X$ e
$X\longrightarrow Y'$; será um isomorfismo sempre que $X'$ seja da forma
$X\times_Y Y''$, onde $Y''$ é um recobrimento não ramificado de $Y$, e
então $Y'$ não será outro senão $Y''$ e (*) será a identidade. Queremos dar
precisamente condições sob as quais $X'$ seja da forma que acabamos de
indicar, isto é, que $Y'$ seja não ramificado e (*) um isomorfismo. Para isso,
introduzamos a fibra $F$ de $X$ em $y$; é um esquema próprio sobre
$\kappa(y)$ do qual $F'$ é um recobrimento (não ramificado se $X'$ o for).
Seja $F'_1$ uma componente conexa de $F'$ correspondente a um ponto $y'_1$ de
$Y'$ sobre $y$. Suponhamos além disso
\begin{enumerate}
  \item[\textit{(ii)}] $X'$ é não ramificado sobre $X$ nos pontos de $F'_1$
    (portanto, $F'_1$ é um recobrimento não ramificado de $F$), e
  \item[\textit{(iii)}] $F'_1$ é um recobrimento geometricamente trivial de
    $F$ (cf. lema 3).
\end{enumerate}

\result{TEOREMA}{11} Nessas condições, existe uma vizinhança aberta $U'$ de
$y'_1$ em $Y'$ tal que (*) é um isomorfismo sobre $U'$. Além disso, $Y'$ é não
ramificado em $y'_1$ sobre $Y$ (mas pode ser ramificado nos demais pontos
$y'$ de $Y'$ sobre $y$).

Naturalmente, as condições \textit{(ii)} e \textit{(iii)} são também
necessárias para que seja válida a conclusão do teorema. A demonstração do
teorema é fácil, com auxílio do lema 3 e do teorema 2.

\src{214}
\result{COROLÁRIO}{1} Suponhamos sempre satisfeita \textit{(i)}. Para que um
recobrimento não ramificado sobre $X$ seja isomorfo à imagem inversa de um
recobrimento não ramificado $Y'$ de $Y$, é necessário e suficiente que $X'$
induza sobre cada fibra $f^{-1}(y)$ um recobrimento geometricamente trivial.

Como, segundo o teorema 11, o conjunto dos pontos de $Y$ para os quais se
cumpre essa condição é aberto, basta verificá-la nos pontos $y$ que são
fechados\ldots{} Assinalemos o enunciado seguinte, equivalente ao corolário 1:
o núcleo do homomorfismo (sobrejetivo segundo o lema 4)
$\pi_1(X)\longrightarrow\pi_1(Y)$ é o subgrupo invariante fechado gerado
pelas imagens em $\pi_1(X)$ dos $\pi_1(\overline{f^{-1}(y)})$, onde
\[
  \overline{f^{-1}(y)}
  =f^{-1}(y)\otimes_{\kappa(y)}\overline{\kappa(y)},
\]
$\overline{\kappa(y)}$ designa um fecho algébrico de $\kappa(y)$. Observar-se-á
que, como não se pode escolher o mesmo ponto base para todas as fibras,
os homomorfismos $\pi_1(\overline{f^{-1}(y)})\longrightarrow\pi_1(X)$ só
ficam determinados em qualquer caso (uma vez escolhido um ponto base para $X$
e, portanto, para $Y$) módulo composição por um automorfismo interior em
$\pi_1(X)$.

\result{COROLÁRIO}{2} Nas condições gerais do teorema 11, suponhamos
além disso que $Y,X,X'$ são íntegros, e sejam $K,L,L'$ seus corpos. Então existe
uma subextensão separável $K'$ de $K$ em $L'$, linearmente disjunta de $L$, tal
que $L'=LK'$ (de onde $L'=L\otimes_K K'$).

(Aplica-se a última parte do lema 3 à fibra geométrica de $X$). O caso de
aplicação mais interessante do teorema 11 obtém-se quando $f$ é um morfismo
separável. Então $X'$ é também separável sobre $Y$; portanto, em virtude
da proposição 4, $Y'$ é não ramificado sobre $Y$, e o segundo membro
$X\times_Y Y'$ de (*) é não ramificado sobre $X$. Conclui-se facilmente:

\result{COROLÁRIO}{3} Suponhamos, além de \textit{(i)}, que $f$ é separável.
Seja $X'$ um recobrimento não ramificado e conexo de $X$; para que
$X$ seja imagem inversa de um
recobrimento não ramificado $Y'$ de $Y$, é necessário e suficiente que o
recobrimento $\overline F'$ induzido sobre uma fibra geométrica
$\overline F=f^{-1}(y)$ admita uma seção regular.

Observar-se-á que não era necessário supor que $\overline F'$ fosse
geometricamente trivial sobre $\overline F$ (o que será verdadeiro \textit{a
posteriori}, embora \textit{a priori} esta condição seja muito mais forte). O
corolário 3 é, de fato, equivalente ao enunciado seguinte:

\src{215}
\result{COROLÁRIO}{4} Seja $f:X\longrightarrow Y$ um morfismo próprio e separável
tal que $f_*(\Osh_X)=\Osh_Y$, seja $\overline F$ a fibra geométrica de um ponto
$y\in Y$, e escolhamos um ponto geométrico em $\overline F$, do qual se obtêm,
mediante os morfismos $\overline F\longrightarrow X\longrightarrow Y$, pontos
geométricos em $X,Y$, que serão tomados como pontos base para os grupos
fundamentais de $\overline F,X,Y$. Nessas condições tem-se a sequência
exata
\[
  \boxed{
    \pi_1(\overline F)\longrightarrow\pi_1(X)
    \longrightarrow\pi_1(Y)\longrightarrow e
  }.
\]

Daí se deduzem facilmente os dois enunciados seguintes de Serre--Lang,
liberados aqui de toda hipótese de normalidade:

\result{COROLÁRIO}{5} Sejam $X,Y$ esquemas conexos sobre um corpo $k$;
suponhamos que o esquema reduzido $X_{\mathrm{réd}}$ é separável sobre $k$
(o que é automaticamente verdadeiro se $k$ é perfeito).
Escolhamos um ponto geométrico $a$ (resp. $b$) em $X$ (resp. $Y$), de onde se
obtêm um ponto geométrico $c=(a,b)$ em $X\times_kY$ e um morfismo natural
\[
  \pi_1(X\times_kY,c)\longrightarrow
  \pi_1(X,a)\times\pi_1(Y,b)
\]
(deduzido dos morfismos funtoriais de $\pi_1(X\times Y,c)$ em $\pi_1(X,a)$
e $\pi_1(Y,b)$). Este morfismo é injetivo, e inclusive bijetivo se $k$ é
algebricamente fechado.

(A sobrejetividade neste último caso é quase trivial). Daí se deduz,
com Serre--Lang:

\result{COROLÁRIO}{6} Seja $X$ um esquema algébrico conexo sobre um corpo
algebricamente fechado $k$, e seja $K$ uma extensão algebricamente fechada
de $k$; então os grupos fundamentais de $X$ e $X\otimes_kK$ são os
mesmos, isto é, todo recobrimento não ramificado deste último esquema
provém por extensão de escalares de um recobrimento não ramificado (único
salvo isomorfismo) de $X$.

\result{OBSERVAÇÕES}{}

1) Utilizando a proposição 4, vê-se que a hipótese
$f_*(\Osh_X)=\Osh_Y$ do corolário 4 não é essencial. No caso geral, em vez
de colocar o grupo unidade $e$ depois de $\pi_1(Y)$, é preciso continuar com
\[
  \pi_0(\overline F)\longrightarrow\pi_0(X)
  \longrightarrow\pi_0(Y)\longrightarrow e.
\]
como em topologia algébrica.

\src{216}
2) Em geral, por ora não se pode dizer nada sobre o núcleo de
$\pi_1(\overline F)\longrightarrow\pi_1(X)$, no qual deveria intervir um
$\pi_2(Y)$. Parece, no entanto, que se deveria poder demonstrar que
$\pi_1(\overline F)\longrightarrow\pi_1(X)$ é injetivo se $Y$ é o espectro
de um anel local $A$, apoiando-se no teorema 12 mais adiante (que afirma que
isso ocorre se $A$ é completo).

O teorema 11 utilizava somente os teoremas 1 e 2. Vamos aplicar agora o
teorema 3, apoiando-nos no lema elementar seguinte:

\result{LEMA}{6} Seja $X$ um esquema, e $X_0$ o esquema reduzido correspondente
(isto é, no qual se anularam os elementos nilpotentes). Então todo
recobrimento não ramificado $X'_0$ de $X_0$ é induzido por um recobrimento
não ramificado $X'$ de $X$, determinado salvo isomorfismo canônico.

Este lema, de natureza puramente local, desempenha aqui um papel análogo ao do
teorema 8 em teoria de módulos. Combinando-o com o teorema de existência
(teorema 3), deduz-se:

\result{TEOREMA}{12} Seja $A$ um anel local noetheriano e completo de corpo
residual $k$. Seja $X$ um esquema próprio sobre $A$. Então todo recobrimento
não ramificado $X'_0$ de $X_0=X\otimes_A k$ é induzido por um recobrimento não
ramificado $X'$ de $X$, único salvo isomorfismo.

Em outros termos:

\result{COROLÁRIO}{1} Escolhamos um ponto geométrico em $X_0$ como ponto base para
os grupos fundamentais de $X_0$ e $X$. Então o homomorfismo canônico
$\pi_1(X_0)\longrightarrow\pi_1(X)$ é um isomorfismo.

Apliquemos agora a $X_0$ o lema 5 (supondo, para simplificar, que
$H^0(X_0,\Osh_{X_0})=k$), e observemos que, por ser $A$ completo, suas
extensões não ramificadas correspondem às extensões não ramificadas de seu
corpo residual, isto é, tem-se $\pi_1(Y)=\pi_1(k)$ (onde
$Y=\operatorname{Spec}(A)$). Obtém-se a sequência exata:
\[
  e\longrightarrow\pi_1(\overline X_0)
  \longrightarrow\pi_1(X)\longrightarrow\pi_1(Y)\longrightarrow e.
\]

\result{COROLÁRIO}{2} Seja $f:X\longrightarrow Y$ um morfismo próprio e plano,
sejam $y_1$ um ponto de $Y$ e $y_0$ uma especialização de $y_1$, e consideremos
as fibras «geométricas» correspondentes $\overline X_1$ e $\overline X_0$;
supõe-se que $\overline X_0$ é separável e conexa (o que implica que
$\overline X_1$ satisfaz as mesmas condições). Então pode-se encontrar um
homomorfismo de grupos
$\pi_1(\overline X_1)\longrightarrow\pi_1(\overline X_0)$, definido salvo um
automorfismo interior, e este homomorfismo é sobrejetivo.

\src{217}
Poder-se-ia esperar que este homomorfismo fosse sempre bijetivo; infelizmente,
em geral não é assim se $\kappa(y_0)$ tem característica $>0$. Obteremos,
no entanto, mais adiante uma cota para o núcleo deste homomorfismo (ao menos se
$\overline X_0$ é liso), da qual se segue que, se $\kappa(y_0)$ tem
característica $0$, o homomorfismo anterior é bijetivo (resultado que também
poderia ser demonstrado por via transcendente). Em todo caso, obtemos já uma cota
para o $\pi_1$ de uma fibra especial por meio daquele de uma fibra genérica.
Utilizando, por exemplo, o fato de que uma curva algébrica em característica
$p$ se pode levantar a uma curva em característica $0$ (teorema 9, corolário
4), obtém-se por via transcendente:

\result{COROLÁRIO}{3} Seja $X_0$ o esquema de uma curva completa lisa sobre
um corpo algebricamente fechado de característica arbitrária, e seja $g$ o
gênero de $X_0$; então $\pi_1(X_0)$ admite $2g$ geradores topológicos,
ligados pela relação bem conhecida.

Daí se deduz, mediante uma técnica bem conhecida de seções
hiperplanas:

\result{COROLÁRIO}{4} Seja $X$ um esquema projetivo liso sobre um corpo
algebricamente fechado de característica arbitrária; então $\pi_1(X)$
admite um número finito de geradores topológicos.

Determinemos o núcleo do homomorfismo
$\pi_1(\overline X_1)\longrightarrow\pi_1(\overline X_0)$. Para isso, pode-se
supor que $Y$ é o espectro de um anel de valoração discreta completo
$V=\Osh_y$ ($y=y_0$). A questão equivale à seguinte: dado um recobrimento
não ramificado $\overline X'_1$ de $\overline X_1$ (que, se se quiser, pode
ser suposto galoisiano), determinar em que condições provém de um
recobrimento não ramificado de $X_0$. \textit{A priori}, o recobrimento dado
provém, por extensão de escalares, de um recobrimento não ramificado $X'_1$
de $X_1\otimes_K K'$, onde $K'$ é uma subextensão finita do fecho
algébrico $\overline K$ do corpo de frações $K$ de $V$; se
$\overline X'_1$ fosse galoisiano de grupo $G$, pode-se escolher $X'_1$
galoisiano de grupo $G$. Tem-se então: para que
$X'_1\otimes_{K'}\overline K=\overline X'_1$ provenha de um recobrimento não
ramificado de $X_0$, é necessário e suficiente que se possa encontrar uma extensão
finita $K''$ de $K'$ em $\overline K$ tal que
$X''_1=X'_1\otimes_{K'}K''$ seja da forma $X''\otimes_{V''}K''$, onde $V''$
é o fecho normal de $V$ em $K''$ e $X''$ é um recobrimento não ramificado
de $X\otimes_VV''$. Suponhamos, por exemplo, que $X_0$ é absolutamente normal;
daí resulta que $X\otimes_VV''$ é normal (pois é plano sobre $V''$ e tem
fibra especial normal), e seu corpo de funções é idêntico ao $K''(X_1)$ de
\[
  X_1\otimes_KK''
  =(X\otimes_VK)\otimes_KK''
  =X\otimes_VK''
  =(X\otimes_VV'')\otimes_{V''}K''.
\]
Seja $L''=K''(X'_1)$ o corpo de funções de $X'_1\otimes_{K'}K''$; é uma
extensão finita separável de $K''(X_1)$, e a condição anterior significa
também que $L''$ é uma extensão não ramificada do corpo de funções de
$X\otimes_VV''$ (isto é, a normalização de $X\otimes_VV''$ em $L''$ é não
ramificada sobre $X\otimes_VV''$). Basta, de resto, verificar que $L''$ é
não ramificado nos pontos da fibra especial de $X\otimes_VV''$ (pois é não
ramificado sobre a fibra genérica $X_1\otimes_KK''$). Se agora $X_0$ é liso,
do «teorema de pureza» de Nagata--Zariski resulta inclusive que basta verificar
que $L''$ é não ramificado sobre o anel local $O''$ do ponto genérico da
fibra especial de $X\otimes_VV''$, que é um anel de valoração discreta,
normalização em $K''(X_1)$ do anel local $O\subset K(X_1)$ do ponto
genérico da fibra especial de $X$. Portanto, obteve-se:

\src{218}
\result{COROLÁRIO}{5} Nas condições e com as notações anteriores, para
que o recobrimento não ramificado
$X'_1\otimes_{K'}\overline K$ de
$\overline X_1=X_1\otimes_K\overline K$ provenha de um recobrimento não
ramificado de $\overline X_0$, é necessário e suficiente que exista uma
subextensão finita $K''$ de $\overline K/K'$ tal que $K''(X'_1)$ seja não
ramificado sobre o anel de valoração discreta $O''\subset K''(X_1)$.

Observemos agora que $O''$ é a normalização em $K''(X_1)$ do anel de
valoração discreta $O'\subset K'(X_1)$ (normalização de $O$ em $K'(X_1)$), e
que $O'$ contém a normalização $V'$ de $V$ em $K'$; um uniformizante $u$ de
$V'$ é também um uniformizante de $O'$. Suponhamos agora que $X'_1$ é
galoisiano de grupo de Galois $G$, de ordem $n$ primo com a característica $p$
de $\kappa(y_0)$ (que é também a característica do corpo residual de
$O'$). Portanto, $K'(X'_1)$ está «moderadamente ramificado» sobre $O'$, de onde
resulta facilmente («lema de Abhyankar») que, se se lhe adjunta uma raiz
$n$-ésima $v$ de um uniformizante de $O'$, torna-se não ramificado sobre a
normalização de $O'$ em $K'(X_1)(v)$. Ora, pode-se tomar como $v$ uma
raiz $n$-ésima de um uniformizante de $V'$, o que prova que se satisfaz a
condição do corolário 5. (Essa possibilidade de utilizar o lema de Abhyankar e
o teorema de pureza foi-me vendida por Serre). Para exprimir o resultado
obtido, introduzamos, para todo grupo compacto totalmente descontínuo $\pi$,
o grupo quociente $\widetilde\pi$ de $\pi$ pelo subgrupo fechado gerado
por seus $p$-subgrupos de Sylow, isto é, o limite projetivo dos grupos
quocientes discretos de $\pi$ que são de ordem prima com $p$. Com esta notação,
obtém-se:

\result{TEOREMA}{13} Seja $f:X\longrightarrow Y$ um morfismo próprio e plano,
sejam $y_1$ um ponto de $Y$ e $y_0$ uma especialização de $y_1$, e suponhamos
$\overline X_0$ conexo e liso. Então o homomorfismo
$\widetilde\pi_1(\overline X_1)\longrightarrow
\widetilde\pi_1(\overline X_0)$ deduzido do homomorfismo sobrejetivo do
teorema 12, corolário 2, é um isomorfismo.

Em outros termos:

\src{219}
\result{COROLÁRIO}{1} A classificação dos recobrimentos galoisianos não
ramificados, de grupo de Galois de ordem prima com a característica $p$ de
$\kappa(y_0)$, é a mesma para $\overline X_0$ e para $\overline X_1$.

Em particular, se $\kappa(y_0)$ tem característica zero, obtém-se por via
algébrica que $\pi_1(\overline X_1)\longrightarrow\pi_1(\overline X_0)$ é
bijetivo.

Assinalemos, por fim, que as técnicas utilizadas proporcionam também o
resultado seguinte, mais geral que o teorema 13:

\result{TEOREMA}{14} Seja $f:X\longrightarrow Y$ um morfismo próprio e liso,
seja $D$ um subesquema fechado de $X$, liso sobre $Y$, de codimensão 1 em
todos os seus pontos. Para uma fibra $Z=f^{-1}(z)$ de $f$, seja $Z'=Z-Z\cap D$ e seja
$\pi_1^t(Z')$ o quociente do grupo fundamental $\pi_1(Z')$ que classifica os
recobrimentos não ramificados de $Z'$ que estão «moderadamente ramificados»
sobre $\overline{Z\cap D}$. Sejam $y_0,y_1$ como no teorema 13. Então se
tem um homomorfismo sobrejetivo (definido salvo um automorfismo interior)
\[
  \pi_1^t(\overline X'_1)\longrightarrow\pi_1^t(\overline X'_0),
\]
e o homomorfismo correspondente
\[
  \widetilde\pi_1^t(\overline X'_1)
  \longrightarrow\widetilde\pi_1^t(\overline X'_0)
\]
é um isomorfismo.

Daí se deduzem as variantes correspondentes dos corolários do
teorema 13 e do corolário 4 do teorema 12. Do mesmo modo, utilizando o
teorema 9, corolário 3, obtém-se por via transcendente:

\result{COROLÁRIO}{} Seja $X_0$ o esquema de uma curva completa lisa sobre um
corpo algebricamente fechado de característica arbitrária, e seja
$S=(s_i)_{1\le i\le n}$ uma parte finita com $n$ elementos de $X_0$. Então
$\pi_1^t(X_0-S)$ admite $2g+n$ geradores topológicos
$x_i,y_i$ ($1\le i\le g$) e $\sigma_j$ ($1\le j\le n$), ligados pela relação
\[
  \left(\prod_i x_i y_i x_i^{-1}y_i^{-1}\right)
  \sigma_1\cdots\sigma_n=1,
\]
onde os $\sigma_j$ são geradores dos grupos de inércia correspondentes
aos $s_j$. Para todo grupo finito $G$ de ordem prima com a característica,
gerado por elementos $x_i,y_i,\sigma_j$ que satisfazem a relação anterior,
existe um recobrimento galoisiano não ramificado de $X_0-S$, de grupo $G$, cujos
grupos de inércia nos pontos $s_j$ estão gerados pelos $\sigma_j$.

Quando $X_0$ é de gênero 0 e $n=3$, tem-se uma solução do «problema dos
três pontos», ao menos para os recobrimentos galoisianos de ordem prima com
a característica. (Aqui, de resto, o teorema 9 é inútil, e parece, por
outra parte, que o corolário anterior pode ser deduzido do caso particular
considerado no problema dos três pontos).

\src{220}
\result{OBSERVAÇÕES}{}

1) Um estudo mais completo, no qual sem dúvida interviriam recobrimentos
galoisianos generalizados de $X,X_0,X_1$ (com grupo de Galois eventualmente
infinitesimal), deveria permitir recuperar o núcleo do teorema 12, corolário
2. Em contrapartida, o estudo dos recobrimentos que admitem ramificação não
«moderada» parece ter de ser muito mais difícil.

2) O lema 6, junto com um resultado de Grauert relativo ao completamento
formal de um esquema projetivo não singular ao longo de uma seção
hiperplana (ou com o teorema, ainda não demonstrado, mencionado na
observação 2 depois do teorema 11), permite também demonstrar em geometria
algébrica «abstrata» o teorema clássico de Lefschetz sobre o grupo
fundamental.

\begin{center}
  \textbf{BIBLIOGRAFIA}
\end{center}

\begin{description}
  \item[{[1]}] Dieudonné (J.) e Grothendieck (A.). --- \emph{Éléments de
    géométrie algébrique}, Publications mathématiques de l'Institut des Hautes
    Études Scientifiques (de próxima publicação).
  \item[{[2]}] Grothendieck (Alexander). --- \emph{The cohomology theory of
    abstract algebraic varieties}. International Congress of Mathematicians
    [1958. Edinburgh] (de próxima publicação).
  \item[{[3]}] Serre (Jean-Pierre). --- \emph{Faisceaux algébriques
    cohérents}, Annals of Math., Series 2, t. 61, 1955, p. 197--278.
  \item[{[4]}] Serre (Jean-Pierre). --- \emph{Géométrie algébrique et
    géométrie analytique}, Ann. Institut Fourier Grenoble, t. 6, 1955--1956,
    p. 1--42.
\end{description}

\par

\def\FgaUnitKey{e182}\clearpage
\chapter*{Errata da exposição n\textsuperscript{o} 182}
\addcontentsline{toc}{chapter}{Errata da exposição n\textsuperscript{o} 182}
\setcounter{footnote}{1}
\markboth{Errata da exposição 182}{Errata da exposição 182}
\sid{e182}
\src{298}

\noindent\textbf{2. Geometria formal e geometria algébrica.}

\begin{center}
[Seminário Bourbaki, t. 11, 1958/59, n\textsuperscript{o} 182, 28 p.]
\end{center}

\noindent\textbf{Página 8, teorema 6.} --- A hipótese «$f$
quase-projetivo» pode ser substituída pela hipótese mais fraca «$f$
separado», em virtude do resultado seguinte (cf. SGA VIII, 6.2): todo morfismo
$f:X\longrightarrow Y$ que seja quase-finito e separado é quase-projetivo.

\src{299}
\noindent\textbf{Página 12, e página 14, observação 1.} --- Assinalemos que J.-P.
Serre construiu uma variedade projetiva não singular, de dimensão 3, sobre
um corpo algebricamente fechado $k$, de característica $p>0$, que não provém
por redução de um esquema próprio sobre um anel local que seja domínio de
integridade, com corpo residual $k$ e corpo de frações de característica
zero.\footnote{Serre (Jean-Pierre). --- Exemples de variétés projectives en
caractéristique $p$ non relevables en caractéristique zéro, \emph{Proc. Nat.
Acad. Sc. U.S.A.}, t. 47, 1961, p. 108--109.} Mumford teria encontrado um exemplo
análogo, com uma superfície projetiva não singular.

\noindent\textbf{Página 15, observações 2 e 3.} --- Atualmente sou menos
otimista quanto aos resultados aqui conjecturados. No entanto, o
problema relativo às variações da estrutura do espaço projetivo,
assinalado ao final da observação 3, foi resolvido afirmativamente por
Hironaka, e o problema análogo para as variedades abelianas foi resolvido
por Koizumi.

\noindent\textbf{Página 16, linha 3.} --- Assinalemos que Mumford construiu
recentemente os esquemas de módulos para as curvas de gênero $g$ (cf. SHMI).
Além disso, o teorema 10 prova que os «esquemas de Jacobi de nível $n$»
da teoria dos módulos são não singulares (e até mesmo lisos sobre $\mathbf Z$).

\noindent\textbf{Página 23, linha 12, primeira linha do corolário
5.} --- Acrescentar: «sendo $X$ ou $Y$ próprio sobre $k$».

A substância dos §§ 1 a 5 está contida na parte publicada de EGA III, e
a dos §§ 6 e 7 está contida em SGA III. Para o estudo do grupo
fundamental, vejam-se SGA V, IX, X, XI, bem como SGA 1962 (exposições X, XII e
XIII), para os teoremas do tipo de Lefschetz e numerosas questões em aberto.
Somente a teoria dos recobrimentos moderadamente ramificados (cf. teorema 14)
ainda não foi objeto de uma redação formal. O corolário do teorema 14,
que fornece a determinação completa dos recobrimentos galoisianos de ordem
coprima à característica de uma curva algébrica sobre um corpo
algebricamente fechado, foi utilizado de maneira essencial em três ocasiões:

1\textsuperscript{o} na demonstração, por Igusa, da desigualdade de Picard
para as superfícies projetivas não singulares em característica arbitrária;

2\textsuperscript{o} no estudo (desenvolvido independentemente por Ogg e
Šafarevič) do grupo dos fibrados principais homogêneos sob uma variedade
abeliana definida sobre um corpo de funções de uma variável, em
característica arbitrária;

3\textsuperscript{o} na demonstração recente, por Artin, de certos
teoremas fundamentais sobre a «cohomologia de Weil» das variedades algébricas.

\par

\def\FgaUnitKey{190}\unit{Exposição 190 -- Técnica de descida I}
\sid{190}
\markboth{Exposição 190}{Exposição 190}
\src{299}
\seminar{12.\textsuperscript{o} ano, 1959/60, n.\textsuperscript{o} 190\\Dezembro de 1959}
\paperhead{TÉCNICA DE DESCIDA E TEOREMAS DE EXISTÊNCIA EM GEOMETRIA ALGÉBRICA\\
I. GENERALIDADES. DESCIDA POR MORFISMOS FIELMENTE PLANOS}%
  {por Alexander GROTHENDIECK}

Do ponto de vista técnico, a presente exposição e as que se seguirão podem
ser consideradas variações sobre o célebre «teorema 90» de Hilbert.
Parece que A. Weil introduziu o método de descida em geometria algébrica sob
o nome de «descida do corpo base». Aliás, Weil restringiu-se ao caso de
extensões finitas separáveis de corpos. P. Cartier tratou depois o caso de
extensões radiciais de altura 1. Por falta da linguagem dos esquemas e,
mais concretamente, por não dispor de elementos nilpotentes nos anéis que
lhe era permitido considerar, Cartier não pudera formular a identidade
essencial entre esses dois casos.

Atualmente, parece que a técnica geral de descida que será exposta (unida,
quando couber, aos teoremas fundamentais da «geometria formal», cf. [3])
está na base da maioria dos teoremas de existência em geometria
algébrica. Convém assinalar, além disso, que essa técnica pode, sem dúvida, ser
transposta à «geometria analítica», e cabe esperar que, num futuro próximo,
os especialistas saibam demonstrar os análogos «analíticos» dos teoremas de
existência em geometria formal que serão apresentados na exposição II. Em
todo caso, os trabalhos recentes de Kodaira--Spencer, cujos métodos parecem
inadequados para definir e estudar «variedades de módulos» numa vizinhança
de seus pontos singulares, mostram com bastante clareza a necessidade de
métodos mais próximos da teoria dos esquemas (que deverão completar,
naturalmente, as técnicas transcendentes).

Na presente exposição I trataremos o caso de descida mais elementar (indicado
no título). No entanto, as aplicações dos teoremas 1, 2 e 3 que se seguem já
são muito numerosas. Limitamo-nos a apresentar algumas a título de exemplo,
sem procurar conferir-lhes aqui a máxima generalidade possível.

Utilizaremos livremente a linguagem dos esquemas, para a qual remetemos à
exposição já citada, bem como a [2]. Assinalemos, contudo, expressamente que os
pré-esquemas considerados na presente exposição não são necessariamente
noetherianos, e os

\src{300}
morfismos não são necessariamente de tipo finito. Assim, se $A$ é um anel local
noetheriano, com completamento $\widehat A$, será preciso considerar o anel não
noetheriano $\widehat A\otimes_A\widehat A$, e os morfismos de esquemas afins
correspondentes às inclusões entre os anéis considerados.

\fsection{A}{Preliminares sobre as categorias}

\subsection*{1. Categorias fibradas, dados de descida, morfismos de
$\mathcal F$-descida}
\addcontentsline{toc}{subsection}{1. Categorias fibradas e dados de descida}

\result{a. DEFINIÇÃO}{1.1} Uma \emph{categoria fibrada} $\mathcal F$ de base
$\mathcal C$ consta de uma categoria $\mathcal C$; para cada
$X\in\mathcal C$, de uma categoria $\mathcal F_X$; para cada
$\mathcal C$-morfismo $f:X\longrightarrow Y$, de um funtor
$f^*:\mathcal F_Y\longrightarrow\mathcal F_X$, denotado também por
\[
  f^*(\xi)=\xi\times_Y X
\]
para $\xi\in\mathcal F_Y$ (considerando $X$ como um «objeto de
$\mathcal C$ sobre $Y$», isto é, provido do morfismo $f$); e, por fim,
para dois morfismos componíveis
$X\xrightarrow{f}Y\xrightarrow{g}Z$, de um isomorfismo de funtores
\[
  c_{f,g}:(gf)^*\longrightarrow f^*g^*,
\]
estando esses dados sujeitos às condições seguintes:

\begin{enumerate}
\item[(i)] $\operatorname{id}_X^*=\operatorname{id}$;
\item[(ii)] $c_{f,g}$ é o isomorfismo identidade se $f$ ou $g$ é um
isomorfismo identidade;
\item[(iii)] dados três morfismos componíveis
$X\xrightarrow{f}Y\xrightarrow{g}Z\xrightarrow{h}T$, o seguinte diagrama de
isomorfismos, formado mediante isomorfismos da forma $c_{u,v}$, é
comutativo:
\[
\begin{array}{ccc}
  (h(gf))^*=((hg)f)^* & \longrightarrow & f^*(hg)^* \\
  \big\downarrow && \big\downarrow \\
  (gf)^*h^* & \longrightarrow & (f^*g^*)h^*=f^*(g^*h^*).
\end{array}
\]
\end{enumerate}

\result{EXEMPLO}{1} Seja $\mathcal C$ uma categoria na qual existem os
produtos fibrados. Define-se então uma categoria fibrada $\mathcal F$ de
base $\mathcal C$ tomando $\mathcal F_X=$ categoria dos objetos de
$\mathcal C$ sobre $X$; o funtor
$f^*:\mathcal F_Y\longrightarrow\mathcal F_X$ correspondente a um morfismo
$f:X\longrightarrow Y$ define-se por meio do produto fibrado
$Z\longmapsto Z\times_YX$.

\result{EXEMPLO}{2} Seja $\mathcal C$ a categoria dos pré-esquemas e, para
$X\in\mathcal C$, seja $\mathcal F_X$ a categoria dos feixes
quase-coerentes de módulos sobre $X$. Se $f:X\longrightarrow Y$ é um morfismo
de pré-esquemas, $f^*:\mathcal F_Y\longrightarrow\mathcal F_X$ é o funtor
imagem inversa

\src{301}
de feixes de módulos. Obtém-se assim uma categoria fibrada de base
$\mathcal C$.

\result{b. DEFINIÇÃO}{1.2} Um diagrama
\[
  E\xrightarrow{u}E'
  \overset{v_1}{\underset{v_2}{\rightrightarrows}}E''
\]
de aplicações de conjuntos chama-se \emph{exato} se $u$ é uma bijeção de
$E$ sobre a parte de $E'$ formada pelos $x'\in E'$ tais que
$v_1(x')=v_2(x')$.

\result{DEFINIÇÃO}{1.3} Seja $\mathcal F$ uma categoria fibrada de base
$\mathcal C$; consideremos um diagrama
\[
  \tag{+}
  S\xleftarrow{\alpha}S'
  \overset{\beta_1}{\underset{\beta_2}{\rightrightarrows}}S''
\]
de morfismos em $\mathcal C$ tais que
$\alpha\beta_1=\alpha\beta_2$; esse diagrama chama-se
$\mathcal F$-\emph{exato} se, para cada par $\xi,\eta$ de elementos de
$\mathcal F_S$, o seguinte diagrama de aplicações de conjuntos
\[
  \tag{+'}
  \operatorname{Hom}(\xi,\eta)
  \xrightarrow{\alpha^*}
  \operatorname{Hom}(\alpha^*\xi,\alpha^*\eta)
  \overset{\beta_1^*}{\underset{\beta_2^*}{\rightrightarrows}}
  \operatorname{Hom}(\gamma^*\xi,\gamma^*\eta),
\]
onde $\gamma=\alpha\beta_1=\alpha\beta_2$, é exato. Neste último diagrama
temos identificado $\beta_i^*\alpha^*$ com
$(\alpha\beta_i)^*=\gamma^*$ mediante $c_{\beta_i,\alpha}$, para simplificar.

\result{DEFINIÇÃO}{1.4} Seja $\mathcal F$ uma categoria fibrada de base
$\mathcal C$; consideremos dois morfismos $\beta_1,\beta_2$ de $S''$ em $S'$,
em $\mathcal C$. Seja
$\xi'\in\mathcal F_{S'}$. Chama-se \emph{dado de colagem} sobre $\xi'$
(relativo ao par $(\beta_1,\beta_2)$) a um isomorfismo de
$\beta_1^*(\xi')$ sobre $\beta_2^*(\xi')$. Se $\xi',\eta'\in\mathcal F_{S'}$
estão munidos, cada um, de um dado de colagem, um morfismo
$u:\xi'\longrightarrow\eta'$ em $\mathcal F_{S'}$ se diz
\emph{compatível com os dados de colagem} se o diagrama seguinte é
comutativo:
\[
\begin{array}{ccc}
  \beta_1^*(\xi') & \longrightarrow & \beta_2^*(\xi') \\
  \big\downarrow && \big\downarrow \\
  \beta_1^*(\eta') & \longrightarrow & \beta_2^*(\eta').
\end{array}
\]
Desse modo, os objetos de $\mathcal F_{S'}$ providos de um dado de colagem
(relativo a $\beta_1,\beta_2$) formam uma categoria. Se
$\alpha:S'\longrightarrow S$ é um morfismo tal que
$\alpha\beta_1=\alpha\beta_2$, então, para todo $\xi\in\mathcal F_S$, o
objeto $\xi'=\alpha^*(\xi)$

\src{302}
de $\mathcal F_{S'}$ está provido de um dado de colagem canônico, visto que
\[
  \beta_i^*\alpha^*(\xi)\simeq(\alpha\beta_i)^*(\xi)=\gamma^*(\xi),
\]
tomando de novo $\gamma=\alpha\beta_1=\alpha\beta_2$; além disso, se
$u:\xi\longrightarrow\eta$ é um morfismo em $\mathcal F_S$, então
\[
  \alpha^*(u):\alpha^*(\xi)\longrightarrow\alpha^*(\eta)
\]
é um morfismo em $\mathcal F_{S'}$ compatível com os dados de colagem
canônicos. Obtém-se assim um funtor canônico da categoria $\mathcal F_S$ para
a categoria dos objetos de $\mathcal F_{S'}$ providos de um dado de colagem
relativo ao par $(\beta_1,\beta_2)$. Dito isto, a definição 1.3 também
pode ser expressa dizendo que o diagrama (+) dessa definição é
$\mathcal F$-exato se o funtor anterior é «plenamente fiel», isto é, se
define uma equivalência da categoria $\mathcal F_S$ com uma subcategoria da
categoria dos objetos de $\mathcal F_{S'}$ providos de um dado de colagem
relativo a $(\beta_1,\beta_2)$.

\result{DEFINIÇÃO}{1.5} Diz-se que um dado de colagem sobre
$\xi'\in\mathcal F_{S'}$ é \emph{efetivo} (com respeito a $\alpha$) se
$\xi'$, provido deste dado, é isomorfo a um $\alpha^*(\xi)$, com
$\xi\in\mathcal F_S$.

Quando o diagrama (+) é $\mathcal F$-exato, o objeto $\xi$ da definição
anterior fica, portanto, determinado salvo isomorfismo único, e a categoria
$\mathcal F_S$ é equivalente à categoria dos objetos de
$\mathcal F_{S'}$ providos de um dado de colagem efetivo.

\medskip\noindent
c. O caso mais importante é aquele em que
\[
  S''=S'\times_S S',
\]
sendo as $\beta_i$ as duas projeções $p_1$ e $p_2$ de
$S'\times_S S'$ sobre seus dois fatores (supomos agora que em
$\mathcal C$ existem os produtos fibrados). Têm-se então duas
condições necessárias imediatas para que seja efetivo um dado de colagem
$\varphi:p_1^*(\xi')\longrightarrow p_2^*(\xi')$ sobre um
$\xi'\in\mathcal F_{S'}$. Tem-se
\[
  \tag{i}
  \Delta^*(\varphi)=\operatorname{id}_{\xi'},
\]
onde $\Delta:S'\longrightarrow S'\times_S S'$ designa o morfismo diagonal,
e onde se identificou $\Delta^*p_i^*(\xi')$ com
$(p_i\Delta)^*(\xi')=\xi'$. Tem-se
\[
  \tag{ii}
  p_{31}^*(\varphi)=p_{32}^*(\varphi)\,p_{21}^*(\varphi),
\]

\src{303}
onde $p_{ij}$ designa a projeção canônica de
$S'\times_S S'\times_S S'$ sobre o produto parcial de seus fatores
$i$-ésimo e $j$-ésimo.

\result{DEFINIÇÃO}{1.6} Chama-se \emph{dado de descida} sobre
$\xi'\in\mathcal F_{S'}$, relativo ao morfismo
$\alpha:S'\longrightarrow S$, a um dado de colagem sobre $\xi'$ relativo ao par
$(p_1,p_2)$ de projeções canônicas
$S'\times_S S'\longrightarrow S'$ que satisfaz as condições (i) e (ii)
anteriores.

\result{DEFINIÇÃO}{1.7} Um morfismo $\alpha:S'\longrightarrow S$ chama-se
\emph{morfismo de $\mathcal F$-descida} se o diagrama de morfismos
\[
  S\xleftarrow{\alpha}S'
  \overset{p_1}{\underset{p_2}{\rightrightarrows}}S'\times_S S'
\]
é $\mathcal F$-exato (definição 1.3); diz-se que $\alpha$ é um
\emph{morfismo de $\mathcal F$-descida estrito} se, além disso, todo dado de
descida (definição 1.6) sobre um objeto de $\mathcal F_{S'}$ é efetivo.

Esta última condição também pode ser enunciada, de modo mais gráfico, assim:
«Dá no mesmo» dar um objeto de $\mathcal F_S$ ou um objeto de
$\mathcal F_{S'}$ provido de um dado de descida.

Observe-se que se um morfismo de $\mathcal F$-descida
$\alpha:S'\longrightarrow S$ admite uma «seção»
$s:S\longrightarrow S'$ (isto é, um morfismo $s$ tal que
$\alpha s=\operatorname{id}_S$), então é um morfismo de
$\mathcal F$-descida estrito. Com efeito, se $\xi'\in\mathcal F_{S'}$ está provido de
um dado de descida relativo a $\alpha$, «provém» de
$\xi=s^*(\xi')$.

\medskip\noindent
d. As noções anteriores podem ser apresentadas de maneira mais intuitiva
introduzindo, para um objeto parâmetro $T$ de $\mathcal C$ sobre $S$, o
conjunto
\[
  \operatorname{Hom}_S(T,S')=S'(T),
\]
cujos elementos serão denotados por $t,t'$, etc. Dado um objeto
$\xi'\in\mathcal F_{S'}$, a cada $t\in S'(T)$ corresponde então um objeto
$t^*(\xi')$ de $\mathcal F_T$, que também será denotado por $\xi'_t$. Um dado de
colagem sobre $\xi'$ (relativo a $p_1,p_2$) fica definido então dando, para
todo $T$ sobre $S$ e todo par de pontos $t,t'\in S'(T)$, um isomorfismo
\[
  \varphi_{t',t}:\xi'_t\longrightarrow\xi'_{t'}
\]
(que satisfaz condições evidentes de funtorialidade ao variar $T$). As
condições (i) e (ii) de 1.c escrevem-se então
\[
  \tag{i bis}
  \varphi_{t,t}=\operatorname{id}_{\xi'_t}
  \qquad\text{para todo $T$ e todo $t\in S'(T)$},
\]

\src{304}
\[
  \tag{ii bis}
  \varphi_{t,t''}=\varphi_{t,t'}\varphi_{t',t''}
  \qquad\text{para todo $T$ e todos $t,t',t''\in S'(T)$}.
\]
Vê-se, além disso, que (ii bis), tomando $t=t'=t''$, implica
$\varphi_{t,t}=\varphi_{t,t}^2$; visto que $\varphi_{t,t}$ é um isomorfismo
por hipótese, daí se obtém a relação (i bis), que é de fato uma consequência
de (ii bis) (portanto, (i) é consequência de (ii)). Mas, se já não se supõe
a priori que os $\varphi_{t,t'}$ sejam isomorfismos (isto é, que
$\varphi:p_1^*(\xi')\longrightarrow p_2^*(\xi')$ seja um isomorfismo), então
(ii bis) não implica necessariamente (i bis); no entanto, a conjunção de
(ii bis) e (i bis) implica que os $\varphi_{t,t'}$ são isomorfismos (pois se
terá
$\varphi_{t,t'}\varphi_{t',t}=\varphi_{t,t}=\operatorname{id}_{\xi'_t}$).

\subsection*{2. Diagramas exatos e epimorfismos estritos, morfismos de
descida. Exemplos}
\addcontentsline{toc}{subsection}{2. Diagramas exatos e epimorfismos estritos}

\result{a. DEFINIÇÃO}{2.1} Seja $\mathcal C$ uma categoria. Um diagrama
\[
  T\xrightarrow{\alpha}T'
  \overset{\beta_1}{\underset{\beta_2}{\rightrightarrows}}T''
\]
de morfismos em $\mathcal C$ chama-se \emph{exato} se, para todo
$Z\in\mathcal C$, o diagrama correspondente de aplicações de conjuntos
\[
  \operatorname{Hom}(Z,T)\longrightarrow
  \operatorname{Hom}(Z,T')
  \rightrightarrows\operatorname{Hom}(Z,T'')
\]
é exato (definição 1.2). Diz-se então que $(T,\alpha)$ (ou, por abuso de
linguagem, que $T$) é um \emph{núcleo} do par de morfismos
$(\beta_1,\beta_2)$.

Evidentemente, esse núcleo está determinado salvo isomorfismo único. Se
$\mathcal C$ é a categoria dos conjuntos, a definição anterior é
compatível com a definição 1.2. Define-se de maneira dual a exatidão de um
diagrama
\[
  S\xleftarrow{\alpha}S'
  \overset{\beta_1}{\underset{\beta_2}{\rightrightarrows}}S''
\]
de morfismos em $\mathcal C$; diz-se então que $(S,\alpha)$ é um
\emph{conúcleo} do par de morfismos $(\beta_1,\beta_2)$.

\result{DEFINIÇÃO}{2.2} Um morfismo $\alpha:S'\longrightarrow S$ chama-se
\emph{epimorfismo estrito} se é um epimorfismo e se, para todo morfismo
$u:S'\longrightarrow Z$, a seguinte condição necessária é também
suficiente para que $u$ se fatore como
$S'\longrightarrow S\longrightarrow Z$: para todo $S''\in\mathcal C$ e todo
par de morfismos $\beta_1,\beta_2:S''\longrightarrow S'$ tal que
$\alpha\beta_1=\alpha\beta_2$, tem-se também $u\beta_1=u\beta_2$.

Se existe o produto fibrado $S'\times_S S'$, equivale a dizer que o diagrama
\[
  S\xleftarrow{\alpha}S'
  \overset{p_1}{\underset{p_2}{\rightrightarrows}}S'\times_S S'
\]

\src{305}
é exato, isto é, que $S$ é um conúcleo do par $(p_1,p_2)$. Em todo
caso, um morfismo conúcleo é um epimorfismo estrito. Assinalemos também que um
epimorfismo estrito que seja um monomorfismo é um isomorfismo. Deixamos ao
leitor desenvolver a noção dual de \emph{monomorfismo estrito}.

Para precisar as relações entre a noção de morfismo de
$\mathcal F$-descida e a de epimorfismo estrito, introduzimos além disso as
definições seguintes:

\result{DEFINIÇÃO}{2.3} Um morfismo $\alpha:S'\longrightarrow S$ chama-se
\emph{epimorfismo universal} (resp. \emph{epimorfismo estrito universal}) se,
para todo $T$ sobre $S$, existe o produto fibrado $T'=S'\times_ST$ e a
projeção $T'\longrightarrow T$ é um epimorfismo (resp. um epimorfismo
estrito).

Nas categorias muito boas (como a categoria dos conjuntos, a
categoria dos feixes de conjuntos sobre um espaço topológico, as categorias
abelianas, etc.), as quatro noções de epimorfismo assim introduzidas
coincidem; ao contrário, todas elas são distintas numa categoria como
a dos pré-esquemas, ou a dos pré-esquemas sobre um pré-esquema não vazio
dado $S$, inclusive se nos limitamos a $S$-esquemas finitos sobre $S$.

\result{DEFINIÇÃO}{2.4} Um morfismo $\alpha:S'\longrightarrow S$ chama-se
\emph{morfismo de descida} (resp. \emph{morfismo de descida estrito}) se é
um morfismo de $\mathcal F$-descida (resp. um morfismo de
$\mathcal F$-descida estrito) (cf. definição 1.7), onde $\mathcal F$ denota
a categoria fibrada de base $\mathcal C$ dos objetos de $\mathcal C$ sobre
objetos de $\mathcal C$ (n\textsuperscript{o} 1, exemplo 1).

\result{PROPOSIÇÃO}{2.1} Se em $\mathcal C$ existem os produtos finitos e
os produtos fibrados (finitos), então, em $\mathcal C$, os morfismos de
descida coincidem com os epimorfismos estritos universais.

\result{b. EXEMPLOS}{} Seja $\mathcal C$ a categoria dos pré-esquemas. Seja
$S\in\mathcal C$, e sejam $S'$ e $S''$ dois pré-esquemas finitos sobre $S$,
isto é, correspondentes a feixes de álgebras $\mathcal A'$, $\mathcal A''$
sobre $S$ que, como feixes de módulos, são quase-coerentes e de tipo finito
(isto é, coerentes se $S$ é localmente noetheriano). Seja
$\alpha:S'\longrightarrow S$ o morfismo estrutural de $S'$, e sejam
$\beta_1,\beta_2$ dois $S$-morfismos de $S''$ em $S'$, definidos por
homomorfismos de álgebras $\mathcal A'\longrightarrow\mathcal A''$, denotados
também por $\beta_1,\beta_2$. Utilizando o fato de que um morfismo finito é
fechado (primeiro teorema de Cohen--Seidenberg), demonstra-se facilmente que o
diagrama
\[
  \tag{+}
  S\xleftarrow{\alpha}S'
  \overset{\beta_1}{\underset{\beta_2}{\rightrightarrows}}S''
\]
em $\mathcal C$ é exato se e somente se também o é o diagrama de feixes

\src{306}
\[
  \mathcal O_S=\mathcal A\longrightarrow\mathcal A'
  \overset{\beta_1}{\underset{\beta_2}{\rightrightarrows}}\mathcal A''
\]
sobre $S$. Em particular, se $\alpha:S'\longrightarrow S$ é um morfismo
finito correspondente a um feixe $\mathcal A'$ de álgebras sobre $S$, então
$\alpha$ é um epimorfismo estrito se e somente se o diagrama de feixes
\[
  \mathcal O_S=\mathcal A\longrightarrow\mathcal A'
  \overset{p_1}{\underset{p_2}{\rightrightarrows}}
  \mathcal A'\otimes_{\mathcal A}\mathcal A'
\]
é exato (é um epimorfismo se e somente se
$\mathcal A\longrightarrow\mathcal A'$ é injetivo). Se $S$ é afim de anel
$A$, e portanto $S'$ é afim de anel $A'$ finito sobre $A$, então
$S'\longrightarrow S$ é um epimorfismo estrito se e somente se
$A\longrightarrow A'$ é um isomorfismo de $A$ sobre o subanel de $A'$
formado pelos $x'\in A'$ tais que
\[
  1_{A'}\otimes_Ax'-x'\otimes_A1_{A'}=0
\]
(é um epimorfismo se e somente se $A\longrightarrow A'$ é injetivo). Como já
assinalamos, mesmo se $S$ é o esquema de um anel local artiniano, um
morfismo finito $S'\longrightarrow S$ que seja um epimorfismo não é
necessariamente um epimorfismo estrito. No entanto, pode-se demonstrar que, se
$S$ é um pré-esquema noetheriano, todo morfismo finito
$S'\longrightarrow S$ que seja um epimorfismo é a composição de uma sucessão
finita de epimorfismos estritos (também finitos). Isso mostra, além disso, que
a composição de dois epimorfismos estritos não é necessariamente um
epimorfismo estrito.

\result{c}{} Se $(+)$ é um diagrama exato de morfismos finitos, então,
para todo morfismo plano $T\longrightarrow S$ de pré-esquemas, o diagrama
obtido de $(+)$ mediante a mudança de base $T\longrightarrow S$ continua
exato. Disso se deduz que, se $X,Y$ são dois $S$-pré-esquemas, sendo $X$
plano sobre $S$, então o seguinte diagrama de aplicações de conjuntos
(onde $X',Y'$ são as imagens inversas de $X,Y$ sobre $S'$, e $X'',Y''$ suas
imagens inversas sobre $S''$)
\[
  \operatorname{Hom}_S(X,Y)\longrightarrow
  \operatorname{Hom}_{S'}(X',Y')
  \rightrightarrows
  \operatorname{Hom}_{S''}(X'',Y'')
\]
é exato. Em particular, se $\mathcal F$ denota a categoria fibrada cuja base
é a categoria $\mathcal C$ dos pré-esquemas e tal que, para
$X\in\mathcal C$, $\mathcal F_X$ é a categoria dos $X$-pré-esquemas
planos, então o diagrama $(+)$ é $\mathcal F$-exato. (Esse resultado deixa
de ser verdadeiro se não se impõe a hipótese de planicidade; em particular, um
epimorfismo estrito finito não é necessariamente um morfismo de descida.)
Vê-se igualmente que $(+)$ é $\mathcal F$-exato se $\mathcal F$ denota a
categoria fibrada para a qual $\mathcal F_X$ é a categoria dos feixes
quase-coerentes e planos sobre o pré-esquema $X$ (também aqui a hipótese

\src{307}
de planicidade é essencial). Num e noutro caso, a questão da efetividade de
um dado de colagem (e, mais concretamente, de um dado de descida quando
$S''=S'\times_SS'$) sobre um objeto plano sobre $S'$ é delicada (e sua
resposta em diversos casos particulares é um dos objetos principais das
presentes exposições). O conferenciante ignora se, para todo epimorfismo
estrito finito $S'\longrightarrow S$, todo dado de descida sobre um feixe
quase-coerente plano sobre $S'$ é efetivo (mesmo supondo que $S$ seja o
espectro de um anel local artiniano e limitando-nos aos feixes localmente
livres de posto 1). Mais geralmente, sejam $A$ um anel e $A'$ uma
$A$-álgebra (tudo é comutativo) tal que o diagrama de aplicações
\[
  A\longrightarrow A'\rightrightarrows A'\otimes_AA'
\]
seja exato, o que equivale também a que o morfismo
$S'\longrightarrow S$ correspondente entre os espectros de $A,A'$ seja um
morfismo de $\mathcal F$-descida, onde $\mathcal F$ é a categoria fibrada
dos feixes quase-coerentes planos. Seja $M'$ um $A'$-módulo plano provido de
um dado de descida a $A$, isto é, de um isomorfismo
\[
  \varphi:M'\otimes_AA'\xrightarrow{\sim}A'\otimes_AM'
\]
de $A'\otimes_AA'$-módulos que satisfaz as condições (i) e (ii) de 1.c
(cuja explicitação em termos de módulos deixamos ao leitor). É efetivo
este dado (com respeito à categoria fibrada dos feixes quase-coerentes
planos)? Seja $M$ o subconjunto de $M'$ formado pelos $x'\in M'$ tais que
\[
  \varphi(x'\otimes_A1_{A'})=1_{A'}\otimes_Ax',
\]
que é um sub-$A$-módulo de $M'$. A injeção canônica
$M\longrightarrow M'$ define um homomorfismo de $A'$-módulos
$M\otimes_AA'\longrightarrow M'$. A efetividade de $\varphi$ significa
então o seguinte: $M$ é um $A$-módulo plano e o homomorfismo anterior é
um isomorfismo.

\result{OBSERVAÇÃO}{} Nas considerações anteriores não havíamos imposto
nenhuma hipótese de planicidade aos morfismos do diagrama $(+)$, o que nos
obrigava, para dispor de uma técnica de descida, a impor hipóteses de
planicidade aos objetos sobre $S,S'$ considerados. No parágrafo 2 imporemos
uma hipótese de planicidade a $\alpha:S'\longrightarrow S$, o que nos permitirá
dispor de uma técnica de descida para objetos sobre $S,S'$ que já não estão
submetidos a condição alguma de planicidade. Em todos os casos intervém uma
hipótese de planicidade. Essa é uma das principais razões da importância da
noção de planicidade em geometria algébrica (cujo papel não podia manifestar-se
enquanto nos limitávamos a corpos base, sobre os quais, com efeito, tudo é
plano!).

\src{308}
\subsection*{3. Aplicação aos morfismos étale}
\addcontentsline{toc}{subsection}{3. Aplicação aos morfismos étale}

Sejam $A$ um anel local e $B$ um álgebra local sobre $A$ cujo ideal maximal
induz o de $A$. Diremos que $B$ é \emph{étale} sobre $A$ (em lugar de
«não ramificada», expressão usada em outros lugares) se satisfaz as condições
seguintes:
\begin{enumerate}
\item[(i)] $B$ é plana sobre $A$;
\item[(ii)] $B/\mathfrak mB$ é uma extensão finita separável de
$A/\mathfrak m=k$ (onde $\mathfrak m$ denota o ideal maximal de $A$).
\end{enumerate}
Quando $A$ e $B$ são noetherianos e $k$ é algebricamente fechado, isso
significa que o homomorfismo $\widehat A\longrightarrow\widehat B$ entre os
completamentos que prolonga $A\longrightarrow B$ é um isomorfismo. Um morfismo de
tipo finito $f:T\longrightarrow S$ diz-se \emph{étale em $x\in T$}, ou também
que $T$ é \emph{étale sobre $S$ em $x$}, se $\mathcal O_x$ é étale sobre
$\mathcal O_{f(x)}$; diz-se que $f$ é \emph{étale}, ou chama-se $f$ um
\emph{morfismo étale}, ou diz-se que $T$ é \emph{étale sobre $S$}, se $f$ é
étale em todo $x\in T$. Assinalemos, além disso, que se $S$ é localmente
noetheriano, o conjunto de pontos de $T$ nos quais $f$ é étale é aberto, e
o uso do «main theorem» de Zariski permite precisar a estrutura de $T/S$
numa vizinhança de tal ponto (por meio de uma equação de tipo bem conhecido).

Quando $S$ é um esquema de tipo finito sobre o corpo dos números
complexos, corresponde-lhe um espaço analítico $\overline S$ no sentido de
Serre [5], exceto que $\overline S$ pode conter elementos nilpotentes em seu
feixe estrutural, o que em nada altera o essencial de [5]. Vê-se então facilmente
que $f$ é étale se e somente se
$\overline f:\overline T\longrightarrow\overline S$ também o é, isto é, se todo
ponto de $\overline T$ admite uma vizinhança na qual $\overline f$ induz um
isomorfismo sobre um aberto de $\overline S$. Em particular, a todo
\emph{recobrimento étale} $T$ de $S$ (isto é, um morfismo finito étale
$f:T\longrightarrow S$) corresponde um recobrimento étale $\overline T$ de
$\overline S$, que é conexo se e somente se $T$ também o é [5]. Vê-se, além
disso, facilmente que, se $T,T'$ são dois esquemas étale sobre $S$, então a
aplicação natural
\[
  \operatorname{Hom}_S(T,T')\longrightarrow
  \operatorname{Hom}_{\overline S}(\overline T,\overline{T'})
\]
é bijetiva, isto é, o funtor $T\longmapsto\overline T$ da categoria dos
esquemas étale sobre $S$ à categoria dos espaços analíticos étale
sobre $\overline S$ é «plenamente fiel» e, portanto, define uma equivalência
da primeira categoria com uma subcategoria da segunda. Um teorema de
Grauert--Remmert [2] implica que, se $S$ é normal, obtém-se assim uma
equivalência entre a categoria dos \emph{recobrimentos étale} de $S$ e a
categoria dos recobrimentos étale (finitos) de $\overline S$, isto é, que
todo recobrimento étale de $\overline S$ é $\overline S$-isomorfo a um
$\overline T$, onde $T$ é um recobrimento étale de $S$. Demonstremos que o
\emph{teorema de Grauert--Remmert continua válido sem hipótese de
normalidade sobre $S$}. Seja, com efeito, primeiro $S'\longrightarrow S$ um
epimorfismo estrito finito; suponhamos demonstrado o

\src{309}
teorema para $S'$ e demonstremos que será verdadeiro para $S$. Seja
$\mathfrak T$ um recobrimento étale de $\overline S$ e consideremos sua imagem
inversa $\mathfrak T'$ sobre $\overline{S'}$, que corresponde a um feixe
analítico coerente $\mathcal A'$ de álgebras sobre $\overline{S'}$, imagem
inversa do feixe de álgebras $\mathcal A$ sobre $\overline S$ que define
$\mathfrak T$. Por hipótese, sobre $\overline{S'}$, $\mathfrak T'$ provém de
um recobrimento étale $T'$ de $S'$, isto é, $\mathcal A'$ provém de um feixe
coerente de álgebras $A'$ sobre $S'$. Por outro lado, $\mathcal A'$ está
provido de um dado de descida canônico relativo a
$\overline{S'}\longrightarrow\overline S$, isto é, de um isomorfismo entre
suas duas imagens inversas sobre
\[
  \overline{S'}\times_{\overline S}\overline{S'}
  =\overline{(S'\times_SS')},
\]
(que satisfaz as condições (i), (ii)), e esse isomorfismo provém, conforme
dito, de um isomorfismo entre os feixes algébricos correspondentes, isto é,
de um dado de descida sobre $A'$ relativo a
$S'\longrightarrow S$. Verifica-se facilmente que este último é efetivo
(pois o de $\mathcal A'$ também o é, e a efetividade de um dado de descida,
tal como se explicitou no número anterior, reconhece-se localmente nos
\emph{completamentos} dos módulos que intervêm). Obtém-se assim um feixe
coerente de álgebras $A$ sobre $S$, que define um recobrimento $T$ de $S$,
o recobrimento procurado. O resultado anterior continua então manifestamente
válido se $S'\longrightarrow S$ for apenas uma composição de um número finito
de epimorfismos estritos finitos, isto é, se for um epimorfismo finito
qualquer (pelo resultado de fatoração assinalado no parágrafo 2). Segue-se que
o teorema de Grauert--Remmert continua válido se $S$ é um esquema
\emph{reduzido}, isto é, se $\mathcal O_S$ não tem elementos nilpotentes, como
se vê introduzindo sua normalização $S'$. Daí se passa facilmente ao caso geral.

Uma demonstração inteiramente análoga, que utiliza de novo o resultado de
fatoração para os epimorfismos estritos finitos e a natureza «formal» da
efetividade dos dados de descida, permite demonstrar o resultado seguinte:
seja $S$ um pré-esquema localmente noetheriano e seja
$S'\longrightarrow S$ um morfismo \emph{finito, sobrejetivo e radicial} (ou,
o que dá no mesmo, um morfismo de tipo finito tal que, para todo $T$
sobre $S$, o morfismo $T'=S'\times_ST\longrightarrow T$ seja um homeomorfismo,
o que também se exprime dizendo que $S'\longrightarrow S$ é um
\emph{homeomorfismo universal}). Para todo $T$ étale sobre $S$, consideremos
sua imagem inversa $T'=T\times_SS'$, que é étale sobre $S'$. Então o
funtor $T\longmapsto T'$ é uma \emph{equivalência de categorias entre a
categoria dos pré-esquemas $T$ étale sobre $S$ e a categoria dos
pré-esquemas $T'$ étale sobre $S'$}. (Utiliza-se a bijetividade de
\[
  \operatorname{Hom}_S(T_1,T_2)\longrightarrow
  \operatorname{Hom}_{S'}(T'_1,T'_2)
\]
para dois pré-esquemas $T_1,T_2$ étale sobre $S$, fato cuja verificação direta
é fácil, e o fato de que o teorema enunciado é verdadeiro se
$S'=(S,\mathcal O_S/\mathcal J)$,

\src{310}
onde $\mathcal J$ é um feixe coerente nilpotente de ideais de $\mathcal O_S$
([4], lema 6)). Observe-se, além disso, que não supomos aqui que os $T$
considerados sejam finitos sobre $S$. Este resultado implica, em particular,
que o morfismo $S'\longrightarrow S$ induz um isomorfismo do grupo
fundamental $[\,]$ de $S'$ sobre o de $S$ («\emph{invariância topológica do
grupo fundamental de um pré-esquema}»).

\subsection*{4. Relações com a 1-cohomologia}
\addcontentsline{toc}{subsection}{4. Relações com a 1-cohomologia}

\medskip\noindent
a. Seja $\mathcal C$ uma categoria na qual sempre existe o produto de dois
objetos, e seja $T\in\mathcal C$. Para todo conjunto finito
$I\ne\varnothing$, pode-se considerar $T^I$; ao variar $I$, obtém-se assim um
funtor covariante da categoria dos conjuntos finitos não vazios para
$\mathcal C$, isto é, o que se pode chamar um \emph{objeto simplicial} de
$\mathcal C$, denotado por $K_T$. Este último depende covariantemente de $T$;
além disso, se $u,v$ são dois morfismos $T\longrightarrow T'$, os morfismos
correspondentes $K_T\longrightarrow K_{T'}$ são homotópicos. Diremos que $T$
\emph{domina} a $T'$ se $\operatorname{Hom}(T,T')\ne\varnothing$; esta é uma
relação de pré-ordem filtrante crescente em $\mathcal C$. Do anterior se
deduz que, se $T$ domina $T'$, existe uma classe canônica (salvo homotopia)
de homomorfismos de objetos simpliciais $K_T\longrightarrow K_{T'}$; em
particular, se $K_T$ e $K_{T'}$ são tais que cada um domina ao outro, então
$K_T$ e $K_{T'}$ são homotopicamente equivalentes. Seja agora $F$ um funtor
(contravariante, para fixar ideias) de $\mathcal C$ para uma categoria abeliana
$\mathcal C'$; então
\[
  C^*(T,F)=F(K_T)
\]
é um objeto cosimplicial de $\mathcal C'$ e, portanto, define da maneira
habitual um complexo (de cocadeias) em $\mathcal C'$, do qual se pode tomar a
cohomologia
\[
  H^*(T,F)=H^*(C^*(T,F))=H^*(F(K_T))
\]
(poder-se-á escrever uma $\mathcal C$ como subíndice de $H^*$ se houver
possibilidade de confusão). Esse é um funtor cohomológico em $F$, cuja
variância ao variar $T$ se deduz do que foi dito sobre os $K_T$; mais
precisamente, para $F$ fixo e $T$ variável em $\mathcal C$ (pré-ordenada pela
relação de dominação), os $H^*(T,F)$ formam um sistema indutivo de objetos
graduados de $\mathcal C'$;
em particular, se $T$ e $T'$ são tais que cada um domina ao outro, então
$H^*(T,F)$ e $H^*(T',F)$ são canonicamente isomorfos.

Suponhamos que em $\mathcal C$ existem os produtos fibrados. Para
$S\in\mathcal C$ fixo, pode-se aplicar o anterior à categoria
$\mathcal C_S$ dos objetos de $\mathcal C$ sobre $S$; escreveremos
$C^*(T/S,F)$ e $H^*(T/S,F)$ em lugar de $C^*(T,F)$ e $H^*(T,F)$ quando
quisermos precisar que trabalhamos na categoria $\mathcal C_S$; assim,

\src{311}
$C^*(T/S,F)$ é um complexo de cocadeias em $\mathcal C'$ que, em dimensão
$n$, é igual a
\[
  F(T\times_ST\times_S\cdots\times_ST)
\]
(onde o parêntese contém $n+1$ fatores).

Observe-se que, como de costume, pode-se definir $H^0(T/S,F)$ sem supor
abeliana a categoria $\mathcal C'$: é o núcleo (definição 2.1), se existe,
do par de morfismos $F(p_i)$ ($i=1,2$)
\[
  F(T)\rightrightarrows F(T\times_ST)
\]
correspondentes às duas projeções
$p_1,p_2:T\times_ST\rightrightarrows T$. Em particular, ter-se-á um morfismo
natural (chamado de \emph{aumento})
\[
  F(S)\longrightarrow H^0(T/S,F)
\]
que será um isomorfismo nos casos favoráveis (em particular, se
$T\longrightarrow S$ é um epimorfismo estrito e $F$ é «exato pela
esquerda»). Do mesmo modo, quando $F$ toma valores na categoria dos
grupos numa categoria $\mathcal C''$, também se pode definir $H^1(T/S,F)$;
quando $\mathcal C''$ é a categoria dos conjuntos (isto é, $F$ toma
valores na categoria dos grupos ordinários, não necessariamente
comutativos), $H^1(T/S,F)$ é o quociente do subgrupo $Z^1(T/S,F)$ de
$C^1(T/S,F)=F(T\times_ST)$ formado pelos $g$ tais que
\[
  F(p_{31})(g)=F(p_{32})(g)\,F(p_{21})(g)
\]
pelo grupo de operadores $F(T)$, que atua sobre $C^1(T/S,F)$ e, em
particular, sobre o subconjunto $Z^1(T/S,F)$ mediante
\[
  \rho(g')\cdot g=F(p_2)(g')\,g\,F(p_1)(g')^{-1}.
\]

\medskip\noindent
b. Seja, por exemplo, $\mathcal F$ uma categoria fibrada de base $\mathcal C$.
Sejam $\xi,\eta\in\mathcal F_S$ e, para todo $S'$ sobre $S$, seja
\[
  F_{\xi,\eta}(S')=
  \operatorname{Hom}(\xi\times_SS',\eta\times_SS').
\]
Assim, $F_{\xi,\eta}$ é um funtor contravariante de $\mathcal C_S$ para a
categoria dos conjuntos. Dito isto, \emph{afirmar que o morfismo de
aumento}
\[
  F_{\xi,\eta}(S)\longrightarrow H^0(S'/S,F_{\xi,\eta})
\]
\emph{é um isomorfismo para todo par de elementos
$\xi,\eta\in\mathcal F_S$ significa que
$\alpha:S'\longrightarrow S$ é um morfismo de $\mathcal F$-descida}
(definição 1.7).

\medskip\noindent
c. Ponhamos igualmente, para $\xi\in\mathcal F_S$ e todo objeto $S'$ de
$\mathcal C$ sobre $S$,
\[
  G_\xi(S')=\operatorname{Aut}(\xi\times_SS') ;
\]
definiu-se assim um funtor contravariante $G_\xi$ de $\mathcal C_S$ para a
categoria dos

\src{312}
grupos. Dito isto, constata-se que $Z^1(S'/S,G_\xi)$ \emph{se identifica
canonicamente com o conjunto dos dados de descida sobre
$\xi'=\xi\times_SS'$ relativos a $S'\longrightarrow S$} (definição 1.6), e
$H^1(S'/S,G_\xi)$ \emph{se identifica com o conjunto das classes (salvo
isomorfismo) de objetos de $\mathcal F_{S'}$ providos de um dado de descida
relativo a $\alpha:S'\longrightarrow S$ que, como objetos de
$\mathcal F_{S'}$, são isomorfos a $\xi'=\xi\times_SS'$}. Portanto, se
$\alpha:S'\longrightarrow S$ é um morfismo de $\mathcal F$-descida
(cf. (b)), então $H^1(S'/S,G_\xi)$ \emph{contém como subconjunto o
conjunto das classes (salvo isomorfismo) de objetos $\eta$ de
$\mathcal F_S$ tais que $\eta\times_SS'$ seja isomorfo em
$\mathcal F_{S'}$ a $\xi\times_SS'$; e essa inclusão é uma identidade se e
somente se todo dado de descida sobre $\xi'=\xi\times_SS'$ relativo a
$\alpha:S'\longrightarrow S$ é efetivo}. (Esse será o caso, em particular,
se $\alpha:S'\longrightarrow S$ é um morfismo de $\mathcal F$-descida
estrito.)

\result{OBSERVAÇÃO}{} Os complexos de cocadeias do tipo $C^*(T/S,F)$
contêm como casos particulares a maioria dos complexos padrão
conhecidos (cohomologia de Čech, cohomologia de grupos, etc.) e desempenham um
papel importante em geometria algébrica, em particular na «cohomologia de
Weil» dos pré-esquemas.

\medskip\noindent
d. \result{EXEMPLO}{1} Seja $S'$ um objeto sobre $S\in\mathcal C$, e seja
$\Gamma$ um grupo de automorfismos de $S'$ tal que $S'$ seja «formalmente
principal sobre $S$, de grupo $\Gamma$», isto é, tal que o morfismo natural
\[
  \Gamma\times S'\longrightarrow S'\times_SS',
\]
onde $\Gamma\times S'$ denota a soma direta de $\Gamma$ cópias de $S'$, seja
um isomorfismo. (Supõe-se que em $\mathcal C$ existem as somas diretas que
intervêm aqui.) Seja $F$ um funtor contravariante de $\mathcal C$ para a
categoria dos grupos abelianos. Então $C^*(S'/S,F)$ é canonicamente
isomorfo ao grupo simplicial padrão de cocadeias homogêneas
$C^*(\Gamma,F(S'))$ e, portanto, $H^*(S'/S,F)$ é canonicamente isomorfo a
$H^*(\Gamma,F(S'))$.

\medskip\noindent
e. \result{EXEMPLO}{2} Seja $\mathcal C$ a categoria dos pré-esquemas. Denota-se
por $G_a$ («grupo aditivo») o funtor contravariante de $\mathcal C$ para a
categoria dos grupos abelianos definido por
\[
  G_a(X)=H^0(X,\mathcal O_X).
\]
Define-se de igual modo o funtor $G_m$ («grupo multiplicativo») por meio de
\[
  G_m(X)=H^0(X,\mathcal O_X)^*
\]
(grupo dos elementos invertíveis do anel $H^0(X,\mathcal O_X)$), e, mais
geralmente, o funtor $\operatorname{Gl}(n)$ («grupo linear de ordem $n$»)
mediante

\src{313}
\[
  \operatorname{Gl}(n)(X)=
  \operatorname{Gl}\bigl(n,H^0(X,\mathcal O_X)\bigr),
\]
que é um funtor de $\mathcal C$ à categoria dos grupos (não
necessariamente comutativos se $n>1$; para $n=1$ se recupera $G_m$). Além disso,
pode-se interpretar $\operatorname{Gl}(n)$ como um funtor de automorfismos
(cf. (c)) considerando a categoria fibrada $\mathcal F$ de base $\mathcal C$
tal que, para $X\in\mathcal C$, $\mathcal F_X$ seja a categoria dos feixes
localmente livres sobre $X$: com efeito, tem-se
\[
  \operatorname{Gl}(n)(X)=
  \operatorname{Aut}_{\mathcal F_X}(\mathcal O_X^n).
\]
De (b) segue-se que, se $\alpha:S'\longrightarrow S$ é um morfismo de
$\mathcal F$-descida (cf. parágrafo 2.c),
$H^1(S'/S,\operatorname{Gl}(n))$ contém o conjunto das classes (salvo
isomorfismo) de feixes localmente livres sobre $S$ cuja imagem inversa sobre
$S'$ é isomorfa a $\mathcal O_{S'}^n$, e essa inclusão é uma igualdade se e
somente se todo dado de descida sobre $\mathcal O_{S'}^n$ (relativo a
$\alpha:S'\longrightarrow S$) é efetivo. Quando $S$ é o espectro de um
anel local, isto significa, portanto,
$H^1(S'/S,\operatorname{Gl}(n))=(e)$, pois todo feixe localmente livre sobre $S$
é então trivial.

Assinalemos a equivalência das condições seguintes para um morfismo
$\alpha:S'\longrightarrow S$:
\begin{enumerate}
\item[(i)] o homomorfismo de aumento
\[
  H^0(S,\mathcal O_S)=G_a(S)\longrightarrow H^0(S'/S,G_a)
\]
é um isomorfismo;
\item[(ii)] $S'\longrightarrow S$ é um morfismo de $\mathcal F$-descida
(sendo $\mathcal F$ a categoria fibrada de base $\mathcal C$ considerada
acima).
\end{enumerate}
Se $S'\longrightarrow S$ é finito, essas condições equivalem também a:
\begin{enumerate}
\item[(iii)] $S'\longrightarrow S$ é um epimorfismo estrito
(cf. parágrafo 2.c).
\end{enumerate}

Suponhamos agora que $S=\operatorname{Spec}(A)$,
$S'=\operatorname{Spec}(A')$; então
\[
  C^n(S'/S,G_a)=C^n(A'/A,G_a)=\bigotimes_A^{n+1}A',
\]
sendo o operador coborde
$C^n(A'/A,G_a)\longrightarrow C^{n+1}(A'/A,G_a)$ a soma alternada dos
operadores de face
\[
  \partial_i(x_0\otimes x_1\otimes\cdots\otimes x_n)
  =x_0\otimes\cdots\otimes x_{i-1}\otimes1_{A'}\otimes
   x_i\otimes\cdots\otimes x_n.
\]
Do mesmo modo,
\[
  C^n(S'/S,G_m)=C^n(A'/A,G_m)
  \simeq\left(\bigotimes_A^{n+1}A'\right)^*,
\]
e as operações simpliciais em $C^*(A'/A,G_m)$ são induzidas pelas de
$C^*(S'/S,G_a)$. Explicitam-se analogamente as operações simpliciais em
$C^*(A'/A,\operatorname{Gl}(n))$. \emph{Em todos os casos que o
conferenciante conhece, tem-se}
\[
  H^i(A'/A,G_a)=0\qquad\text{para $i>0$},
\]
\emph{e, se $A$ é local, tem-se}
\[
  H^1(A'/A,G_m)=0
  \quad\text{e, mais geralmente,}\quad
  H^1(A'/A,\operatorname{Gl}(n))=(e)
\]
(quando $S'\longrightarrow S$ é um morfismo de $\mathcal F$-descida,
isto é, quando o diagrama
$A\longrightarrow A'\rightrightarrows A'\otimes_AA'$ é exato; compare-se com
o parágrafo 2.c). Observe-se que \emph{o «teorema 90» de Hilbert nada mais é
que a relação}
\src{314}
\[
  H^1(S'/S,G_m)=0
\]
\emph{quando $A$ é um corpo e $A'$ uma extensão de Galois finita deste
(cf. exemplo 1), e também se pode exprimir dizendo que, no caso
considerado, $S'\longrightarrow S$ é um morfismo de descida estrito para a
categoria fibrada dos feixes localmente livres de posto 1.} É sob esta
última forma que convém generalizar o teorema de Hilbert, variando tanto
as hipóteses sobre o morfismo $S'\longrightarrow S$ como as relativas aos
feixes quase-coerentes considerados.

Assinalemos, por fim, a equivalência das propriedades seguintes quando
$A$ é um anel local artiniano de ideal maximal $\mathfrak m$ e $A'$ uma
$A$-álgebra (denotando, para todo inteiro $k>0$, por $A_k$ (resp. $A'_k$) os
anéis $A/\mathfrak m^{k+1}$ (resp. $A'/\mathfrak m^{k+1}A'$)):
\begin{enumerate}
\item[(i)] $H^1(A'_k/A_k,G_a)=0$ para todo $k$.
\item[(ii)] $H^1(A'_k/A_k,G_m)=0$ para todo $k$.
\item[(iii)] $H^1(A'_k/A_k,\operatorname{Gl}(n))=(e)$ para todo $k$ e todo
$n$.
\end{enumerate}
Se $S'\longrightarrow S$ é um epimorfismo estrito, as condições
anteriores implicam inclusive que é um morfismo de descida estrito para os
módulos livres (de tipo finito ou não) sobre $A'$.

\result{OBSERVAÇÃO}{} A definição dos grupos $H^i(S'/S,G_m)$, quando
$S,S'$ são esquemas de corpos $A,A'$, se deve a Amitsur. O grupo
$H^2(S'/S,G_m)$ é particularmente interessante como variante «global» do
grupo de Brauer; para essa variante pode-se consultar [1], capítulo VII.

\fsection{B}{Descida por morfismos fielmente planos}

\subsection*{1. Enunciado dos teoremas de descida}
\addcontentsline{toc}{subsection}{1. Enunciado dos teoremas de descida}

\result{DEFINIÇÃO}{1.1} Um morfismo
$\alpha:S'\longrightarrow S$ de pré-esquemas diz-se \emph{plano} se, para todo
$x'\in S'$, $\mathcal O_{x'}$ é um módulo plano sobre o anel
$\mathcal O_{\alpha(x')}$ (isto é, se
$\mathcal O_{x'}\otimes_{\mathcal O_{\alpha(x')}}M$ é um funtor exato no
$\mathcal O_{\alpha(x')}$-módulo $M$). Um morfismo se diz
\emph{fielmente plano} se é plano e sobrejetivo.

Por exemplo, se $S=\operatorname{Spec}(A)$ e
$S'=\operatorname{Spec}(A')$, então $S'$ é plano sobre $S$ se e somente se
$A'$ é um $A$-módulo plano, e $S'$ é fielmente plano sobre $S$ se e somente se
$A'$ é um $A$-módulo fielmente plano (isto é, o funtor
$A'\otimes_A M$ no $A$-módulo $M$ é exato e fiel); isto significa também,
na terminologia de Serre [5], que o par $(A,A')$ é plano. Se $S'$ é
fielmente plano sobre $S$, então o funtor imagem inversa de feixes
quase-coerentes sobre $S$ é exato e fiel; em outras palavras, para que uma
sucessão de homomorfismos de feixes

\src{315}
quase-coerentes sobre $S$ seja exata, é necessário e suficiente que também o seja a sua
imagem inversa sobre $S'$ (em particular, para que um homomorfismo de feixes
quase-coerentes sobre $S$ seja um monomorfismo, resp. um epimorfismo, resp. um
isomorfismo, é necessário e suficiente que também o seja a sua imagem inversa sobre
$S'$). Essa propriedade continua válida se nos restringirmos a um aberto
qualquer de $S$ e, sob essa forma, caracteriza os morfismos fielmente
planos.

\result{DEFINIÇÃO}{1.2} Um morfismo $\alpha:S'\longrightarrow S$ diz-se
\emph{quase-compacto} se a imagem inversa de toda parte aberta quase-compacta
$U$ de $S$ é quase-compacta (isto é, união finita de abertos afins).

Evidentemente, basta verificar esta propriedade para os abertos afins de $S$.
Por exemplo, um morfismo afim (isto é, tal que a imagem inversa de um aberto
afim seja afim) é quase-compacto.

A classe dos morfismos planos, resp. fielmente planos, resp. quase-compactos,
é estável por composição e por «mudança de base», e contém, naturalmente,
os isomorfismos.

\result{TEOREMA}{1} Seja $\alpha:S'\longrightarrow S$ um morfismo de
pré-esquemas, fielmente plano e quase-compacto. Então $\alpha$ é um
\emph{morfismo de descida estrito} (A, definição 1.7) para a categoria
fibrada $\mathcal F$ dos feixes quase-coerentes (A, parágrafo 1, exemplo 2).

Esse enunciado significa duas coisas:
\begin{enumerate}
\item[(i)] Se $F$ e $G$ são dois feixes quase-coerentes sobre $S$, e $F'$ e
$G'$ suas imagens inversas sobre $S'$, então o homomorfismo natural
\[
  \operatorname{Hom}(F,G)\longrightarrow\operatorname{Hom}(F',G')
\]
é uma bijeção do primeiro membro sobre o subgrupo do segundo formado pelos
homomorfismos $F'\longrightarrow G'$ compatíveis com os dados de descida
canônicos sobre esses feixes, isto é, aqueles cujas imagens inversas pelas
duas projeções de $S''=S'\times_SS'$ sobre $S'$ dão um mesmo
homomorfismo $F''\longrightarrow G''$.
\item[(ii)] Todo feixe quase-coerente $F'$ sobre $S'$, provido de um dado de
descida relativo ao morfismo $\alpha:S'\longrightarrow S$
(A, definição 1.6), é isomorfo (provido deste dado) à imagem inversa de
um feixe quase-coerente $F$ sobre $S$.
\end{enumerate}
Tomando $F=\mathcal O_S$ em (i), obtém-se:

\result{COROLÁRIO}{1} Seja $G$ um feixe quase-coerente sobre $S$; sejam
$G'$ e $G''$ suas imagens inversas sobre $S'$ e sobre
$S''=S'\times_SS'$, e sejam $p_1,p_2$ as duas projeções de $S''$ sobre
$S'$. Então o seguinte diagrama de aplicações de conjuntos

\src{316}
\[
  \Gamma(G)\xrightarrow{\alpha^*}\Gamma(G')
  \underset{p_2^*}{\overset{p_1^*}{\rightrightarrows}}\Gamma(G'')
\]
é exato (A, definição 1.(a)).

Por outro lado, a conjunção de (i), (ii) e da definição 1.1 dá o

\result{COROLÁRIO}{2} Seja $G$ como no corolário 1. Então existe uma
correspondência biunívoca entre os subfeixes quase-coerentes de $G$ e os
subfeixes quase-coerentes de $G'$ cujas imagens inversas sobre $S''$, pelas
duas projeções $p_1,p_2$, dão o mesmo subfeixe de $G''$.

Naturalmente, tem-se um enunciado equivalente em termos de feixes quociente.
Como se viu (A, parágrafo 4, (e)), o teorema 1 deve ser considerado uma
generalização do «teorema 90» de Hilbert, e implica como casos particulares
diversas formulações em termos de 1-cohomologia. Para a demonstração, reduz-se
facilmente ao caso em que $S=\operatorname{Spec}(A)$,
$S'=\operatorname{Spec}(A')$, e para (i) reduz-se facilmente a demonstrar o
corolário 1, isto é, a exatidão do diagrama
\[
  M=A\otimes_A M\longrightarrow A'\otimes_A M
  \rightrightarrows A'\otimes_AA'\otimes_A M
\]
para todo $A$-módulo $M$, o que se deduz do lema mais geral:

\result{LEMA}{1.1} Seja $A'$ uma $A$-álgebra fielmente plana. Então, para
todo $A$-módulo $M$, o complexo aumentado por $M$
$C^*(A'/A,G_a)\otimes_A M$ (cf. A, parágrafo 4, (e)) é uma
\emph{resolução} de $M$.

Basta demonstrar que o complexo aumentado deduzido do anterior mediante a
mudança de base de $A$ para $A'$ satisfaz as mesmas conclusões. Isso conduz a
verificar o enunciado quando se substituem $A$ por $A'$ e $A'$ por
$A'\otimes_AA'$, o que nos reduz a um caso no qual existe um
homomorfismo de $A$-álgebras $A'\longrightarrow A$ (ou, em termos
geométricos, ao caso em que $S'$ sobre $S$ admite uma seção). Nesse caso,
isso se deduz das generalidades de A, parágrafo 4, (a). Assinalemos de passagem
o seguinte corolário, que generaliza um enunciado bem conhecido em
cohomologia de Galois (compare-se com A, parágrafo 4, (e)):

\result{COROLÁRIO}{} Se $A'$ é fielmente plana sobre $A$, tem-se
$H^0(A'/A,G_a)=A$ e
\[
  H^i(A'/A,G_a)=0\qquad\text{para $i\geq1$}.
\]

Para demonstrar a parte (ii) do teorema 1, procede-se como para (i),
reduzindo-se ao caso em que $S'$ sobre $S$ admite uma seção, onde se deduz
de (i) (cf. A, parágrafo 1, (c)).

Evidentemente, pode-se variar ad libitum o teorema 1 e seus corolários,
introduzindo diversas estruturas adicionais nos feixes (ou sistemas de
feixes) quase-coerentes considerados. Por exemplo, dar sobre $S$ um feixe
\src{317}
quase-coerente de álgebras comutativas «equivale» a dar sobre $S'$ um feixe desse
tipo provido de um dado de descida relativo a
$\alpha:S'\longrightarrow S$. Tendo em conta a correspondência funtorial
entre tais feixes quase-coerentes sobre $S$ e os pré-esquemas afins sobre
$S$, obtém-se a segunda afirmação do teorema seguinte:

\result{TEOREMA}{2} Seja $\alpha:S'\longrightarrow S$ como no teorema 1.
Então $\alpha$ é um \emph{morfismo de descida} (não estrito em geral)
(A, definição 2.4), e é um \emph{morfismo de descida estrito} para a
categoria fibrada dos esquemas afins sobre pré-esquemas
(A, definição 1.7).

A primeira afirmação do teorema significa o seguinte: sejam $X,Y$ dois
pré-esquemas sobre $S$; sejam $X',Y'$ suas imagens inversas sobre $S'$, e
$X'',Y''$ suas imagens inversas sobre $S''=S'\times_SS'$. Então o seguinte
diagrama de aplicações naturais
\[
  \operatorname{Hom}_S(X,Y)\xrightarrow{\alpha^*}
  \operatorname{Hom}_{S'}(X',Y')
  \underset{p_2^*}{\overset{p_1^*}{\rightrightarrows}}
  \operatorname{Hom}_{S''}(X'',Y'')
\]
é exato; isto é, $\alpha^*$ é uma bijeção de
$\operatorname{Hom}_S(X,Y)$ sobre a parte de
$\operatorname{Hom}_{S'}(X',Y')$ formada pelos homomorfismos compatíveis com
os dados de descida canônicos sobre $X',Y'$ (isto é, aqueles cujas
imagens inversas pelas duas projeções de $S''$ sobre $S'$ são iguais).
Isso se deduz facilmente do teorema 1, corolário 1, quando se restringe
a $Y$ afim sobre $S$; no caso geral, é preciso combinar o teorema 1 com o
resultado seguinte:

\result{LEMA}{1.2} Seja $\alpha:S'\longrightarrow S$ um morfismo fielmente plano
e quase-compacto. Então $S$ identifica-se com um \emph{espaço topológico
cociente} de $S'$; isto é, toda parte $U$ de $S$ tal que
$\alpha^{-1}(U)$ seja aberta é aberta.

Para completar o teorema 2, é preciso dar critérios de efetividade para um dado
de descida sobre um $S'$-pré-esquema $X'$ (quando não se supõe que $X'$ seja
afim sobre $S'$). Assinalemos primeiro que \emph{tal dado de descida não é
necessariamente efetivo}, mesmo se $S$ é o espectro de um corpo $k$, $S'$
o espectro de uma extensão quadrática $k'$ deste e $X'$ um esquema
algébrico próprio de dimensão 2 sobre $S'$ (como se pode ver, segundo Serre,
utilizando a superfície não projetiva de Nagata). \emph{Para que um dado de
descida sobre $X'/S'$ relativo a $\alpha:S'\longrightarrow S$ (fielmente
plano e quase-compacto) seja efetivo, é necessário e suficiente que $X'$ seja
união de abertos $X'_i$, afins sobre $S'$, que sejam «estáveis» pelo dado de
descida sobre $X'$.} Isso certamente se cumpre (quaisquer que sejam $X'/S'$
e o dado de descida sobre $X'$) se o morfismo
$\alpha:S'\longrightarrow S$ é \emph{radicial} (isto é, injetivo e com
extensões residuais que

\src{318}
são radiciais). Pode-se demonstrar também que continua a cumprir-se se
$\alpha:S'\longrightarrow S$ é \emph{finito} e toda parte finita de $X'$,
contida numa fibra de $X'$ sobre $S$, está contida num aberto de $X'$
afim sobre $S$ (esse é o \emph{critério de Weil}). Em particular, isso se
cumpre se $X'/S'$ é quase-projetivo, e nesse caso pode-se demonstrar que o
pré-esquema «descido» $X/S$ também é quase-projetivo (e projetivo se
$X'/S'$ também o é). Em resumo:

\result{TEOREMA}{3} Seja $\alpha:S'\longrightarrow S$ um morfismo de
pré-esquemas, fielmente plano e quase-compacto. Se $\alpha$ é \emph{radicial},
é um \emph{morfismo de descida estrito}. Se $\alpha$ é finito, é um
morfismo de descida estrito com respeito à categoria fibrada dos
pré-esquemas quase-projetivos (ou projetivos) sobre pré-esquemas.

\result{OBSERVAÇÕES}{} Ignoro se, na segunda afirmação anterior, a
hipótese de que $\alpha$ seja um morfismo \emph{finito} é realmente necessária;
em todo caso, verifica-se de maneira «formal» que pode ser substituída pela
seguinte hipótese, aparentemente mais fraca: \emph{para todo ponto de $S$
existe uma vizinhança aberta $U$, um $U'$ finito e fielmente plano sobre $U$, e um
$S$-morfismo de $U'$ em $S'$.} Um caso típico que não entra no anterior é
aquele em que $S=\operatorname{Spec}(A)$,
$S'=\operatorname{Spec}(\widehat A)$, onde $A$ é um anel local noetheriano
e $\widehat A$ seu completamento; ou também aquele em que $S'$ é quase-finito sobre
$S$ (isto é, localmente isomorfo a um aberto de um $S$-esquema finito) e não
finito. Nesses dois casos, o conferenciante também ignora a resposta à
pergunta seguinte: seja $X$ um $S$-esquema tal que $X'=X\times_SS'$ seja
projetivo sobre $S'$; é verdade que $X$ é projetivo sobre $S$?

\subsection*{2. Aplicação à descida de certas propriedades dos morfismos}
\addcontentsline{toc}{subsection}{2. Descida de propriedades dos morfismos}

Seja $P$ uma classe de morfismos de pré-esquemas. Seja
$\alpha:S'\longrightarrow S$ um morfismo de pré-esquemas, seja
$f:X\longrightarrow Y$ um morfismo de $S$-pré-esquemas e seja
$f':X'\longrightarrow Y'$ a imagem inversa de $f$ por $\alpha$. Cabe
perguntar-se então se a relação «$f'\in P$» implica «$f\in P$». A resposta é
afirmativa em muitos casos importantes quando se supõe que $\alpha$ é
\emph{fielmente plano e quase-compacto} (sendo às vezes supérflua essa última
hipótese). Isso se vê diretamente e sem dificuldade se $P$ é a classe dos
morfismos sobrejetivos, resp. radiciais (esses casos se deduzem da
sobrejetividade de $\alpha$), resp. planos, resp. fielmente planos, resp. lisos
(esses casos se deduzem da planicidade fiel de $\alpha$), resp. de tipo finito.
Utilizando os teoremas 1 e 2 e o lema 1.2, vê-se também que ocorre o mesmo se
$P$ é uma das classes seguintes: isomorfismos, imersões abertas, imersões
fechadas, imersões (se $f$ é de tipo finito e $Y$
localmente noetheriano), morfismos afins, morfismos finitos, morfismos
quase-finitos, morfismos abertos, morfismos fechados, homeomorfismos, morfismos
separados,

\src{319}
morfismos próprios. O único caso importante não elucidado é o dos
morfismos projetivos e quase-projetivos, já assinalado na observação do
parágrafo 1.

\subsection*{3. Descida por morfismos finitos fielmente planos}
\addcontentsline{toc}{subsection}{3. Descida por morfismos finitos fielmente planos}

Seja $\alpha:S'\longrightarrow S$ um morfismo finito correspondente a um feixe de
álgebras $\mathcal A'$ sobre $S$ que, como feixe de módulos, é localmente livre
de tipo finito e é $\ne0$ em toda parte; então $\alpha$ é um morfismo
fielmente plano e quase-compacto, ao qual se podem aplicar os resultados
anteriores. Dar um feixe quase-coerente $F'$ sobre $S'$ equivale a dar o feixe
quase-coerente $\alpha_*(F')$ sobre $S$, provido de sua estrutura de
$\mathcal A'$-módulo (observando que
$\mathcal A'=\alpha_*(\mathcal O_{S'})$). Para simplificar, este feixe sobre $S$
também será denotado por $F'$. As duas imagens inversas $p_i^*(F')$ de $F'$
sobre $S'\times_SS'$ correspondem igualmente aos feixes quase-coerentes de
$(\mathcal A'\otimes_{\mathcal O_S}\mathcal A')$-módulos
\[
  F'\otimes_{\mathcal O_S}\mathcal A'
  \quad\text{e}\quad
  \mathcal A'\otimes_{\mathcal O_S}F'.
\]
Dar um $(S'\times_SS')$-homomorfismo do primeiro ao segundo equivale a dar um
homomorfismo de $(\mathcal A'\otimes_{\mathcal O_S}\mathcal A')$-módulos e,
tendo em conta que $\mathcal A'$ é localmente livre, isto equivale a dar
um homomorfismo de $(\mathcal A'\otimes\mathcal A')$-módulos:
\[
  \mathcal U=
  \operatorname{Hom}_{\mathcal O_S}(\mathcal A',\mathcal A')
  =\mathcal A'\otimes\check{\mathcal A}'
  \longrightarrow
  \operatorname{Hom}_{\mathcal O_S}(F',F'),
\]
isto é, a dar, para cada seção $\xi$ de $\mathcal U$ sobre um aberto $V$,
um homomorfismo de $\mathcal O_S$-módulos
$T_\xi:F'|_V\longrightarrow F'|_V$ que satisfaz as condições
\begin{equation}
  T_{f\xi}(x)=fT_\xi(x),
  \qquad T_{\xi f}(x)=T_\xi(fx),
  \tag{3.1}
\end{equation}
onde $f,x$ são, respectivamente, seções de $\mathcal A',F'$ sobre um
aberto de $S$ contido em $V$. As condições (i) e (ii) de um dado de
descida (A, parágrafo 1, (c)) escrevem-se então, respectivamente:
\begin{equation}
  T_{1_{\mathcal U}}x=x,
  \qquad\text{isto é,}\qquad
  T_{1_{\mathcal U}}=\operatorname{id}_{F'},
  \tag{3.2}
\end{equation}
\begin{equation}
  T_{\xi\eta}=T_\xi T_\eta.
  \tag{3.3}
\end{equation}
Em outras palavras, \emph{um dado de descida sobre $F'$ equivale a uma
representação do feixe de $\mathcal O_S$-álgebras
$\mathcal U=\operatorname{Hom}_{\mathcal O_S}(\mathcal A',\mathcal A')$ no
feixe de $\mathcal O_S$-álgebras
$\operatorname{Hom}_{\mathcal O_S}(F',F')$ que satisfaz as duas condições de
linearidade (3.1).} Se se tem um emparelhamento de feixes quase-coerentes sobre
$S'$:
\[
  F'_1\times F'_2\longrightarrow F'_3
\]
(que pode ser interpretado como um emparelhamento de feixes de
$\mathcal A'$-módulos sobre $S$),

\src{320}
e dados de colagem sobre os $F'_i$, definidos por homomorfismos
\[
  T_i:\mathcal U\longrightarrow
  \operatorname{Hom}_{\mathcal O_S}(F'_i,F'_i)
  \qquad (i=1,2,3),
\]
então esses dados são compatíveis com o emparelhamento dado, no sentido
evidente do termo, se e somente se vigora a condição seguinte.

Para toda seção $\xi$ de $\mathcal U$ sobre um aberto, denotando por
\[
  \Delta\xi=\sum_i\xi'_i\otimes_{\mathcal A'}\xi''_i
\]
a seção de $\mathcal U\otimes_{\mathcal A'}\mathcal U$ (considerando
$\mathcal U$ como $\mathcal A'$-módulo mediante sua estrutura esquerda)
definida pela fórmula
\[
  \xi\mathbin{\cdot}(fg)=\sum_i\xi'_i(f)\xi''_i(g)
\]
(onde $f$ e $g$ são duas seções de $\mathcal A'$ sobre um aberto menor),
tem-se a fórmula
\begin{equation}
  T^{(3)}_\xi(x\mathbin{\cdot}y)
  =\sum_iT^{(1)}_{\xi'_i}x\mathbin{\cdot}T^{(2)}_{\xi''_i}y
  \tag{3.4}
\end{equation}
para todo par de seções $x,y$ de $\mathcal A'$ sobre um aberto menor.
(Essa propriedade pode ser expressa
dizendo que os homomorfismos $T^{(i)}$ são compatíveis com a aplicação
diagonal de $\mathcal U$ com respeito ao emparelhamento dado.) Em particular, as
fórmulas (3.1) a (3.4) permitem interpretar, em termos de representações de
álgebras com aplicações diagonais, os dados de descida sobre um feixe
quase-coerente de álgebras sobre $S'$ e, portanto, também (restringindo-se às
álgebras comutativas) os dados de descida sobre um $S'$-esquema afim.

Daí se passa a uma interpretação análoga dos dados de descida sobre um
$S'$-pré-esquema $X'$ qualquer: dar tal $X'$ equivale a dar um pré-esquema
$X'$ sobre $S$, provido de um homomorfismo de $\mathcal O_S$-álgebras
\[
  \mathcal A'\longrightarrow\mathcal O_{X'},
\]
e um dado de descida sobre $X'$ equivale a dar um homomorfismo de feixes
(compatível com o morfismo $h:X'\longrightarrow S'$)
\[
  \mathcal U\longrightarrow
  \operatorname{Hom}_{h^{-1}(\mathcal O_S)}
  (\mathcal O_{X'},\mathcal O_{X'}),
\]
que satisfaz condições análogas às condições (3.1) a (3.4) anteriores.

\textsc{Exemplo 1 («Weil»).} -- Suponhamos que $S'/S$ é um recobrimento étale
de Galois de grupo de Galois $\Gamma$ (cf. A, parágrafos 3 e 4.d). Então, um
dado de descida sobre um feixe quase-coerente $F'$ sobre $S'$ (resp. sobre um
$S'$-pré-esquema $X'$) equivale a dar uma representação de $\Gamma$ por
automorfismos de $(S',F')$ (resp. de $(S',X')$) compatível com as operações
de $\Gamma$ sobre $S'$. Este resultado

\src{321}
é «formal», isto é, demonstra-se em termos de categorias, mas, do
ponto de vista deste número, deduz-se também da estrutura explícita de
$\mathcal U$ (provido de sua estrutura de anel, do homomorfismo de anéis
$\mathcal A'\longrightarrow\mathcal U$ e da aplicação diagonal), conhecida
por completo graças ao resultado seguinte: $\mathcal U$ admite, como
$\mathcal A'$-módulo esquerdo, uma base formada pelas seções de
$\mathcal U$ que correspondem aos elementos de $\Gamma$.

\textsc{Exemplo 2 («Cartier»).} -- Seja $p$ um número primo; suponhamos que
$p\mathcal O_S=0$ (isto é, $\mathcal O_S$ é de característica $p$),
$(\mathcal A')^p\subset\mathcal O_S=\mathcal A$ (isto é, $S'/S$ é radicial
de altura 1) e que o feixe de álgebras $\mathcal A'$ sobre $\mathcal A$ admite
localmente uma $p$-base, isto é, uma família $(x_i)$ de seções tal que
$\mathcal A'$ está gerada como álgebra pelos $x_i$, submetidos às únicas
condições $x_i^p=0$. Supomos que o conjunto dos $i$ é finito, de
cardinal $n$. Seja $\mathfrak g$ o feixe das $\mathcal A$-derivações de
$\mathcal A'$; é um feixe de $\mathcal A'$-módulos localmente livre de posto
$n$ e, além disso, um feixe de $p$-álgebras de Lie sobre $\mathcal A$ (mas não sobre
$\mathcal A'$) que satisfaz a condição
\begin{equation}
  [X,fY]=X(f)Y+f[X,Y].
  \tag{3.5}
\end{equation}

\textsc{Lema.} -- $\mathcal U=
\operatorname{Hom}_{\mathcal O_S}(\mathcal A',\mathcal A')$ está gerada,
como $\mathcal O_S$-álgebra provida de um homomorfismo de álgebras
$\mathcal A'\longrightarrow\mathcal U$, pelo sub-$\mathcal A'$-módulo
esquerdo $\mathfrak g$, com as relações adicionais
\begin{equation}
  \left\{
  \begin{aligned}
    Xf-fX&=X(f),\\
    XY-YX&=[X,Y],\\
    X^p&=X^{(p)}.
  \end{aligned}
  \right.
  \tag{3.6}
\end{equation}

Do lema anterior se deduz que um dado de descida sobre o feixe
quase-coerente $F'$ sobre $S'$ equivale a dar, para todo $X\in\mathfrak g$, um
$\mathcal O_S$-endomorfismo $\overline X$ de $F'$ que satisfaz as condições
\begin{align}
  \overline{fX}&=f\overline X, \tag{3.7}\\
  \overline X(fx)&=X(f)x+f\overline X(x), \tag{3.8}\\
  \overline{[X,Y]}&=[\overline X,\overline Y], \tag{3.9}\\
  \overline{X^{(p)}}&=\overline X^p. \tag{3.10}
\end{align}
(Isso é o que se poderia chamar uma conexão linear sobre $F'$, sem curvatura
e compatível com a potência $p$-ésima.) Explicita-se de igual modo a noção
de dado de descida sobre um $S'$-pré-esquema $X'$; a relação (3.4) é
substituída aqui pela condição de que os $\overline X$ sejam derivações de
$\mathcal O_{X'}$. Como o morfismo $S'\longrightarrow S$ é radicial, o
teorema 3 garante que todo dado de descida desse tipo é

\src{322}
efetivo e, portanto, define um $S$-pré-esquema $X$.

Observe-se que não foi necessário impor nenhuma hipótese de planicidade, de
não singularidade nem de finitude de tipo algum sobre $F'$, resp. $X'$.

\subsection*{4. Aplicação a critérios de racionalidade}
\addcontentsline{toc}{subsection}{4. Aplicação a critérios de racionalidade}

Seja $X$ um $S$-pré-esquema tal que a imagem direta de $\mathcal O_X$ sobre
$S$ seja $\mathcal O_S$; essa propriedade continuará válida após toda mudança
plana $S'\longrightarrow S$ da base $S$. Se $F$ é um feixe invertível
(isto é, localmente livre de posto 1) sobre $X$, os automorfismos de $F$, que se
identificam com as seções invertíveis de $\mathcal O_X$, correspondem
biunívocamente às seções invertíveis de $\mathcal O_S$. Seja então $s$
uma seção de $X$ sobre $S$; chamaremos, de maneira gráfica, seção de $F$
sobre $s$ a uma seção do feixe invertível $s^*(F)$ sobre $S$. Do anterior
deduz-se que, se $F_i$ ($i=1,2$) são dois feixes invertíveis sobre $X$, cada um
provido de uma seção sobre $s$, e se $F_1$ e $F_2$ são isomorfos, existe um
único isomorfismo de $F_1$ sobre $F_2$ compatível com as seções em questão
(isto é, que transforma a primeira na segunda). Por outro lado, e
independentemente da seção $s$, convenhamos em considerar equivalentes
dois feixes invertíveis $F_1$ e $F_2$ sobre $X$ tais que todo ponto de $S$ tenha
uma vizinhança aberta $U$ para a qual as restrições de $F_1$ e $F_2$ a
$X|U$ sejam isomorfas. Então todo feixe invertível $F$ sobre $X$ é equivalente
a um feixe invertível $F_1$ provido de uma seção distinguida sobre $s$
(toma-se $F_1=F\,s^*(F)^{-1}$), e $F_1$ está determinado salvo isomorfismo. Em
outras palavras, a classificação dos feixes invertíveis sobre $X$ salvo
equivalência é a mesma que a classificação, salvo isomorfismo, dos feixes
invertíveis providos de uma seção distinguida.

Como essas propriedades continuam válidas após uma mudança plana
$\alpha:S'\longrightarrow S$ da base (substituindo a seção $s$ por sua
imagem inversa $s'$ por $\alpha$), se conclui, tendo em conta o teorema 1:

\emph{Seja o pré-esquema $X/S$ como acima e suponhamos que admita uma seção
$s$; seja $\alpha:S'\longrightarrow S$ um morfismo fielmente plano e
quase-compacto; seja $F'$ um feixe invertível sobre $X'=X\times_SS'$. Para que $F'$
seja equivalente à imagem inversa sobre $X'$ de um feixe invertível $F$ sobre
$X$, é necessário e suficiente que suas imagens inversas $p_1^*(F')$ e
$p_2^*(F')$ sobre $X'\times_XX'=X\times_S(S'\times_SS')$ sejam equivalentes. Se
for assim, $F$ está determinado salvo equivalência.} (Dir-se-á então que $F'$
é racional sobre $S$.)

Inspirando-se nesse princípio quando $\alpha:S'\longrightarrow S$ é como no
exemplo 1 ou no exemplo 2 do número anterior, obtêm-se os critérios de
racionalidade de Weil ou de Cartier. (Observe-se que esses autores se limitam ao
caso em que $S$ e $S'$

\src{323}
são espectros de corpos; a fortiori, $S$ é então o espectro de um anel
local, e a relação de equivalência introduzida acima nada mais é que a relação
de isomorfia.) No primeiro caso, $F'$ é racional sobre $S$ se e somente se seus
transformados por $\Gamma$ são equivalentes a $F'$. Para exprimir o critério
de racionalidade no segundo caso, consideram-se, em geral, o morfismo
diagonal
\[
  X'\longrightarrow X''=X'\times_XX'
\]
de $X'/X$, o feixe de ideais $\mathcal I$ correspondente sobre
$X'\times_XX'$ e o feixe $\mathcal I/\mathcal I^2$, que se identifica com sua
imagem inversa $\Omega^1_{X'/X}$ sobre $X'$. Como as restrições dos
$F''_i=p_i^*(F')$ ($i=1,2$) à diagonal são isomorfas (pois são isomorfas a
$F'$), isto é, $F''_1(F''_2)^{-1}=F''$ tem uma restrição à diagonal
que é trivial, segue-se que a restrição de $F''$ a
$(X'',\mathcal O_{X''}/\mathcal I^2)$ é dada, salvo isomorfismo, por um
elemento bem determinado $\xi$ de
\[
  H^1(X'',\mathcal I/\mathcal I^2)
  =H^1(X',\Omega^1_{X'/X}).
\]
Além disso, nesse caso tem-se
\[
  \Omega^1_{X'/X}
  =\Omega^1_{S'/S}\otimes_{\mathcal O_S}\mathcal O_X,
\]
e, portanto, se $\Omega^1_{S'/S}$ é localmente livre sobre $S$ (como no
caso de Cartier), $\xi$ define uma seção de
\[
  R^1f'_*(\mathcal O_{X'})\otimes\Omega^1_{S'/S}
\]
sobre $S'$ (chamada classe de Atiyah--Cartier do feixe invertível $F'$ sobre
$X'/S$), cuja anulação é necessária e suficiente para que sejam equivalentes
as imagens inversas de $F'$ pelas duas projeções de
\[
  (X'',\mathcal O_{X''}/\mathcal I^2)
  =X\times_S(S'',\mathcal O_{S''}/\mathcal J^2)
\]
sobre $X'$ (onde $\mathcal J$ é o feixe de ideais sobre
$S''=S'\times_SS'$ definido pelo morfismo diagonal
$S'\longrightarrow S'\times_SS'$). Essa anulação é, pois, trivialmente
necessária para que sejam equivalentes as imagens inversas de $F'$ sobre o
próprio $X''=X\times_SS''$ e, portanto, também para que $F'$ seja equivalente
à imagem inversa de um feixe invertível $F$ sobre $X$. Além disso, a classe de
Atiyah--Cartier pode ser interpretada também como a obstrução à existência,
localmente sobre $S'$, de uma conexão de $F'$ relativa às derivações de
$X'/X$; tal conexão fica, além disso, determinada quando se conhecem as
derivações de $F'$ correspondentes às extensões naturais a $X'$ das
derivações de $S'/S$. Disso e dos desenvolvimentos do número anterior
conclui-se facilmente que, no caso do exemplo 2 desse número e quando $X/S$
admite uma seção, a anulação da classe de Atiyah--Cartier é também
suficiente para que $F'$ seja racional sobre $S$.

\subsection*{5. Aplicação à restrição da base em um esquema abeliano}
\addcontentsline{toc}{subsection}{5. Restrição da base em um esquema abeliano}

Seja $S$ um pré-esquema. Chama-se esquema abeliano sobre $S$ um esquema $X$
liso e próprio

\src{324}
sobre $S$ cujas fibras nos pontos $x\in S$ são variedades abelianas sobre os
$\kappa(x)$. Suponhamos que $S$ seja noetheriano e regular
(isto é, que seus anéis locais sejam regulares); então, utilizando o
teorema de conexidade de Murre [4] (ao menos no caso de «características
iguais», no qual o teorema citado está atualmente demonstrado), pode-se
demonstrar que toda seção racional de $X$ sobre $S$ está definida em toda
parte (isto é, é uma seção), o que generaliza um teorema clássico de
Weil. Deduz-se, mais geralmente, que, se $X'$ é um esquema liso sobre $S$,
então toda $S$-aplicação racional de $X'$ em $X$ está definida em toda
parte. Disso se obtém o seguinte, que generaliza um resultado de
Igusa--Lang: Seja $S$ noetheriano e regular, e denote $K$ seu
anel de funções racionais (produto direto de corpos); seja $X$ um
esquema abeliano sobre $K$. Se $X$ é isomorfo a um $K$-esquema da forma
$X_0\times_S\operatorname{Spec}(K)$, onde $X_0$ é um esquema abeliano sobre
$S$, então $X_0$ está determinado salvo isomorfismo único.

Utilizando o resultado de unicidade anterior, vê-se que a questão da
restrição da base em $X$ é local sobre $S$ (e, portanto, que basta saber
efetuar a restrição aos $\operatorname{Spec}(\mathcal O_x)$, com
$x\in S$). Vê-se de igual maneira que, se
$\alpha:S'\longrightarrow S$ é um morfismo liso de tipo finito, se $K'$ é o
anel de funções racionais de $S'$, e se $X\otimes_KK'$ é da forma
$X'_0\times_{S'}\operatorname{Spec}(K')$, então $X'_0$ está provido de um
dado de descida canônico relativo a $\alpha$. Tendo em conta o teorema 3,
conclui-se:

\result{PROPOSIÇÃO}{5.1} Seja $S$ um pré-esquema noetheriano, regular e
irredutível, de corpo de funções racionais $K$; seja $K'$ uma extensão
finita de $K$, não ramificada sobre $S$; seja $S'$ a normalização de $S$ em
$K'$ (que é, portanto, um recobrimento étale de $S$); e seja $X$ um esquema
abeliano sobre $K$ tal que $X\otimes_KK'$ seja da forma
$X'_0\times_{S'}\operatorname{Spec}(K')$, onde $X'_0$ é um esquema abeliano
projetivo sobre $S'$. Então $X$ é da forma
$X_0\times_S\operatorname{Spec}(K)$, onde $X_0$ é um esquema abeliano
projetivo sobre $S$.

\result{OBSERVAÇÕES}{} O conferenciante ignora se pode ser substituída a
hipótese de que $S'\longrightarrow S$ seja um recobrimento étale sobrejetivo
(que permite utilizar o teorema 3) pela hipótese de que seja um morfismo
liso, sobrejetivo e de tipo finito (inclusive se se supõe que é um morfismo
étale), ou se a proposição continua válida sem supor que $X'_0$ seja
projetivo sobre $S'$ (condição que talvez se cumpra automaticamente).

\subsection*{6. Aplicação a critérios de trivialidade local e isotrivialidade}
\addcontentsline{toc}{subsection}{6. Critérios de trivialidade local e isotrivialidade}

Sejam $S$ um pré-esquema, $G$ um «pré-esquema em grupos» sobre $S$, e $P$ um
pré-esquema sobre $S$ sobre o qual «atua $G$» (pela direita). Diz-se que $P$
é formalmente principal homogêneo sob $G$ se o morfismo bem conhecido
\[
  G\times_SP\longrightarrow P\times_SP
\]

\src{325}
induzido pela ação de $G$ sobre $P$ é um isomorfismo. Doravante, suporemos que
$G$ seja plano sobre $S$ (e, portanto, fielmente plano sobre $S$), e reservaremos
o nome de fibrado principal homogêneo sob $G$ para um fibrado formalmente
principal homogêneo $P$ que seja fielmente plano e quase-compacto sobre $S$.
É imediato que isso equivale a dizer que se pode encontrar uma mudança
fielmente plana e quase-compacta
$S'\longrightarrow S$ da base $S$, tal que o fibrado formalmente principal
homogêneo
\[
  P'=P\times_SS'
\]
sob $G'=G\times_SS'$ seja trivial, isto é, isomorfo a $G'$ (isto é, que
admita uma seção); em particular, pode-se tomar $S'=P$. Observe-se, além disso,
que, se $S$ é localmente noetheriano, a hipótese de planicidade fiel sobre $P$
equivale à hipótese de que
\[
  \widehat P_s=P\times_S\operatorname{Spec}(\widehat{\mathcal O}_s)
\]
seja fielmente plano sobre $\widehat{\mathcal O}_s$ para todo $s\in S$ (onde
$\widehat{\mathcal O}_s$ denota o completamento do anel local
$\mathcal O_s$), como se deduz do fato de que $\widehat{\mathcal O}_s$ é
fielmente plano sobre $\mathcal O_s$. Além disso, se $P$ é de tipo finito sobre
$S$ localmente noetheriano, o conjunto dos pontos $s$ que satisfazem a
condição anterior é construtível; portanto, se $S$ é um «pré-esquema de
Jacobson» (por exemplo, um esquema de tipo finito sobre um corpo ou, mais
geralmente, um anel de Jacobson), basta verificar a
condição em questão nos pontos fechados de $S$. Isso nos reduz ao caso em
que a base é o espectro de um anel local completo $A$. Quando
$S=\operatorname{Spec}(A)$ ($A$ um anel local noetheriano completo) e $P$ é
de tipo finito sobre $S$, a planicidade fiel de $P/S$ equivale também à
existência de um $S'$ finito e plano sobre $S$ tal que $P'$ seja trivial; e, se
além disso $G$ é liso sobre $S$, pode-se supor que $S'$ seja étale sobre $S$.
Por conseguinte, se além disso o corpo residual de $A$ é algebricamente fechado
(«caso geométrico»), $P$ é fielmente plano sobre $A$ se e somente se é trivial.
Assim, pois, se $S$ é um pré-esquema algébrico sobre um corpo algebricamente
fechado e $G$ é liso e de tipo finito sobre $S$, vê-se que a condição de
planicidade fiel sobre $S$ equivale à condição de trivialidade analítica (SLF)
de Serre ([6], pp. 1--12).

Podem-se introduzir outros tipos, mais fortes, de condições sobre $P$ que
tenham a natureza de uma «trivialidade local». Diremos, em particular, que
$P$ é isotrivial (resp. estritamente isotrivial) se, para todo $s\in S$,
existem uma vizinhança aberta $U$ de $s$ e um morfismo finito e fielmente plano
(resp. um recobrimento étale sobrejetivo) $U'\longrightarrow U$ tais que
$P'=P\times_SU'$ seja trivial. (Nossa terminologia difere da de Serre [6], que
chama localmente isotrivial ao que nós chamamos estritamente isotrivial.) A
isotrivialidade estrita é especialmente útil se $G$ é liso sobre $S$, mas, em
contrapartida, é uma noção inadequada nos demais casos.

Se $G$ é afim sobre $S$, todo fibrado principal homogêneo $P$ sob $G$ é afim
pelo parágrafo 2, de onde, graças ao teorema 2, resulta a possibilidade de
«descer»

\src{326}
tais fibrados por morfismos fielmente planos e quase-compactos. Tomemos, em
particular, $G=\operatorname{Gl}(n)_S$, definido pela condição de que o
funtor de $S$-pré-esquemas
\[
  S'\longmapsto\operatorname{Hom}_S(S',G)
\]
(com valores na categoria dos grupos) se identifique com o funtor
\[
  \operatorname{Gl}(n)(S')
  =\operatorname{Gl}\bigl(n,H^0(S',\mathcal O_{S'})\bigr)
\]
de A, parágrafo 4.e. Utilizando os fatos (i) de que todo fibrado principal
homogêneo sob $G$ (resp. todo feixe localmente livre de posto $n$ sobre $S$) se
torna isomorfo ao objeto «trivial» $G$ (resp. $\mathcal O_S^n$) após uma mudança
fielmente plana e quase-compacta adequada de $S$; (ii) de que podem ser descidos
os objetos do tipo considerado (fibrados principais homogêneos sob $G$,
resp. feixes localmente livres de posto $n$) por tais morfismos; e, por fim,
(iii) de que o grupo de automorfismos do fibrado trivial sobre um $S'/S$ é
funtorialmente isomorfo ao grupo de automorfismos do feixe localmente livre
trivial de posto $n$ sobre $S'$, conclui-se «formalmente» que «dá no mesmo»
dar sobre $S$ (ou sobre um $S'/S$) um fibrado principal homogêneo de grupo $G$
ou dar aí um feixe localmente livre de posto $n$. (Mais precisamente, tem-se
uma equivalência de categorias fibradas.) Conclui-se, em particular:

\result{PROPOSIÇÃO}{6.1} Todo fibrado principal homogêneo de grupo
$\operatorname{Gl}(n)_S$ é localmente trivial.

Mediante argumentos conhecidos obtém-se o mesmo resultado para grupos
estruturais tais como $\operatorname{Sl}(n)_S$, $\operatorname{Sp}(n)_S$ e
produtos de tais grupos. Deduz-se também que, se $F$ é um subgrupo fechado
de $G=\operatorname{Gl}(n)_S$, plano sobre $S$, tal que existe o quociente
$G/F$, e se $G$ é um fibrado principal homogêneo isotrivial (resp.
estritamente isotrivial) sobre $G/F$, de grupo $F\times_S(G/F)$, então todo
fibrado principal homogêneo de grupo $F$ é isotrivial (resp. estritamente
isotrivial). Isso se aplica a todos os «grupos lineares» sobre $S$ utilizados
até agora e, em particular, quando $G=S\times_k\Gamma$, onde $S$ é um
pré-esquema sobre o corpo $k$ e $\Gamma$ um grupo linear no sentido clássico
e, em particular, liso, sobre $k$. Isso resolve, pois, para tais grupos, uma
questão de Serre (loc. cit.).

Assinalemos igualmente que, para a maioria dos grupos lisos sobre $S$ conhecidos,
sejam lineares ou não, e em todo caso para os da forma
$S\times_k\Gamma$ anterior, pode demonstrar-se que todo fibrado principal
homogêneo isotrivial é estritamente isotrivial. Isso resolve, em particular,
outra questão de Serre (loc. cit., pp. 1--14), dado que um fibrado principal
homogêneo obtido por uma descida à Cartier (cf. parágrafo 3, exemplo 2) é
manifestamente isotrivial.

\result{OBSERVAÇÃO}{} Uma das dificuldades essenciais nessas questões
(além da questão da existência dos esquemas quociente $G/F$) é a
falta de critérios de efetividade para um dado de descida por um morfismo
fielmente plano não finito.

\src{327}
\begin{center}
  \textbf{BIBLIOGRAFIA}
\end{center}

\begin{description}
  \item[{[1]}] Dieudonné (J.) e Grothendieck (Alexander). --- \emph{Éléments
    de géométrie algébrique}, a aparecer nas Publications
    mathématiques de l'Institut des Hautes Études Scientifiques.
  \item[{[2]}] Grauert (H.) und Remmert (R.). --- \emph{Komplexe Räume}, Math.
    Annalen, t. 136, 1958, pp. 245--318.
  \item[{[3]}] Grothendieck (Alexander). --- \emph{Géométrie formelle et
    géométrie algébrique}. Séminaire Bourbaki, t. 11, 1958/59, n.º 182.
  \item[{[4]}] Murre (J. P.). --- \emph{On a connectedness theorem for a
    birational transformation at a simple point}, Amer. J. Math., t. 80, 1958,
    pp. 3--15.
  \item[{[5]}] Serre (Jean-Pierre). --- \emph{Géométrie algébrique et géométrie
    analytique}, Ann. Inst. Fourier Grenoble, t. 6, 1955--56, pp. 1--42.
  \item[{[6]}] Serre (Jean-Pierre). --- \emph{Espaces fibrés algébriques},
    Séminaire Chevalley, t. 2, 1958: Anneaux de Chow et applications, n.º 1.
\end{description}

\par

\def\FgaUnitKey{e190}\clearpage
\chapter*{Errata da exposição n\textsuperscript{o} 190}
\addcontentsline{toc}{chapter}{Errata da exposição n\textsuperscript{o} 190}
\setcounter{footnote}{0}
\markboth{Errata da exposição 190}{Errata da exposição 190}
\sid{e190}
\src{300}

\noindent\textbf{3. Técnica de descida e teoremas de existência em geometria
algébrica [TDTE I]. I: Generalidades. Descida por morfismos fielmente planos.}

\begin{center}
[Seminário Bourbaki, t. 12, 1959/60, n\textsuperscript{o} 190, 29 p.]
\end{center}

\noindent\textbf{Página 1, linha 10.} --- Agora parece excessivo dizer que a
técnica de descida está «na base da maioria dos teoremas de
existência em geometria algébrica». Isto é verdadeiro em grande medida para as
técnicas não projetivas tratadas nas duas primeiras exposições
da série TDTE I a VI, mas não para as técnicas projetivas (TDTE IV, V,
VI).

\noindent\textbf{Página 5, linha 16.} --- Não é necessário supor que $\alpha$
seja um morfismo de $\mathcal F$-descida.

\noindent\textbf{Página 20, observação.} --- Um morfismo
$S'\longrightarrow S$ quase-finito, étale e sobrejetivo, ou um morfismo
$\operatorname{Spec}(\widehat A)\longrightarrow\operatorname{Spec}(A)$, não é
sempre um morfismo de descida estrito, mesmo que $A$ seja um anel local de uma
curva algébrica sobre um corpo algebraicamente fechado $k$ e
$S=\operatorname{Spec}(A)$. Assim, pode-se encontrar um morfismo próprio e liso
$f:X\longrightarrow S$, que faz de $X$ um fibrado principal de base $S$ sob
uma curva elíptica $E$, tal que $f':X'\longrightarrow S'$ seja projetivo, mas
$f$ não o seja. É, portanto, também um exemplo de fibrado
principal homogêneo não isotrivial sob um esquema abeliano.

\noindent\textbf{Página 26, linha 9.} --- Leia-se «Chow--Iang» em vez de
«Igusa--Lang».

Para diversos detalhes acerca da teoria da descida, vejam-se SGA VI, VII e
VIII.

\par

\def\FgaUnitKey{195}\unit{Exposição 195 -- Técnica de descida II}
\sid{195}
\markboth{Exposição 195}{Exposição 195}
\src{369}
\seminar{12.\textsuperscript{o} ano, 1959/60, n.\textsuperscript{o} 195\\Fevereiro de 1960}
\paperhead{TÉCNICA DE DESCIDA E TEOREMAS DE EXISTÊNCIA EM GEOMETRIA ALGÉBRICA\\
II. O TEOREMA DE EXISTÊNCIA NA TEORIA FORMAL DOS MÓDULOS}%
  {por Alexander GROTHENDIECK}

\fsection{A}{Funtores representáveis e pro-representáveis}

\subsection*{1. Funtores representáveis}
\addcontentsline{toc}{subsection}{1. Funtores representáveis}

Seja $\mathcal C$ uma categoria. Para todo $X\in\mathcal C$, seja $h_X$ o
funtor contravariante de $\mathcal C$ na categoria $(\operatorname{Ens})$
dos conjuntos,
\[
  h_X:\mathcal C\longrightarrow(\operatorname{Ens}),
\]
definido pela fórmula
\[
  h_X(Y)=\operatorname{Hom}(Y,X).
\]
Um morfismo $X\longrightarrow X'$ em $\mathcal C$ induz de maneira evidente um
homomorfismo de funtores $h_X\longrightarrow h_{X'}$: $h_X$ é um funtor
covariante em $X$; isto é, temos definido um funtor covariante canônico
\[
  h:\mathcal C\longrightarrow
  \operatorname{Hom}\bigl(\mathcal C^o,(\operatorname{Ens})\bigr)
\]
de $\mathcal C$ na categoria dos funtores covariantes da categoria dual
$\mathcal C^o$ de $\mathcal C$ na categoria dos conjuntos. Recordemos
então:

\result{PROPOSIÇÃO}{1.1} Este funtor $h$ é \emph{plenamente fiel}; em outros
termos, para todo par $X,X'$ de objetos de $\mathcal C$, a aplicação
natural
\[
  \operatorname{Hom}(X,X')\longrightarrow
  \operatorname{Hom}(h_X,h_{X'})
\]
é bijetiva.

Em particular, se um funtor
$F\in\operatorname{Hom}(\mathcal C^o,(\operatorname{Ens}))$ é isomorfo a um
funtor da forma $h_X$, então $X$ fica determinado salvo isomorfismo
único. Diz-se nesse caso que o funtor $F$ é \emph{representável}. A
proposição anterior significa que o funtor canônico $h$ define uma
equivalência da categoria $\mathcal C$ com a subcategoria plena de
$\operatorname{Hom}(\mathcal C^o,(\operatorname{Ens}))$ formada pelos
funtores representáveis. Este fato fundamenta a noção de «solução de um
problema universal»: tal problema consiste sempre em examinar se um funtor
dado (contravariante como aqui, ou covariante no caso dual) de $\mathcal C$
em $(\operatorname{Ens})$ é representável.

\src{370}
Observemos além disso que, pela própria definição de produto em uma categoria
[1], o funtor $h:X\longmapsto h_X$ comuta com os produtos sempre que estes
existam (e, mais geralmente, com os limites projetivos finitos ou infinitos;
em particular, com os produtos fibrados, com a formação de «núcleos» [2],
etc., sempre que existam): tem-se um isomorfismo de funtores
\[
  h_{X\times X'}\xrightarrow{\sim}h_X\times h_{X'}
\]
sempre que exista $X\times X'$; isto é, têm-se bijeções funtoriais
em $Y$
\[
  h_{X\times X'}(Y)\xrightarrow{\sim}h_X(Y)\times h_{X'}(Y).
\]
Em particular, dar um morfismo
\[
  X\times X'\longrightarrow X''
\]
em $\mathcal C$ (isto é, uma «lei de composição» em $\mathcal C$ entre
$X,X',X''$) equivale a dar um morfismo
\[
  h_{X\times X'}=h_X\times h_{X'}\longrightarrow h_{X''},
\]
isto é, a dar, para todo $Y\in\mathcal C$, uma lei de composição de
conjuntos
\[
  h_X(Y)\times h_{X'}(Y)\longrightarrow h_{X''}(Y)
\]
de tal modo que, para todo morfismo $Y\longrightarrow Y'$ em $\mathcal C$, o
sistema de aplicações entre conjuntos
\[
  h_{X^{(i)}}(Y)\longleftarrow h_{X^{(i)}}(Y')
  \qquad(\text{para }i=0,1,2)
\]
seja um morfismo com respeito às duas leis de composição correspondentes a
$Y$ e $Y'$. Vê-se assim que a noção de estrutura de «$\mathcal C$-grupo»,
«$\mathcal C$-anel», etc. sobre um objeto $X$ de $\mathcal C$ se exprime do
modo mais cômodo (tanto na teoria como na prática) dizendo que, para
todo $Y\in\mathcal C$, o conjunto $h_X(Y)$ está provido de uma lei de grupo
(resp. de anel, etc.) no sentido usual, e que as aplicações
$h_X(Y)\longleftarrow h_X(Y')$ correspondentes a morfismos
$Y\longrightarrow Y'$ são homomorfismos de grupos (resp. de anéis, etc.).
Esta é, por exemplo, a forma mais intuitiva e cômoda de definir os diversos
grupos clássicos $G_a$, $G_m$, $\operatorname{Gl}(n)$, etc. sobre um pré-esquema
base arbitrário $S$ e de escrever as relações clássicas entre esses grupos;
de dotar o esquema afim $V(\mathcal F)$ sobre $S$, definido por um feixe
quase coerente $\mathcal F$, de uma estrutura de «fibrado vetorial»; e de
definir e estudar as diversas variedades de bandeiras associadas
(grassmannianas, fibrados projetivos), etc. O yoga geral consiste
simplesmente em identificar, mediante o funtor canônico $h$, os objetos de
$\mathcal C$ com certos funtores contravariantes: os funtores representáveis,

\src{371}
de $\mathcal C$ na categoria dos conjuntos.

O procedimento habitual de inverter as setas, necessário, por exemplo, no
caso dos esquemas afins para passar da linguagem geométrica à da álgebra
comutativa, leva a dualizar as considerações anteriores e, em
particular, a introduzir também os funtores covariantes representáveis
$\mathcal C\longrightarrow(\operatorname{Ens})$; isto é, os da forma
\[
  Y\longmapsto\operatorname{Hom}(X,Y)=h'_X(Y).
\]

\subsection*{2. Funtores pro-representáveis, pro-objetos}
\addcontentsline{toc}{subsection}{2. Funtores pro-representáveis, pro-objetos}

Seja $\widetilde X=(X_i)_{i\in I}$ um sistema projetivo de objetos de
$\mathcal C$; a ele corresponde um funtor covariante
\[
  h'_{\widetilde X}=\varinjlim_i h'_{X_i},
\]
mais explicitamente,
\[
  h'_{\widetilde X}(Y)=\varinjlim_i h'_{X_i}(Y)
  =\varinjlim_i\operatorname{Hom}(X_i,Y),
\]
de $\mathcal C$ em $(\operatorname{Ens})$. Um funtor de $\mathcal C$ em
$(\operatorname{Ens})$ que seja isomorfo a um funtor desse tipo com $I$
filtrante chama-se \emph{pro-representável}. Pelo número anterior, estes são
exatamente os funtores isomorfos a limites indutivos filtrantes de funtores
representáveis. Seja $\widetilde X'=(X'_j)_{j\in J}$ um segundo sistema
projetivo filtrante em $\mathcal C$ (construído sobre um segundo conjunto
preordenado filtrante de índices, $J$); verifica-se facilmente que há uma
bijeção canônica (que generaliza a proposição 1.1)
\[
  \operatorname{Hom}(h'_{\widetilde X'},h'_{\widetilde X})
  =\varprojlim_j\varinjlim_i\operatorname{Hom}(X_i,X'_j).
\]
Isso leva a introduzir a categoria $\operatorname{Pro}(\mathcal C)$ dos
pro-objetos de $\mathcal C$; seus objetos são os sistemas projetivos de objetos
de $\mathcal C$ (sobre conjuntos variáveis de índices preordenados filtrantes),
e se $\widetilde X=(X_i)_{i\in I}$ e
$\widetilde X'=(X'_j)_{j\in J}$ são dois objetos dessa classe, põe-se
\[
  \operatorname{ProHom}(\widetilde X,\widetilde X')
  =\varprojlim_j\varinjlim_i\operatorname{Hom}(X_i,X'_j),
\]
sendo evidente a composição dos pro-homomorfismos. Por sua própria
construção, $\widetilde X\longmapsto h'_{\widetilde X}$ pode ser considerado
como um funtor contravariante em $\widetilde X$, que estabelece uma equivalência
entre a categoria dual da categoria $\operatorname{Pro}(\mathcal C)$ dos
pro-objetos de $\mathcal C$ e a categoria dos funtores covariantes
pro-representáveis de $\mathcal C$ em $(\operatorname{Ens})$. Naturalmente, um
objeto $X$ de $\mathcal C$ define canonicamente um pro-objeto, denotado

\src{372}
também por $X$, de modo que $\mathcal C$ é equivalente a uma subcategoria
plena de $\operatorname{Pro}(\mathcal C)$. Se então
$\widetilde X=(X_i)_{i\in I}$ é um pro-objeto qualquer de $\mathcal C$, tem-se
(com a identificação anterior)
\[
  \widetilde X=\varprojlim_i X_i,
\]
tomando-se o limite projetivo em $\operatorname{Pro}(\mathcal C)$
(visto que $h'_{\widetilde X}=\varinjlim_i h'_{X_i}$). Convém observar que,
mesmo quando o limite projetivo dos $X_i$ exista em $\mathcal C$, em
geral não será isomorfo ao limite projetivo $\widetilde X$ em
$\operatorname{Pro}(\mathcal C)$, como já resulta evidente quando $\mathcal C$
é a categoria dos conjuntos. Note-se que, por definição,
$\varprojlim_{\mathcal C}X_i=L$ fica definido pela condição de que o
funtor em $Y\in\mathcal C$
\[
  \varprojlim_i\operatorname{Hom}_{\mathcal C}(Y,X_i)
  =\operatorname{Hom}_{\operatorname{Pro}(\mathcal C)}(Y,\widetilde X)
\]
com valores em $(\operatorname{Ens})$ seja representável mediante $L$, isto é,
seja isomorfo a $\operatorname{Hom}_{\mathcal C}(Y,L)$; por conseguinte,
$\varprojlim_{\mathcal C}X_i$ já está definido em termos do pro-objeto
$\widetilde X$ e, mais precisamente, depende funtorialmente do pro-objeto
$\widetilde X$ sempre que esteja definido; não há, pois, inconveniente em
denotá-lo por $\varprojlim_{\mathcal C}(\widetilde X)$. Se os limites
projetivos existem sempre em $\mathcal C$,
$\varprojlim_{\mathcal C}(\widetilde X)$ é um funtor de
$\operatorname{Pro}(\mathcal C)$ em $\mathcal C$, e há um homomorfismo
canônico de funtores
$\varprojlim_{\mathcal C}(\widetilde X)\longrightarrow\widetilde X$.

Como todo funtor, que suporemos covariante para fixar as ideias,
\[
  F:\mathcal C\longrightarrow\mathcal C'
\]
se prolonga de maneira evidente a um funtor
\[
  \operatorname{Pro}(F):\operatorname{Pro}(\mathcal C)
  \longrightarrow\operatorname{Pro}(\mathcal C'),
\]
segue-se que, se em $\mathcal C'$ sempre existem os limites projetivos, $F$
define também canonicamente um funtor composto
\[
  \overline F:\operatorname{Pro}(\mathcal C)\longrightarrow\mathcal C',
  \qquad
  \overline F=\varprojlim_{\mathcal C'}\circ\operatorname{Pro}(F),
\]
que transforma $\widetilde X=(X_i)_{i\in I}$ em
$\varprojlim_{\mathcal C'}F(X_i)$.

Um pro-objeto $\widetilde X$ chama-se \emph{pro-objeto estrito} se é isomorfo a
um pro-objeto $(X_i)_{i\in I}$ cujos morfismos de transição
$X_i\longrightarrow X_j$ são epimorfismos; um funtor definido por tal objeto
chama-se \emph{estritamente pro-representável}. Pode-se exigir, além disso, que $I$
seja ordenado filtrante e que todo epimorfismo $X_i\longrightarrow X'$ seja
equivalente a um epimorfismo

\src{373}
$X_i\longrightarrow X_j$ para um $j\in I$ adequado (determinado de maneira
única por essa condição). Nessas condições, o sistema projetivo
$(X_i)_{i\in I}$ fica determinado salvo isomorfismo único (no sentido usual
dos isomorfismos de sistemas projetivos). Segue-se que em
$\operatorname{Pro}(\mathcal C)$ existe sempre o limite projetivo de um
sistema projetivo de pro-objetos estritos $\widetilde X^{(\alpha)}$ e que, com
as notações anteriores para $F$, $\overline F$, ter-se-á
\[
  \overline F\left(\varprojlim_\alpha\widetilde X^{(\alpha)}\right)
  =\varprojlim_{\alpha,\mathcal C'}
  \overline F\left(\widetilde X^{(\alpha)}\right).
\]
Em particular, se todo pro-objeto de $\mathcal C$ é estrito (cf. o número
seguinte), o funtor prolongado $\overline F$ comuta com os limites
projetivos.

\subsection*{3. Caracterização dos funtores pro-representáveis}
\addcontentsline{toc}{subsection}{3. Caracterização dos funtores pro-representáveis}

Sejam $\mathcal C$, $\mathcal C'$ duas categorias nas quais existem os limites
projetivos finitos (isto é, relativos a conjuntos preordenados finitos, não
necessariamente filtrantes), ou, o que é equivalente, nas quais existem os
produtos finitos e os produtos fibrados finitos (o que implica, em
particular, a existência de um «objeto final» $e$, tal que
$\operatorname{Hom}(X,e)$ se reduza a um elemento para todo $X$). Seja $F$ um
funtor covariante de $\mathcal C$ em $\mathcal C'$. As condições seguintes
são equivalentes:

\begin{enumerate}
\item[(i)] $F$ comuta com os limites projetivos finitos;
\item[(ii)] $F$ comuta com os produtos finitos e os produtos fibrados finitos;
\item[(iii)] $F$ comuta com os produtos finitos e, para todo diagrama exato
\[
  X\longrightarrow X'\rightrightarrows X''
\]
em $\mathcal C$ ([3], A, definição 2.1), o diagrama transformado por $F$
\[
  F(X)\longrightarrow F(X')\rightrightarrows F(X'')
\]
é exato.
\end{enumerate}

Diz-se então que $F$ é \emph{exato à esquerda}.

De agora em diante, supõe-se que em $\mathcal C$ existem os limites projetivos
finitos. Das definições resulta então imediatamente que um funtor
representável é exato à esquerda e, passando ao limite, que um funtor
pro-representável é exato à esquerda.

Para obter uma recíproca, seja
\[
  F:\mathcal C\longrightarrow(\operatorname{Ens})
\]
um funtor covariante, e sejam $X\in\mathcal C$ e $\xi\in F(X)$. Diz-se que
$\xi$ (ou o par \src{374}$(X,\xi)$) é \emph{minimal} se, para todo
$X'\in\mathcal C$ e $\xi'\in F(X')$, e todo monomorfismo estrito ([3], A, 2)
$u:X'\longrightarrow X$ tal que $\xi=F(u)(\xi')$, $u$ é um isomorfismo. Por
outro lado, diz-se que um par $(X,\xi)$ \emph{domina} $(X'',\xi'')$
($\xi\in F(X)$, $\xi''\in F(X'')$) se existe um morfismo
$v:X\longrightarrow X''$ tal que $\xi''=F(v)(\xi)$; se $\xi$ é minimal e $F$
é exato à esquerda, este morfismo $v$ é então único; se $\xi''$ é
minimal, $v$ é sobrejetivo. Deduz-se facilmente a proposição seguinte:

\result{PROPOSIÇÃO}{3.1} Para que $F$ seja estritamente pro-representável, é
necessário e suficiente que satisfaça as duas condições:
\begin{enumerate}
\item[(i)] $F$ é exato à esquerda;
\item[(ii)] todo par $(X,\xi)$, com $\xi\in F(X)$, é dominado por um par
minimal.
\end{enumerate}

Esta última condição é vazia se todo objeto de $\mathcal C$ é artiniano
(tomando um subobjeto $X'$ de $X$ minimal entre aqueles para os quais existe
$\xi'\in F(X')$ tal que $\xi$ seja a imagem de $\xi'$). Daí:

\result{COROLÁRIO}{} Seja $\mathcal C$ uma categoria cujos objetos são
artinianos e na qual existem os limites projetivos finitos. Então os
funtores pro-representáveis de $\mathcal C$ em $(\operatorname{Ens})$ são
exatamente os funtores exatos à esquerda, e são inclusive estritamente
pro-representáveis.

Este último fato significa também que \emph{todo pro-objeto de $\mathcal C$
é estrito}.

\subsection*{4. Exemplo: grupos de tipo galoisiano, grupos proalgébricos}
\addcontentsline{toc}{subsection}{4. Grupos de tipo galoisiano, grupos proalgébricos}

Se $\mathcal C$ é a categoria dos grupos finitos ordinários,
$\operatorname{Pro}(\mathcal C)$ é equivalente à categoria dos grupos
topológicos compactos totalmente desconexos. São grupos desse tipo e suas
generalizações, obtidas substituindo os grupos finitos ordinários por
esquemas em grupos finitos sobre um pré-esquema base dado (por exemplo, os
grupos algébricos finitos sobre um corpo $k$), que farão as vezes de
grupos fundamentais, de homotopia e de homologia absolutos e relativos para
os pré-esquemas. Em todos estes exemplos se aplica o corolário da
proposição 3.1, e é precisamente por meio do funtor associado que se deve
definir $\pi_1$ [2]. O mesmo ocorre partindo da categoria dos grupos
algébricos ou quase algébricos sobre um corpo (ou, mais geralmente, sobre um
pré-esquema noetheriano): obtêm-se os «grupos proalgébricos» de Serre [4].

\subsection*{5. Exemplo: «variedades formais»}
\addcontentsline{toc}{subsection}{5. Exemplo: variedades formais}

Sejam $\Lambda$ um anel noetheriano e $\mathcal C$ a categoria das
$\Lambda$-álgebras $A$ que são módulos artinianos de tipo finito sobre
$\Lambda$ (ou, mais brevemente,

\src{375}
as $\Lambda$-álgebras artinianas). São satisfeitas as condições do corolário da
proposição 3.1. A categoria $\operatorname{Pro}(\mathcal C)$ é aqui
equivalente à categoria das álgebras topológicas $O$ sobre $\Lambda$
isomorfas a limites projetivos topológicos
\[
  O=\varprojlim_i O_i
\]
de álgebras $O_i\in\mathcal C$, isto é, cuja topologia é linear, separada e
completa, e tal que para todo ideal aberto $J_i$ de $O$, $O/J_i$ seja uma
álgebra artiniana sobre $\Lambda$. O funtor
$\mathcal C\longrightarrow(\operatorname{Ens})$ associado a tal álgebra não é
outro que
\begin{align*}
  F(A)=h'_O(A)
  &=\{\text{homomorfismos contínuos de $\Lambda$-álgebras topológicas }
       O\longrightarrow A\}\\
  &=\varinjlim_i
    \operatorname{Hom}_{\Lambda\text{-álgebras}}(O_i,A).
\end{align*}

Observe-se, além disso, que a categoria $\mathcal C$ é essencialmente o produto
das categorias análogas correspondentes aos completamentos dos
anéis locais $\Lambda_{\mathfrak m}$, para os ideais maximais
$\mathfrak m$ de $\Lambda$; portanto, se desejado, pode-se restringir ao caso
em que $\Lambda$ seja um anel local completo dessa classe. Em qualquer caso,
$O$ se decompõe canonicamente no produto topológico de componentes
locais, correspondentes aos «pontos» do esquema formal [2] definido por
$O$. Um ponto dessa classe é definido por um elemento $\xi$ de algum
$F(K)$, onde $K\in\mathcal C$ é um corpo (por exemplo, o corpo residual
da componente local considerada); dois pares $(\xi,K)$ e $(\xi',K')$
definem o mesmo ponto se, e somente se, são dominados por um mesmo
$(\xi'',K'')$, ou se dominam um mesmo $(\xi''',K''')$. (Se os
$\Lambda/\mathfrak m$ são algebricamente fechados, basta tomar a soma
disjunta dos $F(\Lambda/\mathfrak m)$.)

É importante dar condições para que a componente local $O_\xi$ de $O$
correspondente a um $\xi\in F(K)$ seja um anel noetheriano. Quando $\Lambda$
é um anel local completo (noetheriano, recordemos), isso equivale a dizer
que $O_\xi$ é isomorfo a um anel quociente de um anel de séries formais
$\Lambda[[t_1,\ldots,t_n]]$ sobre $\Lambda$. Para dar um critério desse tipo,
introduzamos (para todo anel $A$) a $A$-álgebra $I_A$ dos «números
duais» de $A$ mediante
\[
  I_A=A[t]/t^2A[t].
\]
Seja $\varepsilon:I_A\longrightarrow A$ o homomorfismo de aumentação; este define
(se $A\in\mathcal C$) uma aplicação
\[
  F(\varepsilon):F(I_A)\longrightarrow F(A),
\]

\src{376}
e, utilizando que $F$ é exato à esquerda, define-se de maneira
intrínseca uma estrutura de $K$-módulo na parte
\[
  F(I_A,\xi)=F(\varepsilon)^{-1}(\xi)
\]
de $F(I_A)$ formada pelos $\xi'\in F(I_A)$ que «se reduzem a $\xi$»;
utilizando a forma explícita de $F$ em termos de $O$, verifica-se que este
$A$-módulo se identifica com
\[
  \operatorname{Hom}_{\Lambda}(\mathfrak m_\xi/\mathfrak m_\xi^2,A),
\]
onde $\mathfrak m_\xi$ é o núcleo do homomorfismo
$\xi:O\longrightarrow A$, isto é, se $A$ é um corpo, o ideal maximal da
componente local $O_\xi$ de $O$. Deduz-se imediatamente a proposição
seguinte:

\result{PROPOSIÇÃO}{5.1} Seja $\xi\in F(K)$, onde $K\in\mathcal C$ é um
corpo. Para que a componente local correspondente $O_\xi$ de $O$ seja um
anel noetheriano, é necessário e suficiente que o conjunto $F(I_K,\xi)$ de
elementos de $F(I_K)$ que se reduzem a $\xi$ seja um espaço vetorial
de dimensão finita sobre $K$. Nessas condições, tem-se um isomorfismo
canônico
\[
  F(I_K,\xi)=
  \operatorname{Hom}\!\left(\mathfrak m_\xi/\mathfrak m_\xi^2,K\right),
\]
onde $\mathfrak n_\xi$ é o ideal maximal de $\Lambda$, núcleo do
homomorfismo $\Lambda\longrightarrow K$.
Em particular, a dimensão do $K$-espaço vetorial $F(I_K,\xi)$ é igual à
dimensão do espaço vetorial $\mathfrak m_\xi/\mathfrak m_\xi^2$ sobre o corpo
$O_\xi/\mathfrak m_\xi=K(\xi)$.

Suponhamos que $O_\xi$ seja noetheriano e, para simplificar a escrita, que
$\Lambda$ seja local completo e $O=O_\xi$. Diz-se que $O$ é \emph{liso sobre
$\Lambda$} se $O$ é uma álgebra finita e étale sobre a álgebra completada do
localizado de $\Lambda[t_1,\ldots,t_n]$ em um ideal maximal desse anel que
induz o ideal maximal de $\Lambda$; quando a extensão residual de $O$ sobre
$\Lambda$ é trivial (por exemplo, se o corpo residual de $\Lambda$ é
algebricamente fechado), isso equivale a dizer que o próprio $O$ é isomorfo a uma
álgebra de séries formais desse tipo.
Finalmente, se já não se supõe necessariamente que $O$ seja noetheriano, dir-se-á
também que $O$ é liso sobre $\Lambda$ se $O$ é isomorfo a um limite
projetivo topológico de $\Lambda$-álgebras quociente que são noetherianas e
lisas sobre $\Lambda$ no sentido anterior. Isso se generaliza imediatamente ao caso
em que $\Lambda$ e $O$ já não são supostos locais. Dito isso, tem-se a
proposição seguinte:

\result{PROPOSIÇÃO}{5.2} Para que $O$ seja liso sobre $\Lambda$, é necessário
e suficiente que o funtor associado $F$ transforme epimorfismos em
epimorfismos.

Isto significa que, para todo homomorfismo sobrejetivo
$A\longrightarrow A'$ em $\mathcal C$, $F(A)\longrightarrow F(A')$ é
sobrejetivo. Naturalmente, basta verificar esta condição quando $A$ é local
e, procedendo passo a passo, quando o ideal de $A$ que é núcleo de
$A\longrightarrow A'$ é anulado pelo ideal maximal de $A$. Na prática,
isso reduz a questão a verificar que certa obstrução, ligada a invariantes
infinitesimais

\src{377}
da situação que dá origem ao funtor $F$, seja nula, problema de natureza
cohomológica.

Para terminar, digamos algumas palavras, no contexto anterior, sobre os
\emph{anéis de definição}. Seja ainda $F$ um funtor de $\mathcal C$ em
$(\operatorname{Ens})$, pro-representável mediante uma $\Lambda$-álgebra
topológica $O$. Então, para todo $A\in\mathcal C$ e todo $\xi\in F(A)$,
existe um menor subanel $A'$ de $A$ tal que $\xi$ seja imagem de um elemento
$\xi'$ de $F(A')$ (que fica então determinado de maneira única): com efeito,
basta interpretar $\xi$ mediante um homomorfismo de $O$ em $A$ e tomar como
$A'$ a imagem de $O$ por este último. Diz-se então que $A'$ é o anel
de definição do objeto $\xi\in F(A)$. Se $u:A\longrightarrow B$ é um
homomorfismo de álgebras e $\eta=F(u)(\xi)$, o anel de definição de $\eta$
é a imagem por $u$ do anel de definição de $\xi$. Quando se parte de um
funtor $F$ de $\mathcal C$ em $(\operatorname{Ens})$, a existência dos
anéis de definição e as propriedades que acabamos de dar são, com pequenas
ressalvas, equivalentes a que $F$ seja pro-representável; isto é, em
geral estão longe de ser triviais.

\fsection{B}{Os dois teoremas de existência}

Conservemos as notações de A, parágrafo 5, e partamos de um funtor covariante
\[
  F:\mathcal C\longrightarrow(\operatorname{Ens});
\]
buscam-se critérios manejáveis para que $F$ seja pro-representável, isto é,
exprimível mediante uma $\Lambda$-álgebra $O$ como acima. Em virtude do
corolário de A, proposição 3.1, para isso é necessário e suficiente que $F$ seja
exato à esquerda. No estado atual da técnica de descida (cf. as
questões formuladas em [3], p. 9), esse critério não pode ser verificado
diretamente sob essa forma nos casos mais importantes, e são necessários
critérios aparentemente menos exigentes.

\result{TEOREMA}{1} Para que o funtor $F$ seja pro-representável, é necessário
e suficiente que satisfaça as duas condições seguintes:
\begin{enumerate}
\item[(i)] $F$ comuta com os produtos finitos;
\item[(ii)] para toda álgebra $A\in\mathcal C$ e todo homomorfismo
$A\longrightarrow A'$ em $\mathcal C$ tal que o diagrama
\[
  A\longrightarrow A'\rightrightarrows A'\otimes_A A'
\]
seja exato ([3], A, definição 1.2), o diagrama transformado
\[
  F(A)\longrightarrow F(A')\rightrightarrows F(A'\otimes_A A')
\]

\src{378}
é exato.
\end{enumerate}
Além disso, basta verificar (ii) quando $A$ é local e, ademais, em um dos
dois casos seguintes:
\begin{enumerate}
\item[(a)] $A'$ é um módulo livre sobre $A$;
\item[(b)] o módulo quociente $A'/A$ é um $A$-módulo de comprimento 1.
\end{enumerate}

A demonstração deste teorema é bastante delicada e não pode ser esboçada aqui.
Limitemo-nos a assinalar que ela repousa essencialmente em um estudo das relações
de equivalência (no sentido das categorias) no espectro de uma álgebra
artiniana (estudo que ainda suscita vários problemas cuja solução parece
indispensável ao desenvolvimento ulterior da teoria).

Nas aplicações, a verificação da condição (i) é sempre trivial.
A verificação de (ii) decompõe-se em duas partes: o caso em que $A'$ é um $A$-módulo livre
corresponde à técnica de descida por morfismos planos ([3], teoremas 1, 2 e
3) e não oferece dificuldade. Para tratar o caso (b), será utilizado o resultado
seguinte:

\result{TEOREMA}{2} Sejam $A$ um anel local artiniano de ideal maximal
$\mathfrak m$ e $A'$ uma $A$-álgebra que contém a $A$, tal que
$\mathfrak mA'\subset A$ e que
\[
  A\longrightarrow A'\rightrightarrows A'\otimes_A A'
\]
seja exato (o que ocorre, em particular, se $A'/A$ é um $A$-módulo de
comprimento 1). Seja $\mathcal F$ a categoria fibrada ([3], A, definição 1.1) de
feixes quase coerentes e planos sobre pré-esquemas variáveis. Então o
morfismo $\operatorname{Spec}(A')\longrightarrow\operatorname{Spec}(A)$ é um
morfismo de $\mathcal F$-descida estrito ([3], A, definição 1.7).

Em outros termos, o dado de um $A$-módulo plano $M$ é completamente
equivalente ao dado de um $A'$-módulo plano $M'$, provido de um
$A'\otimes_AA'$-isomorfismo de $M'\otimes_AA'$ sobre $A'\otimes_AM'$ que
satisfaz a condição usual de transitividade para um dado de descida (loco
citato).

O teorema 2 demonstra-se provando primeiro que
\[
  H^i(A'/A,G_a)=0\qquad\text{para }i\geq 1
\]
([3], A, 4(e)), pois a hipótese $\mathfrak mA'\subset A$ permite reduzir-se
facilmente ao caso em que $A$ é um corpo $k$ (a saber, $A/\mathfrak m$). Se
Aplicam-se então as equivalências assinaladas em [3], página 16.

\fsection{C}{Aplicações a alguns casos particulares}

\subsection*{1. Observações gerais sobre os funtores representados por pré-esquemas}
\addcontentsline{toc}{subsection}{1. Observações gerais sobre os funtores representados por pré-esquemas}

Seja $S$ um pré-esquema localmente noetheriano. Diz-se que um pré-esquema $X$
sobre $S$ é \src{379}localmente de tipo finito sobre $S$ se, para todo
$x\in X$ que se projeta em $y\in Y$, existem
uma vizinhança afim de $y$, de anel $A$, e uma vizinhança afim de $x$ situada sobre
o anterior, de anel $B$, tais que $B$ seja uma álgebra de tipo finito sobre
$A$. Encontram-se muitos exemplos importantes de pré-esquemas localmente de
tipo finito sobre $S$ que não são de tipo finito sobre $S$, como soluções de
problemas universais clássicos; assim, importa poder considerar o esquema de
Picard de uma curva como a reunião de uma infinidade de componentes conexas
(que não devem ser confundidas com a componente conexa do elemento neutro, isto
é, a «variedade de Picard»). Em certas ocasiões, é então conveniente
trabalhar na categoria $\mathcal C$ dos pré-esquemas localmente de tipo
finito sobre $S$, para nela examinar a questão da representabilidade de
um funtor contravariante $F$. O objetivo principal destas exposições é
desenvolver uma técnica geral que permita reconhecer se tal funtor $F$ é
representável e estudar as propriedades do $S$-pré-esquema $X$ correspondente
a partir das de $F$. Assinalemos de passagem que, nesse estudo, aparecem
exemplos não patológicos de pré-esquemas sobre $S$ não separados sobre $S$, em
particular como «pré-esquemas de Picard» de excelentes $S$-esquemas; não devem,
portanto, banir-se da geometria algébrica os pré-esquemas que não são
esquemas.

Seja $X$ um pré-esquema localmente de tipo finito sobre $S$, e seja
\[
  F(Y)=\operatorname{Hom}_S(Y,X)
\]
o funtor contravariante associado. Pode-se considerar a restrição $F_0$ de
$F$ à subcategoria $\mathcal C_0$ de $\mathcal C$ formada pelos
pré-esquemas $Y$ sobre $S$ que são artinianos e finitos sobre $S$: quando
$S=\operatorname{Spec}(\Lambda)$, $\mathcal C_0$ é, pois, a categoria dual
da categoria das $\Lambda$-álgebras artinianas considerada em B. Se
$Y=\operatorname{Spec}(A)$, onde $A$ é um anel local artiniano, $Y$
reduz-se a um único ponto $y$ situado sobre um ponto fechado $s$ de $S$, e um
$S$-homomorfismo de $Y$ em $X$ (isto é, um elemento de $F(Y)$) fica definido
pelo dado de um ponto $x\in X$ sobre $s$ e de um
$\mathcal O_{s,S}$-homomorfismo de $\mathcal O_{x,X}$ em $A$. Se existe tal
homomorfismo, então $x$ é necessariamente um ponto fechado de $X$ (pois seu
corpo residual é algébrico sobre o de $s$). Isso mostra que a restrição
$F_0$ de $F$ às «$\Lambda$-álgebras artinianas» é pro-representável, e está
representada pela álgebra topológica sobre $\Lambda$ cujas componentes
locais são os completamentos $\widehat{\mathcal O}_{x,X}$ dos anéis locais
de $X$ nos pontos $x$ de $X$ que são fechados e estão situados sobre pontos
fechados de $Y$. Isso mostra que apenas o conhecimento de $F_0$ fornece
informação precisa sobre a estrutura de $X$ (a saber, a estrutura dos
completamentos de seus anéis locais nos pontos indicados). Observe-se que,
mesmo quando $S$

\src{380}
é o espectro de um corpo algebricamente fechado, somente mediante a
consideração sistemática de «variedades» $Y$ tais que $\mathcal O_Y$ possa
ter elementos nilpotentes, em particular trabalhando com os espectros de
anéis artinianos locais, pode-se chegar à «formulação correta» dos
problemas universais clássicos e captar seu aspecto «infinitesimal».

Quando se parte de um funtor $F$ dado de antemão e se quer decidir se ele é
representável, o estudo do funtor $F_0$ (com o auxílio dos teoremas 1 e 2)
fornecerá indicações quase completas: ou, como ocorre com frequência
(verificando simplesmente, por exemplo, a natureza dos conjuntos
$F(I_K,\xi)$ e seu comportamento funtorial, cf. A), constata-se que $F_0$ já não
é pro-representável (o que explica o fracasso das tentativas realizadas
até agora para definir de maneira mais ou menos natural variedades de módulos
para a classificação dos fibrados vetoriais de posto $>1$); ou consegue-se
verificar que $F_0$ é representável, mas os espaços vetoriais
$F(I_K,\xi)$ não são de dimensão finita, caso em que será preciso contentar-se com
a solução «formal»; ou, por fim, $F_0$ é representado por um
produto de anéis locais completos noetherianos, o que fornece indícios muito
sólidos de que o próprio $F$ é representável e, unido a propriedades análogas,
porém de natureza mais global, que pretendemos desenvolver mais adiante, bastará
sem dúvida para implicar que assim seja efetivamente. Finalmente, encontram-se
problemas geométricos interessantes (vejam-se os parágrafos 4 e 5 abaixo) nos
quais se dispõe apenas do funtor $F_0$ (que não provém de nenhum funtor «global»
$F$), e nos quais haverá motivos para se dar por satisfeito se se conseguir associar-lhe
um «esquema formal de módulos».

Para terminar estas generalidades, indiquemos como a teoria de esquemas
explica anomalias aparentes, como a superfície de Igusa $V$, cuja «variedade
de Picard» $P$ se reduz a um ponto e para a qual, no entanto, tem-se
\[
  H^1(V,\mathcal O_V)\ne 0 ;
\]
nesse caso, $P$ é um grupo «puramente infinitesimal» que não se reduz ao
grupo unidade, isto é, está definido por uma álgebra local $O$ de posto finito
sobre o corpo base $k$ e dotado de uma aplicação diagonal correspondente à
estrutura multiplicativa de $P$; se $\mathfrak m$ é o ideal maximal de
$O$, o dual de $\mathfrak m/\mathfrak m^2$ é canonicamente isomorfo a
$H^1(V,\mathcal O_V)$ (cf. o parágrafo 3 mais abaixo). Somente quando o grupo de
Picard é um grupo algébrico no sentido clássico, isto é, liso sobre o
corpo base $k$, a dimensão de $H^1(V,\mathcal O_V)$ (sempre igual à de
$\mathfrak m/\mathfrak m^2$) é igual à do grupo de Picard.

\subsection*{2. Esquemas $\underline{\operatorname{Hom}}_S(X,Y)$,
$\Pi_{X/S}Z$, $\underline{\operatorname{Aut}}_S(X)$, etc.}
\addcontentsline{toc}{subsection}{2. Esquemas Hom, Pi e Aut}

Sejam $X,Y$ dois pré-esquemas sobre $S$; para todo pré-esquema $T$ sobre $S$,

\src{381}
sejam $X_T=X\times_ST$, $Y_T=Y\times_ST$, e tomemos o conjunto
\[
  F(T)=\operatorname{Hom}_T(X_T,Y_T)
      =\operatorname{Hom}_S(X_T,Y)
      =\operatorname{Hom}_S(X\times_ST,Y)
\]
como funtor contravariante em $T$. Se ele for representável, denota-se por
$\underline{\operatorname{Hom}}_S(X,Y)$ o pré-esquema sobre $S$ que o
representa; tem-se, portanto, um isomorfismo funtorial
\[
  \operatorname{Hom}_S
  \bigl(T,\underline{\operatorname{Hom}}_S(X,Y)\bigr)
  \xrightarrow{\sim}
  \operatorname{Hom}_S(T\times_SX,Y).
\]
Há muitas variantes deste problema universal cuja solução se reduz a ele:
o pré-esquema dos $S$-automorfismos de um $S$-pré-esquema $X$ (que será um
pré-esquema em grupos), o pré-esquema dos $S$-homomorfismos de um
$S$-pré-esquema em grupos em outro (que será um pré-esquema em grupos comutativos
se o segundo esquema em grupos é comutativo), etc. Também se pode
generalizar a definição de $\underline{\operatorname{Hom}}_S(X,Y)$
considerando um pré-esquema $Z$ sobre o pré-esquema $X$ sobre $S$, e o funtor
\[
  F(T)=\operatorname{Hom}_{X_T}(X_T,Z_T)
\]
(conjunto das «seções» de $Z_T$ sobre $X_T$); se esse funtor for
representável, o $S$-pré-esquema que o representa será denotado por
$\Pi_{X/S}Z$; ter-se-á, pois, por definição, um isomorfismo funtorial
\[
  \operatorname{Hom}_S(T,\Pi_{X/S}Z)
  =\operatorname{Hom}_{X_T}(X_T,Z_T).
\]
Tomando $Z=Y\times_SX$, recupera-se
$\underline{\operatorname{Hom}}_S(X,Y)$. Dessas definições resulta para os
novos pré-esquemas assim introduzidos um formulário tão trivial quanto útil, que não
daremos aqui (visto que é válido em toda categoria na qual existem
produtos e produtos fibrados). Mais séria é a questão da existência de
esquemas do tipo $\underline{\operatorname{Hom}}_S(X,Y)$. Constata-se primeiro
que, para $X$ fixo, $\underline{\operatorname{Hom}}_S(X,Y)$ somente pode existir
para todo $Y$ sobre $S$ (ou $\Pi_{X/S}Z$ para todo $Z$ sobre $X$) se $X$ é plano
sobre $S$. Além disso, percebe-se que só é razoável esperar a existência de
uma solução, para $Y$ suficientemente geral, se $X$ é além disso próprio sobre
$S$. Por outro lado, parece que essas condições são suficientes para a
existência de $\underline{\operatorname{Hom}}_S(X,Y)$ e $\Pi_{X/S}Z$, com a
condição de impor, quando couber, uma hipótese do tipo «quase-projetivo»
sobre $Y/S$, respectivamente sobre $Z/X$; isso pode ser verificado, em qualquer
caso, em bastantes situações (por exemplo, quando $Y$ é afim sobre $S$; ou,
mediante construções diretas elementares, quando $X$ é finito sobre $S$).
Eis o que fornecem os teoremas 1 e 2:

\src{382}
\result{PROPOSIÇÃO}{2.1} Sejam $\Lambda$ um anel noetheriano, $X$ e $Y$ dois
pré-esquemas quaisquer sobre $\Lambda$, e tomemos o funtor
\[
  F(A)=\operatorname{Hom}_A(X_A,Y_A)
\]
na categoria $\mathcal C_0$ das $\Lambda$-álgebras artinianas. Se $X$ é
plano sobre $\Lambda$, este funtor é pro-representável.

Além disso, verifica-se que para todo $A\in\mathcal C_0$ e $\xi\in F(A)$ tem-se
um isomorfismo canônico
\[
  F(I_A,\xi)=H^1\!\left(
    X_A,
    \underline{\operatorname{Hom}}_{\mathcal O_{X_A}}
    \bigl(\xi^*(\Omega^1_{Y_A/A}),\mathcal O_{X_A}\bigr)
  \right),
\]
onde $\Omega^1_{Y_A/A}$ é o feixe de $1$-diferenciais de Kähler de $Y_A$
com respeito a $A$. Tomando $A$ como um corpo, obtém-se, utilizando A.5.1 e o
teorema de finitude de [2], o corolário seguinte:

\result{COROLÁRIO}{} Suponhamos que $X$ seja plano e próprio sobre $S$, e que $Y$
seja localmente de tipo finito sobre $S$. Então $F$ é pro-representável e as
componentes locais da álgebra topológica correspondente sobre $\Lambda$ são
anéis noetherianos.

\result{OBSERVAÇÕES}{} Os problemas considerados nesta seção, e muitos
outros, eram habitualmente estudados, no âmbito da geometria algébrica
clássica, mediante as «coordenadas de Chow» dos ciclos no espaço
projetivo, que permitiam considerar esses ciclos como pontos de variedades
projetivas adequadas. Esse procedimento, e em geral o uso das
coordenadas de Chow, parece irremediavelmente insuficiente do ponto de
vista dos esquemas, pois destrói os elementos nilpotentes nas
variedades de parâmetros e, em particular, não se presta a um estudo
satisfatório das variações infinitesimais de ciclos (sem falar de sua
natureza não intrínseca, ligada ao espaço projetivo). A linguagem das
coordenadas de Chow era, infelizmente, a única utilizada por muitos
geômetras algébricos para estudar famílias de variedades ou famílias de
ciclos, o que parece ter constituído um sério obstáculo para esclarecer essas
noções, apesar de seu indubitável interesse técnico (provavelmente provisório). Se
se quiser obter o análogo das variedades de Chow na teoria dos
esquemas, chega-se ao seguinte problema universal: seja $X$ um pré-esquema sobre
$S$; para todo pré-esquema $T$ sobre $S$, considera-se o conjunto $F(T)$ dos
subesquemas fechados de $X_T=X\times_ST$ que são planos sobre $T$. Tenta-se
representar este funtor em $T$ mediante um pré-esquema sobre $S$. Mais
geralmente, pode-se dar um feixe quase coerente $G$ sobre $X$ e tomar como
$F(T)$ o conjunto dos feixes

\src{383}
quociente de $G_T$ que são planos sobre $T$. Parece que o problema tem uma
solução, mediante um esquema $C$ localmente de tipo finito sobre $S$, quando
$X$ é próprio sobre o pré-esquema localmente noetheriano $S$ e $G$ é além disso
coerente. Em qualquer caso, supondo unicamente que $S$ seja localmente
noetheriano, a restrição de $F$ às «$S$-álgebras artinianas» é
pro-representável; e se, além disso, $X$ é próprio sobre $S$ e $G$ é coerente, as
componentes locais do anel topológico correspondente $O$ são
noetherianas. Naturalmente, mesmo depois de demonstrar a existência do
«esquema de Chow» de $X$ sobre $S$, restará
encontrar uma decomposição em abertos disjuntos $C_i$ (correspondentes à
fixação de invariantes contínuos, como o grau e a dimensão dos
ciclos que se fazem variar) que sejam de tipo finito (e não somente localmente de
tipo finito) sobre $S$, precisar as relações deste esquema com as
variedades de Chow clássicas e precisar quando um $C_i$ é inclusive projetivo,
ou ao menos quase-projetivo, sobre $S$.

\subsection*{3. Esquemas de Picard}
\addcontentsline{toc}{subsection}{3. Esquemas de Picard}

Seja $f:X\longrightarrow S$ um $S$-pré-esquema, e tomemos o feixe
multiplicativo $\mathcal O_X^*$ das unidades do feixe estrutural de $X$, e
o grupo
\[
  P(X/S)=H^0\bigl(S,R^1f_*(\mathcal O_X^*)\bigr),
\]
chamado grupo de Picard relativo de $X/S$. Um elemento deste grupo fica,
pois, definido dando uma cobertura aberta $(U_i)$ de $S$ e um feixe invertível
$L_i$ sobre cada $f^{-1}(U_i)$, de modo que, para todo par $(i,j)$,
$L_i|f^{-1}(U_i\cap U_j)$ seja isomorfo a $L_j|f^{-1}(U_i\cap U_j)$, ao menos
localmente sobre $U_i\cap U_j$ (isto é, esses dois feixes sejam «equivalentes»
no sentido de [3], B, 4). Se $X/S$ admite uma seção, então $P(X/S)$ nada mais
é que o conjunto das classes de feixes invertíveis sobre $X/S$ módulo
«equivalência» (loco citato). Ponhamos agora, para todo $T$ sobre $S$,
\[
  F(T)=P(X_T/T) ;
\]
obtém-se um funtor covariante em $T$, que poderá ser chamado funtor de Picard de
$X/S$; se esse funtor for representável, o pré-esquema sobre $S$ que o
representa será chamado pré-esquema de Picard de $X/S$ e será denotado por
$\underline P(X/S)$. Ter-se-á, pois, um isomorfismo de funtores
\[
  \operatorname{Hom}_S\bigl(T,\underline P(X/S)\bigr)
  \xrightarrow{\sim}P(X_T/T).
\]
A formação de pré-esquemas de Picard é compatível com a extensão da
base; em particular, os pré-esquemas de Picard das fibras de $X$ sobre $S$
(que são pré-esquemas sobre os corpos residuais $K(s)$ dos $s\in S$) são
as fibras de $\underline P(X/S)$. Naturalmente, como $P(X_T/T)=F(T)$ é um
grupo comutativo, os pré-esquemas de Picard são pré-esquemas em grupos.
Observe-se também que as jacobianas generalizadas de Rosenlicht são precisamente
as componentes conexas da unidade nos esquemas de Picard de curvas
completas que podem ter singularidades, o que deveria tornar evidente a
maioria de suas propriedades (uma vez demonstrada a existência).

\result{OBSERVAÇÃO}{} A definição adotada aqui só é razoável quando
todo ponto de $Y$ admite uma vizinhança
aberta $U$ sobre a qual $X$ admite uma seção. No caso geral, é preciso
modificar ligeiramente a definição do funtor de Picard para poder obter
também um teorema de existência.

Aqui, as condições plausíveis para a existência de um pré-esquema de Picard
são as seguintes: $X$ é próprio e plano sobre $S$,
$f_*(\mathcal O_X)=\mathcal O_S$, e $X$ admite localmente uma seção sobre
$S$. Essa condição aparece de maneira natural ao aplicar a técnica de
descida, para eliminar os automorfismos de um feixe invertível $L$ sobre $X$
dotando-o de uma seção marcada sobre a seção $s$ ([3], B, 4). Obtém-se,
em particular:

\src{384}
\result{PROPOSIÇÃO}{3.1} Suponhamos que $X$ seja plano sobre
$S=\operatorname{Spec}(\Lambda)$, com $\Lambda$ noetheriano, e que, para todo
$T$ de tipo finito sobre $S$, tem-se
\[
  f_{T*}(\mathcal O_{X_T})=\mathcal O_T
\]
(se $f$ é próprio, ou se $S$ é o espectro
de um corpo, resulta de Künneth que a última condição equivale a
$f_*(\mathcal O_X)=\mathcal O_S$). Então o funtor de Picard de $X/S$ na
categoria das $\Lambda$-álgebras artinianas é pro-representável.

Além disso, tem-se aqui
\[
  F(I_A,\xi)=H^1(X_A,\mathcal O_{X_A}),
\]
em particular:

\result{COROLÁRIO}{} Se $X$ é próprio sobre $S$, então as componentes
locais da $\Lambda$-álgebra topológica $O$ correspondente ao funtor de
Picard são noetherianas.

\result{OBSERVAÇÕES}{} Podem-se generalizar as definições e os resultados
desta seção à classificação dos fibrados principais sobre $X$, tendo como
grupo estrutural um esquema em grupos $G$ sobre $S$ que seja afim, plano sobre
$S$ e comutativo. Quando $G$ não é comutativo, e portanto o fibrado em
grupos adjunto de um fibrado principal (cujas seções são os automorfismos
do fibrado principal) deixa de ser trivial, a proposição 3.1 já não é válida
tal como está. Não obstante, pode-se modificar o problema universal de modo que
continue a ter solução (ao menos, por enquanto, em geometria formal). A regra de
ouro que se deve recordar, no contexto desta \src{385}seção e das
seguintes, e sempre que se busquem «esquemas de módulos» para classes de
objetos definidos apenas a menos de isomorfismo, é sempre a seguinte: eliminar
os possíveis automorfismos dos objetos que se querem classificar mediante
a introdução, se necessário, de estruturas auxiliares (pontos ou
elementos de seções marcadas, fixação de formas diferenciais, etc.), que
serão escolhidas suficientemente inócuas para não modificar substancialmente o
problema inicial.

\subsection*{4. Módulos formais de uma variedade}
\addcontentsline{toc}{subsection}{4. Módulos formais de uma variedade}

Seja $\Lambda$ um anel local noetheriano de corpo residual $k$ (com maior
frequência, $\Lambda$ será igual a $k$, ou a um $p$-anel de Cohen), e seja
$X_0$ um pré-esquema sobre $k$. Para toda $\Lambda$-álgebra artiniana local $A$,
considera-se o conjunto $F(A)$ das classes, salvo isomorfismo, de
$A$-pré-esquemas $X$ planos sobre $A$, dotados de um isomorfismo
\[
  (*)\qquad
  X\otimes_A\kappa(A)\xleftarrow{\sim}
  X_0\otimes_k\kappa(A),
\]
onde $\kappa(A)$ é o corpo residual de $A$; naturalmente, os isomorfismos
entre tais $A$-pré-esquemas planos devem respeitar o isomorfismo precedente
incluído na estrutura. Se $A$ é uma $\Lambda$-álgebra artiniana não
necessariamente local, de componentes locais $A_i$, toma-se como $F(A)$ o
produto dos $F(A_i)$. Desse modo, $F$ torna-se um funtor
multiplicativo em $A$, que poderá ser chamado funtor de módulos para $X_0$ (e
$\Lambda$). Se esse funtor for representável, corresponde-lhe uma
$\Lambda$-álgebra topológica $O$ local, de corpo residual $k$, e o espectro
formal de $O$ será chamado esquema formal de módulos para $X_0$ (e $\Lambda$);
cf. [2] para alguns detalhes sobre essa noção.

Aqui, para aplicar a técnica de descida, os automorfismos «finitos» de $X_0$
são inofensivos, pois não influem na existência de automorfismos (no
sentido precisado acima) dos $A$-pré-esquemas $X$; quando $A$ não se
reduz a um corpo, a condição necessária e suficiente para que $X$ não tenha
nenhum $A$-automorfismo não trivial é que se tenha
\[
  H^0(X_0,\Theta_{X_0/k})=0,
\]
onde $\Theta_{X_0/k}$ denota o feixe das $k$-derivações (= feixe tangente)
de $X_0$. Por outro lado, obtém-se facilmente (ao menos se $X_0$ for liso
sobre $k$)
\[
  F(I_A,\xi)=H^1(X_A,\Theta_{X_A/A}).
\]
Conclui-se então, como de costume:

\src{386}
\result{PROPOSIÇÃO}{4.1} Suponhamos que
$H^0(X_0,\Theta_{X_0/k})=0$. Então existe o esquema formal de módulos para
$X_0$. Se, além disso, $X_0$ for próprio sobre $k$, o esquema formal de módulos é
noetheriano.

\result{OBSERVAÇÕES}{} 1\textsuperscript{o} Quando não se supõe que $X_0$ seja
liso sobre $k$, $F(I_A,\xi)$ identifica-se com um sub-$A$-módulo de
\[
  \operatorname{Ext}^1_{\mathcal O_{P_A}}
  \bigl(P_A;\mathcal I_{X_A},\mathcal O_{X_A}\bigr),
\]
onde se põe $P_A=X_A\times_\Lambda X_A$, onde $\mathcal O_{X_A}$ se
considera como um feixe coerente sobre $P_A$ mediante o morfismo diagonal
$X_A\longrightarrow P_A$, e onde $\mathcal I_{X_A}$ denota o feixe coerente
de ideais sobre $P_A$ definido pelo morfismo diagonal. Mais precisamente,
uma globalização simples da teoria de Hochschild mostra que o
$\operatorname{Ext}^1$ anterior identifica-se com o conjunto das classes,
salvo isomorfismo, de feixes de $I_A$-álgebras planas $\mathcal O$ sobre $X_A$,
dotados de um isomorfismo de aumentação
\[
  \mathcal O\otimes_{I_A}A\longrightarrow\mathcal O_{X_A}
\]
(recorde-se que se põe $I_A=A[t]/(t^2)$). O submódulo $F(I_A,\xi)$ é o
que corresponde aos feixes de álgebras comutativas. As hipóteses de lisura
não são, pois, essenciais na teoria de módulos, contrariamente ao que [2]
dava a entender.

2\textsuperscript{o} Recordemos (loco citato) que, em particular, toda curva
algébrica $X_0$ lisa e própria sobre $k$ admite um esquema formal de módulos,
liso sobre $\Lambda$, de dimensão relativa $3g-3$ se o gênero $g$ é
$\geq2$, e $g$ se $g=0$ ou $1$. Esses dois últimos casos já não entram diretamente
na proposição 4.1. No entanto, no caso das curvas elípticas
($g=1$), poderiam ser reduzidos a ela mediante as considerações que seguem.

Naturalmente, pode-se variar ad libitum a proposição 4.1 considerando
sistemas de esquemas sobre $k$ dotados de diversas estruturas. Suponhamos,
por exemplo, que $X_0$ seja um esquema abeliano sobre $k$, com ponto de origem
marcado (isto é, que $X_0$ se considere como esquema em grupos sobre $k$), e
seja $F(A)$ o conjunto das classes, salvo isomorfismo, de esquemas abelianos
sobre $A$ (isto é, de esquemas em grupos próprios e lisos sobre $A$) dotados
de um isomorfismo de esquemas abelianos $(*)$. Verifica-se que a imposição
de uma estrutura multiplicativa (e inclusive de uma mera «seção unidade»)
elimina os automorfismos infinitesimais e que, por conseguinte, existe um
esquema formal de módulos, correspondente a um anel local completo
noetheriano $O$. Por outro lado, pode-se demonstrar que, se $X$ é um esquema
próprio e liso com fibras «absolutamente conexas» sobre um pré-esquema localmente
noetheriano $S$, toda estrutura multiplicativa sobre $X$ que admita uma
seção unidade é necessariamente associativa e comutativa (sempre que o seja,
ao menos, sobre uma fibra e que $S$ seja conexo), e além disso fica determinada de
maneira única pelo conhecimento

\src{387}
da seção unidade. Além disso, suponhamos que $S$ seja o espectro de um anel
local artiniano $A$ de corpo residual $k$, que $X$ seja próprio sobre $A$ e esteja
dotado de uma seção $s$, e que, por fim, $X\otimes_Ak$ esteja dotado de uma
estrutura de esquema abeliano sobre $k$ cujo elemento neutro é o ponto de
$X\otimes_Ak$ correspondente a $s$; um cálculo simples de obstruções,
unido a um raciocínio devido a Serre, permite demonstrar que existe
efetivamente sobre $X$ uma estrutura multiplicativa que admite a seção
$s$ como seção unidade. (A partir daí, utilizando o «teorema de
existência» de [2] para passar ao caso em que $A$ é noetheriano local completo,
e depois a técnica de descida de [3] para o caso geral, pode-se demonstrar
o enunciado análogo para todo $S$ localmente noetheriano conexo.) Isso prova
que o funtor $F(A)$ considerado aqui é isomorfo ao funtor análogo definido ao
princípio desta seção, prescindindo da estrutura multiplicativa sobre
$X_0$. Segue-se, em particular, que, se $\mathfrak m$ é o ideal maximal de
$O$, $\mathfrak m/\mathfrak m^2$ é canonicamente isomorfo ao dual de
$H^1(X_0,\Theta_{X_0/k})$ e tem, portanto, dimensão $n^2$, onde
$n=\dim X_0$. Seria muito interessante determinar se $O$ é liso sobre $\Lambda$,
isto é, isomorfo a uma álgebra de séries formais em $n^2$ indeterminadas
sobre $\Lambda$. O parágrafo 1, proposição 5.2, permite dar uma formulação
equivalente deste problema como problema de existência de esquemas abelianos
que se reduzam a um esquema abeliano dado. Em qualquer caso, por métodos
transcendentes, vê-se que a resposta é afirmativa se $k$ tem característica
zero. Em característica $p\ne0$, basta evidentemente limitar-se ao caso em que
$\Lambda$ é o anel de vetores de Witt construído sobre um corpo
algebricamente fechado $k$. Talvez tenha chegado o momento para o «funtor
de Greenberg» demonstrar seu valor\ldots

\subsection*{5. Extensão dos recobrimentos}
\addcontentsline{toc}{subsection}{5. Extensão dos recobrimentos}

Sejam $\mathfrak X$ um pré-esquema formal noetheriano [2] e $U$ uma parte
aberta de $\mathfrak X$, definida localmente pela «não anulação» de uma
seção de $\mathcal O_{\mathfrak X}$ que não seja divisor de zero sobre
$\mathfrak X$, e que seja, portanto, grande o bastante para que toda seção
de $\mathcal O_{\mathfrak X}$ sobre um aberto $V$ que seja nula sobre $U\cap V$
seja nula. Seja $\mathcal J$ um «ideal de definição» para $\mathfrak X$, e seja
\[
  X_n=(\mathfrak X,\mathcal O_{\mathfrak X}/\mathcal J^{n+1}),
\]
que é, portanto, um pré-esquema noetheriano ordinário. Se então
$\mathfrak X',\mathfrak X''$ são dois recobrimentos planos de $\mathfrak X$ (isto
é, dois pré-esquemas sobre $\mathfrak X$ definidos por feixes de álgebras que
são coerentes e localmente livres como feixes de módulos), não ramificados sobre
$U$, a aplicação evidente
\[
  \operatorname{Hom}_{\mathfrak X}(\mathfrak X',\mathfrak X'')
  \longrightarrow
  \operatorname{Hom}_{X_0}(X'_0,X''_0)
\]
é injetiva; em particular, um automorfismo de $\mathfrak X'$ que induz a
identidade sobre $X'_0$ é a identidade. Isso permite aplicar à situação a
técnica de descida. Partamos, em particular, de um recobrimento plano $X'_0$
de $X_0$, não ramificado

\src{388}
sobre $U_0$; seja $G(\mathfrak X)$ o conjunto das classes, salvo isomorfismo
(que induzam a identidade sobre $X'_0$), de recobrimentos planos $\mathfrak X'$
de $\mathfrak X$ que induzem $X'_0$ sobre $X_0$ (e são, portanto,
necessariamente não ramificados sobre $U$). Define-se do mesmo modo $G(V)$ para
toda parte aberta $V$ de $\mathfrak X$, e mais geralmente $G(\mathfrak Y)$
para todo pré-esquema formal $\mathfrak Y$ sobre $\mathfrak X$. Dito isto, os
resultados de [2] e [3] implicam em primeiro lugar os resultados seguintes:

\noindent a. Quando $V$ é um aberto variável de $\mathfrak X$, os $G(V)$
formam um feixe sobre $\mathfrak X$, seja $G_{\mathfrak X}=G$. A restrição desse
feixe a $U$ é o feixe constante cujas fibras se reduzem a um elemento.

Mais geralmente, a determinação das fibras de $G_{\mathfrak X}$ é uma
questão de anéis locais completos; mais precisamente:

\noindent b. Para todo $x\in\mathfrak X$, tem-se
\[
  G_x=G\bigl(\operatorname{Spec}(\widehat{\mathcal O}_{\mathfrak X,x})\bigr),
\]
(= classes, salvo isomorfismo, de álgebras finitas e livres $B$ sobre
$\widehat{\mathcal O}_{\mathfrak X,x}$ dotadas de um isomorfismo de
$B\otimes_{\widehat{\mathcal O}_{\mathfrak X,x}}\mathcal O_{X_0,x}$ com
$\mathcal O'_{0x}$, onde $\mathcal O'_0$ é o feixe de álgebras sobre
$X_0$ que define $X'_0$).

\noindent c. Tem-se um isomorfismo canônico
\[
  G_{\mathfrak X}=\varprojlim_nG_{X_n} ;
\]
em outros termos, para todo aberto $V$ de $\mathfrak X$, tem-se
\[
  G(V)=\varprojlim_nG(V_n).
\]

\noindent d. Suponhamos que $\mathfrak X$ seja obtido de um esquema ordinário
próprio $X$ sobre um anel local noetheriano completo $\Lambda$ que possui um
ideal de definição $\mathfrak m$, tomando o completamento $\mathcal J$-ádico de
$\mathcal O_X$, onde $\mathcal J=\mathfrak m\mathcal O_X$. Então
$G(\mathfrak X)$ é canonicamente isomorfo ao conjunto das classes de
recobrimentos planos do esquema ordinário $X$ que «se reduzem a $X'_0$».

De maneira gráfica, pode-se dizer que (a) e (b) estabelecem as relações
fundamentais entre o aspecto local e o aspecto global do problema, (c) dá
as relações entre o aspecto «finito» e o aspecto «infinitesimal» e, por
último, (d) recorda (sob as condições especificadas) a identidade entre o
aspecto «formal» e o aspecto «algébrico».

Suponhamos agora que $\mathfrak X$ esteja definido sobre um anel local
noetheriano completo $\Lambda$, com
$\mathcal J=\mathfrak m\mathcal O_{\mathfrak X}$ (portanto, $X_0$ é um
pré-esquema sobre $\Lambda/\mathfrak m$). Para toda álgebra $A$ finita sobre
$\Lambda$, ponhamos
\[
  F(A)=G(\mathfrak X\times_\Lambda A).
\]
É um funtor covariante em $A$, com valores na categoria dos conjuntos,
e, por (c), esse funtor fica completamente determinado se for conhecido para $A$
artiniana; dizer que esse funtor é pro-representável, isto é, da forma
\[
  F(A)=\operatorname{Hom}_{\Lambda\text{-algèbres top.}}(O,A),
\]

\src{389}
onde $O$ é uma álgebra topológica sobre $\Lambda$ do tipo considerado em A,
parágrafo 5, equivale a dizer que o mesmo ocorre ao se restringir aos $A$
artinianos. A conjunção dos teoremas 1 e 2 implica então efetivamente:

\result{PROPOSIÇÃO}{5.1} O funtor precedente é pro-representável.

Naturalmente, segundo (a), se $U=\mathfrak X$, $G(\mathfrak Y)$ se reduz a um
único elemento para todo $\mathfrak Y$ sobre $\mathfrak X$, e o funtor $F$
tem então escasso interesse (ter-se-á $O=\Lambda$). Parece que, em quase
qualquer outro caso, o anel local topológico $O$ não é noetheriano. Não
obstante, sua existência põe em evidência, de maneira surpreendente, a
natureza «contínua» do conjunto $G(\mathfrak X)$ das soluções (que
corresponde intuitivamente ao fato de que há uma escolha «contínua» no
modo como a ramificação se propaga ao estender $X'_0$). Compare-se esse
resultado com o ponto de vista de J.-P. Serre [4] na teoria do corpo de
classes local, que põe igualmente em relevo o caráter contínuo do grupo de
Galois topológico da extensão abeliana maximal de um corpo local
«geométrico»; ali, o grupo dual (no sentido de Pontryagin) aparece como um
limite indutivo de grupos algébricos (ou ao menos quase algébricos); também
aqui, a classificação das extensões é dada por «variedades» de
dimensão infinita. Por outro lado, no que precede, pode-se tomar como
$\mathfrak X$ o espectro formal de um anel local completo (do qual
$\Lambda$ será, por exemplo, um subanel de Cohen), e cabe esperar que os
resultados desta seção possam ser utilizados no estudo das extensões
de um anel local completo de dimensão $>1$. Tanto no caso local quanto no
global, talvez permitam formular relações precisas entre os fenômenos de
ramificação superior e fenômenos em característica zero (abordáveis por
métodos transcendentes). Em qualquer caso, a análise preliminar da
proposição 5.1 é o que permite estender ao caso «moderadamente ramificado»
os métodos expostos em [2] para estudar o grupo fundamental e resolver por
métodos transcendentes o «problema dos três pontos».

Para terminar, assinalemos que a situação se simplifica quando $X_0$ tem dimensão 1: então, em virtude de (a) e (b),
\[
  G(\mathfrak X)\simeq
  \prod_iG\bigl(\operatorname{Spec}(\widehat{\mathcal O}_{\mathfrak X,x_i})\bigr),
\]
onde os $x_i$ são os pontos de $X_0-U$: podem-se dar arbitrariamente as
extensões «locais» nos pontos de ramificação. Além disso, se $X_0$ é normal,
observa-se que o esquema formal de módulos cuja existência é garantida por 5.1 é
liso sobre $\operatorname{Spec}(\Lambda)$.

\src{390}
\begin{center}
  \textbf{BIBLIOGRAFIA}
\end{center}

\begin{description}
  \item[{[1]}] Grothendieck (Alexander). --- \emph{Sur quelques points
    d'algèbre homologique}, Tohoku Math. J., t. 9, 1957, p. 119--221.
  \item[{[2]}] Grothendieck (Alexander). --- \emph{Géométrie formelle et
    géométrie algébrique}, Séminaire Bourbaki, t. 11, 1958/59, fasc. 3,
    n\textsuperscript{o} 182.
  \item[{[3]}] Grothendieck (Alexander). --- \emph{Technique de descente et
    théorèmes d'existence en géométrie algébrique, I : Généralités, Descente
    par morphismes fidèlement plats}, Séminaire Bourbaki, t. 12, 1959/60,
    fasc. 1, n\textsuperscript{o} 190.
  \item[{[4]}] Serre (Jean-Pierre). --- \emph{Corps locaux et isogénies},
    Séminaire Bourbaki, t. 11, 1958/59, fasc. 3, n\textsuperscript{o} 185.
\end{description}

\par

\def\FgaUnitKey{e195}\clearpage
\chapter*{Errata da exposição n.\textsuperscript{o} 195}
\addcontentsline{toc}{chapter}{Errata da exposição n.\textsuperscript{o} 195}
\setcounter{footnote}{0}
\markboth{Errata da exposição 195}{Errata da exposição 195}
\sid{e195}
\src{300}

\noindent\textbf{4. TDTE II: O teorema de existência na teoria formal dos
módulos.}

\begin{center}
[Seminário Bourbaki, t. 12, 1959/60, n.\textsuperscript{o} 195, 22 p.]
\end{center}

\noindent\textbf{Página 8.} --- A fórmula escrita na proposição 5.1 só é
correta quando $\Lambda$ é um corpo; no caso geral, deve-se substituir
$\mathfrak m_\xi/\mathfrak m_\xi^2$ pelo quociente deste espaço pela
imagem de $\mathfrak n_\xi/\mathfrak n_\xi^2$, onde $\mathfrak n$ é o ideal
maximal de $\Lambda$. Além disso, a definição dada para que $O$ seja liso sobre
$\Lambda$ só é correta quando a extensão residual $k'/k$ é separável.
No caso geral, cf. SGA III, 1.1.

\noindent\textbf{Página 14, observação.} --- Os problemas formulados aqui
ficam completamente resolvidos no caso projetivo mediante os «esquemas de
Hilbert» (cf. TDTE IV). Exemplos de Nagata e Hironaka mostram, em contrapartida, que
os funtores considerados não são necessariamente representáveis se não se impõe
uma hipótese de projetividade, mesmo limitando-se à classificação das
subvariedades de dimensão zero de uma variedade completa não singular de
dimensão 3; isso está ligado ao fato de que o quadrado simétrico de uma
tal variedade pode não existir.

\src{301}
\noindent\textbf{Página 15, n.\textsuperscript{o} 3.} --- Para um estudo mais
completo, veja-se TDTE V.

\noindent\textbf{Página 16, proposição 3.1.} --- Em vez de «se $f$ é
próprio», leia-se «se $f$ é próprio e separável».

\noindent\textbf{Páginas 18 e 19, observações, 2.\textsuperscript{o}} --- Demonstrei
recentemente que o esquema formal de módulos de uma variedade
abeliana sobre um corpo é efetivamente liso sobre o anel de Witt; em
outros termos, todo esquema abeliano sobre um anel artiniano local que seja
quociente de outro provém, por redução, de um esquema abeliano sobre este
último. A demonstração utiliza unicamente as propriedades de variância da
classe de obstrução ao levantamento introduzida na exposição
n.\textsuperscript{o} 182, página 12. Recordemos igualmente que os esquemas de
módulos das curvas de gênero $g$ e dos esquemas abelianos polarizados foram
construídos por Mumford (cf. SHMT).

\noindent\textbf{Página 19, parágrafo 5.} --- Deve-se supor que a seção
considerada para definir $U$ seja um não divisor de zero não somente sobre
$\mathfrak X$, mas também sobre todo $X_n$.

\noindent\textbf{Página 20, linhas 12 e 13.} --- Em vez de
\[
  G_x=G\bigl(\operatorname{Spec}(\widehat{\mathcal O}_{\mathfrak X,x})\bigr),
\]
leia-se
\[
  G_x=G\bigl(\operatorname{Spec}(\mathcal O_{\mathfrak X,x})\bigr)
  \subset
  G\bigl(\operatorname{Spec}(\widehat{\mathcal O}_{\mathfrak X,x})\bigr),
\]
e, em vez de «$\mathcal O'_{0x}$», leia-se
«$\widehat{\mathcal O}'_{0x}$».

\par

\def\FgaUnitKey{212}\unit{Exposição 212 -- Pré-esquemas quocientes}
\sid{212}
\markboth{Exposição 212}{Exposição 212}
\src{99}
\seminares{13.º ano, 1960/61, n.º 212\\Fevereiro de 1961}
\paperhead{TÉCNICAS DE CONSTRUÇÃO E TEOREMAS DE EXISTÊNCIA EM GEOMETRIA ALGÉBRICA\\
III. PRÉ-ESQUEMAS QUOCIENTES}%
  {por Alexander GROTHENDIECK}

\subsection*{Introdução}
\addcontentsline{toc}{subsection}{Introdução}

Os problemas tratados na presente exposição diferem dos considerados nas duas
anteriores, pois se procura representar certos funtores covariantes, e não mais
contravariantes, de esquemas variáveis. O procedimento de passagem ao quociente
é, contudo, essencial em muitos problemas de construção em geometria algébrica,
inclusive os das Exposições I e II ([1], [2]). Assim, a questão da
\emph{efetividade de um dado de descida} sobre um $T$-pré-esquema $X$,
relativamente a um morfismo fielmente plano e quase-compacto
$T\longrightarrow S$, equivale à questão da existência de um quociente de $X$
(que satisfaça as propriedades razoáveis examinadas adiante) pela relação de
equivalência plana em $X$ definida pelo dado de descida; as questões levantadas
em [1], A, \S~2 c serão, sem dúvida, resolvidas ao mesmo tempo que as questões
formuladas no parágrafo 2 da presente exposição. Do mesmo modo, o
\emph{esquema de Picard} de um $S$-esquema $X$ (para a definição, ver [2], C,
\S~3) pode ser definido de várias maneiras como quociente de certos outros
esquemas (de divisores efetivos, ou de imersões em um esquema projetivo) por
relações de equivalência planas, sendo tecnicamente mais simples a definição e
a construção desses esquemas auxiliares: são esquemas da forma
$\underline{\operatorname{Hom}}_S(X,Y)$ e variantes definidos em [2], C,
\S~2, cuja construção será objeto da exposição seguinte (sob hipóteses
adequadas de projetividade). Combinando, portanto, os resultados da presente
exposição e da seguinte, obtém-se a construção dos esquemas de Picard sob
hipóteses adequadas.

O problema da passagem ao quociente na categoria dos pré-esquemas ainda
apresenta várias questões não resolvidas. A mais importante é mencionada no
parágrafo 8. Atualmente, ela é o único obstáculo à construção de
\emph{esquemas de moduli sobre os inteiros para curvas de gênero arbitrário,
variedades abelianas polarizadas}, etc. Isso mostra que sua solução merece os
esforços dos especialistas em grupos algébricos.
\src{100}
\subsection*{1. Relações de equivalência e relações de equivalência efetivas}
\addcontentsline{toc}{subsection}{1. Relações de equivalência e relações de equivalência efetivas}

Sejam $\mathcal C$ uma categoria e $X$ um objeto de $\mathcal C$. Um par de
morfismos
\[
  p_1,p_2:R\rightrightarrows X
\]
é chamado de \emph{par de equivalência} em $\mathcal C$, com \emph{alvo} $X$ e
\emph{fonte} $R$, se, para todo objeto $T$ de $\mathcal C$, as duas aplicações
correspondentes
\[
  p_1(T),p_2(T):R(T)\rightrightarrows X(T)
\]
(onde, para todo objeto $Y$ de $\mathcal C$, se põe
$Y(T)=\operatorname{Hom}(T,Y)$) definem uma aplicação
\[
  R(T)\longrightarrow X(T)\times X(T)
\]
que induz uma bijeção de $R(T)$ sobre o gráfico de uma relação de equivalência
no conjunto $X(T)$. Introduz-se, entre os pares de equivalência com alvo $X$,
uma relação de equivalência evidente; uma classe de equivalência para esta
última é chamada \emph{$\mathcal C$-relação de equivalência em $X$}, ou
simplesmente uma relação de equivalência quando não houver risco de confusão.

Se $X\times X$ existe, dar uma relação de equivalência em $X$ equivale a dar um
subobjeto $R$ de $X\times X$ tal que, para todo objeto $T$ de $\mathcal C$, o
subconjunto de
\[
  (X\times X)(T)=X(T)\times X(T)
\]
correspondente a $R(T)$ seja o gráfico de uma relação de equivalência em
$X(T)$. Denotando por $p_1$ e $p_2$ os morfismos de $R$ em $X$ induzidos pelas
projeções $\operatorname{pr}_1$ e $\operatorname{pr}_2$, a condição anterior
exprime que $(p_1,p_2)$ é um par de equivalência. Também se podem exprimir
diagramaticamente em $\mathcal C$ (desde que existam $X\times X$ e o produto
fibrado $(R,p_2)\times_X(R,p_1)$) os axiomas conjuntistas de uma relação de
equivalência para $R(T)$ em $X(T)$, de acordo com o princípio geral
explicado em [2], A, \S~1. Não precisaremos disso.

Sempre que se tem um par de morfismos $(p_1,p_2)$ com a mesma fonte $R$ e o
mesmo alvo $X$, pode-se definir o \emph{conúcleo} do par como o objeto $Y$ de
$\mathcal C$ que representa o funtor covariante em $Z$
\[
  \operatorname{Hom}_{p_1,p_2}(X,Z),
\]
conjunto dos morfismos $u$ de $X$ em $Z$ tais que $up_1=up_2$. Se $Y$ existe,
fica determinado a menos de um isomorfismo único. Será denotado por
$X/(p_1,p_2)$ ou, por abuso de notação, por $X/R$, sendo esta última notação
empregada sobretudo quando
\src{101}$(p_1,p_2)$ é um par de equivalência: é então usual identificar,
na notação, a relação de equivalência definida pelo par com $R$. Observe-se
que, se $Y$ for considerado um quociente de $X$, ele depende, de fato, apenas
da relação de equivalência definida pelo par $(p_1,p_2)$.

Partamos agora de um morfismo
\[
  f:X\longrightarrow Y
\]
que permite, portanto, considerar $X$ como um «objeto sobre $Y$», e suponhamos
que exista o produto fibrado
\[
  R(f)=X\times_Y X
\]
Sejam $p_1$ e $p_2$ suas duas projeções. Então $(p_1,p_2)$ é um par de
equivalência, chamado par \emph{associado} ao morfismo $f$. Ele define, assim,
uma relação de equivalência, chamada relação associada a $f$.

Um par de morfismos $(p_1,p_2)$ com alvo $X$ e fonte comum $R$ é chamado
\emph{par de equivalência efetivo} se
\begin{enumerate}
  \item[(i)] o conúcleo $X/(p_1,p_2)=Y$ existe;
  \item[(ii)] o produto fibrado $X\times_Y X$ existe;
  \item[(iii)] o morfismo $R\longrightarrow X\times_Y X$ cujas componentes são
  $p_1$ e $p_2$ é um isomorfismo.
\end{enumerate}
Então $(p_1,p_2)$ é de fato um par de equivalência. A relação de equivalência
que ele define é também chamada \emph{relação de equivalência efetiva}.

Um morfismo $f:X\longrightarrow Y$ é chamado \emph{epimorfismo efetivo} se
\begin{enumerate}
  \item[(i)] o produto fibrado $X\times_Y X=R$ existe;
  \item[(ii)] o quociente $X/(p_1,p_2)$ existe, onde $p_1$ e $p_2$ são as duas
  projeções de $R$ em $X$;
  \item[(iii)] o morfismo $X/(p_1,p_2)\longrightarrow Y$ induzido por $f$ é um
  isomorfismo.
\end{enumerate}
Então $f$ é de fato um epimorfismo, e até um epimorfismo estrito (cf. [1], A,
\S~2.3), sendo a recíproca verdadeira se o produto fibrado $X\times_Y X$
existir. Diz-se também que o objeto quociente de $X$ definido pelo epimorfismo
$f$ é um \emph{quociente efetivo} de $X$.

As definições precedentes implicam a seguinte «correspondência de Galois»:
\src{102}
\result{PROPOSIÇÃO}{1.1} Há uma correspondência biunívoca, que preserva as
ordens naturais, entre o conjunto das relações de equivalência efetivas $R$
em $X$ e o conjunto dos quocientes efetivos $Y$ de $X$: a $R$ corresponde o
quociente efetivo $X/R$, e a $Y$ corresponde a relação de equivalência efetiva
definida pela projeção canônica $X\longrightarrow Y$ (que é definida pelo
produto fibrado $X\times_Y X$ munido de suas duas projeções).

Nas categorias muito boas (conjuntos, feixes de conjuntos etc.), todo
quociente é efetivo e toda relação de equivalência é efetiva. Isso já não
ocorre em categorias como a categoria dos pré-esquemas sobre um pré-esquema
$S$ dado, mesmo quando $S$ é o espectro de um corpo e mesmo quando nos
restringimos aos esquemas finitos sobre $S$. As questões de efetividade e até
mesmo (no caso de pré-esquemas não finitos sobre $S$) as questões relativas à
existência de quocientes revelam-se quase sempre delicadas.

\subsection*{2. Exemplo: pré-esquemas finitos sobre $S$}
\addcontentsline{toc}{subsection}{2. Exemplo: pré-esquemas finitos sobre S}

Seja $\mathcal C$ a categoria dos pré-esquemas finitos sobre $S$, supondo-se
$S$ localmente Noetheriano. Ela é equivalente à categoria oposta da categoria
dos feixes coerentes de álgebras comutativas sobre $S$, ou ainda, se $S$ é
afim com anel $A$, à categoria oposta da categoria das $A$-álgebras finitas
sobre $A$ (isto é, finitamente geradas como $A$-módulos). Conclui-se
imediatamente que existem em $\mathcal C$ os limites projetivos finitos e os
limites indutivos finitos. Isso é bem conhecido (sem qualquer hipótese de
finitude\ldots) quanto aos primeiros. Assim, o produto fibrado de pré-esquemas
$X,Y$ sobre $S$ corresponde ao produto tensorial $B\otimes_A C$ das álgebras
correspondentes; o núcleo de dois morfismos $X\rightrightarrows Y$,
correspondentes a dois homomorfismos de $A$-álgebras
$u,v:C\rightrightarrows B$, corresponde ao quociente de $B$ pelo ideal gerado
pelos elementos $u(c)-v(c)$ etc. Para os limites indutivos finitos, basta
considerar, de um lado, as somas finitas, que correspondem ao produto ordinário
de $A$-álgebras, e, de outro, os conúcleos de pares de morfismos
$X\rightrightarrows Y$, que correspondem de fato (como se verifica
imediatamente) ao subanel de $C$ formado pelos elementos nos quais os
homomorfismos $u,v:C\rightrightarrows B$ coincidem (este último é finito sobre
$A$ graças à hipótese Noetheriana). Observe-se, aliás, que se pode mostrar,
utilizando a hipótese Noetheriana, que os limites indutivos finitos e, em
particular, os quocientes assim construídos na categoria $\mathcal C$ dos
pré-esquemas finitos sobre $S$ são, na realidade, quocientes na categoria de
\emph{todos} os pré-esquemas.
\src{103}
Como dissemos em [1], há em $\mathcal C$ epimorfismos não efetivos (ou
estritos;
dá no mesmo, pois existem produtos fibrados). \emph{Ignoro se as relações de
equivalência são sempre efetivas} quando não se impõe nenhuma hipótese de
planicidade. Nessa direção, obtive apenas resultados positivos muito parciais,
indispensáveis à demonstração do teorema fundamental da teoria formal de
moduli (ver [2], B, teo. 1). Assinalemos que, no problema proposto, é fácil
reduzir-se ao caso em que $S$ é o espectro de um anel artiniano local $A$, com
corpo residual algebricamente fechado. Mas, mesmo se $A$ for um corpo, a
resposta não é conhecida.

Também se pode considerar um pré-esquema $X$ sobre $S$ que já não se suponha
finito sobre $S$, mas que seja munido de uma relação de equivalência $R$ em
$X$ tal que $p_1:R\longrightarrow X$ seja um morfismo finito. Diz-se então que
$R$ é uma relação de equivalência finita. Suponhamos, para simplificar, que $S$
e $X$ sejam afins (logo, $R$ é afim, de modo que a situação se reduz a uma
questão de álgebra comutativa pura). \emph{Mesmo nesse caso, ignora-se se
existe um quociente $X/R=Y$ e se o morfismo canônico $X\longrightarrow Y$ é
finito}. (O caso mais simples é aquele em que $S$ é o espectro de um corpo $k$
e $X$ é o espectro de $k[t]$, isto é, a reta afim.) Naturalmente, se os dois
problemas precedentes tiverem soluções afirmativas, pode-se concluir, na
situação presente, que $R$ é efetiva. Observemos que o problema da
\emph{existência} de um quociente $Y$ e da finitude de
$f:X\longrightarrow Y$ se apresenta exatamente nos mesmos termos se, em vez
do grafo de uma relação de equivalência em $X$, houver apenas um pré-grafo de
equivalência em $X$, no sentido do parágrafo 4.

A questão da passagem ao quociente por uma relação de equivalência finita mais
ou menos arbitrária surge na construção de pré-esquemas por «colagem» de
pré-esquemas dados $X_i$ ao longo de certos subpré-esquemas fechados; a lei de
colagem exprime-se precisamente por uma relação de equivalência finita no
pré-esquema soma $X$ dos $X_i$. Deve-se esperar também que a solução dos
problemas aqui propostos, e de diversas variantes, seja uma condição prévia ao
desenvolvimento de uma técnica geral de construções não projetivas, na direção
iniciada em [2].

O único fato positivo geral conhecido pelo autor é o seguinte:

\result{PROPOSIÇÃO}{2.1} Sejam $S$ um pré-esquema localmente noetheriano, $s$
um ponto de $S$ e $\Omega$ uma extensão algebricamente fechada de $k(s)$.
Considere o
\src{104}«funtor fibra» correspondente $F$, que associa a todo $S$-esquema
$X$ finito sobre $S$ o conjunto dos pontos de $X_s$ com valores em $\Omega$.
Esse funtor (trivialmente exato à esquerda) é \emph{exato à direita}, isto é,
comuta com os limites indutivos finitos e, em particular, com os conúcleos de
pares de morfismos.

Usando esse resultado para todos os «pontos geométricos» de $S$, deduz-se que
a categoria «quociente» $\mathcal C'$ de $\mathcal C$, obtida raciocinando
«módulo morfismos radiciais sobrejetivos» (isto é, adjungindo formalmente
inversos aos morfismos em questão), é uma categoria «geométrica», isto é,
satisfaz as mesmas propriedades «de natureza finita» que a categoria dos
conjuntos. Em particular, toda relação de equivalência nela é efetiva. Isso
implica que, se $R$ for uma relação de equivalência em um $X$ finito sobre
$S$, então o morfismo canônico
\[
  R\longrightarrow X\times_Y X\qquad (\text{onde }Y=X/R)
\]
é \emph{radicial e sobrejetivo} (de fato, uma imersão fechada sobrejetiva,
pois é um monomorfismo).

\subsection*{3. O caso de um grupo de operadores}
\addcontentsline{toc}{subsection}{3. O caso de um grupo de operadores}

Suponhamos novamente que $\mathcal C$ seja uma categoria qualquer. Sejam
$G,X$ objetos de $\mathcal C$, e suponhamos que $G$ seja um
$\mathcal C$-grupo de operadores sobre o objeto $X$ de $\mathcal C$. Isso
significa (cf. [2], A, \S~1) que, para todo objeto $T$ de $\mathcal C$, foram
dadas uma estrutura de grupo em $G(T)$ e uma estrutura de $G(T)$-conjunto em
$X(T)$, de tal modo que, variando $T$, as estruturas em questão «variam
funtorialmente» com $T$. Se os produtos $G\times G$ e $G\times X$ existirem
em $\mathcal C$, tal estrutura também poderá ser definida por um par de
morfismos
\[
  G\times G\longrightarrow G,
  \qquad \pi:G\times X\longrightarrow X,
\]
sujeito à condição de que, para todo objeto $T$ de $\mathcal C$, as leis de
composição correspondentes para os conjuntos $G(T)$ e $X(T)$ façam de $G(T)$
um grupo que age sobre $X(T)$. A tradução desse axioma pela comutatividade de
certos diagramas em $\mathcal C$ é fácil, mas fastidiosa e, de fato,
perfeitamente inútil em todos os casos que conheço.

Suponhamos que $G\times X$ exista e consideremos os dois morfismos
\[
  p_1,p_2:G\times X\rightrightarrows X
\]
dados por
\src{105}
\[
  p_1=\operatorname{pr}_1,\qquad p_2=\pi.
\]
Verifica-se imediatamente que $(p_1,p_2)$ é um par de equivalência se, e
somente se, para todo objeto $T$ de $\mathcal C$, a aplicação
\[
  G(T)\times X(T)\simeq(G\times X)(T)
  \longrightarrow X(T)\times X(T),
\]
definida por esse par é injetiva, isto é, se o grupo $G(T)$ age
\emph{livremente} sobre o conjunto $X(T)$, ou seja, se $g\in G(T)$,
$x\in X(T)$ e $g\cdot x=x$ implicam que $g$ é o elemento identidade de
$G(T)$. Diz-se então que $G$ \emph{age livremente sobre $X$} (ou que $X$ é um
\emph{$\mathcal C$-espaço principal sob $G$}). A relação de equivalência
associada ao par $(p_1,p_2)$ chama-se então \emph{relação de equivalência
definida pelo grupo $G$ que age livremente sobre $X$}. Quando $X\times X$
também existe e se considera o morfismo
\[
  p:G\times X\longrightarrow X\times X
\]
definido pelo par $(p_1,p_2)$, a condição de que $G$ aja livremente significa
que $p$ é um \emph{monomorfismo}.

Naturalmente, mesmo quando $G$ não age livremente sobre $X$, procuram-se
critérios para a existência de um quociente de $X$ por $G$, isto é, do
conúcleo do par anterior $(p_1,p_2)$.

O conúcleo em questão será muitas vezes denotado por $X/G$, ou,
preferencialmente, por $G\backslash X$ se $G$ agir à esquerda (reservando-se a
notação anterior para o caso em que $G$ age à direita). Observe-se que, mesmo
quando existe a «imagem» de $G\times X$ por $p$ (definida, por exemplo, como o
menor subobjeto de $X\times X$ pelo qual $p$ se fatora), digamos $R$, ela quase
nunca é uma relação de equivalência em $X$. Se então se tenta passar
diretamente ao quociente por $R$ (ou, mais precisamente, pelo par de morfismos
de $R$ em $X$ induzidos pelas duas projeções $\operatorname{pr}_i$), perdem-se
irremediavelmente as características particulares do par original
$(p_1,p_2)$. É, portanto, importante encontrar uma generalização da noção de
relação de equivalência que se aplique diretamente ao par definido por um
grupo de operadores de $\mathcal C$.

\subsection*{4. Pré-relações de equivalência}
\addcontentsline{toc}{subsection}{4. Pré-relações de equivalência}

Recordemos que uma categoria na qual todo morfismo é um isomorfismo se chama
\emph{grupoide}. Por outro lado, uma categoria deve ser definida como formada
por dois conjuntos subjacentes $(X,R)$, o conjunto de \emph{objetos} e o
conjunto de \emph{setas}, munidos das seguintes estruturas:
\src{106}
\begin{enumerate}
\item[(i)] Um par de aplicações
\[
  p_1,p_2:R\rightrightarrows X
\]
chamadas \emph{aplicação fonte} e \emph{aplicação alvo},

\item[(ii)] Uma aplicação
\[
  \pi:(R,p_2)\times_X(R,p_1)\longrightarrow R
\]
chamada \emph{aplicação de composição}.
\end{enumerate}
Esses dados devem satisfazer axiomas bem conhecidos, que não repetiremos aqui
e que poderiam ser explicitados pela comutatividade de certos diagramas e pela
existência de uma aplicação $D:X\to R$ (necessariamente única) que torna
comutativos outros dois diagramas. $D$ corresponde à passagem de um objeto à
seta identidade correspondente e satisfaz
\[
  p_1D=p_2D=\operatorname{id}_X.
\]
Dizer que a categoria é um grupoide significa então que existe uma aplicação
\[
  s:R\longrightarrow R
\]
(necessariamente única), chamada a \emph{inversão} de $R$, que transforma cada
seta em uma seta inversa, o que poderia ser explicitado pela comutatividade de
outros quatro diagramas, formados com $s$, $\Delta$ e os dados precedentes, e
dos quais os dois primeiros se escrevem:
\[
  p_1s=p_2,\qquad p_2s=p_1.
\]

Recordados esses pontos, as definições gerais em ([2], A, \S~1) mostram, em
particular, o que se deve entender por uma estrutura de
$\mathcal C$-\emph{categoria}, respectivamente de
$\mathcal C$-\emph{grupoide}, sobre um par de objetos $X,R$ de uma categoria
qualquer $\mathcal C$: trata-se, por definição, do dado, para todo objeto $T$
de $\mathcal C$, de uma estrutura de categoria, respectivamente de grupoide,
no sentido conjuntista, cujo conjunto de objetos é $X(T)$ e o conjunto de
setas é $R(T)$, essas estruturas
\guillemotleft variando funtorialmente\guillemotright{} com $T$ variável.
Isso implica, portanto, a definição de dois morfismos
\[
  p_1,p_2:R\rightrightarrows X
\]
chamados \emph{morfismo fonte} e \emph{morfismo alvo} e, quando existe o
produto fibrado desejado, de um morfismo
\src{107}
\[
  \pi:(R,p_2)\times_X(R,p_1)\longrightarrow R
\]
chamado \emph{morfismo de composição}. Esses três dados bastam então para
determinar a estrutura de categoria sobre $X,R$, sendo o axioma a impor a
esses dados o seguinte: para todo $T$, os três dados correspondentes para
$X(T),R(T)$ definem sobre esse par de conjuntos uma estrutura de categoria
(respectivamente, de grupoide). Quando for o caso, isso pode ser expresso pela
comutatividade de certos diagramas que envolvem um morfismo univocamente
determinado
\[
  D:X\longrightarrow R
\]
e, no caso dos grupoides, um morfismo univocamente determinado
\[
  s:R\longrightarrow R;
\]
diagramas que se explicitam como no caso «conjuntista». Essa interpretação
trabalhosa dos axiomas é, felizmente, inútil na prática. O único interesse
teórico da possibilidade de expressar os dados e os axiomas a partir de
morfismos e igualdades de morfismos entre certos produtos fibrados é o
seguinte: se temos um funtor $F:\mathcal C\to\mathcal C'$ exato à esquerda
(isto é, que comuta com produtos finitos e produtos fibrados), ele transforma
uma $\mathcal C$-categoria (respectivamente, um $\mathcal C$-grupoide) em uma
$\mathcal C'$-categoria (respectivamente, um $\mathcal C'$-grupoide), desde
que existam produtos finitos e produtos fibrados em $\mathcal C$.

Na prática, é importante saber interpretar os morfismos
$p_1,p_2,\pi,D,s$ como \emph{operações simpliciais} em um objeto
semissimplicial ou simplicial adequado de $\mathcal C$, ao menos quando
existem produtos fibrados em $\mathcal C$. Para fixar a terminologia,
introduzamos a \emph{categoria $\mathcal S$ dos simplexos padrão} como a
categoria cujos objetos são os conjuntos finitos da forma
\[
  \Delta_n=[0,n],\qquad\text{onde }n\geq 0
\]
(intervalo dos inteiros de $0$ a $n$), e cujos morfismos ou setas são as
\emph{aplicações arbitrárias} entre esses conjuntos finitos. Observe-se que a
categoria $\mathcal S$ é equivalente à categoria dos conjuntos finitos
\emph{não vazios}, tomando-se como morfismos as aplicações entre conjuntos
finitos. Em $\mathcal S$, existe evidentemente o coproduto de uma família
finita \emph{não vazia} de objetos, bem como o coproduto fibrado de dois
objetos sob um terceiro (operação dual ao produto fibrado). Denotemos por
$\mathcal S'$ a subcategoria de $\mathcal S$ que tem os mesmos objetos, mas
cujos morfismos são as \emph{aplicações crescentes} entre os $\Delta_n$. Essa
categoria é equivalente à categoria dos conjuntos finitos não vazios
totalmente ordenados. Nesta categoria, o
\src{108}
coproduto de dois objetos nunca existe, e o coproduto fibrado de dois objetos
$A,B$ sob um terceiro $C$ não existe em geral (tome-se, por exemplo,
$C=\Delta_0$, $A=B=\Delta_1$, com as duas aplicações estruturais
$u:C\to A$ e $v:C\to B$ iguais). Contudo, ele existe em certos casos; por
exemplo,
\[
  C=\Delta_0,\qquad A=\Delta_m,\qquad B=\Delta_n,\qquad
  u(0)=m,\qquad v(0)=0,
\]
e, nesse caso,
\[
  \Delta_m\amalg_{\Delta_0}\Delta_n=\Delta_{m+n}.
\]

Um \emph{objeto simplicial} (respectivamente, um \emph{objeto
semissimplicial}) em uma categoria $\mathcal C$ é, por definição, um funtor
contravariante $K$ de $\mathcal S$ (respectivamente, de $\mathcal S'$) em
$\mathcal C$. Portanto, um objeto simplicial define por restrição um objeto
semissimplicial, mas a primeira noção difere essencialmente da segunda pela
presença de \emph{operações de simetria} nos $K_n=K(\Delta_n)$,
correspondentes às imagens, pelo funtor $K$, dos elementos do grupo simétrico
de $n+1$ letras (considerado como o grupo de automorfismos de $\Delta_n$ em
$\mathcal S$).

Para todo $n$, seja $\Delta'_n$ (respectivamente, $\Delta''_n$) a categoria
finita cujo conjunto de objetos é $\Delta_n$, definida pela relação de ordem
caótica (respectivamente, pela ordem total natural) em $\Delta_n$ (sendo o
conjunto de setas o grafo dessa relação de ordem). É evidente que $\Delta'_n$
(respectivamente, $\Delta''_n$) depende funtorialmente do objeto $\Delta_n$ de
$\mathcal S$ (respectivamente, de $\mathcal S'$). Se $Z$ for uma categoria,
então $\operatorname{Hom}(\Delta'_n,Z)$ (respectivamente,
$\operatorname{Hom}(\Delta''_n,Z)$), quando $\Delta_n$ varia, é um funtor de
$\mathcal S$ (respectivamente, de $\mathcal S'$) para a categoria dos
conjuntos, isto é, um conjunto simplicial (respectivamente, um conjunto
semissimplicial), dito \emph{associado à categoria $Z$}; denotemo-los por $Z'$
e $Z''$. Há, além disso, um homomorfismo natural evidente do conjunto
semissimplicial associado a $Z'$ para $Z''$, e este último é um isomorfismo se,
e somente se, $Z$ for um grupoide. Posto isso:

\result{PROPOSIÇÃO}{4.1} O funtor $Z\longmapsto Z''$ da categoria das
categorias para a categoria dos conjuntos semissimpliciais é plenamente fiel
e define uma equivalência entre a categoria das \emph{categorias} e a
categoria dos conjuntos semissimpliciais, isto é, dos funtores contravariantes
$K$ de $\mathcal S'$ para $(\mathrm{Set})$ que transformam os coprodutos
fibrados $A\amalg_C B$ do tipo especificado acima em produtos fibrados de
conjuntos. Do mesmo modo, o funtor $Z\longmapsto Z'$ da categoria dos
grupoides para a categoria dos conjuntos simpliciais é plenamente fiel e
define uma equivalência entre a categoria dos \emph{grupoides} e a categoria
dos
\src{109}
conjuntos simpliciais, isto é, funtores contravariantes $K$ de $\mathcal S$
em $(\mathrm{Ens})$, que transformam somas amalgamadas em produtos fibrados.

Pode-se, portanto, considerar as categorias como conjuntos semissimpliciais
particulares, e os grupoides como conjuntos simpliciais particulares, com a
condição, naturalmente, de raciocinar «a menos de isomorfismo», como é de rigor
quando se interpretam certas estruturas em termos de outras. O procedimento
usual de redução ao caso conjuntista implica então:

\result{COROLÁRIO}{4.2} O enunciado precedente continua válido quando se
substituem as categorias, os grupoides e os conjuntos simpliciais por
$\mathcal C$-categorias, $\mathcal C$-grupoides e $\mathcal C$-objetos
simpliciais, contanto que existam produtos fibrados em $\mathcal C$.

O objeto semissimplicial $K$ em $\mathcal C$ associado a uma categoria
$(X,R,\ldots)$ em $\mathcal C$ pode ser explicitado interpretando-se a
componente $K_n=K(\Delta_n)$ de $K$ como a $(n+1)$-ésima potência fibrada de
$(R,p_1)$ sobre $X$, ou melhor, pela fórmula de recorrência
\[
  K_0=R,\qquad
  K_n=(K_{n-1},p_n^{(n-1)})\times_X(R,p_1),
\]
onde os $p_i^{(n-1)}$ ($0\leq i\leq n-1$) são as projeções naturais de
$K_{n-1}$ em $X$ (que também se definem por recorrência). Dessa forma,
$p_1,p_2,\pi,D,s$ são interpretados como as operações simpliciais que
correspondem respectivamente aos morfismos em $\mathcal S$: a face $0$ de
$\Delta_1$, a face $1$ de $\Delta_1$, a face $(0,2)$ de $\Delta_2$, a
degenerescência $\Delta_1\to\Delta_0$ e a simetria de $\Delta_1$. Todas as
demais operações semissimpliciais (respectivamente simpliciais) deduzem-se
formalmente das quatro (respectivamente cinco) precedentes por composição e
produtos fibrados.

Chama-se agora \emph{pré-relação de equivalência} em um objeto $X$ de uma
categoria o dado de um grupoide cujo objeto dos objetos (se assim se pode
dizer) é $X$. Tal dado compreende, portanto, entre outras coisas, um objeto $R$
e dois morfismos
\[
  p_1,p_2:R\rightrightarrows X.
\]
Observe-se, porém, que esses dados, por si sós, não determinam a estrutura
considerada, ao contrário do que ocorre com os pares de equivalência. Nesta
exposição, interessamo-nos por essa noção com o objetivo de obter critérios de
possibilidade de passagem ao quociente, isto é, de formação do conúcleo do par
$(p_1,p_2)$. O enunciado do problema, portanto, não faz uso dos dados
adicionais inerentes a um grupoide. Na demonstração dos resultados que se
seguem,
\src{110}
utilizaremos, contudo, esses dados adicionais e, em particular, as operações
simpliciais (incluindo as operações de simetria) até a dimensão $3$ (nas quais
intervém a quarta potência fibrada de $R$ sobre $X$).

Uma relação de equivalência em um objeto $X$ de $\mathcal C$ define uma
pré-relação de equivalência: basta verificá-lo no caso conjuntista, no qual se
associa a uma relação de equivalência em um conjunto $X$ o grupoide cujo
conjunto de objetos é $X$ e cujo conjunto de setas é o grafo da relação de
equivalência.

Um $\mathcal C$-monoide $G$ que opera sobre um objeto $X$ de $\mathcal C$
define uma $\mathcal C$-categoria cujos objetos subjacentes são $R=G\times X$
e $X$ (desde que $G\times X$ exista), e que é um $\mathcal C$-grupoide se, e
somente se, $G$ for um grupo. Basta, ainda, fazer a verificação no caso
conjuntista. Define-se então a composição das setas $(g,a)$ e
$(g',g\cdot a)$ por
\[
  (g',g\cdot a)(g,a)=(g'g,a),
\]
isto é, se $a,b\in X$, então $\operatorname{Hom}(a,b)$ é, por definição, o
transportador de $a$ em $b$, e os morfismos se compõem mediante a composição
dos elementos de $G$.

\medskip
\noindent\textbf{OBSERVAÇÃO. ---} Podem-se evitar as dificuldades lógicas
levantadas por um enunciado como (4.1), subentendendo-se que todos os objetos
considerados se encontram em um «Universo» fixado (que é ele próprio um
conjunto).

\fsection{5}{Quociente por uma relação de equivalência finita e plana}

\result{TEOREMA}{5.1} Sejam $X=\operatorname{Spec}(B)$ um esquema afim e $R$
uma pré-relação de equivalência em $X$, cuja componente $R_1$ é afim, digamos
$R_1=\operatorname{Spec}(C)$. Suponha-se que a primeira projeção
$p_1:R_1\to X$ seja um morfismo finito e localmente livre; em outras palavras,
que o homomorfismo de anéis correspondente $p_1^*:B\to C$ faça de $C$ um
$B$-módulo projetivo de tipo finito. Seja $A$ o subanel de $B$ que é o núcleo
do par de homomorfismos
\[
  p_1^*,p_2^*:B\rightrightarrows C
\]
(isto é, o conjunto dos elementos $b$ tais que
$p_1^*(b)=p_2^*(b)$). Seja $Y=\operatorname{Spec}(A)$, e seja $f:X\to Y$ o
morfismo definido pela inclusão de $A$ em $B$. Nessas condições:
\begin{enumerate}
\item[(i)] $B$ é integral sobre $A$, isto é, $f$ é um morfismo integral.

\item[(ii)] O morfismo $f$ é sobrejetivo, e suas fibras são as classes de
equivalência
\src{111}
conjuntistas $p_2\allowbreak(p_1^{-1}(x))$ em $X$ módulo $R$, e a topologia
de $Y$ é a topologia quociente da de $X$.

\item[(iii)] $Y$ é o quociente de $X$ por $R$ na categoria dos pré-esquemas.

\item[(iv)] Se $R$ provém de uma \emph{relação de equivalência}, então o
morfismo $f:X\to Y$ é finito e localmente livre (isto é, $B$ é um $A$-módulo
projetivo finitamente gerado), e a relação de equivalência é efetiva, isto é,
$R_1\to X\times_Y X$ é um isomorfismo.
\end{enumerate}

Este teorema generaliza o conhecido teorema relativo ao caso de um grupo
finito $G$ que atua por automorfismos no anel $B$ e ao anel de invariantes
$A$, e sua demonstração é análoga à demonstração conhecida. Pode-se precisar
(iii) da seguinte maneira:

\result{COROLÁRIO}{5.2} O morfismo canônico
$R_1\to X\times_Y X$ é \emph{sobrejetivo}.

Seja ainda $R$ uma pré-relação de equivalência «finita e localmente livre» em
$X$, sendo agora $X$ um pré-esquema arbitrário. Suponhamos que se possa
encontrar um pré-esquema $Y$ e um morfismo $f:X\to Y$ tais que $fp_1=fp_2$ e,
além disso, que seja exata a sequência de homomorfismos de feixes de anéis
sobre $Y$:
\[
  \mathcal O_Y\longrightarrow f_*(\mathcal O_X)
  \rightrightarrows g_*(\mathcal O_R)
\]
(onde $g=fp_i$). Resulta então do teorema que se têm conclusões (i) a (iv)
análogas às do teorema; em particular, por (iii), $Y$ é o quociente de $X$ por
$R$ e, portanto, fica determinado a menos de isomorfismo único. Diz-se, nessas
condições, que a pré-relação de equivalência $R$ em $X$ é \emph{admissível}.
Com essa definição:

\result{TEOREMA}{5.3} Sejam $X$ um pré-esquema e $R$ uma pré-relação de
equivalência em $X$ tal que $p_1:R_1\to X$ seja um morfismo finito e localmente
livre. Para que $R$ seja admissível, é necessário e suficiente que toda classe
de equivalência conjuntista $p_2(p_1^{-1}(x))$ em $X$ módulo $R$ esteja
contida em um aberto afim (condição sempre verificada se todo subconjunto
finito de $X$ estiver contido em um aberto afim, por exemplo, se $X$ for
quase-projetivo sobre um esquema afim).

Mostra-se, com efeito, sem dificuldade que toda classe de equivalência módulo
$R$ em $X$ está então contida em um aberto afim \emph{estável por $R$}, e
constrói-se o quociente $Y$ pela colagem das partes obtidas com o auxílio do
Teorema 5.1.
\src{112}

\result{COROLÁRIO}{5.4} Suponhamos satisfeita essa condição e que, além disso,
$R$ provenha de uma relação de equivalência. Então esta última é efetiva, isto
é, $R_1\to X\times_Y X$ é um isomorfismo, e $f:X\to Y$ é um morfismo finito
localmente livre.

Conclui-se imediatamente, por «descida»:

\result{COROLÁRIO}{5.5} Sob as condições de (5.4), $X$ tem posto $n$ em todo
ponto sobre $Y$ se e somente se $(R_1,p_1)$ tem posto $n$ em todo ponto sobre
$X$. Se $X,R_1$ forem $Z$-pré-esquemas e $p_1,p_2$ forem $Z$-morfismos, sendo
portanto $Y$ um $Z$-pré-esquema, então $X$ é plano sobre $Z$ se e somente se
$Y$ o for.

Em resumo:

\result{ESCÓLIO}{} O dado de um morfismo finito, localmente livre e
sobrejetivo $f:X\to Y$ de pré-esquemas equivale ao dado de um pré-esquema $X$
munido de uma relação de equivalência $R$ tal que $p_1:R\to X$ seja finito e
localmente livre, e que toda classe $p_2(p_1^{-1}(x))$ esteja contida em um
aberto afim.

\medskip
\noindent\textbf{OBSERVAÇÕES 5.6.}

\smallskip
\noindent $1^\circ$ Não foi necessário fazer nenhuma hipótese Noetheriana.

\smallskip
\noindent $2^\circ$ Esta noção de passagem ao quociente inclui, como caso
particular, a «descida inseparável» de CARTIER, que corresponde à determinação
dos morfismos finitos localmente livres $f:X\to Y$ tais que
$f_*(\mathcal O_X)$ admita uma $p$-base em relação a $\mathcal O_Y$ (sendo $X$
um pré-esquema dado cujo feixe $\mathcal O_X$ é anulado pelo número primo
$p>0$). Observe-se que o resultado se exprime facilmente sem hipótese de
regularidade sobre os anéis locais e sem supor que $X$ seja um esquema
algébrico sobre um corpo. A teoria de JACOBSON--BOURBAKI é obtida tomando-se
como $X$ o espectro de um corpo de característica $p$.

\smallskip
\noindent $3^\circ$ GABRIEL havia obtido anteriormente um caso particular do
teorema (5.3) na teoria da passagem ao quociente para grupos comutativos
finitos sobre um corpo $k$. [Compare-se (7.4).]

\fsection{6}{Quociente por uma relação de equivalência própria e plana}

\result{TEOREMA}{6.1} Sejam $S$ um pré-esquema localmente Noetheriano, $X$
um $S$-esquema quasi-projetivo, $R$ uma pré-relação de equivalência em $X$,
tal que:
\src{113}
\begin{enumerate}
\item[(a)] $p_1:R_1\to X$ é próprio e plano;
\item[(b)] $R_1\to X\times_S X$ é um morfismo finito ou, o que equivale em
virtude de (a), tem fibras finitas (condição automaticamente satisfeita se
$R$ provém de uma relação de equivalência).
\end{enumerate}
Nessas condições:
\begin{enumerate}
\item[(i)] $X/R=Y$ existe e (se $S$ é Noetheriano) é quase-projetivo sobre $S$.

\item[(ii)] O morfismo canônico $f:X\to Y$ é sobrejetivo, próprio e aberto;
suas fibras são as classes de equivalência $p_2(p_1^{-1}(x))$ em $X$ módulo
$R$; portanto, $Y$ identifica-se com o espaço topológico quociente de $X$ pela
relação de equivalência conjuntista definida por $R$. Por fim,
$R_1\to X\times_Y X$ é sobrejetivo.

\item[(iii)] Se $R$ provém de uma relação de equivalência, esta última é
efetiva, isto é, $R_1\to X\times_Y X$ é um isomorfismo e, além disso,
$f:X\to Y$ é plano (portanto, fielmente plano).
\end{enumerate}

Para a demonstração, reduz-se a (5.1) por meio de quase-seções adequadas de $X$
para $R$, sendo a demonstração análoga à construção dos grupos algébricos
quocientes no Seminário Chevalley.

Em resumo:

\result{ESCÓLIO}{} Seja $X$ quase-projetivo sobre $S$ localmente Noetheriano. O
dado de um morfismo próprio, fielmente plano e sobrejetivo $f:X\to Y$ de $X$
para um $S$-pré-esquema $Y$ equivale ao dado de uma relação de equivalência $R$
em $X$, tal que $p_1:R\to X$ seja próprio e plano.

O mesmo método fornece o resultado seguinte:

\result{TEOREMA}{6.2} Sejam $S$ um pré-esquema Noetheriano, $X$ um
pré-esquema de tipo finito sobre $S$ e $R$ uma pré-relação de equivalência no
$S$-pré-esquema $X$; supõe-se que:
\begin{enumerate}
\item[(a)] $p_1:R_1\to X$ é plano e de tipo finito;
\item[(b)] o morfismo $R_1\to X\times_S X$ é quase-finito (isto é, tem fibras
finitas).
\end{enumerate}
Então existe um aberto $U$ denso em $X$, \emph{saturado} para $R$, tal que:
\begin{enumerate}
\item[(i)] Se $R_U$ é a pré-relação de equivalência induzida por $R$ em $U$,
então $U/R_U$ existe e é de tipo finito sobre $S$.

\item[(ii)] O morfismo canônico $U\to U/R_U$ é sobrejetivo e aberto; suas
fibras são as classes de equivalência conjuntistas para $R_U$ (portanto,
$U/R_U$ é o espaço topológico quociente de $U$ pela relação de equivalência
conjuntista definida por $R_U$); por fim, o morfismo
\[
  (R_U)_1\longrightarrow U\times_{U/R_U}U
\]
é sobrejetivo.
\src{114}

\item[(iii)] Se $R$ provém de uma relação de equivalência, pode-se supor que
$U\to U/R_U$ seja fielmente plano e que $R_U$ seja efetiva.
\end{enumerate}

Este é um resultado de natureza essencialmente «birracional».

\medskip
\noindent\textbf{OBSERVAÇÕES 6.3.}

\smallskip
\noindent $1^\circ$ Ignoro se, em (6.1) e (6.2), a hipótese (b) é inútil. Na
passagem ao quociente por grupos, ela nos obriga praticamente a nos
restringirmos ao caso em que os estabilizadores são todos grupos finitos.

\smallskip
\noindent $2^\circ$ Pode-se perguntar se não há resultados análogos a (6.1) e
(6.2) sem hipótese de planicidade. Não conheço nenhum contraexemplo nessa
direção. Em contrapartida, mesmo conservando a hipótese de planicidade e
restringindo-nos a relações de equivalência tais que $p_1:R\to X$ seja plano e
quase-finito (mas não finito), sendo $X$ afim, pode ocorrer que $R$ não seja
efetiva: tomem-se as relações de equivalência induzidas nos abertos afins que
recobrem a variedade de Nagata (na qual um grupo de ordem dois opera de modo
«não admissível»).

\fsection{7}{Aplicações}

Como dissemos na introdução, a aplicação mais importante de (6.1) é a
construção de esquemas de Picard, bem como de soluções para vários outros
problemas de moduli, aos quais voltaremos posteriormente.

Obtém-se uma demonstração simples do seguinte resultado de SHIMURA:

\result{PROPOSIÇÃO}{7.1} Seja $A$ um esquema abeliano definido sobre um anel de
valorização discreta $V$, com corpo de frações $K$. Então todo esquema
abeliano $B'$ sobre $K$, isógeno a um quociente de $A\otimes_V K$, «tem boa
redução sobre $V$», isto é, é isomorfo a algum $B\otimes_V A$,
 onde $B$ é um esquema
abeliano sobre $V$ (essencialmente único, recordemo-lo).

Pode-se supor que $B'$ seja o quociente de $A_K$ por um subesquema em grupos
$C'$. (N.B.: $C'$ não será, em geral, «reduzido», isto é, seus anéis locais
terão elementos nilpotentes.) Considere-se o subesquema fechado $C$ de $A$ que
é o «fecho» de $C'$, isto é, o menor subesquema fechado de $A$ tal que $C_K$
contenha $C'$. Teremos $C_K=C'$, e do fato de que $V$ é um anel de valorização
discreta deduz-se facilmente que $C$ é um subesquema em grupos de $A$ sobre
$V$. Como $A$ é próprio sobre $\operatorname{Spec}(V)=S$, o mesmo vale para
$C$; por outro lado, $A$ é até projetivo sobre $S$. Pode-se então aplicar
(6.1) para construir $A/C=B$,
\src{115}
que é o $B$ procurado.

Por fim, argumentos essencialmente conhecidos permitem deduzir de (6.2) o
seguinte resultado:

\result{TEOREMA}{7.2} Sejam $S$ o espectro de um anel Artiniano, $F$ e $G$
esquemas em grupos de tipo finito sobre $S$, e $u:F\to G$ um homomorfismo de
esquemas em grupos sobre $S$. Suponhamos:
\begin{enumerate}
\item[(i)] $F$ é plano sobre $S$;
\item[(ii)] o núcleo de $u$ é finito.
\end{enumerate}
Nessas condições, existe o esquema quociente $G/F$; o morfismo canônico
$G\to G/F$ é sobrejetivo e aberto, e suas fibras são as classes de equivalência
conjuntistas definidas pelas ações à direita de $F$ sobre $G$. Por fim, se $u$
é um monomorfismo, então $G\to G/F$ é plano e o morfismo
\[
  G\times F\longrightarrow G\times_{G/F}G
\]
é um isomorfismo; em outras palavras, $G$ é um espaço homogêneo principal
sobre $G/F$, com grupo estrutural (atuando à direita) $F$, ou melhor,
$F\times_S(G/F)$ considerado como esquema em grupos sobre $G/F$ (cf. [1], B,
\S~6).

\result{COROLÁRIO}{7.3} Nessas condições, para que $G$ seja plano sobre $S$,
é necessário e suficiente que $G/F$ o seja. Se essa condição for satisfeita,
a passagem ao quociente por $F$ comuta com toda extensão da base $S$ e, se
$F$ for um subgrupo normal de $G$, $G/F$ poderá ser munido de uma estrutura
de \emph{grupo quociente} de $G$ por $F$.

A situação é particularmente simples quando $S$ é o espectro de um corpo,
pois então todo $S$-pré-esquema é automaticamente plano sobre $S$. Obtém-se:

\result{COROLÁRIO}{7.4} Sejam $F,G$ esquemas em grupos de tipo finito sobre um
corpo $k$, e seja $u:F\to G$ um homomorfismo de $k$-grupos. Então $u$ se
fatora como $F\to F'\to G$, onde $F\to F'$ é o homomorfismo de passagem ao
quociente pelo subgrupo fechado $\operatorname{Ker}u$ de $F$, e $F'\to G$ é
um homomorfismo de grupos que é uma imersão fechada. Existe o quociente
$G/F=G/F'$. O formalismo usual (do tipo dos teoremas de Noether) é válido
entre os grupos algébricos sobre $k$.

Este resultado permite tratar de maneira uniforme a passagem ao quociente nos
grupos algébricos no sentido clássico (isto é, os que são irredutíveis sobre
$k$ e suaves sobre $k$) e a passagem ao quociente por subgrupos
«infinitesimais»
\src{116}
considerados por CARTIER. É vantajoso considerar as «hiperálgebras»
introduzidas por CARTIER, na esteira dos trabalhos de DIEUDONNÉ sobre os grupos
formais, como grupos na categoria dos esquemas formais sobre $k$ e, quando for
o caso (se corresponderem a hiperálgebras de posto finito sobre $k$), como
grupos algébricos finitos sobre $k$.

\fsection{8}{Uma conjectura}

Ela responde à necessidade de saber passar ao quociente pelo grupo projetivo
que atua em certos subesquemas dos «esquemas de Hilbert» (estes últimos
substituindo, na teoria dos esquemas, as variedades de Chow).

Sejam $S$ um pré-esquema e $n$ um inteiro. A todo pré-esquema $S'$ sobre $S$,
façamos corresponder o grupo
$\operatorname{Gl}(n,\Gamma(S',\mathcal O_{S'}))$ das matrizes invertíveis de
$n$ linhas e $n$ colunas com entradas no anel das seções de
$\mathcal O_{S'}$. Obtém-se assim um funtor contravariante em $S'$, que se
mostra facilmente representável; ele corresponde, portanto, a um esquema em
grupos sobre $S$, de fato afim sobre $S$, denotado por $\operatorname{Gl}(n)_S$.
Sua formação é compatível com a mudança de base, de modo que, na realidade,
tudo provém de um esquema em grupos sobre $\mathbf Z$, denotado por
$\operatorname{Gl}(n)$. O grupo $\operatorname{Gl}(1)$, chamado grupo
multiplicativo e frequentemente denotado por $\mathbf G_m$, corresponde ao
funtor $S\longmapsto\Gamma(S,\mathcal O_S)^*$, grupo das «unidades» sobre $S$.
Há um homomorfismo evidente
$\operatorname{Gl}(1)\to\operatorname{Gl}(n)$, e constrói-se facilmente o
grupo quociente, denotado por $\operatorname{GP}(n-1)$ e chamado
\emph{grupo projetivo de grau $n-1$ sobre $\mathbf Z$}. Ele representa o
funtor que associa a $S$ o grupo das seções do feixe
$\widetilde{\operatorname{Gl}(n)}_S/
\widetilde{\operatorname{Gl}(1)}_S$, onde
$\widetilde{\operatorname{Gl}(n)}_S$ denota o feixe dos germes de seções de
$\operatorname{Gl}(n)_S$ sobre $S$. (É preciso observar que as seções de
$\operatorname{GP}(n-1)_S$ sobre $S$ não provêm, em geral, de seções de
$\operatorname{Gl}(n)_S$ sobre $S$!) Note-se que se pode provar que
$\operatorname{GP}(n-1)$ representa também o funtor
$S\longmapsto\operatorname{Aut}_S(\mathbf P_S^{n-1})$, onde
$\mathbf P_S^{n-1}$ é o espaço projetivo padrão de dimensão relativa $n-1$
sobre $S$, ao menos quando $S$ é Noetheriano. É dessa maneira que ele intervém
na teoria de módulos.

Seja $S$ um esquema Noetheriano, que se pode, se desejado, supor afim, e seja
$X$ um $S$-pré-esquema quase-projetivo, munido de um feixe invertível
$\mathcal L$ muito amplo relativamente a $S$. Suponhamos que o grupo
$G=\operatorname{GP}(n)_S$ atue em $X$ e, ao mesmo tempo, em $\mathcal L$ (de
modo compatível com sua ação em $X$), e que atue \emph{livremente} em
$S$.
\src{117}

\result{CONJECTURA}{8.1} Nas condições precedentes:

\smallskip
\noindent $1^\circ$ \emph{A relação de equivalência definida por $G$ é
efetiva, o quociente $Y=X/G$ é de tipo finito sobre $S$ e o morfismo canônico
$f:X\to Y$ é plano e sobrejetivo} (portanto, $X$ torna-se um fibrado principal
homogêneo sobre $Y$, de grupo
$G\times_S Y=\operatorname{GP}(n)_Y$).

\smallskip
\noindent $2^\circ$ Seja $\mathcal L'$ o feixe invertível sobre $Y$ deduzido
de $\mathcal L$ por «descida fielmente plana» ao longo de $f$ (cf. [1], B,
\S~th. 1). Então $\mathcal L'$ é \emph{pré-amplo sobre $Y$ relativamente a
$S$}, isto é, existem um inteiro $m$ e um morfismo quase-finito de $Y$ para um
esquema projetivo padrão conveniente $\mathbf P_S^N$, tais que
$(\mathcal L')^{\otimes m}$ seja isomorfo à imagem inversa de
$\mathcal O_{\mathbf P_S^N}(1)$.

Observe-se que, mesmo se $X$ é separado sobre $S$, pode ocorrer que $Y$ não
seja separado sobre $S$ (situação que ocorre em problemas de módulos nada
patológicos). Se $1^\circ$ se verifica, então $Y$ é separado se e somente se a
relação de equivalência definida por $G$ tem gráfico fechado, isto é, se
$G\times X\to X\times X$ tem imagem fechada (trata-se então de uma imersão
fechada). Se $Y$ é separado, então $\mathcal L'$ é pré-amplo sobre $Y$
relativamente a $S$ se e somente se é amplo, isto é, se uma potência tensorial
adequada define uma imersão projetiva. Nos problemas de módulos mencionados na
introdução, pode-se mostrar que a relação de equivalência obtida tem de fato um
gráfico fechado.

\medskip
\noindent\textbf{OBSERVAÇÕES 8.1. ---} Supusemos
$G=\operatorname{GP}(n)_S$ para fixar as ideias e porque este é atualmente o
caso mais importante na prática. A hipótese razoável a fazer sobre $G$
pareceria antes que $G$ fosse uma das «formas» sobre $S$ de um dos grupos de
Tohokû (cuja construção sobre os inteiros também foi realizada por CHEVALLEY).
O único fato positivo que conheço na direção da conjectura precedente é o
seguinte: seja $X$ \emph{um esquema afim sobre um corpo $k$ de característica
$0$, no qual o grupo $\operatorname{Gl}(n)_k$ ou
$\operatorname{GP}(n-1)_k$ atua livremente. Então a relação de equivalência
definida por $G$ é efetiva, o quociente $X/G$ também é afim e o morfismo
$X\to X/G$ é plano e sobrejetivo}. A demonstração utiliza o fato seguinte (que
por enquanto só está demonstrado em característica $0$): quando se faz $G$
atuar sobre o anel de coordenadas afim $A$ de $G$, considerado como espaço
vetorial sobre $k$, a representação trivial de $G$ ocorre apenas uma vez (em
uma série de composição de um subespaço vetorial de dimensão finita sobre $k$,
estável sob $G$). Parece possível que um uso sistemático da
\src{118}
a teoria das representações lineares de $G$ acabará por fornecer uma
demonstração da conjectura, ao menos quando se está sobre um corpo de base.
Quando já não se está sobre um corpo de base, o conferencista nada sabe.

\section*{BIBLIOGRAFIA}
\addcontentsline{toc}{section}{Bibliografia}

\begin{description}
\item[{[1]}] GROTHENDIECK (Alexander). --- Techniques de descente et théorèmes
d'existence en géométrie algébrique, I: Généralités, Séminaire Bourbaki,
1959/60, n\textsuperscript{o} 190, 29 p.

\item[{[2]}] GROTHENDIECK (Alexander). --- Techniques de descente et théorèmes
d'existence en géométrie algébrique, II: Le théorème d'existence en théorie
formelle des modules, Séminaire Bourbaki, 1959/60,
n\textsuperscript{o} 195, 22 p.

\item[{[3]}] GROTHENDIECK (A.) e DIEUDONNÉ (J.). --- Éléments de géométrie
algébrique, I: Le langage des schémas. --- Paris, Presses universitaires de
France, 1960 (Institut des hautes Études scientifiques, Publications
mathématiques, 4).
\end{description}

\par

\def\FgaUnitKey{e212}\clearpage
\chapter*{Errata da exposição n.º 212}
\addcontentsline{toc}{chapter}{Errata da exposição n.º 212}
\setcounter{footnote}{2}
\markboth{Errata da exposição 212}{Errata da exposição 212}
\sid{e212}
\src{301}

\noindent\textbf{5. TDTE III: Pré-esquemas quocientes.}

\begin{center}
[Seminário Bourbaki, t. 13, 1960/61, n.º 212, 20 p.]
\end{center}

\noindent\textbf{Página 15, asserção (i).} --- O fato de que $Y=X/R$ seja
quase-projetivo sobre $S$ parece, por ora, ter sido demonstrado apenas quando
$R$ provém de uma relação de equivalência.

\noindent\textbf{Páginas 18 e 19.} --- Como assinalamos ao final da exposição
seguinte, a conjectura considerada é decididamente falsa. O «fato positivo»
assinalado na Observação 8.1 parece ter sido demonstrado simultaneamente por
diversos autores (NAGATA, ROSENLICHT, GROTHENDIECK, \ldots).

Assinalemos que a aplicação feita em TDTE V da teoria da passagem ao quociente
aqui desenvolvida à construção dos esquemas de Picard também pode ser
substituída por uma utilização adequada dos esquemas de Hilbert (cf. SHMT).
Como foi dito no § 8, a lacuna mais importante da teoria apresentada
\src{302}
aqui é a ausência de um critério de existência de quocientes por uma relação de
equivalência não própria, tais como as relações de equivalência provenientes de
certas ações do grupo projetivo. Um teorema importante nessa direção foi obtido
por MUMFORD\footnote{MUMFORD (David). --- An elementary theorem in geometric
invariant theory, \emph{Bull. Amer. Math. Soc.}, t. 67, 1961, pp. 483--486.}.
Para um refinamento de seu resultado e diversas aplicações à teoria, ver SHMT.

\par

\def\FgaUnitKey{221}\unit{Exposição 221 -- Os esquemas de Hilbert}
\sid{221}
\markboth{Exposição 221}{Exposição 221}
\src{249}
\seminares{13.º ano, 1960/61, n.º 221\\Maio de 1961}
\paperhead{TÉCNICAS DE CONSTRUÇÃO E TEOREMAS DE EXISTÊNCIA EM GEOMETRIA ALGÉBRICA\\
IV. OS ESQUEMAS DE HILBERT}%
  {por Alexander GROTHENDIECK}

\subsection*{Introdução}
\addcontentsline{toc}{subsection}{Introdução}

As técnicas expostas em [2], I e II, eram essencialmente independentes de
qualquer hipótese de projetividade sobre os esquemas considerados.
Infelizmente, elas não permitem atualmente resolver os problemas de existência
formulados na Exposição II de [2]. Na presente exposição e na seguinte,
resolveremos esses problemas sob hipóteses de projetividade. As técnicas
empregadas são tipicamente projetivas e praticamente não utilizam os
resultados de I e II de [2]. Aqui construiremos os «esquemas de Hilbert»,
destinados a substituir o uso das coordenadas de Chow, como foi dito em [2],
II, n.º 2. Na exposição seguinte, a teoria da passagem ao quociente em
esquemas desenvolvida em III, juntamente com a teoria dos esquemas de Hilbert,
nos permitirá, por exemplo, construir os esquemas de Picard (definidos em [2],
II, n.º 3) sob condições bastante gerais.

Em resumo, pode-se dizer que agora se dispõe de uma técnica de construções
projetivas bastante satisfatória, salvo que ainda falta\footnote{Ver o adendo
ao fim da exposição.} um teorema de passagem ao quociente por grupos como o
grupo projetivo, atuando «sem pontos fixos» (cf. [2], III, n.º 8). A situação
parece até um pouco melhor em geometria analítica (quando nos limitamos ao
estudo dos espaços analíticos «projetivos» sobre um espaço analítico dado),
pois, para os espaços analíticos, desaparece a dificuldade da passagem ao
quociente por um grupo que atua de maneira adequada. Além disso, tanto em
geometria algébrica como em geometria analítica, restaria desenvolver uma
técnica de construção válida sem hipóteses de projetividade.

\fsection{1}{Famílias limitadas de feixes: propriedades de permanência}

Sejam $k$ um corpo e $X$ um $k$-pré-esquema, de tipo finito, para simplificar.
Para toda extensão $K/k$, obtém-se, portanto, um $K$-pré-esquema
$X_K=X\otimes_k K$. Se $\mathcal F$ é um feixe coerente em $X_K$ e $K'$ é
uma extensão de $K$, então $\mathcal F\otimes_K K'=\mathcal F_{K'}$ é um
feixe quase-coerente em $X_K\otimes_K K'=X_{K'}$. Posto isso, se $K$ e $K'$
são duas extensões quaisquer de $k$, $\mathcal F$ é um feixe quase-coerente em
$X_K$ e $\mathcal F'$
\src{250}
um feixe quase-coerente sobre $X_{K'}$, diremos que $\mathcal F$ e
$\mathcal F'$ são \emph{equivalentes} se existem $k$-homomorfismos de $K$,
$K'$ em uma mesma extensão $K''$ de $k$, tais que $\mathcal F_{K''}$ e
$\mathcal F'_{K''}$ sejam isomorfos sobre $X_{K''}$. Trata-se de uma relação de
equivalência, e interessar-nos-emos pelas classes de equivalência de feixes
para essa relação e pelos conjuntos de classes de equivalência. Observe-se que,
se $X_0$ é de tipo finito sobre $k$, então toda classe de feixes
coerentes pode ser definida por um feixe coerente sobre algum $X_K$, onde $K$
é uma extensão de tipo finito de $k$. Pode-se então, na definição das classes
de feixes coerentes, restringir-se às extensões algebricamente fechadas de $k$,
e pode-se também restringir a uma extensão algebricamente fechada fixa $\Omega$
de $k$, de grau de transcendência infinito, sendo dois feixes coerentes
$\mathcal F$ e $\mathcal F'$ sobre $X_\Omega$ equivalentes se e somente se
existe um $k$-automorfismo $\sigma$ de $\Omega$ tal que
$\mathcal F\otimes_\Omega(\Omega,\sigma)$ seja isomorfo a $\mathcal F'$.
Verifica-se que há uma correspondência biunívoca entre as classes de feixes
coerentes para uma ou outra definição.

Sejam $E$ e $E'$ dois conjuntos de classes de feixes coerentes sobre $X$.
Consideremos as classes de todos os feixes da forma
$\mathcal F\otimes\mathcal F'$, onde $\mathcal F$ e $\mathcal F'$ são feixes
coerentes sobre um mesmo $X_K$, com a classe de $\mathcal F$ em $E$ e a de
$\mathcal F'$ em $E'$. Obtém-se assim um conjunto de classes de feixes
coerentes, que será denotado por $E\otimes E'$. Definem-se do mesmo modo os
$\operatorname{Tor}_i(E,E')$, etc. De modo geral, a toda função $U$ que associa
a toda sequência $\mathcal F_1,\ldots,\mathcal F_n$ de $n$ feixes coerentes
sobre um mesmo $X_K$ um conjunto $U(\mathcal F_1,\ldots,\mathcal F_n)$ de
feixes coerentes sobre $X_K$, e que tem uma propriedade evidente de
compatibilidade com os isomorfismos de feixes e as imagens inversas por mudança
de base, associa-se uma função, denotada pelo mesmo símbolo $U$, que associa a
toda sequência de $n$ conjuntos $E_1,\ldots,E_n$ de classes de feixes
coerentes um conjunto $U(E_1,\ldots,E_n)$ de classes de feixes coerentes.

Nosso objetivo neste item é definir certos conjuntos de classes de feixes,
qualificados de \emph{limitados}, e mostrar que as operações $U$ mais usuais,
aplicadas a conjuntos limitados, ainda produzem conjuntos limitados.

Seja $X$ um pré-esquema de tipo finito sobre $S$ Noetheriano. Para todo
$s\in S$, a fibra $X_s$ é um pré-esquema de tipo finito sobre $k(s)$, e
consideraremos classes de feixes coerentes sobre $X_s$ no sentido precedente.
Isso dá sentido à expressão «classe de feixes coerentes sobre uma fibra de
$X/S$» e às expressões análogas. Do mesmo modo, procedendo separadamente sobre
cada fibra, podem-se ainda considerar operações tais como $E\otimes E'$, etc.,
que fazem corresponder a sistemas de conjuntos de
\src{251}
classes de feixes coerentes sobre as fibras de $X/S$, outro conjunto de classes
de feixes coerentes sobre as fibras de $X/S$.

\result{DEFINIÇÃO}{1.1} Seja $E$ um conjunto de classes de feixes coerentes
sobre as fibras de $X/S$. Diz-se que $E$ é \emph{limitado} se existem um
pré-esquema $S'$ de tipo finito sobre $S$ e um feixe coerente $\mathcal F'$ sobre
$X'=X\times_S S'$ tais que $E$ esteja contido no conjunto de classes de feixes
sobre as fibras de $X/S$ definido por $\mathcal F'$.

Por definição, este último faz corresponder a cada $s\in S$ as classes dos
feixes $\mathcal F'\otimes_{S'}k(s')$, em que $s'$ percorre os pontos de $S'$
situados sobre $s$ (de modo que $k(s')$ é uma extensão de $k(s)$ e
$X_s\otimes_{k(s)}k(s')$ identifica-se com a fibra
$X'\otimes_{S'}k(s')=X'_{s'}$ de $X'$ em $s'$). Pode-se dizer que as famílias
limitadas são aquelas contidas em uma \emph{família algébrica} de feixes
coerentes parametrizada por um $S'$ de tipo finito sobre $S$.

Uma união finita de famílias limitadas é limitada (tome-se como pré-esquema a
soma disjunta dos pré-esquemas de parâmetros $S_i$ que definem as famílias
algébricas majorantes). Uma mudança de base $T\longrightarrow S$ transforma uma
família limitada relativa a $X/S$ em uma família limitada relativa a $X_T/T$, e
a recíproca é verdadeira se $T\longrightarrow S$ é sobrejetivo (ou, mais
geralmente, se sua imagem contém os $s$ que intervêm efetivamente na família
dada $E$ para $X/S$). Isso reduz, em princípio, a determinação das famílias
limitadas ao caso em que $S$ é o espectro de uma álgebra de tipo finito sobre o
anel dos inteiros $\mathbf Z$.

Se $E$ e $E'$ são famílias limitadas de classes de feixes relativas a $X/S$,
então $E\otimes E'$ também é limitada. Com efeito, se $E$ e $E'$ são majoradas,
respectivamente, pelas famílias algébricas definidas por $T\longrightarrow S$ e
$\mathcal F$ sobre $X_T$, e por $T'\longrightarrow S$ e $\mathcal F'$ sobre
$X_{T'}$, vê-se que $E\otimes E'$ é majorada pela família algébrica definida por
$T''=T\times_S T'\longrightarrow S$ e pelo feixe $\mathcal F''$ sobre
$X_{T''}$, produto tensorial das imagens inversas de $\mathcal F$ e
$\mathcal F'$ sobre $X_{T''}$. Esse raciocínio só é válido porque o funtor
$\mathcal F\otimes\mathcal F'$ nas variáveis $\mathcal F$, $\mathcal F'$ é exato
à direita e, portanto, comuta com a extensão da base (e, em particular, com a
passagem às fibras). Ele não se aplica, tal como está, às operações locais como
$\operatorname{Tor}_i(E,E')$, $\operatorname{Hom}(E,E')$ e
$\operatorname{Ext}^i(E,E')$. Pode-se, contudo, mostrar que essas operações
continuam transformando conjuntos limitados em conjuntos limitados, procedendo
como para $E\otimes E'$, mas utilizando também resultados do tipo seguinte
(todos contidos em [3], IV, 6.11): uma família limitada $E$ é sempre majorada
por uma família algébrica definida por um feixe coerente $\mathcal F$ sobre um
$X_T$ ($T$ de tipo finito sobre $S$) que é plano relativamente a $T$. (O espaço
de parâmetros inicial é «cortado em partes».) Tais
\src{252}
propriedades de planicidade em feixes adequados asseguram, com efeito, a
comutação de operações como $\operatorname{Tor}_i(\mathcal F,\mathcal F')$
com uma mudança de base arbitrária. O mesmo método aplica-se a operações de
natureza global: imagens diretas e imagens diretas derivadas de feixes
coerentes por morfismos próprios, $\operatorname{Ext}$ globais relativos a
morfismos próprios (cf. [5], III, \S~6), etc.; todas essas operações
transformam famílias limitadas de feixes em famílias limitadas de feixes.
(N. B. -- Aqui os pré-esquemas sobre os quais se tomam os diversos feixes
podem mudar pelas operações do tipo considerado.)

Os dois enunciados seguintes demonstram-se essencialmente pela mesma técnica
de planicidade; para a decomposição primária nas fibras de um morfismo de tipo
finito, veja-se em particular [5], IV.

\result{PROPOSIÇÃO}{1.2} Sejam $E$ e $E'$ conjuntos limitados de classes de
feixes nas fibras de $X/S$, supondo-se $X$ próprio sobre $S$. Então:

\medskip
\noindent (i) A família dos núcleos, conúcleos e imagens de homomorfismos
$\mathcal F\longrightarrow\mathcal F'$, em que a classe de $\mathcal F$
pertence a $E$ e a de $\mathcal F'$ pertence a $E'$, é limitada.

\smallskip
\noindent (ii) A família dos feixes $\mathcal F''$ que são extensões de um
$\mathcal F$ por um $\mathcal F'$, em que a classe de $\mathcal F$ pertence
a $E$ e a de $\mathcal F'$ pertence a $E'$, é limitada.

Mediante uma mudança de base adequada, pode-se supor que $E$ e $E'$ sejam
definidos, respectivamente, por um feixe coerente $\mathcal G$ e um feixe
coerente $\mathcal G'$ sobre um $X_T/T$, sendo $T$ de tipo finito sobre $S$.
Além disso, podem-se supor satisfeitas certas hipóteses de planicidade, que
implicam que a formação dos feixes
$f_{T*}(\operatorname{Hom}_{X_T}(\mathcal G,\mathcal G'))$ e
$\operatorname{Ext}^{1}_{f_T}(\mathcal G,\mathcal G')$ comuta com a mudança
de base por um morfismo arbitrário $T'\longrightarrow T$. Pode-se também
supor que os feixes coerentes precedentes sobre $T$ sejam localmente livres.
Sejam então $T_0$ e $T_1$ os fibrados vetoriais sobre $T$ cujos feixes de
germes de seções são, respectivamente, os feixes precedentes. Define-se então
canonicamente um homomorfismo
$\mathcal G_{T_0}\longrightarrow\mathcal G'_{T_0}$ de feixes coerentes sobre
$X_{T_0}$ e uma extensão
\[
0\longrightarrow\mathcal G_{T_1}\longrightarrow\mathcal G''
\longrightarrow\mathcal G'_{T_1}\longrightarrow0
\]
de feixes coerentes sobre $X_{T_1}$, dotada de uma propriedade universal
evidente. Esse segundo feixe define uma família algébrica que majora a família
considerada em (ii). O mesmo vale para o núcleo, o conúcleo e a imagem do
homomorfismo
\src{253}
precedente, cuja consideração estabelece (i) (desde que se suponha o conúcleo
plano relativamente a $T_0$, caso ao qual ainda se reduz decompondo $T_0$ em
partes\ldots).

\result{PROPOSIÇÃO}{1.3} Seja $E$ uma família limitada de classes de feixes
nas fibras de $X/S$. Então as classes dos feixes estruturais dos
$(\operatorname{supp}\mathcal F)_{\mathrm{red}}$, em que $\mathcal F$ é um
feixe coerente sobre um $X_K$, com $K$ algebricamente fechado, cuja classe
pertence a $E$, formam uma família limitada.

Aqui, $(\operatorname{supp}\mathcal F)_{\mathrm{red}}$ designa o suporte de
$\mathcal F$, munido da estrutura reduzida induzida, isto é, seu feixe
estrutural é o quociente de $\mathcal O_{X_K}$ pelo maior feixe de ideais que
define $\operatorname{supp}\mathcal F$. Pode-se demonstrar o resultado
análogo a (1.3) para os feixes deduzidos canonicamente de $\mathcal F$ pela
teoria da decomposição primária; por exemplo, os
$\mathcal F/\mathcal F_i$, em que os $\mathcal F_i$ são os subfeixes de
$\mathcal F$ primários para as componentes do suporte de $\mathcal F$ e
mínimos com essa propriedade; ou os $\mathcal O_{X_K}/\mathfrak p$, em que
$\mathfrak p$ é um feixe de ideais primos associado a $\mathcal F$; ou os
$\mathcal O_{X_K}/\mathfrak q$, em que $\mathfrak q$ é um feixe de ideais
primários associado a uma componente do suporte de $\mathcal F$ (sendo o
corpo de referência algebricamente fechado).

\fsection{2}{Famílias limitadas e polinômio de Hilbert}

No que segue, supomos que $X$ seja projetivo sobre $S$ e esteja munido de um
feixe muito amplo, denotado por $\mathcal O_X(1)$. Para toda extensão $K$ do
corpo residual $k(s)$ de um ponto $s$ de $S$, considera-se sobre $X_K$ o feixe
correspondente $\mathcal O_{X_K}(1)$, que será ainda muito amplo.

A todo feixe coerente $\mathcal F$ sobre $X_K$ associa-se a função
\[
P_{\mathcal F}(n)=\text{característica de Euler--Poincaré de $\mathcal F(n)$
sobre $X_K$},
\]
que é um polinômio no inteiro $n$, chamado \emph{polinômio de Hilbert} de
$\mathcal F$. Para valores grandes de $n$, $P(n)$ nada mais é do que a
dimensão de $H^0(X_K,\mathcal F(n))$ sobre $K$, pois os
$H^i(X_K,\mathcal F(n))$ são nulos para $i>0$ e $n$ grande.

Se agora $\mathcal F$ for um feixe coerente sobre $X$, plano relativamente a
$S$, então os polinômios de Hilbert dos feixes $\mathcal F_s$ induzidos nas
fibras $X_s$ relativas a uma mesma componente conexa de $S$ serão todos iguais
[5], III, \S~7. Daí resulta (sem hipótese de planicidade) que o conjunto dos
polinômios de Hilbert dos feixes $\mathcal F_s$, $s\in S$, é finito para
todo feixe coerente $\mathcal F$ sobre $X$.
\src{254}
Recordemos, por outro lado, que, se $\mathcal F$ é um feixe coerente sobre $X$,
este é isomorfo a um feixe quociente de um feixe $\mathcal O_X(-n)^N$, para
$n$ e $N$ suficientemente grandes. Portanto, os feixes $\mathcal F_s$
induzidos nas fibras são igualmente quocientes do feixe $\mathcal O(-n)$ na
fibra.

Das duas observações, deduz-se a parte «somente se» do teorema seguinte:

\result{TEOREMA}{2.1} Seja $X$ projetivo sobre um $S$ Noetheriano, e seja
$\mathcal O_X(1)$ muito amplo sobre $X$ relativamente a $S$. Seja $E$ um
conjunto de classes de feixes nas fibras de $X/S$. Para que $E$ seja limitado,
é necessário que
satisfaça as condições seguintes:

\medskip
\noindent (a) Existe um feixe coerente $\mathcal L$ sobre $X$ (que se pode
supor da forma $\mathcal O_X(-n)^N$) tal que $E$ esteja contido na família das
classes de feixes coerentes quocientes de feixes da forma $\mathcal L_K$.

\smallskip
\noindent (b) Os polinômios de Hilbert $P_{\mathcal F}$ dos feixes
$\mathcal F$ cuja classe pertence a $E$ são elementos de um mesmo conjunto
finito de polinômios.

Resta provar a parte «se», que será um caso particular de um resultado mais
preciso. Para todo Módulo coerente $\mathcal F$ sobre um pré-esquema de tipo
finito sobre um corpo $K$, e todo inteiro $r$, seja $N_r$ o sub-Módulo de
$\mathcal F$ cujas seções sobre um aberto são as seções de $\mathcal F$ sobre
esse aberto cujo suporte tem dimensão $<r$. Tem-se, portanto,
$N_r=\mathcal F$ para $r>\dim\operatorname{supp}\mathcal F$, $N_r=0$ para
$r\leq0$, e obtém-se uma filtração crescente finita de $\mathcal F$ cujos
fatores $N_r/N_{r+1}$ têm como ciclos primos associados exatamente os ciclos primos
associados a $\mathcal F$ que têm dimensão $r$. Poremos
\[
\mathcal F_{(r)}=\mathcal F/N_r.
\]
Assim, $\mathcal F_{(r)}$ tem como ciclos primos associados exatamente os
ciclos primos associados a $\mathcal F$ que têm dimensão $\geq r$ e, em
particular, é igual a $\mathcal F$ se e somente se os ciclos primos associados
a $\mathcal F$ têm dimensão $\geq r$. Posto isso:

\result{TEOREMA}{2.2} Nas condições de (2.1), seja $s$ um inteiro e suponhamos
que $E$ satisfaça a condição (a) e a seguinte forma enfraquecida de (b):

\medskip
\noindent $(b_s)$ Os polinômios de Poincaré $P_{\mathcal F}$ dos feixes
$\mathcal F$ cuja classe pertence a $E$ têm coeficientes em graus
$\leq s-1$ que permanecem limitados.

Nessas condições, os feixes $\mathcal F_{(s)}$ (quando a classe de
$\mathcal F$ percorre $E$) formam uma família limitada. Além disso, os
coeficientes de grau $s-2$ dos $P_{\mathcal F}$ permanecem limitados
inferiormente.
\src{255}
Assim:

\result{COROLÁRIO}{2.3} Suponhamos que os feixes $\mathcal F$ cuja classe
pertence a $E$ sejam tais que todos os ciclos primos associados tenham dimensão
$d$ com $s\leq d\leq r$. Então, em (2.1) (b), pode-se considerar apenas os
coeficientes dos $P_{\mathcal F}$ em graus compreendidos entre $s-1$ e $r$.

O fim deste item é dedicado ao esboço da demonstração de (2.2). Os lemas-chave
são os dois lemas seguintes, o primeiro dos quais é bem conhecido (e resume o
conteúdo matemático útil das coordenadas de Chow):

\result{LEMA 2.4 (CHOW)}{} Consideremos os feixes estruturais de subesquemas
$Y$ de fibras $X_K$ ($K$ percorrendo extensões algebricamente fechadas de
corpos residuais de $S$), com $Y$ reduzido e todas as suas componentes de mesma
dimensão $r$ (considerando-se $\mathcal O_Y$ como um feixe quociente de
$\mathcal O_X$). Se os graus dos $Y$ permanecem limitados, os $Y$ formam uma
família limitada.

Aqui, o grau $a$ de $Y$ pode ser definido mais convenientemente pelo
coeficiente do termo dominante de
\[
P_{\mathcal O_Y}=a\frac{n^r}{r!}+\cdots.
\]

\result{LEMA}{2.5} Sejam $\mathcal L$ um feixe coerente sobre $X$ e $E$ um
conjunto de classes de feixes quocientes $\mathcal F$ de feixes $\mathcal L_K$
($K$ uma extensão de corpo residual de $S$). Suponhamos que as fibras de $X$
sobre $S$ tenham dimensão $\leq r$, e ponhamos
\[
P_{\mathcal F}(n)=a_{\mathcal F}\frac{n^r}{r!}
+b_{\mathcal F}\frac{n^{r-1}}{(r-1)!}
+\text{termos de grau $<r-1$}.
\]
Então o coeficiente $a_{\mathcal F}$ permanece limitado, e $b_{\mathcal F}$
permanece limitado inferiormente. Se $b_{\mathcal F}$ permanece limitado,
então a família dos $\mathcal F_{(r)}$ é limitada.

\emph{Demonstração.} -- Substituindo, se necessário, $S$ por uma união de
subesquemas de $S$ que recobre $S$, pode-se supor que existe um morfismo finito
$f:X\longrightarrow\mathbf P^r_S$, tal que $\mathcal O_X(1)$ seja isomorfo à
imagem inversa de $\mathcal O_{\mathbf P^r_S}(1)$; portanto, para todo feixe
coerente $\mathcal F$ sobre $X$, tem-se
$P_{\mathcal F}=P_{f_*(\mathcal F)}$. Por outro lado, verifica-se facilmente
(pela técnica do item precedente) que um conjunto de feixes $\mathcal F$ sobre
$X$ é limitado se e somente se o conjunto dos $f_*(\mathcal F)$ o é. Por fim,
\[
f_*(\mathcal F)_{(r)}=f_*(\mathcal F_{(r)}).
\]
\src{256}
Isso nos reduz, portanto, ao caso $X=\mathbf P^r_S$. Além disso, pode-se supor
que $\mathcal L=\mathcal O_{\mathbf P^r_S}(k)^s$ para $k$, $s$ convenientes. O
coeficiente $a_{\mathcal F}$ satisfaz
\[
0\leq a_{\mathcal F}\leq s
\]
e, portanto, permanece limitado. Dito isso, é equivalente afirmar que os
$P_{\mathcal F}(n)$ têm um coeficiente de $n^{r-1}$ que permanece limitado
inferiormente (resp. limitado), ou afirmar o mesmo para
$P_{\mathcal F}(n-k)=P_{\mathcal F(-k)}(n)$. Isso nos reduz ao caso em que
\[
\mathcal L=\mathcal O_{\mathbf P^r_S}^{s}.
\]

Consideremos a sequência exata
\[
0\longrightarrow N_r\longrightarrow\mathcal F
\longrightarrow\mathcal F_{(r)}\longrightarrow0,
\]
da qual se obtém
\[
P_{\mathcal F}=P_{\mathcal F_{(r)}}+P_{N_r},
\]
e, como o coeficiente de $n^{r-1}$ em $P_{N_r}$ é não negativo (pois
$\dim\operatorname{supp}N_r\leq r-1$), tem-se
\[
b_{\mathcal F_{(r)}}\leq b_{\mathcal F}.
\]
Isso nos permite, para demonstrar as afirmações do lema, substituir
$\mathcal F$ por $\mathcal F_{(r)}$, isto é, supor que os feixes quocientes
$\mathcal F$ de $\mathcal L$ considerados são livres de torção.

Como $\mathbf P^r_K$ é normal, segue-se que $\mathcal F$ é localmente livre de
posto $a=a_{\mathcal F}$ em um aberto $U=\mathbf P^r_K-Y$, onde $Y$ tem
codimensão $\geq2$. Logo, $\bigwedge^a\mathcal F$ é um feixe sobre
$\mathbf P^r_K$ cuja restrição a $U$ é invertível; portanto (como
$\mathbf P^r_K$ é regular e $Y$ tem codimensão $\geq2$), essa restrição é
isomorfa à restrição de um feixe invertível sobre $\mathbf P^r_K$, determinado
salvo isomorfismo. Este último é da forma $\mathcal O_{\mathbf P^r_K}(d)$ para um
inteiro $d$ bem determinado. Como $\bigwedge^a\mathcal F$ é um quociente de
$\bigwedge^a\mathcal O_{\mathbf P^r_K}^{n}\simeq
\mathcal O_{\mathbf P^r_K}^{N}$, com $N=\binom{s}{a}$, possui $N$ seções canônicas,
que definem, portanto, seções de
\src{257}
$\mathcal O_{\mathbf P^r_K}(d)$ sobre $U$, que são restrições de seções
$s_i$ ($1\leq i\leq N$) de $\mathcal O_{\mathbf P^r_K}(d)$ (pois
$\mathbf P^r_K$ é normal e $Y$ tem codimensão $\geq2$). Estas seções
geram $\mathcal O_{\mathbf P^r_K}(d)$ nos pontos de $U$, portanto não são
todas nulas, o que implica que $d\geq0$. Por outro lado, um cálculo simples
mostra que
\[
b_{\mathcal F}=a_{\mathcal F}(r+1)/2+d.
\]
Isso mostra, em particular, que $b_{\mathcal F}>0$, portanto
$b_{\mathcal F}$ é limitado inferiormente. Esse coeficiente permanece limitado
se e somente se $d$ permanece limitado; mostremos que, nesse caso, $\mathcal F$
pertence a uma família limitada. Podemos supor que $a_{\mathcal F}$ e
$b_{\mathcal F}$ sejam fixos, digamos iguais a $a$ e $b$; portanto, $d$ é fixo.
O conhecimento das $N$ seções
$s_i$ de $\mathcal O_{\mathbf P^r_K}(d)$, isto é, de um homomorfismo
\[
s:\bigwedge^a\mathcal L_K\longrightarrow\mathcal O_{\mathbf P^r_K}(d),
\]
permite recuperar $\mathcal F$ como a coimagem do homomorfismo composto
correspondente:
\[
\mathcal L_K\longrightarrow
\operatorname{Hom}\!\left(\bigwedge^{a-1}\mathcal L_K,
\bigwedge^a\mathcal L_K\right)\longrightarrow
\operatorname{Hom}\!\left(\bigwedge^{a-1}\mathcal L_K,
\mathcal O_{\mathbf P^r_K}(d)\right),
\]
em que o primeiro é o homomorfismo canônico proveniente do produto exterior, e
o segundo é deduzido de $s$. A conclusão decorre então de (1.2) (i).

Combinando os dois lemas precedentes, obtém-se:

\result{LEMA}{2.6} Sob as hipóteses preliminares de (2.1), suponhamos que, para
todo $\mathcal F$, se tenha
\[
P_{\mathcal F}(n)=a_{\mathcal F}\frac{n^r}{r!}
+b_{\mathcal F}\frac{n^{r-1}}{(r-1)!}
+\text{termos de grau $<r-1$},
\]
e que os coeficientes $a_{\mathcal F}$ permaneçam limitados. Então os
coeficientes $b_{\mathcal F}$ permanecem limitados inferiormente. Se os
$b_{\mathcal F}$ forem limitados, então os $\mathcal F_{(r)}$ formam uma família
limitada.

Podemos supor que os corpos de base $K$ dos feixes $\mathcal F$ sejam
algebricamente fechados. Munamos cada
$\operatorname{supp}\mathcal F_{(r)}$, união das componentes de grau
$r$ de
$\operatorname{supp}\mathcal F$, da estrutura reduzida induzida. Então os graus dos
$\operatorname{supp}\mathcal F_{(r)}$ permanecem limitados superiormente por
$a$; portanto, em virtude de (2.4), os
$\operatorname{supp}\mathcal F_{(r)}$ formam um conjunto limitado. Além disso,
para toda componente de $\operatorname{supp}\mathcal F_{(r)}$, o comprimento de
$\mathcal F_{(r)}$ nessa componente é $\leq a$; portanto, se
$\mathcal J_{\mathcal F}$ é o Ideal que define
$\operatorname{supp}\mathcal F_{(r)}$, então $\mathcal F_{(r)}$
\src{258}
pode ser considerado um Módulo sobre o subesquema $Y_{\mathcal F}$ de $X$
definido por $\mathcal J_{\mathcal F}^{a}$. Como no lema precedente, reduzimo-nos também
ao caso em que $\mathcal F=\mathcal F_{(r)}$, de modo que $\mathcal F$
provém de um Módulo sobre $Y_{\mathcal F}$. Os $Y_{\mathcal F}$ correspondem
a uma família limitada de Módulos quocientes dos $\mathcal O_{X_K}$ e,
portanto, provêm de um subesquema fechado $Y$ de um esquema $X\times_S T$.
Pode-se então aplicar (2.5) a $Y/T$ e $\mathcal L\otimes_XY$, o que dá a
conclusão.

Podemos agora demonstrar (2.2) por indução sobre o limite superior $r$ dos
$\dim\operatorname{supp}\mathcal F$. O enunciado é trivial para $r<0$;
suponhamos, pois, que $r\geq0$ e que o enunciado tenha sido demonstrado para
$r'<r$. Em virtude de (2.6), os $\mathcal F_{(r)}$ formam um conjunto
limitado; portanto, também em virtude de (1.2) (i), os núcleos dos
homomorfismos $\mathcal L_K\longrightarrow\mathcal F_{(r)}$ formam um
conjunto limitado. Existe, pois, um Módulo coerente $\mathcal L'$ sobre $X$
tal que os núcleos em questão e, consequentemente, também os
\[
N_r(\mathcal F)=\operatorname{Ker}(\mathcal F\longrightarrow\mathcal F_{(r)}),
\]
sejam quocientes dos Módulos $\mathcal L'_K$. Como os $\mathcal F_{(r)}$ são
limitados, os $P_{\mathcal F_{(r)}}$ permanecem limitados, e a fórmula
\[
P_{\mathcal F}=P_{\mathcal F_{(r)}}+P_{N_r}
\]
mostra então que os $P_{N_r}$ satisfazem a mesma condição $(b_s)$ que os
$P_{\mathcal F}$. Logo, os $N_r$ satisfazem as condições (a) e $(b_s)$ e,
pela hipótese de indução, os $(N_r)_{(s)}$ permanecem limitados. Ora,
$\mathcal F_{(s)}$ é uma extensão de $\mathcal F_{(r)}$ por
$(N_r)_{(s)}$; portanto, por (1.2) (ii), os $\mathcal F_{(s)}$ permanecem
limitados. Para a última afirmação de (2.2), observa-se que os núcleos $N_s$ de
$\mathcal F\longrightarrow\mathcal F_{(s)}$ permanecem limitados em virtude
de (1.2) (i), e que o coeficiente do termo em $n^{s-1}$ em $P_{N_s}$ permanece
limitado; o lema 2.6 demonstra então que o coeficiente do termo seguinte
permanece minorado. Isso conclui a demonstração.

\fsection{3}{Esquemas de Hilbert: definição, teorema de existência}

Sejam $X$ um pré-esquema sobre outro $S$, e $\mathcal F$ um Módulo
quase-coerente sobre $X$. Denotemos por
\[
\operatorname{Quot}(\mathcal F/X/S)
\]
o conjunto dos Módulos quase-coerentes quocientes de $\mathcal F$ que são
planos sobre $S$. Seja agora $S'\longrightarrow S$ um morfismo de mudança de
base; ponhamos $X'=X\times_S S'$ e
$\mathcal F'=\mathcal F\otimes_{\mathcal O_S}\mathcal O_{S'}$; assim, $X'$
é um pré-esquema sobre $S'$ munido de um Módulo
\src{259}
quase-coerente $\mathcal F'$, e pode-se considerar
$\operatorname{Quot}(\mathcal F'/X'/S')$. Ponhamos
\[
\underline{\operatorname{Quot}}_{\mathcal F/X/S}(S')
=\operatorname{Quot}(\mathcal F'/X'/S')
\qquad\text{(onde $X'=X\times_SS'$)}.
\]
Se agora $S''\longrightarrow S'$ for um $S$-morfismo, então
$X''=X\times_SS''$ é isomorfo a $X'\times_{S'}S''$, e $\mathcal F''$ é
isomorfo a $\mathcal F'\otimes_{\mathcal O_{S'}}\mathcal O_{S''}$. Como o
funtor imagem inversa
\[
\mathcal G'\longmapsto
\mathcal G'\otimes_{\mathcal O_{S'}}\mathcal O_{S''}
\]
da categoria dos Módulos quase-coerentes sobre $X'$ para a categoria dos
Módulos quase-coerentes sobre $X''$ é exato à direita e transforma Módulos
$S'$-planos em Módulos $S''$-planos, obtém-se uma aplicação natural
\[
\operatorname{Quot}(\mathcal F'/X'/S')\longrightarrow
\operatorname{Quot}(\mathcal F''/X''/S''),
\]
de modo que $\underline{\operatorname{Quot}}_{\mathcal F/X/S}(S')$ é um
funtor contravariante em $S'$ (pré-esquema sobre $S$), com valores na categoria
dos conjuntos. Doravante, suporemos que $X$ seja projetivo sobre um $S$
noetheriano, que $\mathcal F$ seja coerente e, para simplificar, limitar-nos-emos
aos $S'$ sobre $S$ que sejam localmente noetherianos.

\result{TEOREMA}{3.1} Nessas condições, o funtor contravariante
$\underline{\operatorname{Quot}}_{\mathcal F/X/S}$ sobre a categoria dos
$S$-pré-esquemas localmente noetherianos é representável por um
$S$-pré-esquema $\operatorname{Quot}_{\mathcal F/X/S}$, união disjunta de uma
sequência de $S$-esquemas projetivos (a fortiori,
$\operatorname{Quot}_{\mathcal F/X/S}$ é localmente de tipo finito sobre
$S$).

Obteremos tal decomposição da seguinte maneira. Seja $\mathcal O_X(1)$ um
feixe invertível sobre $X$, muito amplo relativamente a $S$. Para todo
polinômio $P(n)$ com coeficientes racionais, seja
$\operatorname{Quot}^P(\mathcal F/X/S)$ a parte de
$\operatorname{Quot}(\mathcal F/X/S)$ formada pelos quocientes coerentes
$\mathcal G$ de $\mathcal F$ que são planos sobre $S$ e cujo polinômio de
Hilbert em todo $s\in S$ é igual a $P$. Poremos então
\[
\underline{\operatorname{Quot}}^P_{\mathcal F/X/S}(S')
=\operatorname{Quot}^P(\mathcal F'/X'/S')
\]
e obteremos assim um subfuntor de
$\underline{\operatorname{Quot}}_{\mathcal F/X/S}$. A propriedade de
invariância dos polinômios de Hilbert recordada no número 2 implica o seguinte:
para que $\underline{\operatorname{Quot}}_{\mathcal F/X/S}$ seja
representável, é necessário e suficiente que os
$\underline{\operatorname{Quot}}^P_{\mathcal F/X/S}$ o sejam; nesse caso, o
$S$-pré-esquema $\operatorname{Quot}_{\mathcal F/X/S}$ que o representa é
isomorfo ao pré-esquema união disjunta dos
$\operatorname{Quot}^P_{\mathcal F/X/S}$ que representam os funtores
$\underline{\operatorname{Quot}}^P_{\mathcal F/X/S}$. Posto isso, (3.1)
será, portanto, consequência do teorema seguinte:
\src{260}
\result{TEOREMA}{3.2} Com as notações anteriores, o funtor
$\underline{\operatorname{Quot}}^P_{\mathcal F/X/S}$ é representável por um
$S$-pré-esquema projetivo $\operatorname{Quot}^P_{\mathcal F/X/S}$.

O restante deste número é dedicado à demonstração de (3.2).

Seja $\nu$ um inteiro. Para todo $S'$ sobre $S$, denotamos por $A_\nu(S')$ o
conjunto dos quocientes $\mathcal G=\mathcal F'/\mathcal K$ de
$\mathcal F'=\mathcal F\otimes_{\mathcal O_S}\mathcal O_{S'}$, que são
coerentes, planos sobre $S'$ e satisfazem as condições seguintes:
\begin{align*}
\text{a.}\quad &R^if'_*(\mathcal G(n))=0 &&\text{para $i>0$ e $n\geq\nu$},\\
\text{b.}\quad &R^if'_*(\mathcal K(n))=0 &&\text{para $i>0$ e $n\geq\nu$},\\
\text{c.}\quad &f'_*(\mathcal K(\nu+k))
=\mathcal S'_k f'_*(\mathcal K(\nu)) &&\text{para $k\geq0$}.
\end{align*}
Para esta última relação, supõe-se que $X$ esteja escrito como o espectro primo
homogêneo de uma álgebra graduada quase-coerente $\mathcal S_*$ sobre $S$,
graduada em graus positivos e gerada por $\mathcal S_1$, e põe-se
$\mathcal S'_* =\mathcal S_*\otimes_{\mathcal O_S}\mathcal O_{S'}$, de modo
que $X'$ seja o espectro primo homogêneo de $\mathcal S'_*$. Para demonstrar
(3.2), reduz-se facilmente ao caso em que $X=\mathbf P^r_S$ (pois $S$ é união
de abertos $U$ tais que $X|U$ seja um subesquema fechado de $\mathbf P^r_U$ e
$\mathcal O_X(1)$ seja induzido por $\mathcal O_{\mathbf P^r_U}(1)$), e ao
caso em que $\mathcal F$ é da forma
$\mathcal O_{\mathbf P^r_S}^{N}$, sendo, portanto, plano sobre $S$. Então, no
que precede, os feixes $\mathcal K$ também são planos sobre $S'$. Resulta das
relações de Künneth ([5], III, \S~7) que as condições (a) e (b) são estáveis
por mudança de base e implicam que, para $n\geq\nu$, a formação de
$f'_*(\mathcal G(n))$ e de $f'_*(\mathcal K(n))$ comute com a mudança de base
(loc. cit.); portanto, (c) também é estável por mudança de base.

Em outras palavras, $A_\nu(S')$ é um funtor contravariante em $S'$, mais
precisamente um subfuntor de
$A(S')=\underline{\operatorname{Quot}}^P_{\mathcal F/X/S}(S')$. Ao variar
$\nu$, obtém-se assim uma sequência crescente de subconjuntos $A_\nu(S')$ de
$A(S')$, cuja união é $A(S')$ em virtude de um teorema bem conhecido de
Serre [5], III, \S~2. Observemos agora que, se $\mathcal G$ é um quociente
coerente de $\mathcal F'$, plano sobre $S'$, e $s$ é um ponto de $S'$ tal que
a mudança de base $\operatorname{Spec}(k(s))\longrightarrow S'$ dá origem a
um quociente $\mathcal G_s$ de $\mathcal F_s$ que satisfaz as condições (a),
(b) e (c), isto é, que pertence a $A_\nu(\operatorname{Spec}(k(s)))$, então
existe uma vizinhança aberta $U$ de $s$ tal que essas mesmas condições sejam
verificadas para $\mathcal G|f'^{-1}(U)$, isto é, esse quociente pertence a
$A_\nu(U)$; para (a) e (b), isso resulta, com efeito, do «teorema das funções
holomorfas» [5], III, \S~7, e (c) resulta do lema de Nakayama e do fato de que,
em todo caso, sabe-se que
$f'_*(\mathcal K(n+k))=\mathcal S'_k f'_*(\mathcal K(n))$
\src{261}
para $n$ suficientemente grande e $k\geq0$ (cf. [5], III, \S~2). Dessas
observações conclui-se o seguinte (compare-se [4], IV). Para que o funtor $A$
seja representável, é necessário e suficiente que os funtores $A_\nu$ o sejam;
nesse caso, o $S$-pré-esquema $Q$ que representa $A$ é uma união crescente de
subpré-esquemas abertos $Q_\nu$ que representam os $A_\nu$.

Seja
\[
M_* =\bigoplus_{n\geq0}f_*(\mathcal F(n))=\mathcal S_*^{N},
\]
de modo que
\[
M'_* =M_*\otimes_{\mathcal O_S}\mathcal O_{S'}
=\bigoplus_{n\geq0}f'_*(\mathcal F'(n))=\mathcal S'_*{}^{N}.
\]
De (a) resulta que:

\noindent (a') $f'_*(\mathcal G(n))$ é localmente livre de posto $P(n)$ para
$n\geq\nu$, [5], III, \S~7, e, de (b), para $i=1$, que:

\noindent (a'') $f'_*(\mathcal G(n))$ é um Módulo quociente de $M'_n$.

Além disso, o conhecimento desse Módulo quociente, para $n=\nu$, determina,
em virtude de (c), os sub-Módulos $f'_*(\mathcal K(n))$ de $M'_n$ para
$n\geq\nu$ e, portanto, determina $\mathcal K$ e, por conseguinte, $\mathcal G$.
Obtém-se assim uma aplicação injetiva
\[
A_\nu(S')\longrightarrow\operatorname{Grass}_{P(\nu)}(M'_\nu)
\]
de $A_\nu(S')$ no conjunto dos Módulos quocientes localmente livres de posto
$P(\nu)$ de $M'_\nu$ e, portanto, um homomorfismo funtorial
\[
i_\nu:A_\nu(S')\longrightarrow
\underline{\operatorname{Grass}}_{P(\nu)}(M_\nu)(S'),
\]
onde o funtor do segundo membro é representado pelo esquema de Grassmann
$\operatorname{Grass}_{P(\nu)}(M_\nu)$ (compare-se [4], V), que é projetivo
sobre $S$. Afirmo:

\result{LEMA}{3.3} $A_\nu(S')$ é um funtor representável, e o morfismo
$Q_\nu\longrightarrow\operatorname{Grass}_{P(\nu)}(M_\nu)$ que representa o
homomorfismo $i_\nu$ é uma imersão (o que implica que $Q_\nu$ é
quase-projetivo sobre $S$).

Essa afirmação é equivalente à seguinte (compare-se [4], IV): suponha-se dado
um Módulo quociente $N$ de $M'_\nu$, localmente livre de posto $P(\nu)$; então
existe um subpré-esquema $Z$ de $S'$ tal que, para todo pré-esquema localmente
noetheriano $T'$
\src{262}
sobre $S'$, a imagem inversa de $N$ sobre $T'$ pertence a
$\operatorname{Im}A_\nu(T')$ se e somente se $T'\longrightarrow S'$ se
fatoriza pelo subpré-esquema $Z$. Mudando as notações, pode-se supor $S'=S$,
isto é, tem-se um quociente $N_\nu$ de $M_\nu$ por um sub-Módulo $R_\nu$.
Para que esse quociente provenha de um elemento de $A_\nu(S)$, é necessário e
suficiente que satisfaça as duas condições seguintes:

\medskip
\noindent (i) $M_{\nu+k}/\mathcal S_kR_\nu$ é localmente livre de posto
$P(\nu+k)$ para $k\geq0$.

\smallskip
\noindent (ii) O subfeixe $\mathcal K$ de $\mathcal F$ definido pelo
sub-Módulo graduado
\[
R_* =\sum_{k\geq0}\mathcal S_kR_\nu
\]
de $M_*$ ([5], II, \S~3) e o quociente $\mathcal G=\mathcal F/\mathcal K$
satisfazem as condições (a) e (b) acima.

Essas condições são manifestamente necessárias; por outro lado, se forem
satisfeitas, então o feixe $\mathcal G$ definido em (ii), uma vez que é
isomorfo ao feixe associado ao $\mathcal S_*$-Módulo graduado $N_*$ que é a
soma dos $M_{\nu+k}/\mathcal S_kR_\nu$, é plano sobre $S$ (pois suas fibras
são somandos diretos de localizações de $N$ em ideais primos homogêneos de
$\mathcal S_*$), e tem polinômio de Hilbert $P$ em virtude de (i). Levando em
conta (ii), vê-se então que (a') e (a'') são satisfeitas, logo, para
$n\geq\nu$, o homomorfismo natural $N_n\longrightarrow f_*(\mathcal G(n))$ é
um homomorfismo sobrejetivo entre Módulos localmente livres do mesmo posto e,
portanto, um isomorfismo; assim, tem-se
\[
f_*(\mathcal K(\nu+k))=\mathcal S_kR_\nu
\quad\text{para todo $k\geq0$},
\]
o que prova que $\mathcal G\in A_\nu(S)$ e que $R_\nu$ é o elemento de
$\operatorname{Grass}_{P(\nu)}(M_\nu)$ definido por $\mathcal G$.

Esse critério (i)--(ii) aplica-se igualmente à situação obtida após uma
mudança de base $S'\longrightarrow S$. Provaremos primeiro que o fato de a
condição (i) ser satisfeita após a mudança de base $S'\longrightarrow S$ se
exprime pela condição de que $S'\longrightarrow S$ se fatorize por certo
subpré-esquema $Z$ de $S$; obtido esse resultado, ficamos reduzidos
(substituindo $S$ por $Z$) ao caso em que a condição (i) já é satisfeita sobre
$S$ e, como ela é estável por mudança de base, resta exprimir a condição (ii).
Mas então, se $U$ denota o conjunto dos $s\in S$ tais que a cohomologia dos
feixes induzidos na fibra $X_s$ por $\mathcal G(n)$ e $\mathcal K(n)$ seja nula
em dimensão $>0$ para $n\geq\nu$, já se observou que $U$ é aberto, e a condição
(ii) será satisfeita após uma mudança de base $S'\longrightarrow S$ se e somente
se $S'\longrightarrow S$ se fatoriza por $U$, o que prova o Lema 3.3. Resta,
portanto, provar o lema seguinte:

\result{LEMA}{3.4} Sejam $S$ um pré-esquema localmente noetheriano, munido de
uma álgebra graduada quase-coerente $\mathcal S_*$ em graus positivos, gerada
por $\mathcal S_1$, $M_*$ um $\mathcal S_*$-Módulo graduado quase-coerente de tipo
finito, $P$ um polinômio com coeficientes racionais e $\nu$ um inteiro. Então
existe um subpré-esquema $Z$ de $S$ (evidentemente único)
\src{263}
com a seguinte propriedade: para todo pré-esquema $S'$ sobre $S$, para que
$M_n\otimes_{\mathcal O_S}\mathcal O_{S'}$ seja localmente livre de posto
$P(n)$ para todo $n\geq\nu$, é necessário e suficiente que
$S'\longrightarrow S$ se fatorize por $Z$.

Pode-se, evidentemente, supor $S$ afim, logo noetheriano. Então:

\result{LEMA}{3.5} Para todo inteiro $N\geq\nu$, seja $U_N$ o aberto de $S$
formado pelos $s\in S$ tais que
\[
\operatorname{rank}_{k(s)}
M_{ns}\otimes_{\mathcal O_{S,s}}k(s)\leq P(n)
\quad\text{para todo $\nu\leq n\leq N$}.
\]
Então a sequência decrescente dos abertos $U_N$ é estacionária.

Sabe-se [3], IV, que $S$ admite uma partição finita em subpré-esquemas
reduzidos $S_i$, tais que cada
$M\otimes_{\mathcal O_S}\mathcal O_{S_i}$ seja plano sobre $S_i$. Pode-se,
portanto, supor $M$ plano, logo os $M_n$ planos. Enfim, pode-se evidentemente
supor $S$ conexo. Mas então existem um inteiro $n_0$ e um polinômio $Q$ tais
que
\[
\operatorname{rank}_{k(s)}
M_{ns}\otimes_{\mathcal O_{S,s}}k(s)=Q(n)
\quad\text{para $n\geq n_0$},
\]
[5], III, \S~7. Suponhamos primeiro $P\neq Q$; então $P(n)\neq Q(n)$ para $n$
grande e, evidentemente, $U_N=\varnothing$ para $N$ grande; a fortiori, a
sequência dos $U_N$ é estacionária. Se, ao contrário, $P=Q$, então $U_N=U_{n_0}$ para $N\geq n_0$, e
a sequência dos $U_N$ é ainda estacionária.

Em particular, o conjunto $U_\infty$ dos $s\in S$ tais que
\[
\tag{*}
\operatorname{rank}_{k(s)}
M_{ns}\otimes_{\mathcal O_{S,s}}k(s)\leq P(n)
\quad\text{para todo $n\geq\nu$},
\]
sendo a interseção dos $U_N$, é aberto. Pode-se então, para a demonstração de
(3.4), substituir $S$ pelo aberto $U$, o que nos reduz ao caso em que a
desigualdade (*) é satisfeita em todo $s\in S$. Além disso:

\result{LEMA}{3.6} Seja $M$ um Módulo sobre um pré-esquema localmente
noetheriano $S$, e seja $r$ um inteiro. Então existe um subpré-esquema $Z$ de
$S$ (evidentemente único) com a seguinte propriedade: para todo $S'$ sobre
$S$, para que $M\otimes_{\mathcal O_S}\mathcal O_{S'}$ seja localmente
livre de posto $r$, é necessário e suficiente que $S'\longrightarrow S$ se
fatorize por $Z$. Se
\[
\operatorname{rank}_{k(s)}
M_s\otimes_{\mathcal O_{S,s}}k(s)\leq r
\quad\text{para todo $s$},
\]
então $Z$ é um subpré-esquema fechado de $S$ (supondo-se $M$ coerente).
\src{264}
Com efeito, o raciocínio precedente nos reduz ao caso em que a desigualdade
(*) vale para todo $s\in S$ (substituindo, se necessário, $S$ pela parte
aberta formada pelos $s$ em que a desigualdade é verificada). Pode-se então
supor que $M$ se insira em uma sequência exata
\[
\mathcal O_S^q\longrightarrow\mathcal O_S^r\longrightarrow M\longrightarrow0,
\]
e a condição considerada sobre os $S'$ sobre $S$ significa também que, na
sequência exata correspondente
$\mathcal O_{S'}^q\longrightarrow\mathcal O_{S'}^r
\longrightarrow M'\longrightarrow0$, a segunda seta é um isomorfismo, isto
é, a primeira é nula. Vê-se então que o subpré-esquema fechado $Z$ de $S$
definido pelo Ideal gerado pelos coeficientes da matriz que define o
homomorfismo $\mathcal O_S^q\longrightarrow\mathcal O_S^r$ satisfaz a
condição desejada.

Retomando então a demonstração de (3.4) no ponto em que a deixamos, denote-se
por $Z_n$ o subesquema fechado de $S$ associado, em virtude de (3.6), ao
Módulo $M_n$ e ao inteiro $r=P(n)$, e por $Z'_N$ o
$\operatorname{Inf}$ dos $Z_n$ para $\nu\leq n\leq N$. Então os $Z'_N$
formam uma sequência decrescente de subesquemas fechados de $Z$, portanto
necessariamente estacionária. Seja $Z$ o valor constante dos $Z'_N$ para $N$
grande. Esse é o $Z$ procurado em (3.4). Isso conclui a demonstração de (3.4)
e, portanto, também a de (3.3).

Demonstramos, pois, que
$\underline{\operatorname{Quot}}^P_{\mathcal F/X/S}$ é representável por um
$S$-pré-esquema $Q$ que é uma união crescente de subpré-esquemas abertos
$Q_\nu$ quase-projetivos sobre $S$. Para ir mais longe, é preciso invocar o
teorema 2.1, do qual se conclui facilmente que $Q$ é quase-compacto (pois é a
imagem do pré-esquema $S'$ de tipo finito sobre $S$ que parametriza a família
dos feixes quocientes dos $\mathcal F_K$ com polinômio de Hilbert $P$).
Portanto, $Q$ é igual a um dos $Q_\nu$, logo quase-projetivo sobre $S$. Para
demonstrar que é projetivo sobre $S$, resta demonstrar que é próprio sobre
$S$; para isso, basta invocar o critério valuativo de properidade sob a forma
[5], II, 7.3.8. Basta verificar o seguinte:

\result{LEMA}{3.7} Sejam $S$ o espectro de um anel de valorização discreta,
$s$ seu ponto genérico, $X$ um pré-esquema sobre $S$, $\mathcal F$ um Módulo
quase-coerente sobre $X$ e $\mathcal G_s$ um Módulo quociente quase-coerente
de $\mathcal F_s=\mathcal F\otimes_{\mathcal O_S}k(s)$ sobre $X_s$. Então
existe um único Módulo quociente quase-coerente $\mathcal G$ de $\mathcal F$,
plano sobre $S$, cuja restrição a $X_s$ seja $\mathcal G_s$.
\src{265}
Com efeito, se $\mathcal G_s=\mathcal F_s/\mathcal K_s$, basta considerar o
maior subfeixe $\mathcal K$ de $\mathcal F$ que induz $\mathcal K_s$, [5],
I, 9.4.2, e tomar $\mathcal G=\mathcal F/\mathcal K$. Verifica-se facilmente
que esse feixe convém.

O teorema 3.2 e, por conseguinte, (3.1) estão completamente demonstrados.

A demonstração prova, ao mesmo tempo, o seguinte:

\result{PROPOSIÇÃO}{3.8} Nas condições de (3.2), sejam
$Q=\operatorname{Quot}^P_{\mathcal F/X/S}$, $X_Q=X\times_SQ$,
$\mathcal F_Q=\mathcal F\otimes_{\mathcal O_S}\mathcal O_Q$, e seja
$\mathcal G$ o quociente coerente de $\mathcal F_Q$, plano sobre $Q$, com
polinômio de Hilbert relativo $P$, tal que $(Q,\mathcal G)$ represente o
funtor $\underline{\operatorname{Quot}}^P_{\mathcal F/X/S}$. Existe um
inteiro $\nu$ tal que, para $n\geq\nu$,
$(f_Q)_*(\mathcal G(n))$ seja um Módulo sobre $Q$ localmente livre de posto
$P(n)$ e muito amplo relativamente a $S$, isto é, que defina uma imersão de
$Q$ em um esquema grassmanniano $\operatorname{Grass}_{P(n)}(M)$ sobre $S$.
A fortiori, para $n\geq\nu$, o feixe
\[
\bigwedge^{P(n)}(f_Q)_*(\mathcal G(n))
\]
sobre $Q$ é invertível e muito amplo relativamente a $S$.

Pode-se, com efeito, reduzir, como em (3.2), ao caso em que $\mathcal F$ seja
plano sobre $S$; então basta tomar um inteiro $\nu$ tal que $A_\nu=A$ (com as
notações precedentes).

O caso de aplicação mais importante de (3.2) é aquele em que
$\mathcal F=\mathcal O_X$. Escreve-se então
\[
\operatorname{Quot}_{\mathcal O_X/X/S}=\operatorname{Hilb}_{X/S},
\qquad
\operatorname{Quot}^P_{\mathcal O_X/X/S}=\operatorname{Hilb}^P_{X/S},
\]
logo se tem uma decomposição
\[
\operatorname{Hilb}_{X/S}=\coprod_P\operatorname{Hilb}^P_{X/S}.
\]
Por definição, $\operatorname{Hilb}_{X/S}$ representa o funtor
$\underline{\operatorname{Hilb}}_{X/S}(S')$, conjunto dos subpré-esquemas
fechados de $X'=X\times_SS'$ que são planos sobre $S'$, e
$\operatorname{Hilb}^P_{X/S}$ representa o subfuntor correspondente aos
subpré-esquemas fechados que admitem um polinômio de Hilbert dado $P$. Esses
pré-esquemas são também chamados, respectivamente, \emph{pré-esquema de
Hilbert} de $X$ sobre $S$ e \emph{pré-esquema de Hilbert de índice $P$}. A
terminologia é justificada pelo papel desempenhado na teoria pelos polinômios
de Hilbert. Sua diferença de natureza em relação às variedades de Chow
clássicas (destinadas a parametrizar ciclos, e não variedades) é da mesma ordem
que aquela entre
\src{266}
o anel de Chow das classes de ciclos de uma variedade e o anel das classes
de feixes da variedade (tal como aparece no teorema de Riemann--Roch [1]);
com efeito, observe-se que, quando $X=\mathbf P^r_S$, sendo $S$ o espectro
de um corpo, o conhecimento do polinômio de Hilbert de um Módulo coerente
$\mathcal F$ sobre $X$ equivale também ao conhecimento das classes de Chern
de $\mathcal F$, ou ainda ao da classe de $\mathcal F$ no anel das classes de
feixes coerentes sobre $X$.

\result{OBSERVAÇÕES}{3.9} Observe-se também que a construção dos
$\operatorname{Quot}_{\mathcal F/X/S}$ e
$\operatorname{Quot}^P_{\mathcal F/X/S}$ foi reduzida ao caso em que
$X=\mathbf P^r_S$, $\mathcal F=\mathcal O_X^N$, sendo $\mathcal O_X(1)$ o
feixe muito amplo usual; mais precisamente, os
$\operatorname{Quot}_{\mathcal F/X/S}$ gerais realizam-se como
subpré-esquemas fechados dos precedentes. Como a formação dos
$\operatorname{Quot}^P_{\mathcal F/X/S}$ é evidentemente compatível com as
mudanças de base $S'\longrightarrow S$, vê-se que se está reduzido ao caso em
que, além disso, $S=\operatorname{Spec}(\mathbf Z)$. Em última análise, fica-se,
portanto, reduzido a estudar os esquemas projetivos sobre $\mathbf Z$:
\[
\mathfrak Q^P_{r,N}=
\operatorname{Quot}^{P}_{(\mathcal O_{\mathbf P^r_{\mathbf Z}})^N/
\mathbf P^r_{\mathbf Z}/\operatorname{Spec}(\mathbf Z)},
\]
e, mais particularmente, os esquemas de Hilbert absolutos:
\[
\operatorname{Hilb}^P_r=\mathfrak Q^P_{r,1}.
\]
Um estudo mais detalhado desses esquemas, começando pela determinação de suas
componentes conexas (são conexos?), e de suas componentes irredutíveis
(em virtude de Serre [7], pode haver componentes irredutíveis situadas
inteiramente sobre um número primo $p\neq0$), seria muito interessante.
Recordemos a questão de Weil sobre se as componentes irredutíveis das fibras de
$\operatorname{Hilb}^P_r$ sobre os $s\in\operatorname{Spec}(\mathbf Z)$
correspondem a extensões ``regulares'' do corpo primo, isto é, se são
``relativamente conexas''. Não se exclui que essas questões sejam mais
acessíveis para os esquemas de Hilbert do que para as ``variedades de Chow''.

\fsection{4}{Variantes}

\noindent a. Nas condições de (3.1), seja $U$ um aberto de $X$, e designemos
por $A'$ o subfuntor de
$A=\underline{\operatorname{Quot}}_{\mathcal F/X/S}$, tal que $A'(S')$ seja
o conjunto dos Módulos quocientes $\mathcal G$ de $\mathcal F'$, planos sobre
$S'$, cujo suporte está contido em $U'$. Vê-se imediatamente que $A'$ é
representável por uma parte aberta do pré-esquema
$\operatorname{Quot}_{\mathcal F/X/S}$ que representa $A$. Segue-se que os
teoremas (3.1) e (3.2)
\src{267}
continuam válidos supondo-se que $X$ seja quase-projetivo sobre $S$ em vez de
projetivo sobre $S$, quando se substitui também, nas conclusões, a palavra
``projetivo'' por ``quase-projetivo'', e se designa agora por
$\underline{\operatorname{Quot}}_{\mathcal F/X/S}(S')$ o conjunto dos
quocientes coerentes $\mathcal G$ de $\mathcal F'$, planos sobre $S'$, cujo
suporte é próprio sobre $S'$.

\medskip
\noindent b. De modo geral, podem-se impor aos quocientes $\mathcal G$ de
$\mathcal F'=\mathcal F\otimes_{\mathcal O_S}\mathcal O_{S'}$, planos sobre
$S'$, toda espécie de condições naturais adicionais, estáveis por mudança de
base, obtendo-se assim tantos subfuntores de
$\underline{\operatorname{Quot}}_{\mathcal F/X/S}$ quantos se pretende
representar. O critério usual permite, em muitos casos, demonstrar que ainda se
obtêm funtores representáveis por partes abertas de
$\operatorname{Quot}_{\mathcal F/X/S}$. Isso ocorre, em particular, quando se
impõe uma das seguintes propriedades adicionais:

\smallskip
\noindent $1^\circ$ As dimensões dos ciclos primos associados aos Módulos
$\mathcal G_{s'}$ ($s'\in S'$) induzidos nas fibras $X'_{s'}$ pertencem a um
conjunto dado de inteiros.

\smallskip
\noindent $2^\circ$ (Quando $\mathcal F=\mathcal O_X$, portanto $\mathcal G$
corresponde a um subpré-esquema fechado $Y$ de $X'$); $Y$ é um pré-esquema
suave ([3], IV) sobre $S$, resp. normal sobre $S$ (isto é, as fibras $Y_{s'}$
são normais ``sobre $k(s')$'', isto é, são normais após toda extensão do corpo
de base), resp. (quando $X$ é plano sobre $S$) é uma $S$-interseção completa
local em $X$ relativamente a $S$ (isto é, as fibras $Y_{s'}$ são interseções
completas locais nos $X_{s'}$).

Outras condições fariam intervir propriedades de natureza cohomológica nos
Módulos $\mathcal G_{s'}$ induzidos nos $X'_{s'}$, etc. Naturalmente, a
conjunção de condições, cada uma das quais representada por um aberto $U_i$ de
$\operatorname{Quot}_{\mathcal F/X/S}$, é representada pelo aberto que é sua
interseção. Por exemplo, considerando, para todo $S'$ sobre $S$, o conjunto dos
subpré-esquemas fechados $Y$ de $X'=X\times_SS'$ que são recobrimentos étales
([3], I) de posto dado $r$ de $S'$, obtém-se um funtor contravariante
representável em $S'$.

\medskip
\noindent c. Os pré-esquemas $\operatorname{Hom}_S(X,Y)$,
$\Pi_{X/S}(Z/X)$, $\operatorname{Isom}_S(X,Y)$, definidos em [2], II, C,
n\textsuperscript{o} 2, existem sob hipóteses de projetividade convenientes e
realizam-se como partes abertas de pré-esquemas de Hilbert convenientes. Como se
tem
\[
\operatorname{Hom}_S(X,Y)=\Pi_{X/S}((X\times_SY)/X),
\]
o caso de $\operatorname{Hom}_S(X,Y)$ reduz-se ao de $\Pi_{X/S}(Z/X)$.
Observa-se então que, para todo $S'$ sobre $S$, o conjunto das seções de
$Z'=Z\times_SS'$ sobre $X'=X\times_SS'$ está em correspondência biunívoca com
o conjunto dos subpré-esquemas $\Gamma$ de $Z'$ (necessariamente fechados se
$Z$ é separado sobre $X$) tais que o morfismo $\Gamma\longrightarrow X'$
induzido por $Z'\longrightarrow X'$ seja um isomorfismo. Dessa maneira,
\src{268}
quando $X$ é plano e próprio sobre $S$ e $Z$ é quase-projetivo sobre $S$,
$\Pi_{X/S}(Z/X)$ existe e realiza-se como um subpré-esquema aberto de
$\operatorname{Hilb}_{Z/S}$. Portanto, quando $X$ é projetivo e plano sobre
$S$ e $Y$ é quase-projetivo sobre $S$, $\operatorname{Hom}_S(X,Y)$ existe e
realiza-se como um subpré-esquema aberto de
$\operatorname{Hilb}_{(X\times_SY)/S}$. Quando $X$ e $Y$ são ambos projetivos
sobre $S$, segue-se imediatamente que $\operatorname{Isom}_S(X,Y)$ também
existe e é representado por uma parte aberta de
$\operatorname{Hom}_S(X,Y)$. Do mesmo modo, quando $X$ é plano e projetivo
sobre $S$ e $Y$ é quase-projetivo sobre $S$, o $S$-pré-esquema
$\operatorname{Imm}_S(X/Y)$, correspondente ao subfuntor do funtor representado
por $\operatorname{Hom}_S(X,Y)$ que corresponde aos $S'$-homomorfismos
$X'\longrightarrow Y'$ que são imersões, é também representável por uma
parte aberta de $\operatorname{Hom}_S(X,Y)$.

Seja $\mathcal L$ (resp. $\mathcal M$) um feixe invertível sobre $X$ (resp.
$Y$), muito amplo relativamente a $S$, donde um feixe muito amplo
$\mathcal L\otimes_{\mathcal O_S}\mathcal M$ sobre $X\times_SY$ relativamente a
$S$. Por conseguinte, para todo polinômio $P$ com coeficientes racionais,
$\operatorname{Hilb}^P_{(X\times_SY)/S}$ está definido e é um pré-esquema
quase-projetivo sobre $S$. Ele induz, portanto, sobre
$\operatorname{Hom}_S(X,Y)$ uma parte simultaneamente aberta e fechada,
quase-projetiva sobre $S$, que denotaremos por
$\operatorname{Hom}_S(X,Y)^P$. Assim, as seções de
$\operatorname{Hom}_S(X,Y)^P$ sobre $S$ são os $S$-morfismos
$g:X\longrightarrow Y$ tais que, para todo inteiro $n$, se tenha
\[
\chi\!\left(\left(\mathcal L\otimes_{\mathcal O_X}g^*(\mathcal M)\right)^{
\otimes n}\right)=P(n).
\]

Obtêm-se dessa maneira generalizações do teorema de Matsusaka, que afirma
que os automorfismos de uma variedade projetiva «polarizada» formam um grupo
algébrico, afirmação que assume aqui um sentido evidentemente mais preciso,
pois dispomos de uma definição desse grupo como solução de um problema
universal. Note-se, aliás, que, sobre um corpo algebricamente fechado, o grupo
dos automorfismos considerado anteriormente é o que se deduz do «verdadeiro»
definido aqui passando ao quociente pelos elementos nilpotentes; isso explica
por que há pouca chance de que as construções mais antigas possam ser feitas
sobre um corpo de base não perfeito, pois o Ideal dos elementos nilpotentes que
aparece após uma extensão do corpo de base não é necessariamente «definido
sobre $k$». Essa mesma observação aplica-se, aliás, também à maioria das construções
de estilo antigo.

\fsection{5}{Estudo diferencial dos esquemas de Hilbert}

Ele decorre do resultado seguinte:
\src{269}
\result{PROPOSIÇÃO}{5.1} Sejam $S$ um pré-esquema, $S_0$ um subpré-esquema
definido por um Ideal quase-coerente $\mathcal J$ de quadrado nulo, $X$ um
$S$-pré-esquema, $\mathcal F$ um Módulo quase-coerente sobre $X$,
$X_0=X\times_SS_0$ e
$\mathcal F_0=\mathcal F\otimes_{\mathcal O_S}\mathcal O_{S_0}$, e, por fim,
$\mathcal G_0=\mathcal F_0/\mathcal K_0$ um Módulo quociente quase-coerente de
$\mathcal F_0$, plano sobre $S_0$. Para todo aberto $U$ de $X$, seja
$\mathcal E(U)$ o conjunto dos Módulos quocientes quase-coerentes $\mathcal G$ de
$\mathcal F|U$, planos sobre $S$, e tais que
$\mathcal G\otimes_{\mathcal O_S}\mathcal O_{S_0}=\mathcal G_0$; assim, quando
$U$ varia, os $\mathcal E(U)$ são as seções de um feixe $\mathcal E$ sobre
$U$. Posto isso, o feixe de grupos
\[
\mathcal A=\operatorname{Hom}_{\mathcal O_{X_0}}
(\mathcal K_0,\mathcal G_0\otimes_{\mathcal O_{S_0}}\mathcal J)
\]
age naturalmente sobre $\mathcal E$, que se torna assim um feixe
«formalmente homogêneo principal sob $\mathcal A$» (isto é, para todo aberto
$U$ em $X$, $\mathcal E(U)$ é vazio ou um conjunto homogêneo principal sob
$\mathcal A(U)$).

Conclui-se daí:

\result{COROLÁRIO}{5.2} Suponhamos que exista localmente sobre $X$ um
prolongamento $\mathcal G$ de $\mathcal G_0$ a um quociente de $\mathcal F$
plano sobre $S$ (isto é, que as fibras do feixe $\mathcal E$ sejam não vazias).
Então existe uma classe de obstrução canônica
\[
c(\mathcal G_0)\in H^1(X,\mathcal A),
\]
cujo anulamento é necessário e suficiente para a existência de um
prolongamento global $\mathcal G$ de $\mathcal G_0$ a um quociente de
$\mathcal F$ plano sobre $S$. Se essa classe se anula, então o conjunto
$\mathcal E(X)$ de todos os prolongamentos possíveis é um conjunto homogêneo
principal sob
\[
\mathcal A(X)=\operatorname{Hom}_{\mathcal O_X}
(\mathcal K_0,\mathcal G_0\otimes_{\mathcal O_{S_0}}\mathcal J).
\]

A existência do prolongamento global fica então garantida, em particular, se
$H^1(X,\mathcal A)=0$.

\result{COROLÁRIO}{5.3} Suponhamos que
$Q=\operatorname{Quot}_{\mathcal F/X/S}$ exista (cf. 4 (a)) -- por exemplo,
que $X$ seja quase-projetivo sobre um $S$ localmente noetheriano e que
$\mathcal F$ seja coerente. Seja $x\in Q$ correspondente a uma extensão
residual $K=k(x)$ de $k(s)$ ($s\in S$):
\src{270}
portanto, $x$ é definido por um Módulo quociente coerente
$\mathcal G_0=\mathcal F_0/\mathcal K_0$ do Módulo
$\mathcal F_0=\mathcal F_K$ sobre o $K$-pré-esquema $X_K$. Seja $\mathcal A$
o feixe coerente sobre $X_K$ definido por
\[
\mathcal A=\operatorname{Hom}_{\mathcal O_{X_0}}(\mathcal K_0,\mathcal G_0).
\]
Então o espaço tangente de Zariski da fibra $Q_s$ no ponto $x$ (dual sobre
$K$ de $\mathfrak m/\mathfrak m^2$, onde $\mathfrak m$ é o ideal maximal de
$\mathcal O_{Q_K,x}$) é canonicamente isomorfo a $H^0(X_K,\mathcal A)$.

O resultado que fornece o espaço tangente de Zariski pode ser generalizado e
fornece uma caracterização, para um $S$-morfismo dado
$g:S'\longrightarrow Q$, isto é, uma seção $g'$ de $Q'=Q\times_SS'$ sobre
$S'$, do Módulo
\[
\Omega=g^*(\Omega^1_{Q/S})=g'^*(\mathcal J/\mathcal J^2)
\]
(onde $\mathcal J$ é o Ideal sobre $Q'$ definido pela seção $g'$ de $Q'$ sobre
$S'$), pela fórmula funtorial no Módulo coerente $\mathcal M$ sobre $S'$:
\[
\operatorname{Hom}_{\mathcal O_{S'}}(\Omega,\mathcal M)
\simeq H^0(X',\mathcal A),
\]
onde $\mathcal A$ é ainda o Módulo sobre $X'=X\times_SS'$ definido por
\[
\mathcal A=\operatorname{Hom}_{\mathcal O_{X'}}
(\mathcal K,\mathcal G\otimes_{\mathcal O_S}\mathcal M),
\]

($\mathcal G=\mathcal F'/\mathcal K$ sendo o Módulo quociente de
$\mathcal F'=\mathcal F\otimes_{\mathcal O_S}\mathcal O_{S'}$ que corresponde
a $g$). Com efeito, basta aplicar (5.1), substituindo ali $S_0$ por $S'$ e $S$
pelo pré-esquema $D(\mathcal M)=(S',\mathcal O_{S'}+\mathcal M)$, onde
$\mathcal M$ é considerado como um Ideal de quadrado nulo.

Quando, em (5.1), se tem $\mathcal F=\mathcal O_X$, portanto o dado de
$\mathcal G_0$ corresponde ao dado de um subpré-esquema fechado $Y_0$ de
$X_0$, plano sobre $S_0$, definido pelo Ideal
$\mathcal J_0=\mathcal K_0$, então (*) fornece
\[
\mathcal A=\operatorname{Hom}_{\mathcal O_{X_0}}
(\mathcal J_0/\mathcal J_0^2,
\mathcal O_{Y_0}\otimes_{\mathcal O_{S_0}}\mathcal J),
\]
onde $\mathcal J/\mathcal J^2$ se interpreta como o feixe conormal a $Y_0$ em
$X_0$, que se denota também por
\src{271}
$\check{\mathcal N}_{Y_0/X_0}$; é então conveniente considerar $\mathcal A$
como um Módulo sobre $Y_0$ e calcular $H^0$ e $H^1$ sobre $Y_0$. Ademais, se
$Y_0$ é localmente uma interseção completa em $X_0$, sendo $X$ plano sobre
$S$, então, em (5.1), fica garantida a possibilidade de prolongamento local;
por outro lado, $\mathcal J/\mathcal J^2$ é localmente livre sobre $Y_0$ e
pode-se escrever
\[
\mathcal A=\mathcal N_{Y_0/X_0}\otimes_{\mathcal O_{S_0}}\mathcal J,
\]
onde o primeiro fator do segundo membro é o feixe normal a $Y_0$ em $X_0$.
Utilizando o critério fundamental de lisura ([3], III, 3.1), obtém-se, por
exemplo:

\result{COROLÁRIO}{5.4} Nas condições de (5.3), suponhamos que
$\mathcal F=\mathcal O_X$, que $X$ seja plano sobre $S$ e que o
subpré-esquema fechado $Y_0$ de $X_0$ que corresponde a $\mathcal G_0$ seja
localmente uma interseção completa. Então o espaço tangente de Zariski a $Q_s$
em $x$ é canonicamente isomorfo a
$H^0(Y_0,\mathcal N_{Y_0/X_0})$. Se
$H^1(Y_0,\mathcal N_{Y_0/X_0})=0$, então o pré-esquema de Hilbert
$\operatorname{Hilb}_{X/S}$ é liso sobre $S$ no ponto $x$.
($\mathcal N_{Y_0/X_0}$ é o feixe normal a $Y_0$ em $X_0$.)

\result{OBSERVAÇÃO}{5.5} Este enunciado aplica-se, em particular, quando $Y_0$
é uma interseção completa em $X_0$ definida por uma equação, isto é, é um
``divisor de Cartier'' positivo. Então $\mathcal N_{Y_0/X_0}$ é isomorfo ao
feixe sobre $Y_0$ induzido pelo feixe invertível $\mathcal J^{-1}$ sobre $X_0$
definido pelo divisor $Y_0$. É essa a situação encontrada, em particular, no
estudo das famílias de divisores positivos sobre uma variedade projetiva não
singular $X_0$. O isomorfismo entre o espaço tangente de Zariski no ponto $x$
de $Q$ (ou, se se preferir, do aberto $D$ de $Q$ que corresponde aos divisores)
e $H^0(Y_0,\mathcal N_{Y_0/X_0})$ era conhecido em geometria algébrica
clássica sob o nome de ``homomorfismo característico'' do primeiro no segundo.
Ele só era definido quando $x$ era um ponto liso da variedade de parâmetros
$T$ de uma ``família contínua completa'' de divisores, isto é, de nosso ponto
de vista, de uma componente irredutível do esquema $D$, munida da estrutura
reduzida induzida. O espaço tangente a $T$ em $x$ é então um subespaço do
espaço tangente a $D$ em $x$; portanto, o homomorfismo característico clássico
é de fato injetivo, mas só é sobrejetivo sob condições adicionais, por exemplo,
se $D$ é integral em $x$. De fato, Zappa [8] construiu um exemplo
\src{272}
(com $X$ uma superfície projetiva não singular sobre o corpo dos complexos) no
qual, mesmo no ponto genérico de $T$, o homomorfismo característico não é
sobrejetivo. Isso significa, portanto, que $D$ não é integral nem sequer no
ponto genérico da componente irredutível considerada. Isso mostra de maneira
particularmente impressionante como as variedades com elementos nilpotentes
são necessárias para compreender fenômenos da mais clássica teoria das
superfícies.

\subsection*{5.6}

Em (5.4), deu-se um critério de lisura, aplicável em particular aos esquemas de
divisores. Kodaira deu em [6] um critério diferente, a saber, o anulamento de
$H^1(X_0,\mathcal L)$, onde $\mathcal L=\mathcal J_0^{-1}$ é o feixe
invertível sobre $X_0$ definido pelo divisor $Y_0$; critério válido quando $S$
é o espectro de um corpo de característica 0 e demonstrado em [6] por via
transcendente no caso em que o corpo de base é $\mathbf C$. Observemos aqui
que, em geral, sendo $S$ novamente arbitrário, a condição de Kodaira é
suficiente para que o morfismo canônico
\[
D\longrightarrow\operatorname{Pic}_{X/S}
\]
do pré-esquema dos divisores para o pré-esquema de Picard de $X/S$ seja liso
no ponto $x$ considerado (como se verifica facilmente pelo critério usual de
lisura, uma vez estabelecida a existência de $\operatorname{Pic}_{X/S}$). Se,
portanto, além disso, $\operatorname{Pic}_{X/S}$ é liso sobre $S$ no ponto
imagem de $x$ (por exemplo, se $\operatorname{Pic}_{X/S}$ é liso sobre $S$),
então $D$ é liso sobre $S$ em $x$. Por outro lado, Cartier demonstrou que todo
pré-esquema em grupos localmente de tipo finito sobre um corpo $k$ de
característica 0 é liso sobre $k$. Combinando esses dois resultados, recupera-se
o resultado de Kodaira. Observe-se que decorre dessas observações que, sobre um
corpo $K$ de característica $p>0$, quando $\operatorname{Pic}_{X/S}$ não é
liso sobre $k$ (o que ocorre quando $X$ é a superfície de Igusa), a condição
$H^1(X_0,\mathcal L)=0$ implica, ao contrário, que $D$ não é liso em $x$, e
nem sequer é reduzido em $x$ se $K$ é algebricamente fechado.

Para terminar, assinalemos ainda o resultado seguinte, que desempenha um papel
importante no estudo diferencial dos espaços fibrados:

\result{PROPOSIÇÃO}{5.7} Sejam $X$ um pré-esquema finito e plano sobre $S$
localmente noetheriano, e $Z$ um pré-esquema sobre $S$, tal que
$\Pi_{X/S}(Z/X)$ exista (o que ocorre se $Z$ é quase projetivo sobre $X$). Se,
então, $Z$ é liso sobre $X$, $\Pi_{X/S}(Z/X)$ é liso sobre $S$.

É uma consequência imediata da definição e do critério usual de lisura ([3],
III, 3.1). Observemos que, quando $X$ é finito e plano sobre $S$, a questão da
existência de $Z/K$ pode ser tratada de maneira muito elementar, sem
\src{273}
utilizar a teoria dos esquemas de Hilbert. Verifica-se, por exemplo, que, se
$X$ é radicial sobre $S$, então $\Pi_{X/S}(Z/X)$ existe sem nenhuma restrição
sobre $Z$. Por exemplo, seja $T$ um $S$-pré-esquema, e seja $T_n$ a
``vizinhança infinitesimal de ordem $n$'' da diagonal de $T\times_ST$ em
$T\times_ST$, munida dos morfismos $p_1,p_2:T_n\longrightarrow T$ induzidos
pelas duas projeções. Considera-se $T_n$ como um pré-esquema finito sobre $T$
por meio de $p_1$; suponhamo-lo, além disso, plano sobre $T$ (o que ocorre se
$T$ é liso sobre $S$). Para todo pré-esquema $X$ sobre $T$, ponhamos
\[
(X/T/S)^{(n)}=
\Pi_{T_n/S}\bigl(p_2^*(X/T)/T_n\bigr),
\]

este é um pré-esquema sobre $T$ chamado \emph{fibrado dos germes de seções de
ordem $n$ de $X$ sobre $T$} (relativamente a $S$). Ele depende funtorialmente
de $X$ e é liso sobre $T$ se $X$ o é.

\fsection{6}{Relações com a noção de norma e os produtos simétricos}

Sejam $S$ um pré-esquema, $X$ e $Y$ $S$-pré-esquemas,
\[
u:(X/S)^n\longrightarrow Y
\]
um $S$-morfismo simétrico da $n$-ésima potência cartesiana de $X/S$ em $Y$.
Suporemos, para simplificar, $S$ localmente noetheriano, e $X$ e $Y$ de tipo
finito sobre $S$. Pode-se então, a todo Módulo coerente $\mathcal F$ sobre $X$,
com suporte finito sobre $S$, plano sobre $S$ e de posto relativo sobre $S$
igual a $n$ (isto é, tal que $f_*(\mathcal F)$ seja um Módulo localmente livre
de posto $n$ sobre $S$), fazer corresponder de maneira natural uma seção de $Y$
sobre $S$:
\[
\pi^u_{X/S}(\mathcal F)\in\Gamma(Y/S).
\]
Não daremos aqui a definição formal, limitando-nos a assinalar que o formalismo
ao qual se chega é uma generalização natural do formalismo usual das normas e
dos traços. Quando a $n$-ésima potência simétrica de $X$ sobre $S$ existe (por
exemplo, quando as órbitas do grupo simétrico $\mathfrak S_n$ que opera sobre
$(X/S)^n$ estão contidas em abertos afins), pode-se tomar para $Y$ essa
potência simétrica $\operatorname{Symm}^n_S(X)$, e obtém-se um elemento canônico
\[
\pi_{X/S}(\mathcal F)\in\Gamma(\operatorname{Symm}^n_S(X)/S),
\]
\src{274}
que permite recuperar os $\pi^u_{X/S}(\mathcal F)$. Outro caso importante é
aquele em que $X$ é um monoide comutativo sobre $S$, e $X=Y$, provindo o
morfismo $u$ da lei de composição de $X$. Escrever-se-á então simplesmente
$\pi(\mathcal F)$ para a seção de $X$ sobre $S$ associada ao Módulo
$\mathcal F$ sobre $X$.

Suponhamos agora que se tenha um Módulo coerente $\mathcal F$ sobre $X$, tal
que $\operatorname{Quot}_{\mathcal F/X/S}$ exista, ou pelo menos tal que o
funtor $\underline{\operatorname{Quot}}^n_{\mathcal F/X/S}$, que associa a
todo $S'$ sobre $S$ o conjunto dos feixes quocientes coerentes $\mathcal M$ de
$\mathcal F'=\mathcal F\otimes_{\mathcal O_S}\mathcal O_{S'}$ que são planos
sobre $S'$ e de posto relativo $n$, seja representável por um $S$-pré-esquema
$\operatorname{Quot}^n_{\mathcal F/X/S}$. (Quando $X$ é quase projetivo sobre
$S$, então $\operatorname{Quot}^n_{\mathcal F/X/S}$ existe de fato e não é
outro, com as notações do número 3, senão
$\operatorname{Quot}^P_{\mathcal F/X/S}$, onde $P$ é o polinômio reduzido ao
termo constante $n$.) Como a formação dos $\pi^u_{X/S}(\mathcal M)$ é
compatível com a mudança de base, obtém-se então um morfismo canônico
\[
\pi^u_{X/S}:\operatorname{Quot}^n_{\mathcal F/X/S}\longrightarrow Y,
\]
em particular, se a $n$-ésima potência simétrica de $X$ sobre $S$ existe:
\[
\pi_{X/S}:\operatorname{Quot}^n_{\mathcal F/X/S}
\longrightarrow\operatorname{Symm}^n_S(X).
\]
O caso mais importante é aquele em que $\mathcal F=\mathcal O_X$, que fornece
um morfismo:
\[
\pi_{X/S}:\operatorname{Hilb}^n_{X/S}
\longrightarrow\operatorname{Symm}^n_S(X).
\]
Este é evidentemente um isomorfismo para $n=0$ ou $n=1$. Mas, para
$n\geq1$, mesmo se $S$ é o espectro
de um corpo $k$, e quando $X$ é liso sobre $S$, ele não é, em geral, um
isomorfismo, nem sequer um morfismo injetivo, pois um subesquema de dimensão 0
de $X$ (correspondente, por exemplo, a um ideal $\mathfrak i$ primário para o
ideal maximal em um anel local $\mathcal O_{X,x}$, para um ponto fechado $x$ de
$X$) não é conhecido quando se conhece o ciclo que ele define (no caso, quando
se conhece a codimensão sobre $k$ de $\mathfrak i$ em $\mathcal O_{X,x}$).
Pode-se dizer apenas o seguinte (sendo $S$ novamente arbitrário):

\medskip
\noindent a. Se $X$ é liso sobre $S$, então o morfismo-norma define um
isomorfismo do aberto de $\operatorname{Hilb}^n_{X/S}$ que corresponde à
classificação dos recobrimentos étales de posto $n$ contidos em $X$ (cf.
n\textsuperscript{o} 4, (b)) com o aberto de
$\operatorname{Symm}^n_S(X)$ que corresponde aos $n$-ciclos sem componentes
múltiplas.
\src{275}

\medskip
\noindent b. Se, além disso, $X$ tem dimensão relativa 1 sobre $S$, então o
morfismo norma define até mesmo um isomorfismo de
$\operatorname{Hilb}^n_{X/S}$ com $\operatorname{Symm}^n_{X/S}$.

Este último fato deve-se a que um subesquema de dimensão 0 de uma curva
algébrica não singular é conhecido quando se conhece o ciclo correspondente.
A mesma observação aplica-se, aliás, mais geralmente aos divisores de Cartier
positivos sobre um esquema algébrico não singular (e não se exclui que, nesse
caso muito particular, a variedade de Chow forneça a mesma coisa que o esquema
de Hilbert).

\fsection{7}{Complementos e questões}

Como observou J.-P. Serre, resulta de um exemplo bem conhecido de Nagata que se
pode encontrar um esquema $S$, espectro de um corpo $k$, um $S$-esquema $S'$,
espectro de uma extensão quadrática $k'$ de $k$, enfim, um $S'$-esquema $X$
próprio e liso (mas não projetivo) de dimensão 3, tal que
$\Pi_{S'/S}(X/S')$ não existe. Isso implica, a fortiori, que o esquema de
Hilbert $\operatorname{Hilb}^2_{X/S}$ não existe (nem sequer o $k$-esquema que
representaria os recobrimentos étales de posto 2 de $S$ contidos em $X$, nem,
a fortiori, o quadrado simétrico de $X$, cf. número precedente). Isso impõe,
portanto, limitações sérias às possibilidades de construções não projetivas em
geometria algébrica. (É, contudo, plausível que tais limitações não se
apresentem em geometria analítica, assim como não se apresentam em geometria
formal (cf. [2], II).) Em contrapartida, se $X$ é um esquema próprio sobre o
espectro $S$ de um corpo $k$, e se $Z$ é quase projetivo sobre $X$, então
$\Pi_{X/S}(Z/X)$ existe e é um esquema, soma disjunta de uma sequência de
esquemas quase projetivos sobre $S$ (como no caso projetivo (3.1)). Para vê-lo,
reduz-se, com efeito, ao caso em que o próprio $X$ é projetivo, dominando $X$
por um $S$-esquema projetivo $X'$; não daremos aqui os detalhes da demonstração,
que utiliza também o resultado de fatoração de um morfismo finito assinalado em
[2], I, A, n\textsuperscript{o} 2 (b). O êxito do método deve-se ao fato de que,
sendo $S$ o espectro de um corpo, o $X'$ que intervém no lema de Chow será
automaticamente plano sobre $S$.

Ignoro se o resultado continua válido sem hipótese sobre $S$, supondo apenas
$X$ próprio e plano sobre $S$, e $Z$ quase projetivo sobre $X$. Um caso
importante nas aplicações é aquele em que $Z$ é um subesquema fechado de $X$;
se, então, $\Pi_{X/S}(Z/X)$ existe, ele é necessariamente um subesquema fechado
de $S$. Pode-se
\src{276}
construí-lo diretamente de maneira bastante simples quando $X$ é projetivo
sobre $S$, sem utilizar a teoria dos esquemas de Hilbert, e o método empregado
mostra, mais geralmente, que, se $Z$ é afim sobre $X$, então
$\Pi_{X/S}(Z/X)$ existe e é afim sobre $S$. Ele mostra igualmente que, se $X$
é próprio e plano sobre $S$ (não necessariamente projetivo sobre $S$), então,
para todo fibrado vetorial $Z$ localmente trivial sobre $X$,
$\Pi_{X/S}(Z/X)$ existe e é um fibrado vetorial sobre $S$. Seria desejável que
esses resultados fossem retomados e unificados.

\section*{BIBLIOGRAFIA}
\addcontentsline{toc}{section}{Bibliografia}

\begin{description}
\item[{[1]}] BOREL (A.) e SERRE (J.-P.). -- Le théorème de Riemann--Roch,
\emph{Bull. Soc. math. France}, v. 86, 1958, p. 97--136.

\item[{[2]}] GROTHENDIECK (Alexander). -- Technique de construction et
théorèmes d'existence en géométrie algébrique, I, II, III, \emph{Séminaire
Bourbaki}, v. 13, 1959/60, n\textsuperscript{o} 190, 29 p.,
n\textsuperscript{o} 195, 22 p.; e v. 14, 1960/61,
n\textsuperscript{o} 212, 20 p.

\item[{[3]}] GROTHENDIECK (Alexander). -- Séminaire de Géométrie algébrique,
I, II, III, IV. -- Paris, Institut des hautes Études scientifiques, 1960/61
(mimeografado).

\item[{[4]}] GROTHENDIECK (Alexander). -- Technique de construction en
géométrie analytique, IV, V, \emph{Séminaire Cartan}, v. 13, 1960/61,
n\textsuperscript{o} 11 e 12.

\item[{[5]}] GROTHENDIECK (A.) e DIEUDONNÉ (J.). -- \emph{Éléments de
géométrie algébrique}, I: Le langage des schémas. -- Paris, Presses
universitaires de France, 1960 (Institut des hautes Études scientifiques,
Publications mathématiques, 4); II, III, IV (a aparecer).

\item[{[6]}] KODAIRA (K.). -- Characteristic linear systems of complete
continuous systems, \emph{Amer. J. of Math.}, v. 78, 1956, p. 716--744.

\item[{[7]}] SERRE (Jean-Pierre). -- Exemples de variétés projectives en
caractéristique $p$ non relevables en caractéristique 0, \emph{Proc. Nat. Acad.
Sc. U.S.A.}, v. 47, 1961, p. 108--109.

\item[{[8]}] ZAPPA (G.). -- Sull'esistenza, sopra le superficie algebriche, di
sistemi continui completi infiniti, la cui curva è a serie caratteristica
incompleta, \emph{Pont. Acad. Sc. Acta}, v. 9, 1945, p. 91--93.
\end{description}

\section*{ADENDO}
\addcontentsline{toc}{section}{Adendo}
\begin{center}
{\small [acrescentado durante a correção das provas]}
\end{center}

Parece agora que as conjecturas de [2], III, n\textsuperscript{o} 8, são falsas,
mesmo para variedades não singulares sobre um corpo de característica 0, tanto
no que concerne à existência quanto à quase projetividade do quociente, e mesmo
quando $G$ opera com um grafo fechado.

\par

\def\FgaUnitKey{e221}\clearpage
\chapter*{Errata da exposição n\textsuperscript{o} 221}
\addcontentsline{toc}{chapter}{Errata da exposição n\textsuperscript{o} 221}
\setcounter{footnote}{3}
\markboth{Errata da exposição 221}{Errata da exposição 221}
\sid{e221}
\src{302}

\noindent\textbf{6. TDTE IV: Os esquemas de Hilbert.}

\begin{center}
[Seminário Bourbaki, t. 13, 1960/61, n\textsuperscript{o} 221, 28 p.]
\end{center}

\noindent\textbf{Página 6, teorema 2.1, 3\textsuperscript{a} linha.} --- Em vez de
«é necessário», leia-se «é necessário e suficiente».

\noindent\textbf{Página 15, linha 13.} --- Em vez de «$P\neq Q$», leia-se
«$P<Q$».

\noindent\textbf{Linha 15.} --- Em vez de «Se, ao contrário, $P=Q$», leia-se
«No caso contrário, teremos $P(n)\geq Q(n)$ para $n$ grande».

\noindent\textbf{Linha 18, fórmula (*).} --- Em vez de «$<$» e «$>$»,
leia-se «$\leq$» e «$\geq$».

\noindent\textbf{5.ª linha a partir de baixo.} --- Em vez de «$r$», leia-se
«$Z$».

\noindent\textbf{Página 18, observação 3.9.} --- Assinalemos que o estudo das
componentes conexas dos esquemas de Hilbert sobre um corpo algebricamente
fechado foi realizado por HARTSHORNE. Esse autor demonstra que os
$\operatorname{Hilb}^{P}_{r}$ são conexos e determina os pares $(r,P)$ para
os quais $\operatorname{Hilb}^{P}_{r}\neq\varnothing$\footnote{HARTSHORNE
(R.). --- Connectedness of the Hilbert scheme (Tese, Harvard, 1962).}.

\noindent\textbf{Página 23, 5.ª linha a partir de baixo.} --- Em vez de «uma
componente», leia-se «de uma componente».

\noindent\textbf{Páginas 23 e 24, observação 5.5.} --- A propósito do exemplo
de ZAPPA, assinalemos que MUMFORD acaba até mesmo de construir uma componente
irredutível do esquema de Hilbert para $\mathbf P^3_{\mathbf C}$ (cujos pontos
gerais representam curvas não singulares de grau 14 e gênero 24), que não é
reduzida em seu ponto genérico. Realizando o blow-up de uma das curvas obtidas,
ele obtém também um esquema projetivo regular de dimensão 3 sobre $\mathbf C$,
cujo esquema formal de moduli não é reduzido em seu ponto genérico ou, o que
equivale ao mesmo, cuja variedade local de moduli, no sentido da geometria
analítica, sobre $\mathbf C$ (cf. Seminário Cartan, t. 13, 1960/61), não é
reduzida em nenhum de seus pontos.

\par

\def\FgaUnitKey{232}\unit{Exposição 232 -- Os esquemas de Picard: teoremas de existência}
\sid{232}
\markboth{Exposição 232}{Exposição 232}
\src{143}
\seminares{14.º ano, 1961/62, n.º 232\\Fevereiro de 1962}
\paperhead{TÉCNICA DE DESCIDA E TEOREMAS DE EXISTÊNCIA EM GEOMETRIA ALGÉBRICA\\
V. OS ESQUEMAS DE PICARD: TEOREMAS DE EXISTÊNCIA}%
  {por Alexander GROTHENDIECK}

\fsection{1}{Grupos e funtores de Picard relativos}

Para todo pré-esquema (mais geralmente, todo espaço anelado) $X$, chamamos
\emph{grupo de Picard} (absoluto) de $X$, e denotamos por
$\operatorname{Pic}(X)$, o grupo das classes, a menos de isomorfismo, de
Módulos invertíveis (isto é, localmente isomorfos a $\mathcal O_X$) sobre $X$.
Tem-se, portanto, um isomorfismo canônico
\begin{equation}
\operatorname{Pic}(X)\xrightarrow{\sim}H^1(X,\mathcal O_X^*),
\tag{1.1}
\end{equation}
onde $\mathcal O_X^*$ denota o feixe das unidades de $\mathcal O_X$ (que se
identifica, com efeito, com o feixe dos automorfismos do Módulo invertível
padrão $\mathcal O_X$). Observemos que
$X\longmapsto\operatorname{Pic}(X)$ é evidentemente um funtor contravariante
em $X$, e que o isomorfismo (1.1) é funtorial.

Se $X$ é um pré-esquema sobre outro pré-esquema $S$, então, quando $S'$ varia
na categoria $(\mathrm{Sch})/S$ dos pré-esquemas sobre $S$, obtém-se, pelo que
precede, um funtor contravariante
$S'\longmapsto\operatorname{Pic}(X\times_S S')$. Esse funtor não tem qualquer
possibilidade de ser «representável» ([4], II, A), pois, devido à existência de
automorfismos não triviais dos Módulos invertíveis que se pretende classificar,
esse funtor não é de «natureza local» ([5], IV, 5.4). Convém, portanto,
«torná-lo local», introduzindo, de modo geral, para todo pré-esquema relativo
$X/S$, um grupo de natureza relativa,
\begin{equation}
\operatorname{Pic}'(X/S)=H^0\bigl(S,R^1f_*(\mathcal O_X^*)\bigr),
\tag{1.2}
\end{equation}
(onde $f:X\longrightarrow S$ é o morfismo estrutural) (compare-se [4], II,
C 3). No lugar citado, esse grupo é chamado grupo de Picard relativo; será
preferível aqui chamá-lo \emph{grupo de Picard relativo restrito} de $X/S$,
por razões que aparecerão. Quando $S'$ varia em $(\mathrm{Sch})/S$,
$S'\longmapsto\operatorname{Pic}'(X\times_S S'/S')$ é um funtor contravariante
em $S'$, denotado também por $\operatorname{Pic}'_{X/S}$, portanto definido
essencialmente pela fórmula
\begin{equation}
\operatorname{Pic}'_{X/S}(S')=
\operatorname{Pic}'(X\times_S S'/S').
\tag{1.3}
\end{equation}
\src{144}
Este funtor é agora «de natureza local», visto que se fez o necessário para
isso. Intuitivamente, o segundo membro de (1.3) interpreta-se como o conjunto
das «famílias algébricas» de classes de feixes invertíveis sobre (as fibras de)
$X/S$, indexadas pelo pré-esquema de parâmetros $S'/S$. Quando o funtor
$\operatorname{Pic}'$ é representável, o pré-esquema sobre $S$ que o representa
é denotado por $\operatorname{Pic}'_{X/S}$ e chamado o
\emph{pré-esquema de Picard} de $X$ sobre $S$; portanto, teremos
\begin{equation}
\operatorname{Hom}_S(S',\operatorname{Pic}'_{X/S})
\simeq\operatorname{Pic}'_{X/S}(S')
=\operatorname{Pic}'(X\times_S S'/S').
\tag{1.4}
\end{equation}

Há, contudo, casos importantes em que $\operatorname{Pic}'_{X/S}$ não é
representável (exemplo: variedade de «Brauer--Severi» sobre um corpo $k$, sem
ponto racional sobre $k$), e nos quais, contudo, existe uma definição natural
de um pré-esquema de Picard relativo. Isso decorre do fato de que, na definição
do funtor $\operatorname{Pic}'$ a partir dos grupos de Picard absolutos
$\operatorname{Pic}(X\times_S S'/S')$, ainda não se localizou o suficiente;
mais precisamente, $\operatorname{Pic}'$ não é, em geral, «compatível com a
descida fielmente plana». Explicitemos.

Seja $\mathfrak M$ o conjunto dos morfismos de pré-esquemas que são
\emph{fielmente planos e quase-compactos}; esse conjunto é estável por mudança
de base e por composição. Seja $P$ um funtor contravariante de
$(\mathrm{Sch})/S$ na categoria dos conjuntos e, para todo $S$-morfismo
$u:T'\longrightarrow T$, $u\in\mathfrak M$, consideremos o diagrama
\begin{equation}
P(T)\longrightarrow P(T')\rightrightarrows P(T'\times_TT')
\tag{1.5}
\end{equation}
obtido pela aplicação de $P$ ao diagrama
\[
T\longleftarrow T'\mathrel{\substack{\xleftarrow{\,\mathrm{pr}_1\,}\\[-0.35ex]
\xleftarrow[\,\mathrm{pr}_2\,]{}}}T'\times_TT'.
\]
Se $P$ é representável, resulta da teoria da descida ([4], I, B, teor. 2) que
o diagrama (1.5) é exato para todo $u\in\mathfrak M$. Exprime-se esse fato
dizendo que $P$ é compatível com $\mathfrak M$, neste caso que $P$ é
«compatível com a descida fielmente plana», ou ainda que o «pré-feixe» $P$
sobre $(\mathrm{Sch})/S$ é um «feixe» para a noção de localização fornecida
pelo conjunto $\mathfrak M$. Quando $P$ é arbitrário, um procedimento padrão,
bem conhecido no caso da localização topológica usual, permite associar-lhe
um «feixe» $\mathcal P$ e um homomorfismo de funtores
$P\longrightarrow\mathcal P$, que seja universal num sentido evidente. O
cálculo de $\mathcal P$ pode ser explicitado da seguinte maneira: para definir
$\mathcal P(T)$, designa-se, para todo $T'$ sobre $T$ tal que o morfismo
$u:T'\longrightarrow T$ pertença a $\mathfrak M$, por
\src{145}
$\overline H^0(T'/T,P)$ o subconjunto de $P(T')$ formado pelos elementos $\xi$
cujas duas imagens $\xi_1,\xi_2$ em $P(T'\times_TT')$ são tais que existe um
morfismo $v:T''\longrightarrow T'\times_TT'$, $v\in\mathfrak M$, para o qual
$\xi_1$ e $\xi_2$ têm a mesma imagem em $P(T'')$. [N.B. -- O conjunto
$\overline H^0$ assim definido é, portanto, maior que o conjunto
$H^0(T'/T,P)$ introduzido em [4], I, A, 4(a).] Quando $T'$ varia sobre $T$
fixo (sempre com $u\in\mathfrak M$), os $\overline H^0(T'/T,P)$ formam um
sistema indutivo (quando se mune o conjunto dos $T'$ da pré-ordem definida pela
dominação), e põe-se
\begin{equation}
\mathcal P(T)=\varinjlim_{T'}\overline H^0(T'/T,P).
\tag{1.6}
\end{equation}
A lei funtorial em $T$ dessa expressão explicita-se de modo evidente.

Quando
\[
P(T)=\operatorname{Pic}(X\times_ST),
\]
o funtor contravariante sobre $(\mathrm{Sch})/S$ definido por (1.6) denomina-se
\emph{funtor de Picard relativo} de $X$ sobre $S$ e denota-se por
$\operatorname{Pic}_{X/S}$; denomina-se \emph{grupo de Picard relativo} de
$X$ sobre $S$, e denota-se por $\operatorname{Pic}(X/S)$, o grupo
$\operatorname{Pic}_{X/S}(S)$. Obtém-se então uma bijeção evidente
\begin{equation}
\operatorname{Pic}_{X/S}(T)\simeq\operatorname{Pic}(X\times_ST/T).
\tag{1.7}
\end{equation}

Um elemento de $\operatorname{Pic}(X/S)$ é, portanto, definido com o auxílio de
um elemento $\xi'$ de um grupo $\operatorname{Pic}(X\times_SS')$ (onde
$S'\longrightarrow S$ é fielmente plano e quase-compacto), tal que se possa
encontrar um morfismo fielmente plano e quase-compacto
$S''\longrightarrow S'\times_SS'$ para o qual as duas imagens inversas de
$\xi'$ em $\operatorname{Pic}(X\times_SS'')$ sejam iguais. Um elemento $\xi'$
de $\operatorname{Pic}(X\times_SS')$ e um elemento $\xi_1$ de
$\operatorname{Pic}(X\times_SS_1)$ (satisfazendo as condições que acabamos de
explicitar) definem o mesmo elemento de $\operatorname{Pic}(X/S)$ se e somente
se existe um morfismo fielmente plano e quase-compacto
$S'_1\longrightarrow S'\times_SS_1$ tal que as imagens dos dois elementos em
questão em $\operatorname{Pic}(X\times_SS'_1)$ sejam iguais. É frequentemente
conveniente continuar trabalhando com o funtor
$P'=\operatorname{Pic}'_{X/S}$ introduzido acima; constata-se imediatamente que
o morfismo canônico $P\longrightarrow P'$ define um isomorfismo
\begin{equation}
\mathcal P\xrightarrow{\sim}\mathcal P',
\tag{1.8}
\end{equation}
o que fornece uma descrição de $\operatorname{Pic}_{X/S}$ em termos de
$\operatorname{Pic}'_{X/S}=P'$, frequentemente mais cômoda. Em virtude de 2.3
abaixo, quando, na descrição de $\operatorname{Pic}(X/S)$ que acabamos de
explicitar, se substitui $P$ por $P'$, pode-se, com efeito,
\src{146}
tomar $S''=S'\times_SS'$, $S'_1=S'\times_SS_1$, ao menos sob as condições
explicitadas em loco citato.

Quando o funtor $\operatorname{Pic}_{X/S}$ é representável, diz-se que $X/S$
admite um pré-esquema de Picard, e o pré-esquema sobre $S$ que representa o
funtor denomina-se \emph{pré-esquema de Picard} de $X$ sobre $S$ e continua a
ser denotado por $\operatorname{Pic}_{X/S}$. Para isso, basta evidentemente que
$P'=\operatorname{Pic}'_{X/S}$ seja representável, pois então $P'$ já é um
``feixe'', e a fórmula (1.8) demonstra que o morfismo
$P'\longrightarrow\mathcal P'$ se identifica com o morfismo canônico
\begin{equation}
\operatorname{Pic}'_{X/S}\longrightarrow\operatorname{Pic}_{X/S},
\tag{1.9}
\end{equation}
que é então um isomorfismo. Portanto, nossa terminologia é compatível com a
introduzida acima em (1.4). Em geral, quando $\operatorname{Pic}_{X/S}$ existe,
ele é definido pelo isomorfismo funtorial
\begin{equation}
\operatorname{Hom}_S(S',\operatorname{Pic}_{X/S})\xrightarrow{\sim}
\operatorname{Pic}(X\times_SS'/S').
\tag{1.10}
\end{equation}

\fsection{2}{Relações entre os diversos grupos de Picard relativos e absolutos}

\result{PROPOSIÇÃO}{2.1} Seja $f:X\longrightarrow S$ um morfismo tal que
$\mathcal O_S\xrightarrow{\sim}f_*(\mathcal O_X)$. Tem-se então uma sequência
exata
\[
0\longrightarrow\operatorname{Pic}(S)\longrightarrow\operatorname{Pic}(X)
\longrightarrow\operatorname{Pic}'(X/S).
\]
Quando $X$ admite uma seção sobre $S$, o último morfismo é sobrejetivo, isto é,
tem-se um isomorfismo
\[
\operatorname{Pic}'(X/S)\xrightarrow{\sim}
\operatorname{Pic}(X)/\operatorname{Pic}(S).
\]

A sequência exata pode ser considerada como a sequência exata em graus baixos
correspondente à sequência espectral de Leray para $f$ e $\mathcal O_X^*$. A
segunda afirmação também é formal.

\result{PROPOSIÇÃO}{2.2} Seja $f:X\longrightarrow S$ um morfismo quase-compacto
e separado tal que $\mathcal O_S\xrightarrow{\sim}f_*(\mathcal O_X)$, e seja
$S'\longrightarrow S$ um morfismo fielmente plano e quase-compacto. Então:

\noindent (i) $\operatorname{Pic}'(X/S)\longrightarrow
\operatorname{Pic}'(X\times_SS'/S')$ é injetivo;

\noindent (ii) se $X$ admite localmente uma seção sobre $S$ (isto é, todo
$s\in S$ possui uma vizinhança aberta $U$ tal que $X|U$ tenha uma seção sobre
$U$), então o diagrama
\[
\operatorname{Pic}'(X/S)\longrightarrow
\operatorname{Pic}'(X\times_SS'/S')\rightrightarrows
\operatorname{Pic}'(X\times_SS''/S'')
\]
\src{147}
(onde $S''=S'\times_SS'$) é exato.

A primeira afirmação resulta, graças às propriedades elementares da descida
fielmente plana, da observação geral seguinte. Se $f:X\longrightarrow S$ é um
morfismo tal que $\mathcal O_S\xrightarrow{\sim}f_*(\mathcal O_X)$, então o
funtor $\mathcal F\longmapsto f^*(\mathcal F)$, da categoria dos Módulos
localmente livres de tipo finito sobre $S$ para a categoria dos Módulos
localmente livres de tipo finito sobre $X$, é plenamente fiel, e sua imagem
essencial é formada pelos Módulos $\mathcal G$ sobre $X$ tais que
$f_*(\mathcal G)$ seja localmente livre e que o homomorfismo canônico
\[
f^*\bigl(f_*(\mathcal G)\bigr)\longrightarrow\mathcal G
\]
seja um isomorfismo. A segunda afirmação foi demonstrada pela teoria da descida
em [4], I, B, 4.

Os resultados de 2.2 também se enunciam assim:

\result{COROLÁRIO}{2.3} Sob as condições de 2.2, o homomorfismo canônico
(1.9) $\operatorname{Pic}'_{X/S}\longrightarrow\operatorname{Pic}_{X/S}$ é
injetivo, e até bijetivo se $X$ admite uma seção localmente sobre $S$. (Logo,
neste último caso, há identidade entre o grupo de Picard relativo
$\operatorname{Pic}(X/S)$ e o grupo de Picard relativo restrito
$\operatorname{Pic}'(X/S)$.)

Combinando com 2.1, obtém-se, portanto:

\result{COROLÁRIO}{2.4} Sob as condições de 2.2, tem-se uma sequência exata
\[
0\longrightarrow\operatorname{Pic}(S)\longrightarrow\operatorname{Pic}(X)
\longrightarrow\operatorname{Pic}(X/S).
\]
Quando $X$ admite uma seção sobre $S$, o último homomorfismo é sobrejetivo,
isto é, tem-se então um isomorfismo
\[
\operatorname{Pic}(X/S)\xrightarrow{\sim}
\operatorname{Pic}(X)/\operatorname{Pic}(S).
\]

\result{OBSERVAÇÃO}{2.5} Seja $f:X\longrightarrow S$ um morfismo tal que
$\mathcal O_S\xrightarrow{\sim}f_*(\mathcal O_X)$, e seja $g$ uma seção de
$X$ sobre $S$. Seja $\mathcal L$ um Módulo invertível sobre $X$; chamemos
$g$-rigidificação de $\mathcal L$ um isomorfismo
$\mathcal O_S\xrightarrow{\sim}g^*(\mathcal L)$, e chamemos \emph{Módulo
invertível $g$-rigidificado} um Módulo invertível $\mathcal L$ sobre $X$ munido
de uma $g$-rigidificação. Todo automorfismo de tal estrutura é trivial, e
$\operatorname{Pic}'(X/S)$ identifica-se com o conjunto das classes, a menos de
isomorfismo, de Módulos invertíveis $g$-rigidificados sobre $S$. (É esse fato que permitiu
utilizar a teoria da descida para demonstrar 2.2(ii).) Isso fornece uma nova
interpretação de $\operatorname{Pic}(X/S)$, pelo menos quando $f$ é, além disso,
quase-compacto e separado, portanto
$\operatorname{Pic}(X/S)\simeq\operatorname{Pic}'(X/S)$ por 2.3.
\src{148}
\result{OBSERVAÇÃO}{2.6} Seja $f:X\longrightarrow S$ um morfismo como em
2.2, e seja $S'\longrightarrow S$ um morfismo fielmente plano e
quase-compacto tal que exista um $S$-morfismo $S'\longrightarrow X$, isto é,
tal que exista uma seção de $X'=X\times_SS'$ sobre $S'$. Sejam
$S''=S'\times_SS'$ e $X''=X\times_SS''$, e consideremos a sequência exata
\[
\operatorname{Pic}(X/S)\longrightarrow\operatorname{Pic}(X'/S')
\rightrightarrows\operatorname{Pic}(X''/S'').
\]
Aplicando 2.4, obtém-se uma sequência exata
\[
\operatorname{Pic}(X/S)\longrightarrow
\operatorname{Pic}(X')/\operatorname{Pic}(S')
\rightrightarrows
\operatorname{Pic}(X'')/\operatorname{Pic}(S''),
\]
em particular, todo elemento do grupo de Picard relativo já «provém» de um
elemento de $\operatorname{Pic}(X')$. Isso fornece, portanto, uma simplificação
substancial da descrição do grupo de Picard relativo dada no número anterior,
e o mesmo vale para o funtor de Picard de $X$ sobre $S$, pois, para todo $T$
sobre $S$, pode-se aplicar o que precede a $X\times_ST/T$ e ao morfismo
$T'=S'\times_ST\longrightarrow T$. Se, por exemplo, o próprio $f$ é fielmente
plano, pode-se tomar $S'=X$, o que permite, quando $f$ é, além disso, de tipo
finito (resp. suave, etc.), limitar-se, na descrição do funtor de Picard relativo
$T\longmapsto\operatorname{Pic}(X\times_ST/T)$, às mudanças de base
$T'\longrightarrow T$ que são de tipo finito (resp. suaves, etc.). Quando $f$
é projetivo e plano e $S$ é localmente noetheriano, demonstra-se que se pode
tomar acima um $S'\longrightarrow S$ tal que $S'$ seja soma disjunta de
recobrimentos planos $S'_i$ de abertos $S_i$ de $S$ que cobrem $S$; se $f$ é
até mesmo separável, pode-se tomar $S'_i$ étale sobre $S_i$.

\fsection{3}{O teorema principal de existência: enunciado}

Não se dispõe, nem mesmo a título conjectural, de um enunciado de existência
de pré-esquemas de Picard que englobe todos os casos conhecidos. Uma condição
«praticamente necessária», se assim se pode dizer, é que
$f:X\longrightarrow S$ seja próprio (assegurando propriedades essenciais de
finitude) e plano. Essas condições não são suficientes, mesmo se $S$ é o
espectro da álgebra dos números duais $k[t]/(t^2)$ sobre um corpo $k$ (digamos,
o corpo $\mathbf C$ dos números complexos), e $X$ tem dimensão 1. No momento de
escrever esta exposição, os teoremas de existência mais importantes para o
pré-esquema de Picard deduzem-se do teorema seguinte:

\result{TEOREMA}{3.1} Seja $f:X\longrightarrow S$ um morfismo de pré-esquemas
localmente noetherianos. Supõe-se:

\noindent (i) $f$ é projetivo;
\src{149}
\noindent (ii) $f$ é plano;

\noindent (iii) as fibras geométricas de $f$ são íntegras.

Sob essas condições, $\operatorname{Pic}_{X/S}$ existe.

A demonstração, que será esboçada nas duas seções seguintes, mostrará ao mesmo
tempo o seguinte: seja $\xi$ a seção de $\operatorname{Pic}_{X/S}$
correspondente a um feixe muito amplo $\mathcal O_X(1)$ sobre $X/S$ (isto é,
induzido por uma imersão projetiva
$X\longrightarrow\mathbf P(\mathcal E)$); existe então uma parte aberta $U$ de
$\operatorname{Pic}_{X/S}$, união disjunta de partes abertas quase-projetivas
sobre $S$, tal que $U$ seja estável pela translação por $\xi$, e que
$\operatorname{Pic}_{X/S}$ seja a união crescente dos abertos $U-n\xi$ (todos
isomorfos a $U$). Disso resulta, em particular, que, sob as condições de 3.1,
$\operatorname{Pic}_{X/S}$ é separado sobre $S$.

\result{OBSERVAÇÃO}{3.2} Vê-se por exemplos (com $S$ o espectro de um anel de
valorização discreta e $X$ de dimensão relativa 1 sobre $S$, por exemplo) que,
se em 3.1 se omite a hipótese (iii), substituindo-a pela hipótese mais fraca de
que, para todo $s\in S$, o homomorfismo
\[
k(s)\longrightarrow H^0(X_s,\mathcal O_{X_s})
\]
é um isomorfismo, então $\operatorname{Pic}_{X/S}$ não é necessariamente
separado sobre $S$; tanto no caso em que as fibras geométricas de $f$ são
reduzidas, mas uma fibra geométrica genérica íntegra ``explode'' por
especialização em duas componentes irredutíveis, como no caso em que as fibras
geométricas de $f$ são irredutíveis, mas uma fibra geométrica genérica íntegra
se especializa em uma ``fibra múltipla''. O primeiro caso ocorre, por exemplo,
com uma cônica que degenera em duas retas concorrentes; um exemplo do segundo me
foi fornecido por D. Mumford, com uma curva elíptica que degenera em uma curva
elíptica dupla. Esses exemplos são válidos em toda característica.

\result{OBSERVAÇÃO}{3.3} Sob as condições de 3.1, ignoro se
$\operatorname{Pic}_{X/S}$ é uma união disjunta de abertos que são de tipo
finito, portanto quase-projetivos, sobre $S$.
 Observe-se
que a consideração dos polinômios de Hilbert $Q\in\mathbf Q[t]$ permite, como no
caso dos esquemas de Hilbert ([4], IV), decompor
$\operatorname{Pic}_{X/S}$ em uma soma disjunta de abertos
$\operatorname{Pic}^{Q}_{X/S}$, e parece plausível que estes sejam
quase-projetivos sobre $S$; isso é o que se verá, ao menos, na próxima exposição,
quando $f$ for um morfismo suave. Atente-se para o fato de que, se se substituir
a hipótese (i) pela hipótese: $X$ é localmente projetivo sobre $S$ (o que basta
para a validade de 3.1, pois a questão da existência de
$\operatorname{Pic}_{X/S}$ é manifestamente local sobre $S$), é fácil, por
outro lado,
\src{150}
dar exemplos em que $\operatorname{Pic}_{X/S}$ contém componentes conexas que
não são de tipo finito sobre $S$. Seja, por exemplo, $X_0$ uma variedade
algébrica projetiva não singular sobre um corpo $k$ algebricamente fechado,
munida de um automorfismo $u$ e de um elemento $\xi$ do grupo de
Néron--Severi de $X_0$, tais que os $u^n(\xi)$ sejam dois a dois distintos.
Pode-se, por exemplo, tomar para $X_0$ o produto de uma curva elíptica $E$ por
si mesma e, para $u$, o automorfismo
\[
(x,y)\longmapsto(x,y+x)
\]
de $E\times E$. Seja $S$ a união de duas curvas irredutíveis não singulares que
se cortam em dois pontos $a$ e $b$. Existe sobre $S$ um recobrimento principal
conexo $P$ de grupo $\mathbf Z$; utilizando as ações de $\mathbf Z$ sobre $X_0$
definidas por $u$, obtém-se um fibrado associado sobre $S$, de fibra $X_0$
(trivial sobre $S-a$ e $S-b$), na realidade um esquema abeliano sobre $S$ no
caso particular considerado. Vê-se facilmente que
$\operatorname{Pic}_{X/S}$, que é também o fibrado associado a $P$ e às ações
de $\mathbf Z$ sobre $\operatorname{Pic}_{X_0/k}$ por meio de $u$, contém uma
componente conexa isomorfa a
\[
P\times\operatorname{Pic}^{0}_{X_0/k},
\]
(onde $\operatorname{Pic}^{0}$ denota a componente conexa do elemento neutro
em $\operatorname{Pic}$), a qual não é de tipo finito sobre $S$. [Fenômenos
análogos também ocorrem em diversos casos de pré-esquemas de Picard não
separados sobre $S$, como os considerados em 3.3.]

\fsection{4}{Divisores de Cartier relativos e fibrados projetivos}

Teremos de utilizar apenas divisores positivos e omitiremos esse qualificativo
adicional no restante deste item.

Seja $X$ um pré-esquema. Um \emph{divisor de Cartier} sobre $X$, ou simplesmente
divisor sobre $X$, para simplificar, é um subpré-esquema fechado $D$ de $X$,
definido por um Ideal $\mathcal J$ que é um Módulo invertível, isto é, gerado
localmente por uma seção que não é divisor de zero de $\mathcal O_X$. A $D$
associamos um Módulo invertível, a saber,
\[
\mathcal L(D)=\mathcal J^{-1},
\]
e a injeção canônica $\mathcal J\longrightarrow\mathcal O_X$ fornece um
homomorfismo canônico
\[
s_D:\mathcal O_X\longrightarrow\mathcal J^{-1}=\mathcal L(D),
\qquad\text{isto é,}\quad s_D\in\Gamma(X,\mathcal L(D)).
\]
Além disso, o dado de um divisor é essencialmente equivalente ao dado de um
Módulo invertível $\mathcal L$ sobre $X$, munido de uma seção $s$ que não seja
divisor de zero em ponto algum, associando reciprocamente a tal par
$(\mathcal L,s)$ o «divisor» de $s$, denotado por $\operatorname{div}(s)$. Para
um $\mathcal L$ invertível dado sobre $X$, o conjunto dos
\src{151}
divisores $D$ que definem $\mathcal L$ está em correspondência biunívoca com o
conjunto quociente
\[
\Gamma(X,\mathcal L)^*/\Gamma(X,\mathcal O_X^*),
\]
onde $\Gamma(X,\mathcal L)^*$ designa a parte de $\Gamma(X,\mathcal L)$ formada
pelas seções que não são divisores de zero em nenhum ponto.

Suponhamos agora que se tenha um morfismo localmente de tipo finito
$f:X\longrightarrow S$, e suponhamos, para simplificar, $S$ localmente
Noetheriano. Sejam $\mathcal J$ um Ideal coerente sobre $X$, $D$ o subesquema de
$X$ que ele define, $x\in X$, e $s=f(x)$. Demonstra-se que as seguintes
condições são equivalentes:

\noindent (i) $\mathcal J$ é invertível em $x$ (isto é, $\mathcal J_x$ é gerado
por um elemento regular de $\mathcal O_{X,x}$) e $D$ é plano sobre $S$ em $x$;

\noindent (ii) $X$ e $D$ são planos sobre $S$ em $x$, e $D_s$ é um divisor de
Cartier sobre a fibra $X_s$ no ponto $x$;

\noindent (iii) $X$ é plano sobre $S$ em $x$, e $\mathcal J_x$ é gerado por um
elemento $f_x$ que induz sobre $X_s$ um germe não divisor de zero.

Diz-se então que $D$ é um \emph{divisor de Cartier relativo} sobre $X/S$ (ou
simplesmente divisor relativo sobre $X/S$) no ponto considerado. Observe-se, por
(i), que $D$ é também um divisor relativo nos pontos vizinhos de $x$; portanto,
se $X$ e $D$ são planos sobre $S$, e $D$ é próprio sobre $S$, então o conjunto
dos $s\in S$ tais que $D_s$ seja um divisor de Cartier em $X_s$ (isto é, tais
que $D$ seja um divisor de Cartier relativo nos pontos de $X_s$) é uma parte
aberta de $S$. Por outro lado, na definição precedente fez-se o necessário para
que a noção de divisor de Cartier relativo seja estável por uma mudança de base
arbitrária $S'\longrightarrow S$. Consideremos então o conjunto
$\operatorname{Div}(X/S)$ dos divisores relativos sobre $X/S$, e em seguida o
funtor contravariante em $S'$ variável sobre $S$ definido por
\[
\operatorname{Div}_{X/S}(S')=
\operatorname{Div}(X\times_SS'/S').
\]

Suponhamos $X$ plano e próprio sobre $S$. Então, pela caracterização (ii) dos
divisores de Cartier relativos, $\operatorname{Div}_{X/S}$ pode ser considerado
como um subfuntor do funtor $\operatorname{Hilb}_{X/S}$ definido em [4], IV, e o
morfismo de inclusão
\[
\operatorname{Div}_{X/S}\longrightarrow\operatorname{Hilb}_{X/S}
\]
é ``representável por imersões abertas'' (cf. [5], IV, 3.13), em virtude das
observações precedentes. Utilizando o teorema principal de existência de [4],
IV, obtém-se:
\src{152}
\result{PROPOSIÇÃO}{4.1} Suponhamos que $f:X\longrightarrow S$ seja projetivo
e plano. Então o funtor $\operatorname{Div}_{X/S}$ é representável e, mais
precisamente, é representado por um aberto de $\operatorname{Hilb}_{X/S}$.

Utilizando (para um feixe $\mathcal O_X(1)$ muito amplo dado sobre $X/S$) a
decomposição canônica de $\operatorname{Hilb}_{X/S}$ em soma dos abertos
$\operatorname{Hilb}^{Q}_{X/S}$ correspondentes aos polinômios de Hilbert
$Q\in\mathbf Q[t]$, deduz-se, aliás, uma decomposição análoga
\[
\operatorname{Div}_{X/S}=
\coprod_{Q\in\mathbf Q[t]}\operatorname{Div}^{Q}_{X/S}
\]
em soma de abertos disjuntos quase-projetivos sobre $S$.

Utilizando a aplicação $D\longmapsto\mathcal L(D)$, obtém-se, por outro lado,
um homomorfismo funtorial
\begin{equation}
\operatorname{Div}_{X/S}\longrightarrow\operatorname{Pic}_{X/S}
\tag{+}
\end{equation}
que nos propomos estudar; ver-se-á que ele é relativamente representável ([5],
IV, 3) sob condições bastante gerais. Partamos, portanto, de um elemento $\xi$
de $\operatorname{Pic}_{X/S}(S')$, supondo, para simplificar a notação, que
$S'=S$, e mostremos que o subfuntor correspondente de
$\operatorname{Div}_{X/S}$ é representável. Tomemos primeiro o caso em que
$\xi$ é definido por um Módulo invertível $\mathcal L$ sobre $X$. Suporemos
que $X$ seja próprio e plano sobre $S$ e que as fibras geométricas de $X$
sobre $S$ sejam íntegras, o que implica também ([2], III, par. 7) que se tem
$\mathcal O_S\xrightarrow{\sim}f_*(\mathcal O_X)$ e que essa relação permanece
válida após toda mudança de base $S'\longrightarrow S$. Então os divisores de
Cartier relativos $D$ sobre $X/S$ tais que $\mathcal L(D)$ e $\mathcal L$
definem o mesmo elemento de
$\operatorname{Pic}(X/S)=\operatorname{Pic}_{X/S}(S)$, isto é, em virtude de
2.4, tais que $\mathcal L(D)$ e $\mathcal L$ sejam isomorfos localmente sobre
$S$, estão em correspondência biunívoca com as seções do feixe quociente
$f_*(\mathcal L)^*/\mathcal O_S^*$. Essa correspondência é compatível com as
mudanças de base. Por outro lado, considerações gerais do tipo «Künneth» de
loco citato (cf. também [6]) mostram que o fato de $X/S$ ser próprio e de
$\mathcal L$ ser plano sobre $S$ implica a existência de um Módulo coerente
$\mathcal Q$ sobre $S$, definido a menos de isomorfismo único, e de um
isomorfismo de feixes
\[
f_*(\mathcal L)\xrightarrow{\sim}
\operatorname{Hom}_{\mathcal O_X}(\mathcal Q,\mathcal O_S),
\]
 sendo, além disso, a formação de $\mathcal Q$
compatível com as mudanças de base. Aqui $f_*(\mathcal L)^*$ designa o
subfeixe de conjuntos de $f_*(\mathcal L)$ cujas seções sobre $U$ são as
seções de $\mathcal L$ sobre $f^{-1}(U)$ que definem divisores de Cartier
relativos sobre
\src{153}
Cartier relativos sobre $f^{-1}(U)/U$, isto é, que induzem seções que não são
divisores de zero sobre os $X_s$ ($s\in U$). Utilizando a hipótese de que as
fibras $X_s$ são íntegras, isso significa simplesmente que as seções induzidas
sobre as fibras $X_s$ não são identicamente nulas, ou ainda, em termos de
homomorfismos locais $\mathcal Q\longrightarrow\mathcal O_S$, que tais
homomorfismos sejam sobrejetivos (Nakayama). Isso mostra que o conjunto das
seções de $f_*(\mathcal L)^*/\mathcal O_S^*$ está em correspondência biunívoca
com o conjunto dos Módulos quocientes invertíveis de $\mathcal Q$, ou ainda,
pela definição do fibrado projetivo $\mathbf P(\mathcal Q)$ associado ao Módulo
coerente $\mathcal Q$ (cf. [5], V, 2), com o conjunto das seções de
$\mathbf P(\mathcal Q)$ sobre $S$. Essa descrição é compatível com a formação
das imagens inversas, obtendo-se, por conseguinte, o teorema seguinte.

\result{TEOREMA}{4.3} Sejam $f:X\longrightarrow S$ um morfismo próprio e plano
com fibras geométricas íntegras, sendo $S$ localmente noetheriano, e
$\mathcal L$ um Módulo invertível sobre $X$. Para todo $S'$ sobre $S$, seja
$T(S')$ o conjunto dos divisores relativos $D$ sobre $X\times_SS'/S'$ tais que
$\mathcal L(D)$ seja isomorfo a
$\mathcal L\otimes_{\mathcal O_S}\mathcal O_{S'}$ localmente sobre $S'$ (isto
é, tais que $\mathcal L(D)$ e
$\mathcal L\otimes_{\mathcal O_S}\mathcal O_{S'}$ definam o mesmo elemento de
$\operatorname{Pic}(X\times_SS'/S')$). Então existe um Módulo coerente
$\mathcal Q$ sobre $S$, determinado a menos de isomorfismo único, tal que o
funtor $T$ seja representável pelo fibrado projetivo $\mathbf P(\mathcal Q)$.

\result{COROLÁRIO}{4.4} Se supusermos, além disso, que $f$ é projetivo, então o
homomorfismo funtorial
$\operatorname{Div}_{X/S}\longrightarrow\operatorname{Pic}_{X/S}$ é
representável por morfismos projetivos.

Com efeito, quando $X$ admite uma seção (resp. localmente uma seção) sobre $S$,
o homomorfismo precedente é representável por fibrados projetivos (resp. por
fibrados projetivos locais), graças a 4.3 e 2.1. No caso em que $f$ é
quase-projetivo, reduz-se facilmente ao caso precedente por um método de
descida, utilizando quase-seções finitas e planas de $X$ localmente sobre $S$.

\result{OBSERVAÇÃO}{4.5} Sob as condições de 4.3, o Módulo $\mathcal Q$ sobre
$S$ não é, em geral, localmente livre, o que se exprime pelo fato de que a
dimensão das fibras reduzidas de $\mathcal Q$, isto é, a dos
$H^0(X_s,\mathcal L_s)$ para $s\in S$ variável, pode dar saltos. Dado um Módulo
coerente $\mathcal Q$ sobre o pré-esquema localmente noetheriano $S$, verifica-se
facilmente que, para um $s\in S$ dado, $\mathcal Q$ é livre em $s$ se e somente
se $\mathbf P(\mathcal Q)$ é plano sobre $S$ nos pontos acima de $s$ (caso em
que será até liso sobre $S$ nos pontos acima de $s$); no caso em questão,
estando $\mathcal Q$ definido em termos de $\mathcal L$ como acima, isso
significa também que a formação da imagem direta $f_*(\mathcal L)$ «comuta com
a mudança de base na
\src{154}
vizinhança de $s$», ou ainda que
$f_*(\mathcal L)_s\longrightarrow H^0(X_s,\mathcal L_s)$ é sobrejetivo. Isso
ocorre, por exemplo, se $H^1(X_s,\mathcal L_s)=0$. Sob reserva da existência
dos pré-esquemas em questão, esses critérios aplicam-se, em particular, à
situação universal
$\operatorname{Div}_{X/S}\longrightarrow\operatorname{Pic}_{X/S}$, e fornecem
uma condição necessária e suficiente, resp. suficiente, para que esse morfismo
seja liso em um ponto dado de $\operatorname{Div}_{X/S}$.

\fsection{5}{Demonstração do teorema principal de existência}

Sob as condições de 3.1, escolhamos um Módulo $\mathcal O_X(1)$ muito amplo
sobre $X/S$, e seja $\xi$ o elemento correspondente de
$\operatorname{Pic}(X/S)$. Para abreviar, ponhamos
\[
P(S')=\operatorname{Pic}(X\times_SS'/S'),
\]
e suponhamos primeiro, para simplificar, que $X/S$ admita uma seção. Seja
$P^+(S')$ a parte de $P(S')$ formada pelas classes dos $\mathcal L$ invertíveis
sobre $X\times_SS'$ tais que
\[
R^if'_*(\mathcal L(n))=0\quad\text{para $i>0$ e todo $n\geq0$},
\qquad
f'_*(\mathcal L(n))\neq0\quad\text{para todo $n\geq0$}.
\]
Essas condições são estáveis por mudança de base, definindo, portanto, um
subfuntor $P^+$ de $P$, evidentemente estável pela translação por $\xi$.
Utilizando os «teoremas A e B» de Serre ([2], III, 2) e as generalidades
([5], IV, 5), vê-se facilmente que $P$ é representável se e somente se $P^+$ o
é, e, nesse caso, $P^+$ será representado por um aberto $U$ do pré-esquema
$\operatorname{Pic}_{X/S}$ que representa $P$, e este último será a união
crescente dos abertos $U-n\xi$.

Para abreviar, seja $D=\operatorname{Div}_{X/S}$, e seja $D^+$ a imagem inversa
de $P^+$ pelo morfismo canônico $D\longrightarrow P$. Tem-se, portanto, um
morfismo
\begin{equation}
D^+\longrightarrow P^+,
\tag{+}
\end{equation}
e já se sabe que $D^+$ é representável por um aberto $D^+$ do pré-esquema
$D=\operatorname{Div}_{X/S}$ (cuja existência é garantida por 4.1). Resulta
facilmente de 4.3 que o morfismo (+) é representável por morfismos lisos,
projetivos e sobrejetivos (e, mais precisamente, por fibrados projetivos
associados a Módulos localmente livres não nulos em toda parte): isso decorre
do fato de que, quando $\mathcal L$ sobre $X\times_SS'$ é como no início deste
item, então $f'_*(\mathcal L)$ é um Módulo localmente livre $\neq0$, cuja
formação comuta com a mudança de base; com as notações de 4.3, $\mathcal Q$
será então isomorfo ao dual de $f'_*(\mathcal L)$. Utilizando o critério geral
([5], IV, 4.7), concluiremos daí a representabilidade de $P^+$. No lugar citado,
tomaremos para $\mathfrak C$ o conjunto dos morfismos fielmente
\src{155}
planos e quase-compactos de pré-esquemas (que são de fato epimorfismos
efetivos em virtude de [4], I B). A condição (a) de loco citato, a saber, que
(+) é representável por morfismos pertencentes a $\mathfrak C$, está verificada,
como acabamos de ver; o mesmo ocorre com a condição (b), que significa que o
funtor $P^+$ é compatível com a descida fielmente plana quase-compacta, o que
é imediato. Resta provar a condição (c) de loco citato, a saber, que a
equivalência $R$ no pré-esquema $D^+$ deduzida do morfismo
$\mathfrak C$-representável (+) é $\mathfrak C$-efetiva, isto é, é efetiva e
tal que $D^+\longrightarrow D^+/R$ pertença a $\mathfrak C$. Para isso,
observamos primeiro que os abertos $D^{+Q}$ de $D^+$ correspondentes aos
diversos polinômios de Hilbert virtuais $Q\in\mathbf Q[t]$ são estáveis por $R$
(pois as fibras de $R$ são conexas), o que nos reduz a mostrar que, para todo
$Q$, a relação de equivalência $R^Q$ induzida em $D^{+Q}$ é
$\mathfrak C$-efetiva. Ora, $D^{+Q}$ é agora quase-projetivo e, por outro lado,
a relação de equivalência $R^Q$ é projetiva e plana. Estamos, portanto, nas
condições de aplicação de [4], III, 6.1, o que implica o resultado desejado.

No caso geral em que $X/S$ não admite necessariamente uma seção, reduzimo-nos
facilmente ao caso anterior pela técnica de descida, ou retomamos a demonstração
anterior com a modificação necessária na definição de $P^+$.

\result{OBSERVAÇÃO}{5.1} O método seguido é essencialmente o de Matsusaka para
a construção projetiva das variedades de Picard. O resultado de [4], III
invocado para a possibilidade da passagem ao quociente efetivo também pode ser
facilmente deduzido do teorema de existência dos esquemas de Hilbert (cf., por
exemplo, [6]) (classicamente, esses quocientes eram construídos mediante o uso
das coordenadas de Chow). Observe-se que a formação do aberto
$\operatorname{Pic}^{+}_{X/S}$ de $\operatorname{Pic}_{X/S}$ e sua decomposição
em abertos $\operatorname{Pic}^{+Q}_{X/S}$ quase-projetivos sobre $S$, segundo
os polinômios de Hilbert dos divisores que definem os Módulos invertíveis
considerados, é compatível com a mudança de base (o que permite a aplicação da
técnica de descida).

\result{OBSERVAÇÃO}{5.2} Não se exclui que $\operatorname{Pic}_{X/S}$ exista
sempre que $f:X\longrightarrow S$ seja próprio e plano e tal que os
homomorfismos
\[
k(s)\longrightarrow H^0(X_s,\mathcal O_{X_s})\qquad(s\in S)
\]
sejam isomorfismos (significando então esta última condição também que
$\mathcal O_S\xrightarrow{\sim}f_*(\mathcal O_X)$ e que essa relação permaneça
válida após toda mudança de base $S'\longrightarrow S$). É ao menos o que se
prova no âmbito dos espaços analíticos quando $f$ é, além disso, projetivo, por
um método diferente (o de Chow, salvo engano) explicado em [5], IX,
3.1. Nesse método, a passagem ao
\src{156}
quociente por uma relação de equivalência própria e plana em um aberto do
esquema dos divisores é substituída por uma passagem ao quociente pelo grupo
projetivo no esquema das imersões de $X$ em $\mathbf P^r_S$. Esse método poderá
provavelmente adaptar-se ao caso dos esquemas utilizando os resultados de
Mumford sobre a passagem ao quociente pelo grupo projetivo [8]; por ora, só
existe demonstração escrita quando $X$ tem ``muitas seções locais'' sobre $S$,
por exemplo, quando $X$ é separável sobre um anel local completo.
 Em princípio, o método em questão teria alcance mais geral,
pois também fornece a existência de pré-esquemas de Picard nos casos em que
estes não são separados e nos quais, portanto, o primeiro método necessariamente
falha. (Tecnicamente, a dificuldade provém do fato de que, quando as fibras
geométricas de $f$ deixam de ser íntegras, o funtor considerado em 4.3 deixa de
ser representável pelo próprio fibrado projetivo $\mathbf P(\mathcal Q)$ e
passa a sê-lo por um aberto deste, o que conduz à delicada questão da passagem
ao quociente por uma relação de equivalência plana, mas não própria.)

\result{OBSERVAÇÃO}{5.3} Observe-se que a demonstração apresentada aqui não
recorre nem à construção prévia das jacobianas das curvas ou famílias de curvas,
nem à teoria das variedades abelianas ou dos esquemas abelianos, e por isso se
distingue essencialmente de exposições tradicionais como a do livro de Lang [7]
ou a exposição de Chevalley [1], que seguem a via esboçada por A. Weil. Mesmo no
caso das jacobianas das curvas não singulares sobre um corpo algebricamente
fechado (digamos, o corpo dos números complexos), a construção da jacobiana aqui
apresentada é a única conhecida que fornece as propriedades muito fortes que
adotamos como definição no número 1 (essencialmente as de Chevalley, mas levando
em conta ``variedades de parâmetros'' com elementos nilpotentes). Que a
construção dos esquemas de Picard deva preceder, e não seguir, a teoria das
variedades abelianas é claro a priori pelo fato de que, em geral, os esquemas de
Picard não são esquemas abelianos nem se reduzem a eles, como já se vê no caso
das curvas singulares sobre um corpo algebricamente fechado, em que se encontram
as ``jacobianas generalizadas'' de Rosenlicht, que não são variedades abelianas.
Além disso, a teoria das variedades abelianas, e mais geralmente a dos esquemas
abelianos, ganha muito em simplicidade quando se dispõe de uma teoria dos
esquemas de Picard em geral. Em particular, a teoria de dualidade para os
esquemas abelianos, sobretudo os resultados do tipo Cartier, torna-se quase
formal a partir daí (cf., por exemplo, [10]).
\src{157}
\result{OBSERVAÇÕES}{5.4} O «princípio de compatibilidade» de Igusa para a
jacobiana de uma curva que degenera em uma curva singular só pode ser bem
compreendido como um teorema de existência do esquema de Picard de um esquema
relativo em curvas $X/S$, não necessariamente suave sobre $S$. Trata-se,
portanto, de um caso particular do teorema principal de existência 3.1 quando a
curva especializada é íntegra (isto é, em termos clássicos, irredutível de
multiplicidade 1). Observe-se que, por ora, o caso de uma curva especial
redutível (mesmo quando as componentes têm multiplicidade 1, isto é, quando a
curva especial é separável sobre o corpo residual) escapa aos teoremas de
existência conhecidos, salvo no caso em que se está sobre um anel de valorização
discreta completo com corpo residual algebricamente fechado; compare-se 5.2.
Essa questão de existência certamente surgirá por ocasião de uma construção
geométrica, na teoria dos esquemas, das «compactificações» de Baily--Satake dos
esquemas modulares das curvas de gênero $g$. (Essa compactificação é conhecida,
por ora, apenas para $g=1$, graças aos trabalhos de Igusa.)

\fsection{6}{Teoremas de existência relativos}

Esboçaremos aqui alguns casos úteis em que a existência de certos esquemas de
Picard implica a existência de alguns outros, o que permite deduzir do teorema
principal 3.1 diversos outros teoremas de existência.

\result{PROPOSIÇÃO}{6.1} Seja $f:X\longrightarrow S$ um morfismo plano e
projetivo tal que, na fatorização de Stein $f=f''f'$, o morfismo
$f':X\longrightarrow S'$ seja plano e tenha fibras geométricas íntegras
(portanto, satisfaça às hipóteses de 3.1), e o morfismo finito
$f'':S'\longrightarrow S$ seja plano. Então $\operatorname{Pic}_{X/S}$ existe
e (com as notações introduzidas em [4], II, C 2) tem-se um isomorfismo canônico
\[
\operatorname{Pic}_{X/S}\xrightarrow{\sim}
\Pi_{S'/S}\operatorname{Pic}_{X/S'}.
\]

Para demonstrá-lo, estabelece-se primeiro um isomorfismo de funtores
\[
\operatorname{Pic}_{X/S}\xrightarrow{\sim}
\Pi_{S'/S}\operatorname{Pic}_{X/S'},
\]
utiliza-se 3.1, que implica que $\operatorname{Pic}_{X/S'}$ é representável, e
utiliza-se a estrutura assinalada no n\textsuperscript{o} 3 de
$\operatorname{Pic}_{X/S'}$ (que implica que toda parte finita de uma fibra de
$\operatorname{Pic}_{X/S'}$ sobre $S$ está contida em um aberto afim) para a
existência de $\Pi_{S'/S}\operatorname{Pic}_{X/S'}$.
\src{158}
Por exemplo, quando $X$ é um esquema soma de esquemas $X_i$ sobre $S$ que
satisfazem as condições de 3.1, o enunciado 6.1 reduz-se ao enunciado trivial
\[
\operatorname{Pic}_{X/S}\xrightarrow{\sim}
\prod_i\operatorname{Pic}_{X_i/S}.
\]

\result{COROLÁRIO}{6.2} Seja $f:X\longrightarrow S$ um morfismo projetivo e
plano com fibras geométricas localmente íntegras (por exemplo, um morfismo
projetivo e normal); então $\operatorname{Pic}_{X/S}$ existe.

Nesse caso, aliás, $S'$ é um recobrimento étale de $S$ (o que é verdadeiro
sempre que $f$ é separável, isto é, plano com fibras geométricas reduzidas), e
verifica-se que o teorema de estrutura enunciado no n\textsuperscript{o} 3 para
$\operatorname{Pic}_{X/S}$ continua válido graças à estrutura análoga de
$\operatorname{Pic}_{X/S'}$.

Por outro lado, um procedimento de descida fornece um teorema de existência
relativo, cujo alcance depende, aliás, da solução das questões de descida não
plana levantadas em [4], I, A (c), e do qual nos limitamos a explicitar aqui um
caso particular:

\result{PROPOSIÇÃO}{6.3} Seja $f:X\longrightarrow S$ um morfismo próprio, sejam
$X_1$ e $X_2$ dois subpré-esquemas de $X$ planos sobre $S$, definidos por dois
Ideais coerentes $\mathcal I_1$ e $\mathcal I_2$ tais que
$\mathcal I_1\cap\mathcal I_2=(0)$ e que
$\mathcal O_X/(\mathcal I_1+\mathcal I_2)$ seja plano sobre $S$ (isto é, o
subpré-esquema de $X$ que é o supremo de $X_1$ e $X_2$ é $X$, enquanto seu ínfimo $X$
é plano sobre $S$).
 Suponhamos, além disso, que, para todo $s\in S$, os
homomorfismos
\[
k(s)\longrightarrow H^0(X_{i,s},\mathcal O_{X_{i,s}}),\qquad i=1,2,
\]
sejam bijetivos. Então o homomorfismo natural de funtores
\[
\operatorname{Pic}_{X/S}\longrightarrow
\operatorname{Pic}_{X_1/S}\times\operatorname{Pic}_{X_2/S}
\]
é representável por morfismos afins; portanto, se
$\operatorname{Pic}_{X_1/S}$ e $\operatorname{Pic}_{X_2/S}$ existem, o mesmo
ocorre com $\operatorname{Pic}_{X/S}$, e o morfismo canônico
\[
\operatorname{Pic}_{X/S}\longrightarrow
\operatorname{Pic}_{X_1/S}\times\operatorname{Pic}_{X_2/S}
\]
é afim.

Por descida fielmente plana, reduzimo-nos ao caso em que $Z$ admite uma seção
sobre $S$, definindo, portanto, seções de $X$, $X_1$, $X_2$ sobre $S$, o que
permite eliminar os automorfismos nas estruturas consideradas, como explicado
em 2.5. A demonstração consiste então em observar que o dado de um Módulo
invertível
\src{159}
«rigidificado» $\mathcal L$ sobre $X$ equivale ao dado de uma tripla
$(\mathcal L_1,\mathcal L_2,u)$, em que $\mathcal L_i$ é um Módulo invertível
«rigidificado» sobre $X_i$, e $u$ é um isomorfismo de
$\mathcal L_1|Z$ em $\mathcal L_2|Z$, compatível com as rigidificações.
Resta apenas verificar que, para $\mathcal L_1$ e $\mathcal L_2$ fixados, o
dado de $u$ se exprime por uma seção de um esquema conveniente sobre $S$, afim
sobre $S$, o que é fácil. Deduz-se facilmente de 6.3:

\result{COROLÁRIO}{6.4} Seja $X$ um esquema próprio e separável sobre um corpo
$k$, e sejam $X_i$ as componentes irredutíveis de $X$. Se existem os
$\operatorname{Pic}_{X_i/k}$, também existe $\operatorname{Pic}_{X/k}$, e o
morfismo canônico
\[
\operatorname{Pic}_{X/k}\longrightarrow\prod_i\operatorname{Pic}_{X_i/k}
\]
é afim.

Em combinação com 6.2, isso mostra, por exemplo, a existência de
$\operatorname{Pic}_{X/k}$ sempre que $X$ é um esquema projetivo e separável
sobre um corpo $k$. Quando $X$ deixa de ser separável sobre $k$, tem-se também
um resultado de redução, utilizando um raciocínio devido a Oort [11]. O método
aplica-se igualmente sobre um esquema de base arbitrário (caso útil, por
exemplo, para demonstrar, na exposição seguinte, o resultado de finitude
anunciado em 3.3). Para evitar um enunciado excessivamente técnico,
limitemo-nos ao caso de um corpo de base:

\result{PROPOSIÇÃO}{6.5} Seja $X$ um esquema próprio sobre um corpo $k$, e seja
$X_0$ um subesquema com o mesmo conjunto subjacente (portanto, definido por um
Ideal nilpotente sobre $X$). Então o morfismo funtorial
\[
\operatorname{Pic}_{X/k}\longrightarrow\operatorname{Pic}_{X_0/k}
\]
é representável por morfismos afins. Em particular, se
$\operatorname{Pic}_{X_0/k}$ existe, o mesmo ocorre com
$\operatorname{Pic}_{X/k}$, e o morfismo
$\operatorname{Pic}_{X/k}\longrightarrow\operatorname{Pic}_{X_0/k}$ é afim.

Combinando isso com 6.4, conclui-se facilmente:

\result{COROLÁRIO}{6.6} Seja $X$ um pré-esquema projetivo sobre um corpo $k$;
então $\operatorname{Pic}_{X/k}$ existe.

\result{OBSERVAÇÃO}{6.6} É extremamente plausível que, para todo esquema $X$
próprio sobre um corpo $k$, $\operatorname{Pic}_{X/k}$ exista. Os resultados precedentes reduzem essa questão ao caso em que $k$ é
algebricamente fechado e $X$ é íntegro. Sabe-se, então, que existe um esquema
íntegro $X'$ projetivo sobre $k$ e um morfismo dominante
$g:X'\longrightarrow X$ (lema de Chow). Bastaria, portanto, mostrar que o
\src{160}
morfismo funtorial correspondente
\[
\operatorname{Pic}_{X/k}\longrightarrow\operatorname{Pic}_{X'/k}
\]
é representável (e, sem dúvida, representável por morfismos afins), pois se
sabe, por outro lado, que $\operatorname{Pic}_{X'/k}$ é representável. Isso
levanta questões de descida não plana ainda não resolvidas. Observe-se que,
quando se restringe a consideração do funtor $\operatorname{Pic}_{X/k}$ aos
pré-esquemas reduzidos (sendo $X$ próprio e íntegro sobre $k$ algebricamente
fechado), obtém-se de fato um funtor representável, como mostrou Chevalley [1]
no caso em que $X$ é normal, e Seshadri [12], por um método de descida, no caso
geral. Mas, com nossas notações, o esquema construído por esses autores não é
$\operatorname{Pic}_{X/k}$, mas
$(\operatorname{Pic}_{X/k})_{\mathrm{réd}}$, o esquema reduzido correspondente
a $\operatorname{Pic}_{X/k}$.

\result{OBSERVAÇÃO}{6.7} Seja, mais geralmente, $f:X'\longrightarrow X$ um
morfismo sobrejetivo de pré-esquemas próprios sobre $k$. Então, considerações de
descida não plana levam a conjecturar que
\[
\operatorname{Pic}_{X/k}\longrightarrow\operatorname{Pic}_{X'/k}
\]
é um morfismo afim, o que implicaria, em particular (dividindo pelas
componentes conexas dos elementos neutros), que o homomorfismo correspondente
sobre os grupos de Néron--Severi é injetivo módulo torsão. É o que se pode
verificar pela teoria das interseções quando $X$ e $X'$ são não singulares. A
resposta não parece conhecida em nenhum outro caso.

\result{OBSERVAÇÃO}{6.8} Ao contrário do que se poderia pensar, a consideração
dos esquemas de Picard de esquemas algébricos com elementos nilpotentes é útil,
e até indispensável, em diversas questões. Assim, quando $X$ é um esquema
projetivo, digamos, regular, e $Y$ uma seção hiperplana, convém
considerar as «vizinhanças infinitesimais» de todas as ordens $X_n$ de $Y$ e os
esquemas de Picard $\operatorname{Pic}_{X_n/k}$; quando $X$ é irredutível de
dimensão $\geq4$ (resp. $\geq3$), o morfismo canônico
\[
\operatorname{Pic}_{X/k}\longrightarrow\operatorname{Pic}_{X_n/k}
\qquad(n\text{ grande})
\]
é um isomorfismo (resp. induz um isomorfismo para as imagens inversas dos
subgrupos de torsão de Néron--Severi), resultado que nos será útil no estudo
qualitativo dos esquemas de Picard na exposição seguinte. Do mesmo modo, a
consideração dos esquemas de Picard de certas curvas com elementos nilpotentes
e os teoremas fundamentais de geometria formal [3] permitem verificar, no caso
de características iguais, uma conjectura de Mumford, a saber, que, para todo
anel local noetheriano normal completo $A$ de dimensão 2, o grupo das classes
de divisores de $A$ pode ser considerado como o conjunto dos pontos racionais
sobre $k$ de um
\src{161}
grupo algébrico $G$ sobre o corpo residual $k$ ($G$ sendo, além disso,
determinado canonicamente uma vez escolhido um corpo de coeficientes em $A$).
No caso em que $A$ tem dimensão arbitrária, é plausível que exista um grupo
proalgébrico sobre $k$ que desempenhe o mesmo papel que $G$ acima e que se
construa, quando se pode ``dessingularizar'' $\operatorname{Spec}(A)$, como um
limite projetivo de esquemas de Picard de esquemas projetivos convenientes (com
elementos nilpotentes) sobre $k$.

\section*{BIBLIOGRAFIA}
\addcontentsline{toc}{section}{Bibliografia}

\begin{description}
\item[{[1]}] CHEVALLEY (Claude). -- Sur la théorie de Picard,
\emph{Amer. J. of Math.}, v. 82, 1960, p. 435--490.

\item[{[2]}] GROTHENDIECK (A.) e DIEUDONNÉ (J.). -- Éléments de géométrie
algébrique, Chapitre I et suivants. -- Paris, Presses universitaires de France,
1960--1961 (Institut des Hautes Études Scientifiques, Publications
mathématiques 4, 8, \ldots).

\item[{[3]}] GROTHENDIECK (Alexander). -- Géométrie formelle et géométrie
algébrique, \emph{Séminaire Bourbaki}, v. 11, 1958/59,
n\textsuperscript{o} 182, 28 p.

\item[{[4]}] GROTHENDIECK (Alexander). -- Techniques de descente et théorèmes
d'existence en géométrie algébrique, I--IV, \emph{Séminaire Bourbaki}, v. 12,
1959/60, n\textsuperscript{o} 190, 29 p. e n\textsuperscript{o} 195, 22 p.;
v. 13, 1960/61, n\textsuperscript{o} 212, 20 p. e
n\textsuperscript{o} 221, 28 p.

\item[{[5]}] GROTHENDIECK (Alexander). -- Techniques de construction en
géométrie analytique, I--X, \emph{Séminaire Cartan}, v. 13, 1960/61:
Déformation des structures analytiques complexes, n\textsuperscript{o} 7--17.

\item[{[6]}] GROTHENDIECK (Alexander). -- Séminaire de géométrie algébrique,
redigido por Lichtenbaum, Harvard University, 1961 (a aparecer).

\item[{[7]}] LANG (Serge). -- \emph{Abelian varieties}. -- Nova York,
Interscience Publishers, 1959 (Interscience Tracts in pure and applied
Mathematics, 7).

\item[{[8]}] MUMFORD (D.). -- An elementary theorem in geometric invariant
theory, \emph{Bull. Amer. math. Soc.}, v. 67, 1961, p. 483--486.

\item[{[9]}] MUMFORD (D.). -- The topology of normal singularities of an
algebraic surface and a criterion for simplicity. -- Paris, Presses
universitaires de France, 1961 (Institut des Hautes Études Scientifiques,
Publications mathématiques, 9).

\item[{[10]}] MUMFORD (D.) e TATE (J.). -- Séminaire de géométrie algébrique,
Harvard University, Spring Term 1962 (a aparecer).

\item[{[11]}] OORT (F.). -- Sur le schéma de Picard,
\emph{Bull. Soc. math. France}, v. 90, 1962 (a aparecer).

\item[{[12]}] SESHADRI (C. S.). -- Tese a aparecer em
\emph{Annali di Matematica}.
\end{description}

\par

\def\FgaUnitKey{e232}\clearpage
\chapter*{Errata da exposição n.\textsuperscript{o} 232}
\addcontentsline{toc}{chapter}{Errata da exposição n.\textsuperscript{o} 232}
\markboth{Errata da exposição 232}{Errata da exposição 232}
\sid{e232}
\src{303}

\noindent\textbf{7. TDTE V: Esquemas de Picard. Teoremas de existência.}

\begin{center}
[Seminário Bourbaki, vol. 14, 1961/62, n.\textsuperscript{o} 232, 19 p.]
\end{center}

\noindent\textbf{Página 7, observação 3.3.} --- A questão levantada aqui foi
resolvida afirmativamente por MUMFORD (ver os comentários à exposição seguinte).

\noindent\textbf{Página 13, observação 5.2.} --- Como assinalo no início da
exposição seguinte, a conjectura de existência formulada aqui é falsa; contudo,
MUMFORD consegue demonstrar, por seus métodos, um teorema ligeiramente mais
fraco.

\noindent\textbf{Página 14, linha 7.} --- Depois de «completo», acrescentar «com
corpo residual algebricamente fechado». O contraexemplo de MUMFORD que acabamos
de mencionar mostra, aliás, que essa restrição é indispensável.

\noindent\textbf{Página 17, observação 6.6.} --- Como assinalamos no início da
exposição seguinte, a questão levantada aqui acaba de ser resolvida
afirmativamente por MURRE.

\noindent\textbf{Página 18, observação 6.7.} --- A questão levantada aqui tem
resposta afirmativa (cf. o último parágrafo dos comentários à exposição
seguinte).

\noindent\textbf{Página 18, observação 6.8, 4\textsuperscript{a} linha.} --- Em vez de
«regular», leia-se «suave sobre $k$».

\par

\def\FgaUnitKey{236}\unit{Exposição 236 -- Os esquemas de Picard: propriedades gerais}
\sid{236}
\markboth{Exposição 236}{Exposição 236}
\src{221}
\seminares{14.º ano, 1961/62, n.º 236\\Maio de 1962}
\paperhead{TÉCNICA DE DESCIDA E TEOREMAS DE EXISTÊNCIA EM GEOMETRIA ALGÉBRICA\\
VI. OS ESQUEMAS DE PICARD: PROPRIEDADES GERAIS}%
  {por Alexander GROTHENDIECK}

\fsection{0}{Complementos à exposição anterior ([3], V)\textsuperscript{(1)}}

Houve alguns avanços nas questões de existência dos pré-esquemas de Picard
levantadas em [3], V:

\medskip
\noindent a. (Mumford). Não é verdade, em geral, que, quando
$f:X\longrightarrow S$ é um morfismo projetivo e separável (= plano com fibras
separáveis), exista o pré-esquema de Picard $\operatorname{Pic}_{X/S}$, mesmo
que as fibras de $f$ tenham dimensão 1 e que $S$ seja o espectro de um anel de
valorização discreta completo. Um exemplo é obtido tomando
$S=\operatorname{Spec}\mathbf R[[t]]$ e, para $X$, o subesquema de
$\mathbf P^2_S$ definido pela equação (onde $x,y,z$ são variáveis homogêneas)
\[
x^2+y^2=tz^2,
\]
que representa, portanto, uma cônica que degenera geometricamente em duas
retas concorrentes, embora a fibra especial sobre o corpo $\mathbf R$ seja,
não obstante, irredutível (ela é dada pela equação $x^2+y^2=0$ sobre
$\mathbf R$). Vê-se facilmente que, após a extensão étale
$S'\longrightarrow S$, com $S'=\operatorname{Spec}\mathbf C[[t]]$, existe o
pré-esquema de Picard de $X'/S'$, do qual se obtém uma descrição explícita como
união disjunta de cópias de $\widetilde S$, onde $\widetilde S$ é obtido de
$S$ duplicando infinitas vezes a origem. Verifica-se facilmente que o dado de
descida sobre $\operatorname{Pic}_{X'/S'}$ para $S'\longrightarrow S$ (dado
aqui pela ação do grupo de Galois $\mathbf Z/2\mathbf Z$ de $S'$ sobre $S$)
não é efetivo: o grupo permuta entre si certos pontos duplicados (de modo que
há órbitas que não estão contidas em um aberto afim). Mumford pode demonstrar,
contudo, que, se $f:X\longrightarrow S$ é um morfismo projetivo separável tal
que, para todo $s\in S$, as componentes irredutíveis de $X_s$ são
geometricamente irredutíveis sobre $k(s)$, então
$\operatorname{Pic}_{X/S}$ existe; sua demonstração apoia-se em um refinamento
de seu teorema de passagem ao quociente, cf. [7]. Observe-se, por outro lado,
que continua sendo possível que, sem hipóteses sobre as componentes
irredutíveis das fibras $X_s$, o esquema $\operatorname{Pic}^{\tau}_{X/S}$
(que será introduzido mais adiante) exista mesmo assim.
\src{222}
\noindent b. (Murre). Seja $X$ um esquema próprio sobre um corpo. Então
$\operatorname{Pic}_{X/k}$ existe. A demonstração retoma em parte a
demonstração de Chevalley [1] e apoia-se essencialmente na estrutura de grupo
do funtor de Picard.

Para alguns complementos relativos à teoria dos esquemas de Picard,
especialmente em relação com os esquemas abelianos, será útil consultar [7].
Por fim, uma lacuna notável da presente exposição é a ausência de «critérios
de equivalência» que permitam comparar o esquema de Picard de um esquema
projetivo com os de suas seções hiperplanas; os teoremas-chave para desenvolver
tais critérios encontram-se em [5], aos quais devem ser acrescentados os
teoremas de existência dos esquemas de Picard.

\fsection{1}{Propriedades topológicas dos pré-esquemas em grupos comutativos}

Sejam primeiro $k$ um corpo e $G$ um pré-esquema em grupos sobre $k$. Como o
elemento neutro $e$, por ser racional sobre $k$, é necessariamente fechado,
segue-se imediatamente que a diagonal de $G\times_kG$ é fechada e, portanto,
que $G$ é separado: todo pré-esquema em grupos sobre um corpo é separado.
Denotaremos por $G^\circ$ a componente conexa do elemento neutro $e$. Como $e$
é racional sobre $k$, $G^\circ$ é, de fato, geometricamente conexa, isto é, a
formação de $G^\circ$ é compatível com a extensão do corpo de base. Segue-se
também que $G^\circ$ é estável por multiplicação como conjunto e, quando $G$ é
localmente noetheriano, de modo que $G^\circ$ é aberto, pode-se considerar
$G^\circ$ como um subgrupo aberto de $G$. Doravante, suporemos que $G$ é
localmente de tipo finito sobre $k$; então $G^\circ$ é geometricamente
irredutível e de tipo finito sobre $k$. Com efeito, pode-se supor que $k$ é
algebricamente fechado e, além disso, que $G$ é reduzido (pois
$G_{\mathrm{red}}$ será então um subgrupo de $G$, tendo em vista que
$G_{\mathrm{red}}\times_kG_{\mathrm{red}}$ será ainda reduzido), logo liso
sobre $k$ em um aberto não vazio e, portanto, em toda parte, como se vê
transladando esse aberto. Mas então $G$ é localmente irredutível, de modo que
suas componentes irredutíveis coincidem com suas componentes conexas, e
$G^\circ$ é irredutível. Seja então $U$ uma vizinhança afim de $e$ em
$G^\circ$; usando a irredutibilidade de $G^\circ$, vê-se imediatamente que
$U\cdot U=G^\circ$, o que prova que $G^\circ$ é quase-compacto e, portanto, de
tipo finito sobre $k$.

Para simplificar, suponhamos que $G$ seja comutativo. Para todo inteiro $n>0$,
seja $G^{(n)}$ a imagem inversa de $G^\circ$ pelo homomorfismo $\varphi_n$ de
elevação à $n$-ésima potência em $G$; assim, $G^{(n)}$ é um subgrupo aberto de
$G$.
Ponhamos:
\[
G^\tau=\bigcup_nG^{(n)}.
\]
\src{223}
\[
G^\sigma=\bigcup_{(n,p)=1}G^{(n)},
\qquad
G^\rho=\bigcup_hG^{(p^h)},
\]
onde $p$ é o expoente característico do corpo $k$. Obtêm-se, portanto,
subgrupos abertos de $G$ que satisfazem
\[
G^\sigma\cap G^\rho=G^\circ,
\qquad
G^\sigma G^\rho=G^\tau.
\]

\result{OBSERVAÇÃO}{} Pode-se construir o esquema em grupos quociente
$G/G^\circ=N$ (cf. [3], IV) e definir então $G^\tau$, $G^\sigma$ e
$G^\rho$ como as imagens inversas em $G$ do subgrupo de torção de $N$
(respectivamente, de sua componente $p$-primária, respectivamente, do
complemento natural desta última, soma das componentes $q$-primárias para
$q$ primo $\neq p$).
Observe-se, a esse respeito, que $N$ é um esquema em grupos discreto separável
sobre $k$ e, portanto, uma vez escolhida uma clausura algébrica $\overline k$
de $k$, que dá lugar a um grupo de Galois $\pi$, identifica-se com um grupo
discreto ordinário sobre o qual $\pi$ age por automorfismos. É dessa maneira
que se pode interpretar de modo evidente a construção do subgrupo de torção e
a decomposição deste em suas componentes $q$-primárias. Quando $G$ é o esquema
de Picard de um esquema $X$ próprio sobre $k$, $N$ pode ser chamado
\emph{esquema de Néron--Severi} (reduzido) de $X$ sobre $k$. Quando
$G^\circ_{\mathrm{red}}$ é um subesquema em grupos de $G^\circ$, o que
ocorre, em particular, sempre que $k$ é perfeito ou que $G^\circ$ é próprio
sobre $k$ (por exemplo, quando $X$ é geometricamente normal), convém introduzir
também o quociente $N'=G/G^\circ_{\mathrm{red}}$, que tende a comportar-se
melhor do que $N$ quanto à especialização, isto é, quando $X$ varia em uma
família de esquemas algébricos.

Sejam agora $S$ um pré-esquema localmente noetheriano e $G$ um pré-esquema em
grupos sobre $S$, localmente de tipo finito sobre $S$. Não supomos que $G$ seja
de tipo finito sobre $S$, nem separado sobre $S$. Poremos então
\[
G^\circ=\bigcup_{s\in S}(G_s)^\circ,
\]
e, quando $G$ for comutativo,
\[
G^\tau=\bigcup_{s\in S}(G_s)^\tau,
\qquad
G^\sigma=\bigcup_{s\in S}(G_s)^\sigma,
\qquad
G^\rho=\bigcup_{s\in S}(G_s)^\rho.
\]
\src{224}
Trata-se de subconjuntos de $G$, estáveis pela multiplicação em $G$, o que
evidentemente não implica que possam ser definidos mediante subpré-esquemas em
grupos de $G$. Em particular, parece não existir, em geral, um subpré-esquema
em grupos de $G$ cujo conjunto subjacente seja $G^\circ$. Naturalmente, se um
desses conjuntos for aberto, então, munido da estrutura induzida, será um
subpré-esquema aberto em grupos de $G$. Veremos que isso sempre ocorre com
$G^\tau$; assim, do ponto de vista dos funtores representáveis, e em particular
do ponto de vista das «especializações», a equivalência numérica comporta-se
de maneira mais satisfatória do que a equivalência algébrica. Eis as principais
propriedades gerais dos conjuntos que acabamos de definir:

\result{TEOREMA}{1.1} $G^\circ$, $G^\tau$, $G^\sigma$ e $G^\rho$ são
localmente construtíveis. Além disso:

\noindent (i) $G^\circ$ é quase-compacto sobre $S$. Quando os $G_s^\circ$ são
próprios e $G$ é separado sobre $S$, então $G^\circ$ é próprio sobre $S$ e,
portanto, fechado em $G$.

\noindent (ii) $G^\tau$ é aberto. Se $G^\circ$ é fechado, o mesmo vale para
$G^\tau$.

\noindent (iii) Se $G^\circ$ é fechado, o mesmo vale para $G^\sigma$, desde
que se esteja em característica igual, isto é, que todos os corpos residuais de
$S$ tenham a mesma característica. Se $G^\circ$ é fechado e
$G\longrightarrow S$ é universalmente aberto nos pontos de $G^\sigma$ (cf.
Corolário 1.5 abaixo), então $G^\sigma\longrightarrow S$ é universalmente
aberto.

\noindent (iv) Se $G^\circ$ é fechado, o mesmo vale para $G^\rho$. Suponhamos
que se esteja em característica igual e que, para todo inteiro $n>0$ tal que
$(n,p)=1$, o homomorfismo de elevação à $n$-ésima potência em $G$ seja aberto;
então $G^\rho$ é aberto.

Demos algumas indicações sobre a demonstração. O fato de $G^\circ$ ser
localmente construtível e quase-compacto sobre $S$ está contido no lema
seguinte:

\result{LEMA}{1.2} Suponhamos que $S\neq\varnothing$. Então existem um aberto
não vazio $U$ em $S$, um esquema em grupos $H$ de tipo finito sobre $U$ com
fibras conexas e, por fim, um homomorfismo de esquemas em grupos
$H\longrightarrow G|U$ que é uma imersão aberta cuja imagem é
$G^\circ|U$.

Para demonstrar o lema, pode-se supor que $S$ seja irredutível; seja $\eta$ seu
ponto genérico, e façamos a mudança de base
$S'=\operatorname{Spec}(\mathcal O_{S,\eta})\longrightarrow S$. Obtém-se um
esquema em grupos $G'$ sobre um anel artiniano local
$\mathcal O_{S,\eta}=A$, no qual há um subesquema aberto em grupos
$G'{}^\circ$ de tipo finito sobre $A$, como se disse acima. Este provém,
portanto, de um esquema em grupos $H$ de tipo finito sobre uma vizinhança
aberta $U$ de $\eta$, e a imersão canônica
$G'{}^\circ\longrightarrow G'$ provém de uma imersão
\src{225}
aberta $H\longrightarrow G|U$, que será um homomorfismo de esquemas em grupos
para $U$ suficientemente pequeno. Como as fibras de $H$ são conexas quando se
toma $U$ suficientemente pequeno, e têm todas a mesma dimensão, igual à das
fibras de $G$ para $U$ suficientemente pequeno, segue-se que, para todo
$s\in U$, a imagem de $H_s$ em $G_s$ é exatamente $G_s^\circ$ (após reduzir
novamente $U$), o que demonstra 1.2.

A segunda afirmação de (i) está contida no lema seguinte (que se aplica a uma
vizinhança aberta quase-compacta de $G^\circ$ em $G$):

\result{LEMA}{1.3} Sejam $X$ um pré-esquema separado de tipo finito sobre um
$S$ localmente noetheriano, $g$ uma seção de $X$ sobre $S$ e $X^\circ$ a união
das componentes conexas dos $g(s)$ nos $X_s$. Seja $s\in S$ tal que
$X_s^\circ$ seja próprio sobre $k(s)$. Então existe uma vizinhança aberta $U$
de $s$ tal que $X^\circ|U$ seja próprio sobre $U$ e, a fortiori, fechado em
$X|U$.

Por descida fielmente plana da base, reduzimo-nos ao caso em que $S$ é o
espectro de um anel local completo e $s$ é seu ponto fechado. Aplicando [2],
III, 5.5.1, vê-se que $X$ se decompõe na união disjunta de dois abertos $X'$ e
$X''$, o primeiro próprio sobre $S$ e tal que $X'_s=X_s^\circ$. Isso nos reduz
ao caso em que $X=X'$, isto é, em que $X$ é próprio sobre $S$. Nesse caso,
conclui-se por uma demonstração padrão que utiliza o critério valuativo para
que uma parte seja própria (omitido em [2], cap. II).

Demonstremos que $G^\tau$ é aberto ou, o que equivale ao mesmo, tendo em vista
que a formação de $G^\tau$ (assim como a de $G^\circ$, $G^\sigma$ e $G^\rho$)
comuta com a extensão da base: para toda seção $g$ de $G$ sobre $S$,
$g^{-1}(G^\tau)$ é aberto. Isso significa duas coisas:

\medskip
\noindent a. Seja $y\in S$ tal que $g(y)\in G_y^\tau$. Então, para todo
$y'\in\overline{\{y\}}$ suficientemente próximo de $y$, tem-se
$g(y')\in G_{y'}^\tau$.

\medskip
\noindent b. Seja $y'\in S$ tal que $g(y')\in G_{y'}^\tau$. Então, para toda
generalização $y$ de $y'$, tem-se $g(y)\in G_y^\tau$.

Para (a), observe-se que existe um inteiro $n>0$ tal que
$g^n(y)\in G_y^\circ$; como $G^\circ$ é construtível,
$(g^n)^{-1}(G^\circ)$ é construtível e, portanto, tem-se
$g^n(y')\in G_{y'}^\circ$ para $y'\in\overline{\{y\}}$ suficientemente
próximo de $y$. Para (b), observe-se que as potências
$g(y')^n=g^n(y')$ permanecem em um aberto quase-compacto de $G_{y'}$ (pois
estão contidas em um número finito de classes módulo $G_{y'}^\circ$); logo,
existe um aberto quase-compacto $U$ em $G$ que contém os $g^n(y')$ e, portanto,
também suas generalizações $g^n(y)$. Assim, as potências de $g(y)$ permanecem
em um aberto quase-compacto de $G_y$, o que implica facilmente que
$g(y)\in G_y^\tau$.
\src{226}
Suponhamos $G^\circ$ fechado; demonstremos que o mesmo vale para $G^\tau$.
Como já se sabe que $G^\tau$ é aberto e, portanto, localmente construtível,
resta demonstrar que é estável por especialização, o que decorre do fato de ser
uma união de fechados, a saber, as imagens inversas, pelos homomorfismos
$\varphi_n$ de elevação à $n$-ésima potência, do fechado $G^\circ$.

O mesmo argumento demonstrará que $G^\sigma$ e $G^\rho$ são fechados se
$G^\circ$ o for (sob a condição, no primeiro caso, de se estar em
característica igual), uma vez demonstrado que $G^\sigma$ e $G^\rho$ são
localmente construtíveis. Ora, para $x\in G^\tau$, seja $\nu(x)$ o menor
inteiro $n>0$ tal que o homomorfismo $\varphi_n$ de elevação à $n$-ésima
potência envie $x$ a $G^\circ$. Assim, $G^\sigma$ (respectivamente,
$G^\rho$) é formado pelos $x\in X$ tais que $\nu(x)$ seja primo com $p$
(respectivamente, uma potência de $p$), e nossa afirmação de construtibilidade
decorre então da seguinte afirmação mais precisa:

\result{LEMA}{1.4} A função $\nu$ sobre $G^\tau$ é localmente construtível.

Isso significa, com efeito, que, para todo inteiro $n>0$, o conjunto dos
$x\in G^\tau$ tais que $\nu(x)=n$ é localmente construtível; ora, ele é a
diferença entre $\varphi_n^{-1}(G^\circ)$ e a união dos
$\varphi_d^{-1}(G^\circ)$, onde $d$ percorre os divisores próprios de $n$.
Como $G^\circ$ é localmente construtível, o mesmo vale para os
$\varphi_d^{-1}(G^\circ)$ e, portanto, também para a diferença anterior.

Suponhamos $G^\circ$ fechado e $G\longrightarrow S$ universalmente aberto nos
pontos de $G^\sigma$; demonstremos que $G^\sigma\longrightarrow S$ é aberto,
isto é, que transforma uma vizinhança em $G^\sigma$ de um $x\in G^\sigma$
em uma vizinhança de $y=f(x)$. Como $G^\sigma$ é localmente construtível,
reduzimo-nos a demonstrar que, para toda generalização $y'$ de $y$, existe uma
generalização $x'$ de $x$ em $G^\sigma$ sobre $y'$. Isso nos reduz, por
mudança de base, ao caso em que $S$ é o espectro de um anel de valorização
discreta e em que $y,y'$ são, respectivamente, o ponto fechado e o ponto
genérico. Usando o fato de que $G\longrightarrow S$ é aberto em $x$ (logo,
existe uma generalização $x_1$ de $x$ em $G$ sobre $y'$), pode-se supor ainda
que existe uma seção $g$ de $G$ sobre $S$ tal que $x=g(y)$ (após fazer outra
mudança de base). Se $k(y')$ tem característica 0, basta tomar qualquer
generalização $x'$ de $x$ em $G$ sobre $y'$; ela pertence a $G^\tau$ porque
$G^\tau$ é aberto e, portanto, a $G^\sigma$ porque
$G_{y'}^\sigma=G_{y'}^\tau$. Se a característica de $k(y')$ é $p>0$,
ponhamos
\[
\nu(g(y'))=p^hm\qquad\text{com }(m,p)=1 ;
\]
\src{227}
sejam $a$ e $b$ dois inteiros tais que $ap^h+bm=1$; ponhamos
\[
g_1=g^{ap^h},\qquad g_2=g^{bm},\qquad\text{logo}\qquad g=g_1g_2 ;
\]
portanto, por construção, tem-se $g_1(y')\in G^\sigma$ e
$g_2(y')\in G^\rho$. Como $G^\circ$ é fechado, segue-se que
$g_1(y)\in G^\sigma$ e $g_2(y)\in G^\rho$; daí, em virtude de
\[
g(y)=g_1(y)g_2(y)\in G^\sigma,
\]
também $g_2(y)\in G^\sigma$ e, portanto,
\[
g_2(y)\in G_y^\sigma\cap G_y^\rho=G_y^\circ.
\]
Ora, da hipótese e do fato de que $S$ é o espectro de um anel de valorização
discreta, resulta que $G-(G_y-G_y^\circ)$ é um aberto de $G$ sobre o qual
$G\longrightarrow S$ induz um morfismo aberto; logo, por todo ponto de
$G_y^\circ$ passa uma «quase-seção». Por conseguinte, após fazer ainda uma
extensão da base $S$, pode-se supor que existe uma seção $g'_2$ de $G^\circ$
sobre $S$ tal que $g'_2(y)=g_2(y)$. Ponhamos $g'=g_1g'_2$; então, por
construção, $g'(y)=g(y)=x$, e
\[
g'(y')=g_1(y')g'_2(y')\in G^\sigma,
\]
logo $g(y')=x'$ é uma generalização de $x$ em $G^\sigma$ sobre $y'$, o que demonstra
(iii).

Demonstremos, por fim, a última afirmação (iv). Reduzimo-nos a provar que, se
$x\in G^\rho$, então toda generalização $x'$ de $x$ pertence a $G^\rho$.
Pode-se supor (após tomar as imagens por $\varphi_{n^h}$ para um $h$
conveniente) que se tem até mesmo
$x\in G^\circ$. Então, para todo inteiro $n$ primo com a característica, $x$
pertence à imagem de $\varphi_n$ (pois a elevação à $n$-ésima potência em um
grupo conexo de tipo finito sobre um corpo de característica prima com $n$ é
sobrejetiva). Como $\varphi_n$ é aberta, segue-se que $x'$ também pertence à
imagem de $\varphi_n$. Mais precisamente, seja $U$ um aberto quase-compacto de
$G^\tau$ que contém $G_y^\circ$; então se tem $x'\in\varphi_n(U)$ para todo
$n$ primo com $p$. Tomando para $n$ um múltiplo comum dos fatores primos com
$p$ dos $\nu(z)$, para $z\in U$, obtém-se $x'\in G^\rho$.

A afirmação (1.1), (iii), completa-se assim:

\result{COROLÁRIO}{1.5} Seja $n>1$ um inteiro tal que o homomorfismo de
elevação à $n$-ésima potência $\varphi_n:G\longrightarrow G$ seja
universalmente aberto (por exemplo, étale); ponhamos
\[
G^{(n)}=\varphi_n^{-1}(G^\circ),
\]
e suponhamos que as fibras conexas $G_s^\circ$ não «contenham
\src{228}
componente aditiva» (isto é, o grupo obtido pela extensão do corpo $k(s)$ a
uma clausura algébrica não contém nenhum subgrupo isomorfo a
$\mathbf G_a$). Então $G\longrightarrow S$ é universalmente aberto nos
pontos dos $G^{(n^h)}$. Em particular, se os $G_s^\circ$ não contêm componente
aditiva e se, para todo inteiro $n>1$, o homomorfismo $\varphi_n$ é
universalmente aberto nos pontos de característica residual prima com
$p$, então
$G\longrightarrow S$ é universalmente aberto nos pontos de $G^\sigma$ e,
portanto, por (1.1), (iii), se $G^\circ$ é fechado,
$G^\sigma\longrightarrow S$ é universalmente aberto. Nesses casos,
$s\longmapsto\dim G_s$ é uma função localmente constante sobre $S$.

Com efeito, da hipótese resulta que o núcleo ${}_nG$ de $\varphi_n$ é
universalmente aberto sobre $S$; logo, $G\longrightarrow S$ é universalmente
aberto nos pontos de ${}_nG$ e, portanto (substituindo $n$ por $n^h$), nos
pontos dos ${}_{n^h}G$. Ora, a hipótese sobre as fibras $G_s^\circ$ serve
exatamente para assegurar que os pontos de ordem uma potência de $n$ em $G_s$
são densos na união dos $G^{(n^h)}$; daí resulta facilmente que
$G\longrightarrow S$ é universalmente aberto em todos os pontos dessa união,
em particular ao longo da seção neutra, de onde decorre facilmente que
$s\longmapsto\dim G_s$ é uma função localmente constante.

\result{OBSERVAÇÃO}{1.6} Recordemos que um morfismo $X\longrightarrow S$ é
dito universalmente aberto se transforma todo aberto em um aberto e conserva
essa louvável propriedade após toda mudança de base $S'\longrightarrow S$.
Isso significa também (quando $S$ é localmente noetheriano e
$X\longrightarrow S$ é localmente de tipo finito) que toda componente
irredutível de $X$ domina $S$ e que essa propriedade é preservada após toda
mudança de base $S'\longrightarrow S$. Em ambas as afirmações, basta, aliás
(sob as hipóteses de finitude acima), testar mudanças de base
$S'\longrightarrow S$, em que $S'$ é o espectro de um anel de valorização
discreta (completo, com corpo residual algebricamente fechado, se assim se
quiser~...). A definição estende-se de maneira evidente ao caso de um
subconjunto $Z$ de $X$ (como o subconjunto $G^\sigma$ de $G$). É um fenômeno
perfeitamente normal, mesmo partindo de um morfismo projetivo liso
$X\longrightarrow S$ com fibras geométricas conexas (por exemplo, o produto
fibrado da família modular de curvas elípticas consigo mesma sobre
$S=\operatorname{Spec}\mathbf C[j]$), que $\operatorname{Pic}_{X/S}$ não
seja universalmente aberto sobre $S$, isto é, que possa haver componentes
irredutíveis de $\operatorname{Pic}_{X/S}$ situadas inteiramente sobre um
único ponto de $S$: isso está ligado ao fato de que o posto do grupo de
Néron--Severi das fibras de $X/S$ pode apresentar saltos para cima (fenômenos
de «multiplicação complexa»). Em contrapartida, (1.5) assegura-nos que, nos
casos favoráveis, $\operatorname{Pic}_{X/S}^\sigma$ (e, ao que parece, na
maioria das vezes até mesmo $\operatorname{Pic}_{X/S}^\tau$) é
universalmente aberto sobre $S$.
\src{229}
Por fim, eis um caso útil em que, excepcionalmente, $G^\circ$ consente em ser
aberto:

\result{COROLÁRIO}{1.7} Suponhamos que $G\longrightarrow S$ seja
universalmente aberto nos pontos de $G^\circ$ (cf. (1.5)) e que as fibras
$G_s$ sejam separáveis e, portanto, lisas sobre $k(s)$ (esta última condição é automaticamente satisfeita em
característica residual zero, em virtude de um resultado de Cartier). Então
$G^\circ$ é aberto em $G$. Se, além disso, $S$ é reduzido, $G^\circ$ é liso
(em particular, plano) sobre $S$.

A primeira afirmação pode, aliás, ser precisada observando-se que, se, para um
$s\in S$ dado, $G\longrightarrow S$ é universalmente aberto nos pontos de
$G_s^\circ$ e $G_s^\circ$ é separável sobre $k(s)$, então $G^\circ$ é uma
vizinhança de $G_s^\circ$ (e, aliás, a hipótese feita em $s$ continuará
verificada nos pontos vizinhos). Esse enunciado, independente de toda estrutura
de grupos, será demonstrado em [2], IV, §~7. A última afirmação de (1.7), também
independente de toda estrutura de grupos e de toda questão de componente
conexa, é um caso particular de um critério de planicidade dado em [IV], §~5,
que implica, mais geralmente:

\result{COROLÁRIO}{1.8} Seja $U$ um subconjunto aberto de uma fibra $G_s$ tal
que $G\longrightarrow S$ seja universalmente aberto nos pontos de $U$ (cf.
(1.5)). Se $G_s$ é separável sobre $k(s)$ e $S$ é reduzido em $s$, então $G$ é
plano sobre $S$ nos pontos de $U$ (e, portanto, neste caso, liso sobre $S$ nos
pontos de $U$).

\result{OBSERVAÇÕES}{1.9} Ignoro se, quando $G$ é separado sobre $S$,
$G^\circ$ ou $G^\tau$ é sempre fechado em $G$; isso parece pouco provável.
Encontram-se, em todo caso, contraexemplos evidentes se se abandona a hipótese
de separação, por exemplo tomando a reta afim cuja origem é duplicada uma
infinidade enumerável de vezes, obtendo-se assim um esquema em grupos $G$ sobre
a reta afim cujas fibras se reduzem todas ao grupo unidade, salvo uma cuja
fibra é $\mathbf Z$. Esse esquema em grupos é um subesquema aberto em grupos,
a saber, o fecho da seção unidade, no pré-esquema de Picard do $S$-esquema $X$
correspondente a uma família de cônicas que degenera em duas retas concorrentes.

Além disso, mesmo partindo de um esquema em grupos finito e plano $G$ sobre o
espectro $S$ de um anel de valorização discreta $V$ e, por conseguinte, se
$G^\circ$ se reduz à seção unidade, sendo, portanto, fechado, várias
conclusões tornam-se falsas quando se abandonam certas hipóteses. Suponhamos
que $V$ seja de característica igual $p>0$ e seja $G$ o núcleo do homomorfismo
$\mathbf G_a\longrightarrow\mathbf G_a$ definido pelo homomorfismo de
funtores
\[
f\longmapsto f^p-tf,
\]
onde $t$ é um uniformizante de $V$.
\src{230}
(Lembremos que, por definição, o «grupo aditivo» $\mathbf G_a$ sobre $S$
representa o funtor $S'\longmapsto\Gamma(S',\mathcal O_{S'})$.) Então, como
$S$-esquema,
\[
G=\operatorname{Spec}V[x]/(x^p-tx)
\]
é a união de $p$ «retas» concorrentes, representando um grupo
$\mathbf Z/p\mathbf Z$ que degenera em um grupo infinitesimal de tipo aditivo.
Neste exemplo, $G^\circ=G^\sigma$ (e é uma das $p$ retas), portanto não é
aberto em $G$. No caso de característica mista, de característica residual
$p>0$, partamos do esquema em grupos $H=\boldsymbol\mu_p$, núcleo da aplicação
de elevação à $p$-ésima potência em $\mathbf G_m$, portanto
\[
H=\operatorname{Spec}V[x]/(x^p-1) ;
\]
é ainda a união de $p$ retas concorrentes, representando um grupo
$\mathbf Z/p\mathbf Z$ (em característica zero) que degenera em um grupo
infinitesimal de tipo multiplicativo em característica $p$. Seja $H'$ o
«esquema em grupos constante» definido pelo grupo finito ordinário
$\mathbf Z/p\mathbf Z$, portanto a soma disjunta, indexada por
$\mathbf Z/p\mathbf Z$, de cópias de $S$, ou ainda
$\operatorname{Spec}V^{\mathbf Z/p\mathbf Z}$. Então $G=H\times_S H'$ descreve
um grupo $(\mathbf Z/p\mathbf Z)^2$ em característica zero, que degenera em um
grupo infinitesimal vezes $\mathbf Z/p\mathbf Z$ em característica $p$. Aqui
$G^\rho$ é a união da seção unidade e da fibra especial, portanto $G^\rho$ não
é aberto, ao passo que $G^\sigma$ é a união da seção unidade e da fibra
genérica, portanto não é fechado, contrariamente ao que se afirma nos casos de
característica igual em (1.1), (iii) e (iv). Naturalmente, esses fenômenos estão
ligados à característica $p>0$. Os resultados precedentes dão, com efeito:

\result{COROLÁRIO}{1.10} Suponhamos que as características residuais de $S$
sejam todas nulas, portanto $G^\tau=G^\sigma$ e $G^\circ=G^\rho$. Então
$G^\tau=G^\sigma$ é aberto, e até aberto e fechado se $G$ é separado sobre $S$
e se os $G_s^\circ$ são próprios; sob essa mesma hipótese, $G^\circ=G^\rho$ é
próprio sobre $S$, portanto fechado em $G$. Suponhamos, por fim, que, para todo
inteiro $n>1$, o homomorfismo de elevação à $n$-ésima potência em $G$ seja
universalmente aberto; então $G^\circ$ é aberto; e, se além disso os
$G_s^\circ$ não têm componente aditiva (por exemplo, se os $G_s^\circ$ são
próprios, como acima), então $G^\tau\longrightarrow S$ é universalmente aberto,
e até liso se $S$ é reduzido.

Assinalemos, por fim, o seguinte resultado fácil:

\result{PROPOSIÇÃO}{1.11} Existe uma parte aberta $U$ de $S$ tal que o conjunto
$G'$ dos pontos de $G$ nos quais $G$ é liso (resp. plano) sobre $S$ seja o
conjunto subjacente a um subesquema em grupos aberto induzido de $G|U$. Além
disso, toda seção de $G$ sobre $U$ é uma seção de $G'$ sobre $U$.
\src{231}
\result{COROLÁRIO}{1.12} Se $G$ é liso (resp. plano) sobre $S$ nos pontos da
seção unidade, também o é nos pontos de toda seção de $G$
sobre $S$ e em todos os pontos de $G^\circ$. Se, além disso, para todo inteiro
$n>0$, o homomorfismo de elevação à $n$-ésima potência
$\varphi_n:G\longrightarrow G$ é étale nos pontos de característica prima com
$n$, então $G$ é liso (resp. plano) sobre $S$ em todos os pontos de $G^\sigma$.

\fsection{2}{Aplicação às propriedades locais dos esquemas de Picard}

\result{TEOREMA}{2.1}

\noindent(i) Seja $f:X\longrightarrow S$ um morfismo próprio e liso, e
suponhamos que $\operatorname{Pic}_{X/S}$ exista (por exemplo, que $f$ seja um
morfismo projetivo). Então $\operatorname{Pic}_{X/S}$ é separado sobre $S$ e,
para toda parte fechada $Z$ de $\operatorname{Pic}_{X/S}$ que seja de tipo
finito sobre $S$, $Z$ é próprio sobre $S$.

\noindent(ii) Seja $X$ um pré-esquema sobre um corpo $k$, próprio e
geometricamente normal. Então $\operatorname{Pic}_{X/S}^\circ$\footnote{O ramo
corrente traz $\operatorname{Pic}_{X/k}^\circ$, pois $S$ não intervém neste
enunciado sobre o corpo $k$.} é próprio sobre $k$.

\begin{proof}
(i) Os critérios valuativos ([2], II, parágrafo 7) reduzem-nos a provar o
seguinte: se $S$ é o espectro de um anel de valorização discreta completo e $U$
o aberto reduzido ao ponto genérico de $S$, então toda seção racional de
$\operatorname{Pic}_{X/S}$ sobre $S$, isto é, toda seção sobre $U$, prolonga-se
de maneira única a uma seção sobre $S$. Tendo em conta a definição de
$\operatorname{Pic}_{X/S}$, isso equivale ao seguinte enunciado: para todo
módulo invertível $\mathcal L$ sobre $V=f^{-1}(U)$, existe um módulo invertível
sobre $X$ que prolonga $\mathcal L$, único a menos de isomorfismo. Ora, isso
decorre facilmente da descrição dos módulos invertíveis sobre $V$, resp. $X$,
em termos de classes de divisores de Cartier, tendo em conta que os anéis
locais de $X$ são regulares ($X$ sendo liso sobre $S$ regular), portanto
fatoriais pelo teorema de Auslander, o que implica que todo divisor sobre
$S$\footnote{O ramo corrente traz aqui «sobre $X$»; a fatorialidade invocada é
a dos anéis locais de $X$.} é um divisor de Cartier. Com efeito, todo divisor
sobre $V$ pode ser prolongado a um divisor sobre $X$ tomando-se o seu fecho.
\end{proof}

\result{OBSERVAÇÃO}{2.2} A demonstração permanece válida supondo-se apenas que
$f$ é plano e que suas fibras $X_s$ são localmente interseções completas e
lisas sobre $k(s)$ em codimensão $\leq2$, tendo em conta o seguinte resultado
demonstrado em [5]: um anel local noetheriano de interseção completa que é
regular em codimensão $\leq3$ é fatorial («conjectura de Samuel»). Note-se que o
\src{232}
resultado torna-se falso se se substitui «codimensão $\leq2$» por «codimensão
$\leq1$», isto é, pela hipótese «normal», como se pode constatar no exemplo de
uma família de quádricas não singulares que degenera em um cone quadrático.

(ii) Utilizando o lema de Chow, reduzimo-nos ao caso em que $X$ é projetivo,
portanto imerso em $\mathbf P_k^n$; podemos supor $X$ conexo. Se $\dim X=1$,
então $X$ é liso sobre $k$, e aplica-se (i). Se
$\dim X\geq2$, pode-se utilizar uma variante dos «critérios de equivalência»
conhecidos, que implica existir um número finito de curvas $Y_i$ lisas em $X$
(obtidas como interseções de $X$ com subespaços lineares convenientes de
$\mathbf P_k^n$), tais que
\[
\operatorname{Pic}_{X/k}^\tau\longrightarrow
\prod_i\operatorname{Pic}_{Y_i/k}^\tau
\]
seja um monomorfismo e induza, a fortiori, um monomorfismo nas componentes
conexas. Como a componente conexa do segundo membro é própria sobre $k$ pelo
que precede, e como se trata de um homomorfismo de esquemas em grupos, que será
necessariamente uma imersão fechada, segue-se que
$\operatorname{Pic}_{X/k}^\circ$ também é próprio sobre $k$. Pode-se evitar o
recurso aos delicados critérios de equivalência utilizando-se a estrutura dos
grupos algébricos comutativos sobre um corpo algebricamente fechado (devida a
Chevalley--Borel); reduzimo-nos a provar que todo morfismo da reta afim

em $\operatorname{Pic}_{X/k}^\tau$ é constante, o que equivale a dizer que todo
módulo invertível sobre $X[t]$ provém de um módulo invertível sobre $X$, resultado bem conhecido e
muito elementar, que nem sequer utiliza o fato de $X$ ser próprio sobre $k$ (a
hipótese de $X$ ser normal permitindo reduzir-nos imediatamente ao caso em que
$X$ é regular).

\result{COROLÁRIO}{2.3} Seja $f:X\longrightarrow S$ um morfismo próprio e
normal (isto é, plano e com fibras geométricas normais), e suponhamos que
$\operatorname{Pic}_{X/S}$ exista; então $\operatorname{Pic}_{X/S}^\circ$ é
próprio sobre $S$, portanto fechado em $\operatorname{Pic}_{X/S}$; além disso,
$\operatorname{Pic}_{X/S}^\tau$ e $\operatorname{Pic}_{X/S}^\rho$ também são
fechados, assim como $\operatorname{Pic}_{X/S}^\sigma$ em «característica
igual».

Aplicam-se (1.1) e (2.1), (ii).

\result{COROLÁRIO}{2.4} Seja $f:X\longrightarrow S$ um morfismo próprio e liso,
tal que $\operatorname{Pic}_{X/S}$ exista e seja uma soma disjunta de esquemas
$P^{(i)}$ de tipo finito sobre $S$ (cf. [3], V, 4.1). Então cada $P^{(i)}$ é
próprio sobre $S$.

Isso decorre de (2.1), (i). Como se assinalou em (2.2), o resultado pode ser
generalizado fazendo-se hipóteses menos restritivas sobre as fibras de $f$, mas
torna-se falso se nos limitamos a supor $f$ normal. Nesse caso, ignoro, aliás,
se $\operatorname{Pic}_{X/S}^\tau$ é, não obstante, próprio sobre $S$, mesmo
admitindo que seja de tipo finito sobre $S$.
\src{233}
\result{TEOREMA}{2.5} Seja $f:X\longrightarrow S$ um morfismo próprio e plano,
tal que $\operatorname{Pic}_{X/S}$ exista, e, para todo inteiro $n$, seja
$\varphi_n$ o homomorfismo de elevação à $n$-ésima potência nesse pré-esquema
em grupos. Então $\varphi_n$ é étale em todo ponto
$x\in X$\footnote{O ramo corrente traz $x\in\operatorname{Pic}_{X/S}$, pois
$\varphi_n$ é o endomorfismo de potência do pré-esquema de Picard.} de
característica residual prima com $n$.

Pela caracterização infinitesimal dos morfismos étales, essa afirmação equivale
à seguinte:

\result{LEMA}{2.6} Suponhamos que $S$ seja o espectro de um anel artiniano local
$A$ cujo ideal maximal $\mathfrak m$ tenha potência $(\nu+1)$-ésima nula, e
sejam
\[
A_{\nu-1}=A/\mathfrak m^\nu,
\qquad X_{\nu-1}=X\otimes_A A_{\nu-1},
\]
$\mathcal L$ um módulo invertível sobre $X$ e $\mathcal L'_{\nu-1}$ um módulo
invertível sobre $X_{\nu-1}$ cuja potência tensorial $n$-ésima seja isomorfa a
\[
\mathcal L_{\nu-1}=\mathcal L\otimes_A A_{\nu-1}.
\]
Então existe um módulo invertível $\mathcal L'$ sobre $X$ cuja potência
tensorial $n$-ésima é isomorfa a $\mathcal L$ (se $n$ é primo com a
característica residual de $k=k(A)$).

Demonstremos o lema. Ponhamos
\[
V=\mathfrak m^\nu=\mathfrak m^\nu/\mathfrak m^{\nu+1} ;
\]
este é um espaço vetorial sobre $k=k(A)$. Começamos por prolongar
$\mathcal L'_{\nu-1}$ a um módulo invertível qualquer $\mathcal L'$ sobre $X$.
A obstrução a fazê-lo encontra-se em
\[
H^2(X_0,\mathcal O_{X_0})\otimes_k V,
\]
mas, em virtude da hipótese feita sobre $\mathcal L'_{\nu-1}$ e do fato de
$\mathcal L_{\nu-1}$ poder ser prolongado, vê-se que o produto dessa obstrução
por $n$ é nulo, portanto ela própria é nula, pois $n$ é primo com a
característica. Além disso, a arbitrariedade na extensão efetuada encontra-se
em
\[
H^1(X_0,\mathcal O_{X_0})\otimes_k V,
\]
e o desvio $\xi$ de $\mathcal L'^{\otimes n}$ em relação a $\mathcal L$
encontra-se no mesmo módulo; se se procura corrigir $\mathcal L'$ de modo a
tornar nulo esse desvio, reduzimo-nos a encontrar um $\eta$ no referido módulo
tal que $n\eta=\xi$. Isso é novamente possível porque $n$ é primo com a
característica.

\result{COROLÁRIO}{} Sob as condições de (2.5), suponhamos, além disso, que os
esquemas de Picard das fibras $X_s$ não contenham componente aditiva (por
exemplo, que os $X_s$ sejam geometricamente normais, cf. (2.1), (ii)). Então
$\operatorname{Pic}_{X/S}\longrightarrow S$ é universalmente aberto nos pontos
de $\operatorname{Pic}_{X/S}^\circ$. Se
$\operatorname{Pic}_{X/S}^\circ$ é fechado (por exemplo, se os $X_s$ são
geometricamente normais, cf. (2.3)), então
$\operatorname{Pic}_{X/S}^\sigma$ é ele próprio universalmente aberto sobre
$S$. Por fim, no caso de característica igual,
$\operatorname{Pic}_{X/S}^\rho\longrightarrow S$ é universalmente aberto.

Aplicam-se (1.5) e (1.1).
\src{234}
\result{COROLÁRIO}{2.7} Seja $f:X\longrightarrow Y$\footnote{O ramo corrente
traz $f:X\longrightarrow S$, pois tanto o pré-esquema de Picard quanto a função
seguinte têm base $S$.} um morfismo próprio e plano, tal que
$\operatorname{Pic}_{X/S}$ exista. Então a função
\[
s\longmapsto\dim\operatorname{Pic}_{X_s/k(s)}
\]
sobre $S$ é semicontínua superiormente (isto é, pode dar saltos para cima, mas
não para baixo), e é até contínua (isto é, localmente constante) se os
$\operatorname{Pic}_{X_s/k(s)}$ não contêm componente aditiva.

A primeira afirmação é trivial, ou quase, para todo pré-esquema em grupos
localmente de tipo finito sobre uma base localmente noetheriana, pois basta
examinar ao longo da seção unidade. A segunda afirmação decorre de (2.5).

\result{OBSERVAÇÃO}{2.8} Sejam $s,s'\in S$ tais que $s$ seja uma especialização
de $s'$. Então (2.7) equivale a uma desigualdade (resp. igualdade) entre as
dimensões dos esquemas de Picard de $X_{s'}$ e de sua «especialização» $X_s$.
Aliás, Serre havia observado, ainda antes da construção dos esquemas de Picard,
que a invariância da dimensão dos Picard dos $X_s$, no caso de um morfismo liso
$f:X\longrightarrow S$, era uma consequência formal da teoria da
especialização do grupo fundamental ([4], X) e das relações clássicas de Kummer
entre os pontos de ordem finita na variedade de Picard e o grupo fundamental
abelianizado ([4], XI). Se designamos, de modo geral, por $\alpha,\mu,\lambda$
as dimensões das partes abeliana, multiplicativa e aditiva de
$\operatorname{Pic}_{X_s/k(s)}^\circ$, e definimos do mesmo modo
$\alpha',\mu',\lambda'$, as relações conhecidas exprimem-se pelas seguintes
desigualdades:
\[
\tag{*}\alpha+\mu+\lambda\geq\alpha'+\mu'+\lambda'
\]
(válida contanto que $\operatorname{Pic}_{X/S}$ exista, portanto provavelmente
em todos os casos), desigualdade que, para $\lambda=\lambda'=0$, se reduz a uma
igualdade, válida sob as mesmas hipóteses de existência:
\[
\alpha+\mu=\alpha'+\mu' ;
\]
por outro lado, tem-se
\[
\tag{**}2\alpha+\mu\leq2\alpha'+\mu'
\]
\src{235}
pelo argumento de Serre, se os $X_s$ são separáveis (sem sequer supor a
existência de $\operatorname{Pic}_{X/S}^\tau$), ou se os
${}_n\operatorname{Pic}_{X/S}$ (núcleos de $\varphi_n$ no funtor de Picard) são
separados sobre $S$ (tendo em conta que são étales sobre $S$ graças a (2.5)).
Tender-se-ia a conjecturar que (*) é uma igualdade em todos os casos, ou pelo
menos se os $X_s$ são separáveis, e que se têm as desigualdades
\[
\tag{***}\alpha\leq\alpha',\qquad\lambda\geq\lambda',
\]
que deveriam ser válidas sempre que se tenha um pré-esquema em grupos,
localmente de tipo finito sobre $S$ localmente noetheriano, no qual a dimensão
das fibras seja constante (ver (1.3) para um resultado positivo nessa direção).

\result{OBSERVAÇÃO}{2.9} Em todos os casos conhecidos,
$\operatorname{Pic}_{X/S}^\tau$ é universalmente aberto sobre $S$, mas sem
dúvida não se deve tirar daí ilusões excessivas, mesmo que
$f:X\longrightarrow S$ seja liso; em todo caso, Mumford
construiu um exemplo (é verdade que com $S$ não reduzido, sendo de fato $S$ o
espectro de um anel artiniano) no qual $\operatorname{Pic}_{X/S}^\tau$ não é
plano sobre $S$, fazendo variar infinitesimalmente a superfície de Igusa. O
ponto considerado por Mumford encontra-se, aliás, em
$\operatorname{Pic}_{X/S}^\rho$, e continua possível (para
$f:X\longrightarrow S$ liso) que $\operatorname{Pic}_{X/S}^\tau$ seja plano
sobre $S$ nos pontos de $\operatorname{Pic}_{X/S}^\sigma$; o conferencista
duvida, contudo, que isso seja sempre assim, mesmo restringindo-se aos pontos de
$\operatorname{Pic}_{X/S}^\circ$ e supondo que $S$ seja o espectro de um anel
de valorização discreta. A questão está, aliás, ligada ao estudo dos pontos
fixos de um esquema abeliano sob um grupo finito de automorfismos, situação
para a qual parecem faltar exemplos.

Parece que, mesmo limitando-nos a $f:X\longrightarrow S$ liso e projetivo, os
resultados de regularidade local sobre $\operatorname{Pic}_{X/S}$ enunciados no
presente número e as conjecturas assinaladas em (2.8) esgotam aproximadamente o
que se pode dizer a esse respeito sem hipóteses mais particulares sobre a
natureza das fibras de $f$. Lembremos, contudo, que, se as fibras geométricas
de $\operatorname{Pic}_{X/S}$ são reduzidas e sem componente aditiva, então
decorre de (1.8) e (2.5) que $\operatorname{Pic}_{X/S}$ é liso sobre $S$ nos
pontos de $\operatorname{Pic}_{X/S}^\sigma$ quando $S$ é reduzido; esse
resultado vale até mesmo, se $f:X\longrightarrow S$ é normal, sem hipótese
sobre $S$, como veremos em (3.5). Assinalemos, a esse respeito:

\result{PROPOSIÇÃO}{2.10}

\noindent(i) Se $\operatorname{Pic}_{X/S}$ é liso (resp. plano) sobre $S$ nos
pontos da seção unidade, também o é em todos os pontos de
$\operatorname{Pic}_{X/S}^\sigma$ e nos pontos de toda seção de
$\operatorname{Pic}_{X/S}$ sobre $S$.
\src{236}
\noindent(ii) Seja $s\in S$ tal que
\[
H^2(X_s,\mathcal O_{X_s})=0,
\]
então existe uma vizinhança aberta $U$ de $s$ tal que
$\operatorname{Pic}_{X/S}|U$ seja liso sobre $U$.

\noindent(iii) Seja $X$ um esquema próprio sobre um corpo; então tem-se
\[
\dim\operatorname{Pic}_{X/k}\leq\dim H^1(X,\mathcal O_X),
\]
com igualdade se e somente se $\operatorname{Pic}_{X/k}$ é liso sobre $k$;
isso ocorre sempre se $k$ tem característica zero.

(i) decorre de (2.5) e (1.12), (ii) do critério infinitesimal de lisura e de um
cálculo de obstruções bem conhecido, tendo em conta que (pelo «teorema de
semicontinuidade») a hipótese feita em $s$ continuará verificada nos pontos
vizinhos. Por fim, (iii) decorre do fato de que $H^1(X,\mathcal O_X)$ é
isomorfo ao espaço tangente de Zariski no elemento neutro de
$\operatorname{Pic}_{X/k}$; a última afirmação é um caso particular de um
teorema de Cartier segundo o qual um «grupo formal» em característica zero é
formalmente liso sobre $k$.

\fsection{3}{O subesquema abeliano canônico de
\texorpdfstring{$\operatorname{Pic}_{X/S}$}{Pic(X/S)} e o esquema de Albanese}

\result{PROPOSIÇÃO}{3.1} Sejam $k$ um corpo e $G$ um esquema em grupos de tipo
finito sobre $k$, comutativo e «sem componente aditiva». Então
$G_{\mathrm{red}}^\circ$ é separável sobre $k$ e, por conseguinte, é um
subesquema em grupos liso sobre $k$.

Sendo isso trivial se $k$ é perfeito, em particular para $G_{\bar k}$, onde
$\bar k$ é o fecho algébrico de $k$, basta ver que
$(G_{\bar k}^\circ)_{\mathrm{red}}$ provém de um subesquema de $G$. Ora, a
hipótese sobre $G_{\bar k}$ (ausência de componente aditiva) resulta facilmente que existe um inteiro
$m$ tal que $(G_{\bar k}^\circ)_{\mathrm{red}}$ seja a imagem esquemática do
homomorfismo de elevação à $m$-ésima potência em $G_{\bar k}$. Como esse último
homomorfismo provém do homomorfismo análogo para $G$, a imagem esquemática deste
último responde à questão.

\result{COROLÁRIO}{3.2} Seja $X$ um esquema normal e próprio sobre $k$ tal que
$\operatorname{Pic}_{X/k}$ exista; então existe um subesquema abeliano $A$ de
$\operatorname{Pic}_{X/k}$ cujo conjunto subjacente é
$\operatorname{Pic}_{X/k}^\circ$.

Com efeito, em virtude de (2.1), (ii), sendo
$\operatorname{Pic}_{X/k}^\circ$ próprio sobre $k$, ele satisfaz a condição de
(3.1).
\src{237}
O resultado precedente mostra, portanto, que, em certos casos, a clássica
«variedade de Picard» (que, na teoria atual, seria denotada por
$(\operatorname{Pic}_{X/k})_{\mathrm{red}}^\circ$) «é definida sobre $k$»,
sem supor que o corpo $k$ seja perfeito.

Seja agora $f:X\longrightarrow S$ um esquema relativo próprio e plano, com
\[
\mathcal O_S\xrightarrow{\sim}f_*(\mathcal O_X)
\]
para simplificar, tal que $\operatorname{Pic}_{X/S}$ exista e que
$\operatorname{Pic}_{X/S}^\circ$ seja próprio sobre $S$. Suporemos, além
disso, para (3.3), (ii), que existe um aberto de $\operatorname{Pic}_{X/S}$
contendo $\operatorname{Pic}_{X/S}^\circ$ que seja quase-projetivo sobre $S$;
essa condição é satisfeita, como vimos, se $f$ é projetivo e tem fibras
geométricas separáveis e irredutíveis. Recordemos que se chama esquema
abeliano sobre $S$ um esquema em grupos sobre $S$, próprio e liso sobre
$S$, com fibras geométricas conexas. Propomo-nos
examinar se existe um subesquema em grupos $A$ de
$\operatorname{Pic}_{X/S}$ que seja um esquema abeliano e cujo conjunto
subjacente seja $\operatorname{Pic}_{X/S}^\circ$. Acabamos de ver que sempre
existe um se $S$ é o espectro de um corpo. Eis o que se sabe dizer no caso
geral considerado aqui:

\result{TEOREMA}{3.3} Sob as condições precedentes:

\noindent(i) Se existe um subesquema abeliano de
$\operatorname{Pic}_{X/S}$ cujo conjunto subjacente é
$\operatorname{Pic}_{X/S}^\circ$, ele é único. Sua formação é, por
conseguinte, compatível com as mudanças de base.

\noindent(ii) Para que exista tal subesquema abeliano, basta que exista um
após toda mudança de base $S'\longrightarrow S$, em que $S'$ é artiniano
local; se $S$ é o espectro de um anel local, basta até mesmo testar com os
$S'=\operatorname{Spec}(A_n)$, em que $A_n=A/\mathfrak m^{n+1}$. Se $S$ é
reduzido, basta igualmente testar com as mudanças de base
$S'\longrightarrow S$, em que $S'$ é o espectro de um anel de valorização
discreta (completo com corpo residual algebricamente fechado, se assim se
desejar).

\noindent(iii) Suponhamos que $A$ exista e seja
$B=\operatorname{Alb}^0(X/S)$ o esquema abeliano dual (isto é,
$B=\operatorname{Pic}_{A/S}^\circ$ [7]). Pode-se então construir
canonicamente um espaço homogêneo principal
$P=\operatorname{Alb}^1(X/S)$ sob $B$ e um $S$-morfismo
$X\longrightarrow P$ que seja universal para os $S$-morfismos de $X$ em
esquemas para-abelianos (isto é, espaços homogêneos principais sob esquemas
abelianos). A formação de $\operatorname{Alb}^0(X/S)$,
$\operatorname{Alb}^1(X/S)$ e
$X\longrightarrow\operatorname{Alb}^1(X/S)$ comuta com a mudança de base.

Esbocemos a demonstração:

(i) é uma propriedade geral de rigidez dos subesquemas abelianos dos esquemas
em grupos comutativos: se dois desses subesquemas coincidem como conjuntos em
um ponto $s\in S$, eles coincidem sobre toda a componente
\src{238}
conexa de $s$ [7]. (Este resultado generaliza um teorema clássico de Chow.)

(ii) Utilizando os esquemas de Hilbert, vê-se que o funtor que, a todo $S'$
sobre $S$, associa o conjunto (reduzido a 1 ou 0 elementos) dos subesquemas
abelianos canônicos de $(\operatorname{Pic}_{X/S})\times_S S'$ é
representável por um esquema $T$ de tipo finito sobre $S$. Em virtude de (i),
$T\longrightarrow S$ é um monomorfismo e, em virtude de (3.2), é
sobrejetivo. Dizer que existe um subesquema abeliano canônico de
$\operatorname{Pic}_{X/S}$ significa que $T$ tem uma seção sobre $S$, ou ainda
que $T\longrightarrow S$ é um isomorfismo. Isso equivale também a dizer que
$T\longrightarrow S$ é étale ou, no caso em que $S$ é reduzido, que
$T\longrightarrow S$ é próprio. Daí resulta imediatamente (ii).

(iii) Utilizando simplesmente a definição de $\operatorname{Pic}_{X/S}$,
constata-se que, para todo esquema abeliano $C$ sobre $S$, o dado de um
$S$-morfismo de $X$ em um espaço homogêneo principal sob $C$ equivale ao dado
de um homomorfismo de grupos
$C'\longrightarrow\operatorname{Pic}_{X/S}$, em que $C'$ é o esquema abeliano
dual de $C$. Ora, se existe o subesquema abeliano canônico $A$ de
$\operatorname{Pic}_{X/S}$, esses homomorfismos fatoram-se necessariamente por
$A$ (como se pode ver ainda utilizando os pontos de ordem finita). Daí resulta
imediatamente (iii).

\result{OBSERVAÇÕES}{3.4} Denotaremos por
$\operatorname{Pic}_{X/S}^{\circ\circ}$ o subesquema abeliano canônico de
$\operatorname{Pic}_{X/S}$, caso exista. Infelizmente, nem sempre é esse o
caso, como se pode ver fazendo variar infinitesimalmente a superfície de Igusa
(pela deformação modular de primeira ordem). É possível, contudo, que
$\operatorname{Pic}_{X/S}^{\circ\circ}$ exista ao menos se $S$ é reduzido ou,
o que equivale ao mesmo segundo (ii), se $S$ é o espectro de um anel de
valorização discreta. Sejam então $X_0,X_1$ a fibra especial e
genérica de $X$, e seja
\[
A_1=\operatorname{Pic}_{X_1/K}^\circ,
\]
em que $K$ é o corpo de frações do anel de valorização $V$. Em virtude de
Koizumi, existe um esquema abeliano $A$ sobre $S$, essencialmente único, cuja
fibra genérica é $A_1$, e vê-se facilmente, como em (2.1), (i) (supondo
doravante $X$ liso sobre $S$), que o morfismo identidade de $A_1$ se prolonga em um
morfismo
\[
A\longrightarrow\operatorname{Pic}_{X/S}.
\]
Daí um homomorfismo
\[
\tag{*}A_0\longrightarrow\operatorname{Pic}_{X_0/k}^{\circ\circ},
\]
que se demonstra facilmente ser um homomorfismo sobrejetivo cujo núcleo é um
$p$-grupo finito, em que $p$ é o expoente característico do corpo residual $k$
\src{239}
(ainda utilizando os pontos de ordem finita). Isto posto, as condições
seguintes sobre $X/S$ são equivalentes:

\medskip
\noindent a. O homomorfismo precedente (*) é um isomorfismo (o que Shimura
exprimiria dizendo que a formação da «variedade de Picard» é «compatível com a
especialização»).

\noindent b. $\operatorname{Pic}_{X/S}^{\circ\circ}$ existe (então não será
outro senão $A$).

\noindent c. (Para lembrar) Os
$\operatorname{Pic}_{X_n/S}^{\circ\circ}$ existem.

Segundo a observação feita sobre o núcleo de (*), a condição (a) é satisfeita
se a característica residual é nula, mas esse resultado será generalizado de
modo notável em (3.5).

Naturalmente, se $\operatorname{Pic}_{X/S}$ é liso sobre $S$ nos pontos de
$\operatorname{Pic}_{X/S}^\circ$, este último é
aberto em $\operatorname{Pic}_{X/S}$ (cf. (1.7)) e, portanto, munido da
estrutura induzida, é um subesquema abeliano de
$\operatorname{Pic}_{X/S}$, logo igual a
$\operatorname{Pic}_{X/S}^{\circ\circ}$, que existe nesse caso. Mas há muito
mais:

\result{TEOREMA}{3.5} Sob as condições de (3.3), seja $s\in S$ tal que
$\operatorname{Pic}_{X_s/k(s)}$ seja liso sobre $k(s)$ (ou, o que equivale ao
mesmo,
\[
\dim\operatorname{Pic}_{X_s/k(s)}=\dim H^1(X_s,\mathcal O_{X_s})).
\]
Então existe uma vizinhança aberta $U$ de $s$ tal que
$\operatorname{Pic}_{X/S}$ seja liso sobre $S$ em os pontos de
$\operatorname{Pic}_{X/S}^\circ|U$, que é, portanto, um subesquema abeliano aberto em
$\operatorname{Pic}_{X/S}|U$. A fortiori,
$\operatorname{Pic}_{X|U/U}^{\circ\circ}$ existe.

Indiquemos o princípio da demonstração. O que precede reduz-nos ao caso em que
$S$ é o espectro de um anel artiniano local $A$, e raciocina-se por indução
sobre a ordem infinitesimal de $A$. Pode-se, portanto, supor que
$\operatorname{Pic}_{X_n/A_n}^\circ$ seja liso sobre $A_n$ e limitar-se a
provar que $\operatorname{Pic}_{X_{n+1}/A_{n+1}}^\circ$ é liso sobre
$A_{n+1}$. Constata-se que, para isso, basta construir um esquema abeliano
$P_{n+1}$ sobre $A_{n+1}$ que prolongue
\[
P_n=\operatorname{Pic}_{X_n/A_n}^\circ,
\]
e, ao mesmo tempo, um Módulo invertível $\mathcal L_{n+1}$ sobre
$X_{n+1}\times_{A_{n+1}}P_{n+1}$ que prolongue o Módulo invertível
$\mathcal L_n$ sobre $X_n\times_{A_n}P_n$ que intervém na definição do esquema
de Picard $\operatorname{Pic}_{X_n/A_n}$ como solução de um problema universal.
(N.B. --- Podemos supor que $X$ seja munido de uma seção sobre $S$~...) Para
realizar a construção, deve-se utilizar o seguinte resultado-chave: todo
esquema abeliano definido sobre um quociente de anel artiniano local pode ser
prolongado (em outros termos, o «esquema formal dos Módulos» absoluto ([3], II)
para um esquema abeliano sobre um corpo algebricamente fechado é liso sobre o
anel dos
\src{240}
vetores de Witt sobre $k$); esse próprio resultado é obtido utilizando as
propriedades formais gerais da obstrução ao levantamento e as operações do
grupo. Admitido esse resultado, começa-se por prolongar $P_n$ de uma maneira
qualquer a $P_{n+1}$; encontra-se então uma obstrução ao levantamento de
$\mathcal L_n$, obstrução que pertence a
\[
H^2(X_0\times P_0,\mathcal O_{X_0\times P_0})\otimes_k V
\qquad(\text{onde }V=\mathfrak m^{n+1}/\mathfrak m^{n+2}),
\]
e até mesmo ao subespaço
\[
H^1(X_0,\mathcal O_{X_0})\otimes
H^1(P_0,\mathcal O_{P_0})\otimes_k V
\]
(tendo em conta que a restrição de $\mathcal L_n$ aos dois fatores $X_n$ e
$P_n$ é trivial). Ora, este último espaço não é outro senão
\[
H^1(P_0,\mathcal T_{P_0/k})\otimes V,
\]
em que $\mathcal T_{P_0/k}$ é o feixe tangente a $P_0/k$; portanto, é também o
espaço que exprime a indeterminação que havia no levantamento de $P_n$ a
$P_{n+1}$ ([3], II). Por conseguinte, pode-se corrigir esse levantamento (de
uma e uma só maneira, como convém) de modo a anular a obstrução ao levantamento
de $\mathcal L_n$.

\result{COROLÁRIO}{3.6} Sob as condições de (3.5),
$R^1f_*(\mathcal O_X)$ é um Módulo localmente livre sobre $S$ em uma
vizinhança de $s$, e sua formação comuta com a mudança de base.

Com efeito, esse Módulo não é outro senão o Módulo tangente a
$\operatorname{Pic}_{X/S}$ ao longo da seção unidade.

\result{OBSERVAÇÃO}{3.7} Utilizando o mesmo argumento que em (3.5), pode-se
mostrar que, se $S'$ é um subesquema de $S$ definido por um Ideal coerente
nilpotente e se se supõe apenas que $\operatorname{Pic}_{X'/S'}^\circ$ existe e
tem fibras lisas, então $\operatorname{Pic}_{X/S}^\circ$ existe
necessariamente e é um esquema abeliano sobre $S$. Isso permite construir
esquemas de Picard em certos casos apesar da ausência de uma hipótese
projetiva; por exemplo, o esquema abeliano dual de um esquema abeliano sobre um
anel artiniano sempre existe. Por outro lado, utilizando (3.6) no caso em que
$X$ é um esquema abeliano sobre $S$ e utilizando a estrutura conhecida de
$H^*(X_s,\mathcal O_{X_s})$ como álgebra exterior sobre
$H^1(X_s,\mathcal O_{X_s})$ (Rosenlicht--Serre), conclui-se que, para todo
$i$, $R^if_*(\mathcal O_X)$ é localmente livre e, mais precisamente, isomorfo
à $i$-ésima potência exterior de $R^1f_*(\mathcal O_X)$.

\result{OBSERVAÇÃO}{3.8} No caso de um morfismo projetivo liso
$f:X\longrightarrow S$, com $S$ reduzido e de características residuais
nulas, o resultado (3.6) era conhecido por métodos transcendentes, como
consequência da teoria de Hodge. De fato, todos os
$R^pf_*(\Omega^q_{X/S})$ são então localmente livres. Há, em contrapartida,
contraexemplos
\src{241}
no caso de características mistas para «$R^1f_*(\mathcal O_X)$ localmente
livre», com o auxílio das variedades de Serre ([8], p.~50). Parece não haver
nenhum contraexemplo em característica igual.

\result{COROLÁRIO}{3.8} Sob as condições de (3.5),
$\operatorname{Pic}_{X/S}|U$ é liso sobre $U$ em todos os pontos de
$\operatorname{Pic}_{X/S}^\sigma|U$.

Aplica-se (2.10), (i). Em particular, tendo em conta (2.10), (iii):

\result{COROLÁRIO}{3.9} Suponhamos que $S$ tenha características residuais
nulas. Então $\operatorname{Pic}_{X/S}^\tau$ é liso sobre $S$.

\result{OBSERVAÇÃO}{3.10} Deduz-se daí, por exemplo, se
$\operatorname{Pic}_{X/S}$ é também próprio sobre $S$, que o grupo de torção
dos grupos de Néron--Severi das fibras geométricas de $f$ permanece constante
sobre toda componente conexa de $S$ (o que, aliás, é evidente por métodos
transcendentes quando $f$ é liso e projetivo). Assinalemos que a utilização
direta de (2.5) permite mostrar, mais geralmente, que, se
$\operatorname{Pic}_{X/S}$ é próprio sobre $S$ (por exemplo,
$f:X\longrightarrow S$ liso e projetivo) e se $q$ é um número primo distinto
das características residuais de $S$, então a componente $q$-primária dos
grupos de torção de Néron--Severi das fibras geométricas de $X/S$ permanece
constante sobre toda componente conexa de $S$. Contudo, já não ocorre o mesmo
em característica $p>0$ para a componente $p$-primária da torção. Continua
possível, porém, que o posto sobre $k(s)$ de
\[
\operatorname{Pic}_{X_s/k(s)}^\tau/
\operatorname{Pic}_{X_s/k(s)}^{\circ\circ}=T_{X_s/k(s)}
\]
permaneça localmente constante; quando $S$ é reduzido, pode-se mostrar que isso
equivale a que $\operatorname{Pic}_{X/S}^{\circ\circ}$ exista e que
$\operatorname{Pic}_{X/S}^\tau$ seja plano sobre $S$, e basta testar a questão
para $S$ espectro de um anel de valorização discreta. Foi isso que verifiquei
nos poucos exemplos que examinei; mas, como o enunciado correspondente com $S$
artiniano é falso (cf. observações (2.9) e (3.4)), sem dúvida não convém
alimentar ilusões excessivas.

\fsection{4}{O teorema de finitude para o esquema de Picard}

Seja $f:X\longrightarrow S$ um morfismo projetivo e plano tal que
$\operatorname{Pic}_{X/S}$ exista; então a consideração dos «polinômios de
Hilbert» $Q\in\mathbf Q[t]$ permite decompor $\operatorname{Pic}_{X/S}$ em uma
soma de abertos $\operatorname{Pic}_{X/S}^Q$. Quando não se fazem hipóteses
adicionais, que assegurem, por exemplo, que $\operatorname{Pic}_{X/S}$ seja
separado sobre $S$, não será verdade, em geral, que esses abertos sejam de tipo
finito sobre $S$; obtém-se um contraexemplo quando $X$ é «uma cônica que
degenera em duas retas». É
\src{242}
possível, contudo, que assim seja se $f$ é separável e tem fibras
geometricamente irredutíveis. A questão está ligada à de saber se
$\operatorname{Pic}_{X/S}^\tau$ é de tipo finito sobre $S$, o que talvez seja
verdade sem hipótese alguma sobre as fibras de $X/S$. Quando
$f:X\longrightarrow S$ é liso, note-se que
$\operatorname{Pic}_{X/S}^\tau$ está contido em um dos
$\operatorname{Pic}_{X/S}^Q$ (graças ao fato de que, sobre uma variedade
projetiva não singular, a equivalência «por torção» é mais fina que (na
realidade, «idêntica à», segundo Matsusaka) a equivalência numérica para os
divisores), portanto é de tipo finito sobre $S$ se os
$\operatorname{Pic}_{X/S}^Q$ o são. Note-se que as questões de finitude que
acabamos de levantar conservam um sentido evidente mesmo sem supor a existência
de $\operatorname{Pic}_{X/S}$, pois se exprimem dizendo que certas famílias de
Módulos invertíveis são «limitadas» no sentido de [3], IV: consideremos, para
todo corpo algebricamente fechado $k$, os subesquemas íntegros de
$\mathbf P_k^r$, de dimensão $n$ e de grau $d$, e os Módulos invertíveis sobre
tais pré-esquemas, que tenham polinômio de Hilbert $Q$ (onde $r,n,d,Q$ são
dados); trata-se de mostrar que a família desses Módulos (considerados como
Módulos coerentes sobre os $\mathbf P_k^r$) é limitada, isto é, pode ser
parametrizada por um esquema de tipo finito sobre $\mathbf Z$.

Utilizando o método de Matsusaka [6], uma demonstração bastante técnica (que
recorre aos «critérios de equivalência» sob uma forma conveniente) permite
responder afirmativamente quando nos restringimos às subvariedades não
singulares de $\mathbf P_k^r$. Mais precisamente, obtém-se o resultado seguinte:

\result{TEOREMA}{4.1} Sejam: $X\longrightarrow S$ um
morfismo liso e projetivo com fibras geometricamente conexas, sendo $S$
noetheriano, $\mathcal O_X(1)$ um Módulo muito amplo sobre $X$ relativamente a
$S$, $E$ um subconjunto de $\operatorname{Pic}_{X/S}$, correspondente a uma
família $(\mathcal L_i)$ de Módulos invertíveis sobre as fibras geométricas
$\overline X_{s_i}$ de $X/S$, $D_i$ um divisor (não necessariamente efetivo)
sobre $\overline X_{s_i}$ que define $\mathcal L_i$,
\[
a_n^{(i)}t^n+\cdots+a_0^{(i)}
\]
o polinômio de Hilbert de $\mathcal L_i$, $\xi=\xi_i$ um divisor que define
$\mathcal O_X(1)$, isto é, uma seção hiperplana. As condições seguintes são
equivalentes:

\noindent a. $Q$ é quase-compacto,
isto é, está contido em um aberto de tipo finito sobre $S$, isto é, a família
dos $(\mathcal L_i)$ é limitada.

\noindent b. $E$ está contido na reunião de um número finito dos conjuntos
$\operatorname{Pic}_{X/S}^Q$, $Q\in\mathbf Q[t]$, isto é, o conjunto dos
polinômios de Hilbert dos $\mathcal L_i$ é finito.

\noindent b'. (Se as fibras de $X/S$ têm todas a mesma dimensão $n$.) Os dois
coeficientes $a_{n-1}^{(i)}$ e $a_{n-2}^{(i)}$ dos polinômios de Hilbert dos
$\mathcal L_i$ permanecem em um conjunto finito.
\src{243}
\noindent b''. (Se as fibras de $X/S$ têm todas a mesma dimensão $n$.) Os
inteiros
\[
\xi^{n-1}D_i\qquad\text{e}\qquad\xi^{n-2}D_i^2
\]
(graus projetivos de $D_i$ e de $D_i^2$) permanecem majorados.

\result{COROLÁRIO}{4.2} Seja $f:X\longrightarrow S$ um morfismo liso e
projetivo com fibras geometricamente conexas; então os esquemas
$\operatorname{Pic}_{X/S}^Q$ e $\operatorname{Pic}_{X/S}^\tau$
($Q\in\mathbf Q[t]$) são projetivos sobre $S$.

Como são de tipo finito sobre $S$ segundo (4.1), pode-se aplicar (2.4).

\section*{BIBLIOGRAFIA}
\addcontentsline{toc}{section}{Bibliografia}

\begin{description}
\item[{[1]}] CHEVALLEY (Claude). -- Sur la théorie de Picard,
\emph{Amer. J. of Math.}, t. 82, 1960, p. 435--490.

\item[{[2]}] GROTHENDIECK (A.) e DIEUDONNÉ (J.). -- Éléments de géométrie
algébrique, Capítulo I e seguintes. -- Paris, Presses universitaires de France,
1960--1961 (Institut des Hautes Études Scientifiques, 4, 8, 11, \ldots).

\item[{[3]}] GROTHENDIECK (Alexander). -- Technique de descente et théorèmes
d'existence en géométrie algébrique, I--V, \emph{Séminaire Bourbaki}, t. 12,
1959/60, n\textsuperscript{o} 190, 29 p. e n\textsuperscript{o} 195, 22 p.;
t. 13, 1960/61, n\textsuperscript{o} 212, 20 p. e
n\textsuperscript{o} 221, 28 p.; t. 14, 1961/62,
n\textsuperscript{o} 232, 19 p.

\item[{[4]}] GROTHENDIECK (Alexander). -- Séminaire de géométrie algébrique,
1960/61. -- Paris, Institut des Hautes Études Scientifiques, 1961.

\item[{[5]}] GROTHENDIECK (Alexander). -- Séminaire de géométrie algébrique,
1961/62, redigido por um grupo de ouvintes. -- Paris, Institut des Hautes Études
Scientifiques (a publicar).

\item[{[6]}] MATSUSAKA (T.). -- The criteria for algebraic equivalence and the
torsion group, \emph{Amer. J. Math.}, t. 79, 1957, p. 53--66.

\item[{[7]}] MUMFORD (D.) e TATE (J.). -- Séminaire de géométrie algébrique,
Harvard University, semestre da primavera de 1962 (a publicar).

\item[{[8]}] SERRE (Jean-Pierre). -- Sur la topologie des variétés algébriques
en caractéristique $p$, \emph{Symposium internacional de topología algebraica}
[1956, México], p. 24--53. -- México, Universidad Nacional Autónoma, 1958.
\end{description}

\par

\def\FgaUnitKey{e236}\clearpage
\chapter*{Errata e complemento à exposição n\textsuperscript{o} 236}
\addcontentsline{toc}{chapter}{Errata e complemento à exposição
n\textsuperscript{o} 236}
\markboth{Errata da exposição 236}{Errata da exposição 236}
\sid{e236}
\src{303}

\noindent\textbf{8. TDTE VI: Os esquemas de Picard. Propriedades gerais.}

\begin{center}
[Seminário Bourbaki, t. 14, 1961/62, n\textsuperscript{o} 236, 23 p.]
\end{center}

\noindent\textbf{Página 12, linha 16.} --- Em vez de «reta afim», leia-se
«reta afim privada da origem».

\noindent\textbf{Linha 18.} --- Em vez de «$X[t]$», leia-se
«$X[t,t^{-1}]$».

\noindent\textbf{Página 16, proposição 3.1, 6\textsuperscript{a} linha.} ---
Em vez de «Ora, a hipótese ...», leia-se «Ora, da hipótese ...».

\noindent\textbf{Página 21, § 4: O teorema de finitude para o esquema de
Picard.} --- As questões de finitude do gênero das levantadas neste parágrafo
foram resolvidas quase totalmente desde a redação da presente exposição.
Indiquemos os principais fatos agora conhecidos nessa direção. Para simplificar
os enunciados, supomos implicitamente que todos os pré-esquemas de Picard que
neles intervêm existem, embora uma modificação evidente desses enunciados
permita eliminar toda hipótese de existência. Doravante, $S$ designa um esquema
noetheriano, e $X,Y$, esquemas próprios sobre $S$.

\medskip
\noindent(i) Seja $f:X\longrightarrow Y$ um morfismo sobrejetivo; então
\[
f^*: \operatorname{Pic}_{Y/S}\longrightarrow\operatorname{Pic}_{X/S}
\]
é um morfismo de tipo finito.
\src{304}
\noindent(ii) Suponhamos $Y$ projetivo sobre $S$, munido de um Módulo
invertível amplo relativamente a $S$, e seja $X$ o esquema dos zeros de uma
seção qualquer desse Módulo, $f:X\longrightarrow Y$ a imersão canônica.
Suponhamos, enfim, que as componentes irredutíveis das fibras de $Y/S$ tenham
dimensão $\geq3$. Então
\[
f^*: \operatorname{Pic}_{Y/S}\longrightarrow\operatorname{Pic}_{X/S}
\]
é de tipo finito.

\noindent(iii) Suponhamos que $X$ seja projetivo sobre $S$ e que todas as suas
fibras geométricas sejam íntegras e de dimensão $n$. Seja
$\mathcal O_X(1)$ um Módulo invertível amplo sobre $X$, que permita definir
polinômios de Hilbert. Seja $M$ um subconjunto de
$\operatorname{Pic}_{X/S}$; então $M$ é quase-compacto se, e somente se, nos
polinômios de Hilbert
\[
a_0x^n+a_1x^{n-1}+a_2x^{n-2}+\cdots
\]
dos elementos de $M$, os coeficientes $a_1$ e $a_2$ permanecem limitados.

\noindent(iv) Para todo inteiro $n\neq0$, o homomorfismo de elevação à
$n$-ésima potência no pré-esquema em grupos $\operatorname{Pic}_{X/S}$ é um
morfismo de tipo finito.

Notemos que (i) e (ii) significam também que (sob as hipóteses consideradas
nesses enunciados (i) e (ii)) um subconjunto $M$ de
$\operatorname{Pic}_{Y/S}$ é quase-compacto se, e somente se, sua imagem em
$\operatorname{Pic}_{X/S}$ o é. Conclui-se daí que um Módulo invertível
$\mathcal L$ sobre $Y$ é $\tau$-equivalente a zero se, e somente se, sua imagem
inversa sobre $X$ o é; em outros termos, $\operatorname{Pic}_{Y/S}^\tau$ é a
imagem inversa de $\operatorname{Pic}_{X/S}^\tau$. Em particular, para mostrar
que o primeiro pré-esquema é de tipo finito sobre $S$, basta prová-lo para o
segundo, pois $\operatorname{Pic}_{Y/S}\longrightarrow
\operatorname{Pic}_{X/S}$ é de tipo finito. Utilizando então (i), o lema de
Chow e (iii), obtemos:

\medskip
\noindent(v) $\operatorname{Pic}_{X/S}^\tau$ é de tipo finito sobre $S$.

De modo geral, a conjunção de (i) para um morfismo finito e de (ii) permite
reduzir, para a maioria das questões de finitude, ao caso em que $X/S$ tem
fibras geométricas íntegras normais de dimensão $\leq2$; muitas vezes, aplicando
(i) a um morfismo sobrejetivo não necessariamente finito e a resolução das
singularidades das superfícies algébricas (devida, em característica qualquer,
a Abhyankar), reduzimos ao caso em que $X/S$ é até mesmo liso, portanto tem
fibras geométricas não singulares de dimensão 2. Isso permite, por exemplo,
tendo em conta (v) e a desigualdade de Picard--Igusa que majora o posto do grupo
de Néron--Severi de uma superfície projetiva não singular, provar a seguinte
generalização do teorema de finitude de Néron:

\medskip
\noindent(vi) Seja $X/S$ próprio sobre $S$ e, de resto, arbitrário; então os
grupos de Néron--Severi
\[
\operatorname{Pic}_{X_i/k_i}/\operatorname{Pic}_{X_i/k_i}^\circ
\]
das fibras geométricas $X_i/k_i$ de $X/S$
\src{305}
são de tipo finito, e seu posto e a ordem de seu subgrupo de torção permanecem
limitados.

O mesmo método de redução ao caso das superfícies não singulares e os teoremas
conhecidos nesse caso (a saber, o teorema de finitude de Néron e o fato de que a
forma de interseção sobre o grupo de Néron--Severi é não degenerada) implicam:

\medskip
\noindent(vii) Seja $X$ um esquema próprio sobre um corpo algebricamente
fechado. Existe um número finito de curvas fechadas íntegras $C_i$
($1\leq i\leq r$) em $X$, com a seguinte propriedade: para todo subconjunto
$M$ de $\operatorname{Pic}_{X/k}$, $M$ é quase-compacto se, e somente se, os
inteiros
\[
\deg_{C_i'}\mathcal L\qquad(\mathcal L\in M)
\]
permanecem limitados (onde $C_i'$ designa a normalização de $C_i$).

Pode-se tomar aqui, para $r$, o posto do grupo de Néron--Severi. Uma vez sabido
que este último é de tipo finito, (vii) reduz-se ao fato de que as formas
lineares sobre o grupo de Néron--Severi definidas por curvas $C$ em $X$ só se
anulam simultaneamente sobre os elementos de torção do grupo de Néron--Severi.
No caso em que $X$ é projetivo e não singular, esse resultado, assim como (v),
era devido a Matsusaka. Utilizando (vii), obtém-se facilmente a seguinte
caracterização dos Módulos invertíveis $\tau$-equivalentes a zero sobre $X$:

\medskip
\noindent(viii) Sejam $X/k$ um esquema próprio sobre um corpo e $\mathcal L$ um
Módulo invertível sobre $X$. Condições equivalentes:

\noindent a. $\mathcal L$ é $\tau$-equivalente a zero;

\noindent b. Para todo Módulo coerente $F$, tem-se
\[
\chi(F\otimes\mathcal L)=\chi(F),
\]
onde $\chi$ designa a característica de Euler--Poincaré;

\noindent b'. Como (b), mas com $F=\mathcal O_Y$, sendo $Y$ um subesquema
fechado íntegro de dimensão 1 em $X$;

\noindent c. Para todo $Y$ como acima, designando por $Y'$ a curva normalizada,
tem-se
\[
\deg\mathcal L_{Y'}=0.
\]

Se $X/k$ é projetivo e munido de um Módulo invertível amplo
$\mathcal O_X(1)$, as condições precedentes são também equivalentes às
seguintes:

\noindent d. Para todo inteiro $m$, $\mathcal L^{\otimes m}(1)$ é amplo.

\noindent e. (Se $X$ é íntegro.) Para todo par de inteiros $m,n$, tem-se
\[
\chi(\mathcal L^{\otimes m}(n))=\chi(\mathcal O_X(n))
\]
\src{306}
(isto é, (b) é verdadeiro tomando $F=\mathcal L^{\otimes m}(n)$).

Para a suficiência desta última condição, note-se que ela significa que os
polinômios de Hilbert dos $\mathcal L^{\otimes m}$ são todos iguais; portanto,
em virtude do critério de Mumford (iii), os $\mathcal L^{\otimes m}$ permanecem
em um subconjunto quase-compacto de $\operatorname{Pic}_{X/k}$, isto é, vale
(a). As condições (b), (b') e (c) devem ser consideradas variantes (sobre um
esquema próprio qualquer) da noção de equivalência numérica, definida
habitualmente sobre variedades projetivas não singulares. Para estas últimas,
a equivalência de (a) e (c) era evidentemente bem conhecida (Matsusaka).

O critério (e) implica também o resultado seguinte:

\medskip
\noindent(ix) Seja $f:X\longrightarrow S$ um morfismo projetivo e plano cujas
fibras geométricas são íntegras. Então $\operatorname{Pic}_{X/S}^\tau$ é aberto
e fechado em $\operatorname{Pic}_{X/S}$.

Limitar-nos-emos a alguns comentários sobre a demonstração dos resultados-chave
(i), (ii), (iii) (o resultado (iv) é um pouco apartado dos demais e se demonstra
utilizando somente (i) para os morfismos finitos sobrejetivos radiciais, mais
precisamente, para um morfismo de Frobenius). Para (i), utilizam-se de modo
essencial as ideias da descida não plana (ver, em particular, TDTE I, página 9).
Ocorre que (tendo em vista apenas resultados de finitude) a falta de critérios
de efetividade para dados de descida é inofensiva. Por outro lado, Mumford
demonstrou recentemente uma forma um pouco menos forte de (iii), a saber, o
critério que envolve todos os coeficientes do polinômio de Hilbert. Seu
argumento é extremamente simples e inspirado pela demonstração de um critério
de amplitude de Nakai (enunciado por este último para as superfícies não
singulares e generalizado por Mumford aos morfismos projetivos quaisquer).
Parece-me, aliás, que esse argumento só é válido mediante uma ligeira restrição
adicional sobre as fibras de $X/S$, a propriedade $(S_2)$ de Serre, verificada
se as fibras geométricas são normais. Utiliza-se então esse critério restrito na
demonstração de (ii): o critério (i) permite, com efeito, reduzir ao caso em que
$Y/S$ é plano com fibras geométricas íntegras e normais e, aplicando o critério
de Mumford, reduz-se facilmente ao caso em que $X/S$ satisfaz as mesmas
condições. Da hipótese de dimensão resulta então que as fibras geométricas de
$Y$ e de $X$ têm profundidade $\geq2$ em seus pontos fechados, o que permite
aplicar os «critérios de equivalência» sob a forma que lhes é dada em SGA 1962
(exposição XII) e concluir a demonstração de (ii). Uma vez que dispomos de (i) e
(ii), não é difícil eliminar do critério de Mumford toda hipótese de normalidade
sobre as fibras e demonstrá-lo também sob a forma mais forte dada em (iii).
\src{307}
Notemos, enfim, que a demonstração de (i) e (ii) mostra igualmente que, no
caso em que $S$ é o espectro de um corpo $k$, o morfismo
\[
\operatorname{Pic}_{Y/k}\longrightarrow\operatorname{Pic}_{X/k}
\]
é afim (e não apenas de tipo finito).

\par

\def\FgaUnitKey{com}\unit{Comentários de 1962}
\sid{com}
\markboth{Comentários de 1962}{Comentários de 1962}
\src{297}
\seminares{14.º ano, 1961/62, Suplemento\\Maio de 1962}
\paperhead{FUNDAMENTOS DA GEOMETRIA ALGÉBRICA.\\COMENTÁRIOS}%
  {por Alexander GROTHENDIECK}

\section*{Advertência}
\addcontentsline{toc}{section}{Advertência}

Apresentamos, no Seminário Bourbaki, entre 1957 e 1962, oito exposições
relativas aos Fundamentos da geometria algébrica. Com exceção da primeira,
essas exposições são redigidas no âmbito dos esquemas. Todos os resultados
enunciados encontrarão seu lugar nos «Éléments de géométrie algébrique» de Jean
DIEUDONNÉ e Alexander GROTHENDIECK. Contudo, o conteúdo de nenhuma dessas
exposições está atualmente coberto nem por um dos capítulos já publicados ou
em preparação, nem por qualquer outro livro ou artigo, e sem dúvida não o
estará ainda por alguns anos. Essa é a principal razão que nos levou a reunir
essas exposições, que colocarão à disposição dos usuários certo número de
noções e resultados-chave da teoria dos esquemas, enquanto se aguarda uma
redação sistemática. Além disso, a leitura desses textos permitirá ao leitor
familiarizar-se rapidamente com os referidos resultados e noções, sem que os
pormenores, necessariamente volumosos, de um tratado sistemático constituam
um obstáculo, e fornecerá as motivações indispensáveis para o estudo de tal
tratado.

Para comodidade do leitor, reunimos aqui certo número de comentários e
erratas, agrupados por exposições, que assinalarão em particular, de passagem,
os progressos realizados desde a redação do texto e indicarão algumas
referências adicionais.

Vários resultados apresentados nessas exposições foram tratados em detalhe no
Seminário de Geometria Algébrica do IHES, bem como em dois seminários
consecutivos realizados em Harvard em 1961/62, conduzidos o primeiro por mim e
o segundo por MUMFORD--TATE, para os quais estão atualmente em preparação
notas mimeografadas (redigidas pelo Sr. LICHTENBAUM).

\section*{Siglas utilizadas}
\addcontentsline{toc}{section}{Siglas utilizadas}

\begin{description}
\item[SGA :] Seminário de Geometria Algébrica, Institut des Hautes Études
Scientifiques, 1960--1962.

\src{298}
\item[SHG :] Seminário GROTHENDIECK, Harvard University, 1961/62.

\item[SHMT :] Seminário MUMFORD--TATE, Harvard University, 1961/62.

\item[EGA :] GROTHENDIECK (A.) e DIEUDONNÉ (J.). --- Éléments de géométrie
algébrique. --- Paris, Presses universitaires de France, 1960, 1961, …
(Institut des Hautes Études Scientifiques, Publications mathématiques, 4, 8,
11, …).

\item[TDTE :] GROTHENDIECK (Alexander). --- Technique de descente et théorèmes
d'existence en géométrie algébrique, I--VI, \emph{Séminaire Bourbaki}, t. 12,
1959/60, n.os 190 e 195; t. 13, 1960/61, n.os 212 e 221; t. 14, 1961/62,
n.os 232 e 236.
\end{description}

\par

\end{document}
